\documentclass[10pt,landscape]{article}
\usepackage[margin=0.6in]{geometry}
\usepackage[T1]{fontenc}
\usepackage{booktabs}
\usepackage[hidelinks]{hyperref}
\title{River LINE Timetable Book\\\large Current (NJ TRANSIT, Sept 2026) and proposed weekday timetables}
\author{ISE Senior Design Team 22}
\date{September 28, 2026}
\begin{document}
\maketitle
\noindent Times are departures, rounded to the minute. ``--'' means the trip does not serve that station.
\begin{itemize}
\item \textbf{NJT current} is transcribed from NJ TRANSIT's GTFS feed (release of Sept.~15, 2026; Mobility Database mirror mdb-509).
\item \textbf{Proposed T8--T3} come from the pattern optimizer described in the project report. Each keeps NJ TRANSIT's 150-minute round trip and spaces trains as evenly as the single track allows. Every scheduled meet falls on double track in the calibrated track model, all day.
\item \textbf{Proposals B8, B7, E5, E6, G8, S8, R8, W7} (revision 3) change one design choice each: meet buffers and layover length (B8, B7, G8, R8), the midday service (E5, E6), holds before single track (S8), and summer running times (W7). All are conflict-free. B8 is the recommended weekday timetable; see Section~7 of the project report for their simulated punctuality.
\item \textbf{Uniform U8--U4} use a single headway in each period. Each period is conflict-free, but the changeovers between peak and off-peak produce the conflicts listed below and need retiming before use.
\item All proposed timetables are planning drafts: running times are NJ TRANSIT's scheduled times, and siding locations are calibrated from NJ TRANSIT's own meets rather than taken from an official track chart.
\end{itemize}
\begin{center}\small
\begin{tabular}{@{}lrrrrrrrr@{}}\toprule
Timetable & Trains peak/mid & Trips & Train-h & AM gap $\to$Trenton & AM gap $\to$Camden & Midday gap & Avg wait & Meet conflicts\\
 & & & & mean/max (min) & mean/max (min) & mean (min) & (min) & \\\midrule
NJT current (Sept 2026) & 8/6 & 87 & 109 & 17.5/30 & 18.2/30 & 30.0 & 12.9 & 1 \\
Proposed T8 (8 trains peak, 6 midday) & 8/6 & 88 & 109 & 18.4/21 & 18.8/24 & 26.2 & 12.9 & 0 \\
Proposed T7 (7 peak, 6 midday) & 7/6 & 80 & 99 & 21.8/26 & 21.3/30 & 25.6 & 14.0 & 0 \\
Proposed T6 (6 peak, 5 midday) & 6/5 & 69 & 86 & 25.0/26 & 24.7/31 & 28.3 & 15.5 & 0 \\
Proposed T5 (5 all day) & 5/5 & 62 & 77 & 30.2/31 & 30.2/31 & 30.0 & 16.0 & 0 \\
Proposed T4 (4 all day) & 4/4 & 48 & 60 & 37.7/38 & 37.5/38 & 37.5 & 21.1 & 0 \\
Proposed T3 (3 all day) & 3/3 & 37 & 46 & 51.5/52 & 46.3/52 & 49.7 & 27.2 & 0 \\
Proposal B8: 8 trains, 160-min cycle, 2.5-min meet buffers, 9-min layovers & 8/6 & 83 & 109 & 19.3/25 & 21.2/31 & 26.7 & 13.5 & 0 \\
Proposal B7: 7 trains, 160-min cycle, 2.5-min meet buffers, 9-min layovers & 7/6 & 76 & 100 & 22.5/27 & 22.6/32 & 27.2 & 14.6 & 0 \\
Proposal E5: 8 peak, 5 midday (frees a car) & 8/5 & 80 & 99 & 17.9/25 & 20.3/28 & 29.2 & 13.3 & 0 \\
Proposal E6: 8 peak, 6 midday, even midday & 8/6 & 87 & 108 & 18.1/23 & 19.8/26 & 23.2 & 12.7 & 0 \\
Proposal G8: 8 trains, 2-min meet buffers & 8/6 & 88 & 109 & 19.4/27 & 16.9/29 & 24.4 & 12.3 & 0 \\
Proposal S8: 8 trains, slack held before single track & 8/6 & 87 & 108 & 18.3/23 & 19.6/26 & 23.3 & 12.8 & 0 \\
Proposal R8: 8 trains, 160-min cycle, 8-min layovers & 8/6 & 83 & 109 & 19.6/25 & 19.5/26 & 26.8 & 13.3 & 0 \\
Proposal W7: summer (+6\% run time), 7 trains & 7/5 & 71 & 93 & 22.6/29 & 22.1/32 & 31.9 & 15.4 & 0 \\
Uniform U8: every 21 min peak, 30 off-peak & 8/5 & 74 & 98 & 21.0/21 & 20.0/21 & 30.0 & 14.2 & 7 \\
Uniform U7: every 23 min peak, 30 off-peak & 8/6 & 71 & 100 & 23.0/23 & 22.6/23 & 30.0 & 14.7 & 6 \\
Uniform U6: every 26 min peak, 30 off-peak & 7/6 & 66 & 96 & 26.0/26 & 24.7/26 & 30.0 & 15.7 & 4 \\
Uniform U5: every 29 min peak, 30 off-peak & 6/5 & 63 & 77 & 29.0/29 & 23.8/29 & 30.0 & 16.8 & 6 \\
Uniform U4: every 38 min peak, 38 off-peak & 4/4 & 50 & 63 & 38.0/38 & 38.0/38 & 38.0 & 19.8 & 0 \\
\bottomrule\end{tabular}\end{center}
\noindent Avg wait: random passenger arrivals at Walter Rand, weighted by an assumed hourly demand profile. Meet conflicts: pairs of opposing trains scheduled on the same single-track section (calibrated track model, 0.5-min buffer).
\clearpage
\tableofcontents
\clearpage
\section{NJT current (Sept 2026)}\label{tt:njt_current}
\subsection*{To Trenton (read down) --- trips 1--13 of 42}
\begin{center}\footnotesize\setlength{\tabcolsep}{3pt}\begin{tabular}{@{}lrrrrrrrrrrrrr@{}}\toprule
Station & 1 & 2 & 3 & 4 & 5 & 6 & 7 & 8 & 9 & 10 & 11 & 12 & 13 \\\midrule
Entertainment Center & 5:45 & -- & -- & 6:00 & 6:13 & 6:43 & 6:58 & 7:13 & 7:43 & 7:58 & 8:13 & 8:28 & 8:58 \\
Aquarium & 5:47 & -- & -- & 6:02 & 6:15 & 6:45 & 7:00 & 7:15 & 7:45 & 8:00 & 8:15 & 8:30 & 9:00 \\
Cooper St/Rutgers & 5:49 & -- & -- & 6:04 & 6:17 & 6:47 & 7:02 & 7:17 & 7:47 & 8:02 & 8:17 & 8:32 & 9:02 \\
Walter Rand TC & 5:55 & 5:51 & -- & 6:08 & 6:21 & 6:51 & 7:06 & 7:21 & 7:51 & 8:06 & 8:21 & 8:36 & 9:06 \\
36th Street & 6:00 & 5:56 & -- & 6:13 & 6:26 & 6:56 & 7:11 & 7:26 & 7:56 & 8:11 & 8:26 & 8:41 & 9:11 \\
Pennsauken TC & 6:02 & -- & -- & 6:15 & 6:28 & 6:58 & 7:13 & 7:28 & 7:58 & 8:13 & 8:28 & 8:43 & 9:13 \\
Route 73/Pennsauken & -- & 5:59 & -- & 6:19 & 6:31 & 7:01 & 7:16 & 7:31 & 8:01 & 8:16 & 8:31 & 8:46 & 9:16 \\
Palmyra & 6:06 & 6:01 & -- & 6:21 & 6:34 & 7:04 & 7:19 & 7:34 & 8:04 & 8:19 & 8:34 & 8:49 & 9:19 \\
Riverton & 6:08 & -- & -- & 6:22 & 6:36 & 7:06 & 7:21 & 7:36 & 8:06 & 8:21 & 8:36 & 8:51 & 9:21 \\
Cinnaminson & 6:10 & -- & -- & 6:24 & 6:38 & 7:08 & 7:23 & 7:38 & 8:08 & 8:23 & 8:38 & 8:53 & 9:23 \\
Riverside & -- & 6:07 & -- & 6:28 & 6:43 & 7:13 & 7:28 & 7:43 & 8:13 & 8:28 & 8:43 & 8:58 & 9:28 \\
Delanco & 6:16 & -- & -- & 6:30 & 6:45 & 7:15 & 7:30 & 7:45 & 8:15 & 8:30 & 8:45 & 9:00 & 9:30 \\
Beverly/Edgewater Park & 6:20 & -- & -- & 6:34 & 6:49 & 7:19 & 7:34 & 7:49 & 8:19 & 8:34 & 8:49 & 9:04 & 9:34 \\
Burlington South & 6:25 & 6:13 & 5:56 & 6:38 & 6:53 & 7:23 & 7:38 & 7:53 & 8:23 & 8:38 & 8:53 & 9:08 & 9:38 \\
Burlington Towne Ctr & 6:27 & 6:15 & 5:58 & 6:40 & 6:55 & 7:25 & 7:40 & 7:55 & 8:25 & 8:40 & 8:55 & 9:10 & 9:40 \\
Florence & 6:33 & 6:20 & 6:04 & 6:46 & 7:01 & 7:31 & 7:46 & 8:01 & 8:31 & 8:46 & 9:01 & 9:16 & 9:46 \\
Roebling & 6:36 & -- & 6:07 & 6:49 & 7:04 & 7:34 & 7:49 & 8:04 & 8:34 & 8:49 & 9:04 & 9:19 & 9:49 \\
Bordentown & 6:42 & 6:27 & 6:12 & 6:54 & 7:09 & 7:39 & 7:54 & 8:09 & 8:39 & 8:54 & 9:09 & 9:24 & 9:54 \\
Cass Street & 6:48 & -- & 6:18 & 7:00 & 7:15 & 7:45 & 8:00 & 8:15 & 8:45 & 9:00 & 9:15 & 9:30 & 10:00 \\
Hamilton Avenue & 6:50 & -- & 6:20 & 7:02 & 7:17 & 7:47 & 8:02 & 8:17 & 8:47 & 9:02 & 9:17 & 9:32 & 10:02 \\
Trenton TC & 6:52 & 6:36 & 6:22 & 7:05 & 7:20 & 7:50 & 8:05 & 8:20 & 8:50 & 9:05 & 9:20 & 9:35 & 10:05 \\
\bottomrule\end{tabular}\end{center}
\subsection*{To Trenton (read down) --- trips 14--26 of 42}
\begin{center}\footnotesize\setlength{\tabcolsep}{3pt}\begin{tabular}{@{}lrrrrrrrrrrrrr@{}}\toprule
Station & 14 & 15 & 16 & 17 & 18 & 19 & 20 & 21 & 22 & 23 & 24 & 25 & 26 \\\midrule
Entertainment Center & 9:28 & 9:58 & 10:28 & 10:58 & 11:28 & 11:58 & 12:28 & 12:58 & 13:28 & 13:58 & 14:28 & 14:58 & 15:28 \\
Aquarium & 9:30 & 10:00 & 10:30 & 11:00 & 11:30 & 12:00 & 12:30 & 13:00 & 13:30 & 14:00 & 14:30 & 15:00 & 15:30 \\
Cooper St/Rutgers & 9:32 & 10:02 & 10:32 & 11:02 & 11:32 & 12:02 & 12:32 & 13:02 & 13:32 & 14:02 & 14:32 & 15:02 & 15:32 \\
Walter Rand TC & 9:36 & 10:06 & 10:36 & 11:06 & 11:36 & 12:06 & 12:36 & 13:06 & 13:36 & 14:06 & 14:36 & 15:06 & 15:36 \\
36th Street & 9:41 & 10:11 & 10:41 & 11:11 & 11:41 & 12:11 & 12:41 & 13:11 & 13:41 & 14:11 & 14:41 & 15:11 & 15:41 \\
Pennsauken TC & 9:43 & 10:13 & 10:43 & 11:13 & 11:43 & 12:13 & 12:43 & 13:13 & 13:43 & 14:13 & 14:43 & 15:13 & 15:43 \\
Route 73/Pennsauken & 9:46 & 10:16 & 10:46 & 11:16 & 11:46 & 12:16 & 12:46 & 13:16 & 13:46 & 14:16 & 14:46 & 15:16 & 15:46 \\
Palmyra & 9:49 & 10:19 & 10:49 & 11:19 & 11:49 & 12:19 & 12:49 & 13:19 & 13:49 & 14:19 & 14:49 & 15:19 & 15:49 \\
Riverton & 9:51 & 10:21 & 10:51 & 11:21 & 11:51 & 12:21 & 12:51 & 13:21 & 13:51 & 14:21 & 14:51 & 15:21 & 15:51 \\
Cinnaminson & 9:53 & 10:23 & 10:53 & 11:23 & 11:53 & 12:23 & 12:53 & 13:23 & 13:53 & 14:23 & 14:53 & 15:23 & 15:53 \\
Riverside & 9:58 & 10:28 & 10:58 & 11:28 & 11:58 & 12:28 & 12:58 & 13:28 & 13:58 & 14:28 & 14:58 & 15:28 & 15:58 \\
Delanco & 10:00 & 10:30 & 11:00 & 11:30 & 12:00 & 12:30 & 13:00 & 13:30 & 14:00 & 14:30 & 15:00 & 15:30 & 16:00 \\
Beverly/Edgewater Park & 10:04 & 10:34 & 11:04 & 11:34 & 12:04 & 12:34 & 13:04 & 13:34 & 14:04 & 14:34 & 15:04 & 15:34 & 16:04 \\
Burlington South & 10:08 & 10:38 & 11:08 & 11:38 & 12:08 & 12:38 & 13:08 & 13:38 & 14:08 & 14:38 & 15:08 & 15:38 & 16:08 \\
Burlington Towne Ctr & 10:10 & 10:40 & 11:10 & 11:40 & 12:10 & 12:40 & 13:10 & 13:40 & 14:10 & 14:40 & 15:10 & 15:40 & 16:10 \\
Florence & 10:16 & 10:46 & 11:16 & 11:46 & 12:16 & 12:46 & 13:16 & 13:46 & 14:16 & 14:46 & 15:16 & 15:46 & 16:16 \\
Roebling & 10:19 & 10:49 & 11:19 & 11:49 & 12:19 & 12:49 & 13:19 & 13:49 & 14:19 & 14:49 & 15:19 & 15:49 & 16:19 \\
Bordentown & 10:24 & 10:54 & 11:24 & 11:54 & 12:24 & 12:54 & 13:24 & 13:54 & 14:24 & 14:54 & 15:24 & 15:54 & 16:24 \\
Cass Street & 10:30 & 11:00 & 11:30 & 12:00 & 12:30 & 13:00 & 13:30 & 14:00 & 14:30 & 15:00 & 15:30 & 16:00 & 16:30 \\
Hamilton Avenue & 10:32 & 11:02 & 11:32 & 12:02 & 12:32 & 13:02 & 13:32 & 14:02 & 14:32 & 15:02 & 15:32 & 16:02 & 16:32 \\
Trenton TC & 10:35 & 11:05 & 11:35 & 12:05 & 12:35 & 13:05 & 13:35 & 14:05 & 14:35 & 15:05 & 15:35 & 16:05 & 16:35 \\
\bottomrule\end{tabular}\end{center}
\subsection*{To Trenton (read down) --- trips 27--39 of 42}
\begin{center}\footnotesize\setlength{\tabcolsep}{3pt}\begin{tabular}{@{}lrrrrrrrrrrrrr@{}}\toprule
Station & 27 & 28 & 29 & 30 & 31 & 32 & 33 & 34 & 35 & 36 & 37 & 38 & 39 \\\midrule
Entertainment Center & 15:58 & 16:13 & 16:28 & 16:43 & 17:13 & 17:28 & 17:58 & 18:13 & 18:28 & 18:43 & 18:58 & 19:28 & 19:58 \\
Aquarium & 16:00 & 16:15 & 16:30 & 16:45 & 17:15 & 17:30 & 18:00 & 18:15 & 18:30 & 18:45 & 19:00 & 19:30 & 20:00 \\
Cooper St/Rutgers & 16:02 & 16:17 & 16:32 & 16:47 & 17:17 & 17:32 & 18:02 & 18:17 & 18:32 & 18:47 & 19:02 & 19:32 & 20:02 \\
Walter Rand TC & 16:06 & 16:21 & 16:36 & 16:51 & 17:21 & 17:36 & 18:06 & 18:21 & 18:36 & 18:51 & 19:06 & 19:36 & 20:06 \\
36th Street & 16:11 & 16:26 & 16:41 & 16:56 & 17:26 & 17:41 & 18:11 & 18:26 & 18:41 & 18:56 & 19:11 & 19:41 & 20:11 \\
Pennsauken TC & 16:13 & 16:28 & 16:43 & 16:58 & 17:28 & 17:43 & 18:13 & 18:28 & 18:43 & 18:58 & 19:13 & 19:43 & 20:13 \\
Route 73/Pennsauken & 16:16 & 16:31 & 16:46 & 17:01 & 17:31 & 17:46 & 18:16 & 18:31 & 18:46 & 19:01 & 19:16 & 19:46 & 20:16 \\
Palmyra & 16:19 & 16:34 & 16:49 & 17:04 & 17:34 & 17:49 & 18:19 & 18:34 & 18:49 & 19:04 & 19:19 & 19:49 & 20:19 \\
Riverton & 16:21 & 16:36 & 16:51 & 17:06 & 17:36 & 17:51 & 18:21 & 18:36 & 18:51 & 19:06 & 19:21 & 19:51 & 20:21 \\
Cinnaminson & 16:23 & 16:38 & 16:53 & 17:08 & 17:38 & 17:53 & 18:23 & 18:38 & 18:53 & 19:08 & 19:23 & 19:53 & 20:23 \\
Riverside & 16:28 & 16:43 & 16:58 & 17:13 & 17:43 & 17:58 & 18:28 & 18:43 & 18:58 & 19:13 & 19:28 & 19:58 & 20:28 \\
Delanco & 16:30 & 16:45 & 17:00 & 17:15 & 17:45 & 18:00 & 18:30 & 18:45 & 19:00 & 19:15 & 19:30 & 20:00 & 20:30 \\
Beverly/Edgewater Park & 16:34 & 16:49 & 17:04 & 17:19 & 17:49 & 18:04 & 18:34 & 18:49 & 19:04 & 19:19 & 19:34 & 20:04 & 20:34 \\
Burlington South & 16:38 & 16:53 & 17:08 & 17:23 & 17:53 & 18:08 & 18:38 & 18:53 & 19:08 & 19:23 & 19:38 & 20:08 & 20:38 \\
Burlington Towne Ctr & 16:40 & 16:55 & 17:10 & 17:25 & 17:55 & 18:10 & 18:40 & 18:55 & 19:10 & 19:25 & 19:40 & 20:10 & 20:40 \\
Florence & 16:46 & 17:01 & 17:16 & 17:31 & 18:01 & 18:16 & 18:46 & 19:01 & 19:16 & 19:31 & 19:46 & 20:16 & 20:46 \\
Roebling & 16:49 & 17:04 & 17:19 & 17:34 & 18:04 & 18:19 & 18:49 & 19:04 & 19:19 & 19:34 & 19:49 & 20:19 & 20:49 \\
Bordentown & 16:54 & 17:09 & 17:24 & 17:39 & 18:09 & 18:24 & 18:54 & 19:09 & 19:24 & 19:39 & 19:54 & 20:24 & 20:54 \\
Cass Street & 17:00 & 17:15 & 17:30 & 17:45 & 18:15 & 18:30 & 19:00 & 19:15 & 19:30 & 19:45 & 20:00 & 20:30 & 21:00 \\
Hamilton Avenue & 17:02 & 17:17 & 17:32 & 17:47 & 18:17 & 18:32 & 19:02 & 19:17 & 19:32 & 19:47 & 20:02 & 20:32 & 21:02 \\
Trenton TC & 17:05 & 17:20 & 17:35 & 17:50 & 18:20 & 18:35 & 19:05 & 19:20 & 19:35 & 19:50 & 20:05 & 20:35 & 21:05 \\
\bottomrule\end{tabular}\end{center}
\subsection*{To Trenton (read down) --- trips 40--42 of 42}
\begin{center}\footnotesize\setlength{\tabcolsep}{3pt}\begin{tabular}{@{}lrrr@{}}\toprule
Station & 40 & 41 & 42 \\\midrule
Entertainment Center & 20:30 & 21:00 & 21:29 \\
Aquarium & 20:32 & 21:02 & 21:31 \\
Cooper St/Rutgers & 20:34 & 21:04 & 21:33 \\
Walter Rand TC & 20:38 & 21:08 & 21:37 \\
36th Street & 20:43 & 21:13 & 21:42 \\
Pennsauken TC & 20:45 & 21:15 & 21:44 \\
Route 73/Pennsauken & 20:48 & 21:18 & 21:47 \\
Palmyra & 20:51 & 21:21 & -- \\
Riverton & 20:53 & 21:23 & -- \\
Cinnaminson & 20:55 & 21:25 & -- \\
Riverside & 21:00 & 21:30 & -- \\
Delanco & 21:02 & 21:32 & -- \\
Beverly/Edgewater Park & 21:06 & 21:36 & -- \\
Burlington South & 21:10 & 21:40 & -- \\
Burlington Towne Ctr & 21:12 & 21:42 & -- \\
Florence & 21:18 & 21:48 & -- \\
Roebling & 21:21 & 21:51 & -- \\
Bordentown & 21:26 & 21:56 & -- \\
Cass Street & 21:32 & 22:02 & -- \\
Hamilton Avenue & 21:34 & 22:04 & -- \\
Trenton TC & 21:37 & 22:07 & -- \\
\bottomrule\end{tabular}\end{center}
\subsection*{To Camden (read down) --- trips 1--13 of 45}
\begin{center}\footnotesize\setlength{\tabcolsep}{3pt}\begin{tabular}{@{}lrrrrrrrrrrrrr@{}}\toprule
Station & 1 & 2 & 3 & 4 & 5 & 6 & 7 & 8 & 9 & 10 & 11 & 12 & 13 \\\midrule
Trenton TC & -- & -- & -- & 5:45 & -- & 5:57 & -- & 6:27 & 6:42 & 6:57 & 7:12 & 7:27 & 7:57 \\
Hamilton Avenue & -- & -- & -- & 5:47 & -- & 5:59 & -- & 6:29 & 6:44 & 6:59 & 7:14 & 7:29 & 7:59 \\
Cass Street & -- & -- & -- & 5:49 & -- & 6:01 & -- & 6:31 & 6:46 & 7:01 & 7:16 & 7:31 & 8:01 \\
Bordentown & -- & -- & -- & 5:56 & -- & 6:08 & -- & 6:38 & 6:53 & 7:08 & 7:23 & 7:38 & 8:08 \\
Roebling & -- & -- & -- & 6:01 & -- & 6:13 & -- & 6:43 & 6:58 & 7:13 & 7:28 & 7:43 & 8:13 \\
Florence & -- & -- & -- & 6:04 & -- & 6:16 & -- & 6:46 & 7:01 & 7:16 & 7:31 & 7:46 & 8:16 \\
Burlington Towne Ctr & -- & -- & -- & 6:09 & -- & 6:21 & 5:59 & 6:51 & 7:06 & 7:21 & 7:36 & 7:51 & 8:21 \\
Burlington South & -- & -- & -- & 6:11 & -- & 6:23 & 6:01 & 6:53 & 7:08 & 7:23 & 7:38 & 7:53 & 8:23 \\
Beverly/Edgewater Park & -- & -- & -- & 6:15 & -- & 6:27 & 6:04 & 6:57 & 7:12 & 7:27 & 7:42 & 7:57 & 8:27 \\
Delanco & -- & -- & -- & 6:19 & -- & 6:31 & 6:07 & 7:01 & 7:16 & 7:31 & 7:46 & 8:01 & 8:31 \\
Riverside & -- & -- & -- & 6:21 & -- & 6:33 & 6:09 & 7:03 & 7:18 & 7:33 & 7:48 & 8:03 & 8:33 \\
Cinnaminson & -- & -- & -- & 6:26 & -- & 6:38 & 6:13 & 7:08 & 7:23 & 7:38 & 7:53 & 8:08 & 8:38 \\
Riverton & -- & -- & -- & 6:28 & -- & 6:40 & 6:15 & 7:10 & 7:25 & 7:40 & 7:55 & 8:10 & 8:40 \\
Palmyra & -- & -- & -- & 6:30 & -- & 6:42 & 6:17 & 7:12 & 7:27 & 7:42 & 7:57 & 8:12 & 8:42 \\
Route 73/Pennsauken & -- & 5:30 & -- & 6:33 & -- & 6:45 & 6:19 & 7:15 & 7:30 & 7:45 & 8:00 & 8:15 & 8:45 \\
Pennsauken TC & -- & 5:32 & -- & 6:35 & -- & 6:47 & 6:21 & 7:17 & 7:32 & 7:47 & 8:02 & 8:17 & 8:47 \\
36th Street & 5:27 & 5:36 & 5:40 & 6:39 & 5:55 & 6:51 & 6:25 & 7:21 & 7:36 & 7:51 & 8:06 & 8:21 & 8:51 \\
Walter Rand TC & 5:32 & 5:41 & 5:45 & 6:44 & 6:00 & 6:56 & 6:30 & 7:26 & 7:41 & 7:56 & 8:11 & 8:26 & 8:56 \\
Cooper St/Rutgers & 5:35 & -- & 5:48 & 6:47 & 6:03 & 6:59 & 6:33 & 7:29 & 7:44 & 7:59 & 8:14 & 8:29 & 8:59 \\
Aquarium & 5:37 & -- & 5:50 & 6:49 & 6:05 & 7:01 & 6:35 & 7:31 & 7:46 & 8:01 & 8:16 & 8:31 & 9:01 \\
Entertainment Center & 5:39 & -- & 5:52 & 6:52 & 6:07 & 7:04 & 6:38 & 7:34 & 7:49 & 8:04 & 8:19 & 8:34 & 9:04 \\
\bottomrule\end{tabular}\end{center}
\subsection*{To Camden (read down) --- trips 14--26 of 45}
\begin{center}\footnotesize\setlength{\tabcolsep}{3pt}\begin{tabular}{@{}lrrrrrrrrrrrrr@{}}\toprule
Station & 14 & 15 & 16 & 17 & 18 & 19 & 20 & 21 & 22 & 23 & 24 & 25 & 26 \\\midrule
Trenton TC & 8:12 & 8:27 & 8:57 & 9:27 & 9:57 & 10:27 & 10:57 & 11:27 & 11:57 & 12:27 & 12:57 & 13:27 & 13:57 \\
Hamilton Avenue & 8:14 & 8:29 & 8:59 & 9:29 & 9:59 & 10:29 & 10:59 & 11:29 & 11:59 & 12:29 & 12:59 & 13:29 & 13:59 \\
Cass Street & 8:16 & 8:31 & 9:01 & 9:31 & 10:01 & 10:31 & 11:01 & 11:31 & 12:01 & 12:31 & 13:01 & 13:31 & 14:01 \\
Bordentown & 8:23 & 8:38 & 9:08 & 9:38 & 10:08 & 10:38 & 11:08 & 11:38 & 12:08 & 12:38 & 13:08 & 13:38 & 14:08 \\
Roebling & 8:28 & 8:43 & 9:13 & 9:43 & 10:13 & 10:43 & 11:13 & 11:43 & 12:13 & 12:43 & 13:13 & 13:43 & 14:13 \\
Florence & 8:31 & 8:46 & 9:16 & 9:46 & 10:16 & 10:46 & 11:16 & 11:46 & 12:16 & 12:46 & 13:16 & 13:46 & 14:16 \\
Burlington Towne Ctr & 8:36 & 8:51 & 9:21 & 9:51 & 10:21 & 10:51 & 11:21 & 11:51 & 12:21 & 12:51 & 13:21 & 13:51 & 14:21 \\
Burlington South & 8:38 & 8:53 & 9:23 & 9:53 & 10:23 & 10:53 & 11:23 & 11:53 & 12:23 & 12:53 & 13:23 & 13:53 & 14:23 \\
Beverly/Edgewater Park & 8:42 & 8:57 & 9:27 & 9:57 & 10:27 & 10:57 & 11:27 & 11:57 & 12:27 & 12:57 & 13:27 & 13:57 & 14:27 \\
Delanco & 8:46 & 9:01 & 9:31 & 10:01 & 10:31 & 11:01 & 11:31 & 12:01 & 12:31 & 13:01 & 13:31 & 14:01 & 14:31 \\
Riverside & 8:48 & 9:03 & 9:33 & 10:03 & 10:33 & 11:03 & 11:33 & 12:03 & 12:33 & 13:03 & 13:33 & 14:03 & 14:33 \\
Cinnaminson & 8:53 & 9:08 & 9:38 & 10:08 & 10:38 & 11:08 & 11:38 & 12:08 & 12:38 & 13:08 & 13:38 & 14:08 & 14:38 \\
Riverton & 8:55 & 9:10 & 9:40 & 10:10 & 10:40 & 11:10 & 11:40 & 12:10 & 12:40 & 13:10 & 13:40 & 14:10 & 14:40 \\
Palmyra & 8:57 & 9:12 & 9:42 & 10:12 & 10:42 & 11:12 & 11:42 & 12:12 & 12:42 & 13:12 & 13:42 & 14:12 & 14:42 \\
Route 73/Pennsauken & 9:00 & 9:15 & 9:45 & 10:15 & 10:45 & 11:15 & 11:45 & 12:15 & 12:45 & 13:15 & 13:45 & 14:15 & 14:45 \\
Pennsauken TC & 9:02 & 9:17 & 9:47 & 10:17 & 10:47 & 11:17 & 11:47 & 12:17 & 12:47 & 13:17 & 13:47 & 14:17 & 14:47 \\
36th Street & 9:06 & 9:21 & 9:51 & 10:21 & 10:51 & 11:21 & 11:51 & 12:21 & 12:51 & 13:21 & 13:51 & 14:21 & 14:51 \\
Walter Rand TC & 9:11 & 9:26 & 9:56 & 10:26 & 10:56 & 11:26 & 11:56 & 12:26 & 12:56 & 13:26 & 13:56 & 14:26 & 14:56 \\
Cooper St/Rutgers & 9:14 & 9:29 & 9:59 & 10:29 & 10:59 & 11:29 & 11:59 & 12:29 & 12:59 & 13:29 & 13:59 & 14:29 & 14:59 \\
Aquarium & 9:16 & 9:31 & 10:01 & 10:31 & 11:01 & 11:31 & 12:01 & 12:31 & 13:01 & 13:31 & 14:01 & 14:31 & 15:01 \\
Entertainment Center & 9:19 & 9:34 & 10:04 & 10:34 & 11:04 & 11:34 & 12:04 & 12:34 & 13:04 & 13:34 & 14:04 & 14:34 & 15:04 \\
\bottomrule\end{tabular}\end{center}
\subsection*{To Camden (read down) --- trips 27--39 of 45}
\begin{center}\footnotesize\setlength{\tabcolsep}{3pt}\begin{tabular}{@{}lrrrrrrrrrrrrr@{}}\toprule
Station & 27 & 28 & 29 & 30 & 31 & 32 & 33 & 34 & 35 & 36 & 37 & 38 & 39 \\\midrule
Trenton TC & 14:27 & 14:57 & 15:27 & 15:57 & 16:12 & 16:42 & 16:57 & 17:12 & 17:27 & 17:42 & 18:12 & 18:27 & 18:57 \\
Hamilton Avenue & 14:29 & 14:59 & 15:29 & 15:59 & 16:14 & 16:44 & 16:59 & 17:14 & 17:29 & 17:44 & 18:14 & 18:29 & 18:59 \\
Cass Street & 14:31 & 15:01 & 15:31 & 16:01 & 16:16 & 16:46 & 17:01 & 17:16 & 17:31 & 17:46 & 18:16 & 18:31 & 19:01 \\
Bordentown & 14:38 & 15:08 & 15:38 & 16:08 & 16:23 & 16:53 & 17:08 & 17:23 & 17:38 & 17:53 & 18:23 & 18:38 & 19:08 \\
Roebling & 14:43 & 15:13 & 15:43 & 16:13 & 16:28 & 16:58 & 17:13 & 17:28 & 17:43 & 17:58 & 18:28 & 18:43 & 19:13 \\
Florence & 14:46 & 15:16 & 15:46 & 16:16 & 16:31 & 17:01 & 17:16 & 17:31 & 17:46 & 18:01 & 18:31 & 18:46 & 19:16 \\
Burlington Towne Ctr & 14:51 & 15:21 & 15:51 & 16:21 & 16:36 & 17:06 & 17:21 & 17:36 & 17:51 & 18:06 & 18:36 & 18:51 & 19:21 \\
Burlington South & 14:53 & 15:23 & 15:53 & 16:23 & 16:38 & 17:08 & 17:23 & 17:38 & 17:53 & 18:08 & 18:38 & 18:53 & 19:23 \\
Beverly/Edgewater Park & 14:57 & 15:27 & 15:57 & 16:27 & 16:42 & 17:12 & 17:27 & 17:42 & 17:57 & 18:12 & 18:42 & 18:57 & 19:27 \\
Delanco & 15:01 & 15:31 & 16:01 & 16:31 & 16:46 & 17:16 & 17:31 & 17:46 & 18:01 & 18:16 & 18:46 & 19:01 & 19:31 \\
Riverside & 15:03 & 15:33 & 16:03 & 16:33 & 16:48 & 17:18 & 17:33 & 17:48 & 18:03 & 18:18 & 18:48 & 19:03 & 19:33 \\
Cinnaminson & 15:08 & 15:38 & 16:08 & 16:38 & 16:53 & 17:23 & 17:38 & 17:53 & 18:08 & 18:23 & 18:53 & 19:08 & 19:38 \\
Riverton & 15:10 & 15:40 & 16:10 & 16:40 & 16:55 & 17:25 & 17:40 & 17:55 & 18:10 & 18:25 & 18:55 & 19:10 & 19:40 \\
Palmyra & 15:12 & 15:42 & 16:12 & 16:42 & 16:57 & 17:27 & 17:42 & 17:57 & 18:12 & 18:27 & 18:57 & 19:12 & 19:42 \\
Route 73/Pennsauken & 15:15 & 15:45 & 16:15 & 16:45 & 17:00 & 17:30 & 17:45 & 18:00 & 18:15 & 18:30 & 19:00 & 19:15 & 19:45 \\
Pennsauken TC & 15:17 & 15:47 & 16:17 & 16:47 & 17:02 & 17:32 & 17:47 & 18:02 & 18:17 & 18:32 & 19:02 & 19:17 & 19:47 \\
36th Street & 15:21 & 15:51 & 16:21 & 16:51 & 17:06 & 17:36 & 17:51 & 18:06 & 18:21 & 18:36 & 19:06 & 19:21 & 19:51 \\
Walter Rand TC & 15:26 & 15:56 & 16:26 & 16:56 & 17:11 & 17:41 & 17:56 & 18:11 & 18:26 & 18:41 & 19:11 & 19:26 & 19:56 \\
Cooper St/Rutgers & 15:29 & 15:59 & 16:29 & 16:59 & 17:14 & 17:44 & 17:59 & 18:14 & 18:29 & 18:44 & 19:14 & 19:29 & 19:59 \\
Aquarium & 15:31 & 16:01 & 16:31 & 17:01 & 17:16 & 17:46 & 18:01 & 18:16 & 18:31 & 18:46 & 19:16 & 19:31 & 20:01 \\
Entertainment Center & 15:34 & 16:04 & 16:34 & 17:04 & 17:19 & 17:49 & 18:04 & 18:19 & 18:34 & 18:49 & 19:19 & 19:34 & 20:04 \\
\bottomrule\end{tabular}\end{center}
\subsection*{To Camden (read down) --- trips 40--45 of 45}
\begin{center}\footnotesize\setlength{\tabcolsep}{3pt}\begin{tabular}{@{}lrrrrrr@{}}\toprule
Station & 40 & 41 & 42 & 43 & 44 & 45 \\\midrule
Trenton TC & 19:27 & 19:57 & 20:27 & 21:00 & 21:15 & 21:28 \\
Hamilton Avenue & 19:29 & 19:59 & 20:29 & 21:02 & -- & 21:30 \\
Cass Street & 19:31 & 20:01 & 20:31 & 21:04 & -- & 21:32 \\
Bordentown & 19:38 & 20:08 & 20:38 & 21:11 & 21:24 & 21:39 \\
Roebling & 19:43 & 20:13 & 20:43 & 21:16 & -- & 21:44 \\
Florence & 19:46 & 20:16 & 20:46 & 21:19 & 21:31 & 21:47 \\
Burlington Towne Ctr & 19:51 & 20:21 & 20:51 & 21:24 & 21:36 & 21:52 \\
Burlington South & 19:53 & 20:23 & 20:53 & 21:26 & 21:38 & 21:55 \\
Beverly/Edgewater Park & 19:57 & 20:27 & 20:57 & 21:30 & 21:42 & -- \\
Delanco & 20:01 & 20:31 & 21:01 & 21:34 & -- & -- \\
Riverside & 20:03 & 20:33 & 21:03 & 21:36 & 21:46 & -- \\
Cinnaminson & 20:08 & 20:38 & 21:08 & 21:41 & 21:51 & -- \\
Riverton & 20:10 & 20:40 & 21:10 & 21:43 & -- & -- \\
Palmyra & 20:12 & 20:42 & 21:12 & 21:45 & 21:54 & -- \\
Route 73/Pennsauken & 20:15 & 20:45 & 21:15 & 21:48 & 21:57 & -- \\
Pennsauken TC & 20:17 & 20:47 & 21:17 & 21:50 & 21:59 & -- \\
36th Street & 20:21 & 20:51 & 21:21 & 21:54 & 22:03 & -- \\
Walter Rand TC & 20:26 & 20:56 & 21:26 & 21:59 & 22:08 & -- \\
Cooper St/Rutgers & 20:29 & 20:59 & 21:29 & 22:02 & -- & -- \\
Aquarium & 20:31 & 21:01 & 21:31 & 22:04 & -- & -- \\
Entertainment Center & 20:34 & 21:04 & 21:34 & 22:07 & -- & -- \\
\bottomrule\end{tabular}\end{center}
\clearpage
\section{Proposed T8 (8 trains peak, 6 midday)}\label{tt:proposed_T8}
\subsection*{To Trenton (read down) --- trips 1--13 of 43}
\begin{center}\footnotesize\setlength{\tabcolsep}{3pt}\begin{tabular}{@{}lrrrrrrrrrrrrr@{}}\toprule
Station & 1 & 2 & 3 & 4 & 5 & 6 & 7 & 8 & 9 & 10 & 11 & 12 & 13 \\\midrule
Entertainment Center & 5:55 & 6:12 & 6:32 & 6:52 & 7:09 & 7:30 & 7:51 & 8:08 & 8:25 & 8:42 & 9:02 & 9:22 & 9:39 \\
Aquarium & 5:57 & 6:14 & 6:34 & 6:54 & 7:11 & 7:32 & 7:53 & 8:10 & 8:27 & 8:44 & 9:04 & 9:24 & 9:41 \\
Cooper St/Rutgers & 5:59 & 6:16 & 6:36 & 6:56 & 7:13 & 7:34 & 7:55 & 8:12 & 8:29 & 8:46 & 9:06 & 9:26 & 9:43 \\
Walter Rand TC & 6:03 & 6:20 & 6:40 & 7:00 & 7:17 & 7:38 & 7:59 & 8:16 & 8:33 & 8:50 & 9:10 & 9:30 & 9:47 \\
36th Street & 6:08 & 6:25 & 6:45 & 7:05 & 7:22 & 7:43 & 8:09 & 8:21 & 8:38 & 8:55 & 9:15 & 9:35 & 9:52 \\
Pennsauken TC & 6:10 & 6:27 & 6:47 & 7:07 & 7:24 & 7:45 & 8:11 & 8:23 & 8:40 & 8:57 & 9:17 & 9:37 & 9:54 \\
Route 73/Pennsauken & 6:13 & 6:30 & 6:50 & 7:10 & 7:27 & 7:48 & 8:14 & 8:26 & 8:43 & 9:00 & 9:20 & 9:40 & 9:57 \\
Palmyra & 6:16 & 6:33 & 6:53 & 7:13 & 7:30 & 7:51 & 8:17 & 8:29 & 8:46 & 9:03 & 9:23 & 9:43 & 10:00 \\
Riverton & 6:18 & 6:35 & 6:55 & 7:15 & 7:32 & 7:53 & 8:19 & 8:31 & 8:48 & 9:05 & 9:25 & 9:45 & 10:02 \\
Cinnaminson & 6:20 & 6:37 & 6:57 & 7:17 & 7:34 & 7:55 & 8:21 & 8:33 & 8:50 & 9:07 & 9:27 & 9:47 & 10:04 \\
Riverside & 6:28 & 6:42 & 7:02 & 7:22 & 7:39 & 8:00 & 8:26 & 8:38 & 8:58 & 9:12 & 9:32 & 9:52 & 10:09 \\
Delanco & 6:30 & 6:44 & 7:04 & 7:24 & 7:41 & 8:02 & 8:28 & 8:40 & 9:00 & 9:14 & 9:34 & 9:54 & 10:11 \\
Beverly/Edgewater Park & 6:34 & 6:48 & 7:08 & 7:28 & 7:45 & 8:06 & 8:32 & 8:44 & 9:04 & 9:18 & 9:38 & 9:58 & 10:15 \\
Burlington South & 6:38 & 6:52 & 7:13 & 7:34 & 7:49 & 8:10 & 8:36 & 8:48 & 9:08 & 9:22 & 9:43 & 10:04 & 10:19 \\
Burlington Towne Ctr & 6:40 & 6:54 & 7:15 & 7:36 & 7:51 & 8:12 & 8:38 & 8:50 & 9:10 & 9:24 & 9:45 & 10:06 & 10:21 \\
Florence & 6:46 & 7:00 & 7:21 & 7:42 & 7:57 & 8:18 & 8:44 & 8:56 & 9:16 & 9:30 & 9:51 & 10:12 & 10:27 \\
Roebling & 6:49 & 7:03 & 7:24 & 7:45 & 8:00 & 8:21 & 8:47 & 8:59 & 9:19 & 9:33 & 9:54 & 10:15 & 10:30 \\
Bordentown & 6:54 & 7:08 & 7:29 & 7:50 & 8:05 & 8:26 & 8:52 & 9:04 & 9:24 & 9:38 & 9:59 & 10:20 & 10:35 \\
Cass Street & 7:00 & 7:14 & 7:35 & 7:56 & 8:11 & 8:32 & 8:58 & 9:10 & 9:30 & 9:44 & 10:05 & 10:26 & 10:41 \\
Hamilton Avenue & 7:02 & 7:16 & 7:37 & 7:58 & 8:13 & 8:34 & 9:00 & 9:12 & 9:32 & 9:46 & 10:07 & 10:28 & 10:43 \\
Trenton TC & 7:05 & 7:19 & 7:40 & 8:01 & 8:16 & 8:37 & 9:03 & 9:15 & 9:35 & 9:49 & 10:10 & 10:31 & 10:46 \\
\bottomrule\end{tabular}\end{center}
\subsection*{To Trenton (read down) --- trips 14--26 of 43}
\begin{center}\footnotesize\setlength{\tabcolsep}{3pt}\begin{tabular}{@{}lrrrrrrrrrrrrr@{}}\toprule
Station & 14 & 15 & 16 & 17 & 18 & 19 & 20 & 21 & 22 & 23 & 24 & 25 & 26 \\\midrule
Entertainment Center & 10:00 & 10:38 & 11:12 & 11:32 & 11:52 & 12:09 & 12:30 & 13:08 & 13:42 & 14:02 & 14:22 & 14:39 & 15:00 \\
Aquarium & 10:02 & 10:40 & 11:14 & 11:34 & 11:54 & 12:11 & 12:32 & 13:10 & 13:44 & 14:04 & 14:24 & 14:41 & 15:02 \\
Cooper St/Rutgers & 10:04 & 10:42 & 11:16 & 11:36 & 11:56 & 12:13 & 12:34 & 13:12 & 13:46 & 14:06 & 14:26 & 14:43 & 15:04 \\
Walter Rand TC & 10:08 & 10:46 & 11:20 & 11:40 & 12:00 & 12:17 & 12:38 & 13:16 & 13:50 & 14:10 & 14:30 & 14:47 & 15:08 \\
36th Street & 10:13 & 10:51 & 11:25 & 11:45 & 12:05 & 12:22 & 12:43 & 13:21 & 13:55 & 14:15 & 14:35 & 14:52 & 15:13 \\
Pennsauken TC & 10:15 & 10:53 & 11:27 & 11:47 & 12:07 & 12:24 & 12:45 & 13:23 & 13:57 & 14:17 & 14:37 & 14:54 & 15:15 \\
Route 73/Pennsauken & 10:18 & 10:56 & 11:30 & 11:50 & 12:10 & 12:27 & 12:48 & 13:26 & 14:00 & 14:20 & 14:40 & 14:57 & 15:18 \\
Palmyra & 10:21 & 10:59 & 11:33 & 11:53 & 12:13 & 12:30 & 12:51 & 13:29 & 14:03 & 14:23 & 14:43 & 15:00 & 15:21 \\
Riverton & 10:23 & 11:01 & 11:35 & 11:55 & 12:15 & 12:32 & 12:53 & 13:31 & 14:05 & 14:25 & 14:45 & 15:02 & 15:23 \\
Cinnaminson & 10:25 & 11:03 & 11:37 & 11:57 & 12:17 & 12:34 & 12:55 & 13:33 & 14:07 & 14:27 & 14:47 & 15:04 & 15:25 \\
Riverside & 10:30 & 11:08 & 11:42 & 12:02 & 12:22 & 12:39 & 13:00 & 13:38 & 14:12 & 14:32 & 14:52 & 15:09 & 15:30 \\
Delanco & 10:32 & 11:10 & 11:44 & 12:04 & 12:24 & 12:41 & 13:02 & 13:40 & 14:14 & 14:34 & 14:54 & 15:11 & 15:32 \\
Beverly/Edgewater Park & 10:36 & 11:14 & 11:48 & 12:08 & 12:28 & 12:45 & 13:06 & 13:44 & 14:18 & 14:38 & 14:58 & 15:15 & 15:36 \\
Burlington South & 10:40 & 11:18 & 11:52 & 12:13 & 12:34 & 12:49 & 13:10 & 13:48 & 14:22 & 14:43 & 15:04 & 15:19 & 15:40 \\
Burlington Towne Ctr & 10:42 & 11:20 & 11:54 & 12:15 & 12:36 & 12:51 & 13:12 & 13:50 & 14:24 & 14:45 & 15:06 & 15:21 & 15:42 \\
Florence & 10:48 & 11:26 & 12:00 & 12:21 & 12:42 & 12:57 & 13:18 & 13:56 & 14:30 & 14:51 & 15:12 & 15:27 & 15:48 \\
Roebling & 10:51 & 11:29 & 12:03 & 12:24 & 12:45 & 13:00 & 13:21 & 13:59 & 14:33 & 14:54 & 15:15 & 15:30 & 15:51 \\
Bordentown & 10:56 & 11:34 & 12:08 & 12:29 & 12:50 & 13:05 & 13:26 & 14:04 & 14:38 & 14:59 & 15:20 & 15:35 & 15:56 \\
Cass Street & 11:02 & 11:40 & 12:14 & 12:35 & 12:56 & 13:11 & 13:32 & 14:10 & 14:44 & 15:05 & 15:26 & 15:41 & 16:02 \\
Hamilton Avenue & 11:04 & 11:42 & 12:16 & 12:37 & 12:58 & 13:13 & 13:34 & 14:12 & 14:46 & 15:07 & 15:28 & 15:43 & 16:04 \\
Trenton TC & 11:07 & 11:45 & 12:19 & 12:40 & 13:01 & 13:16 & 13:37 & 14:15 & 14:49 & 15:10 & 15:31 & 15:46 & 16:07 \\
\bottomrule\end{tabular}\end{center}
\subsection*{To Trenton (read down) --- trips 27--39 of 43}
\begin{center}\footnotesize\setlength{\tabcolsep}{3pt}\begin{tabular}{@{}lrrrrrrrrrrrrr@{}}\toprule
Station & 27 & 28 & 29 & 30 & 31 & 32 & 33 & 34 & 35 & 36 & 37 & 38 & 39 \\\midrule
Entertainment Center & 15:21 & 15:38 & 15:55 & 16:12 & 16:32 & 16:52 & 17:09 & 17:30 & 17:51 & 18:08 & 18:25 & 18:42 & 19:02 \\
Aquarium & 15:23 & 15:40 & 15:57 & 16:14 & 16:34 & 16:54 & 17:11 & 17:32 & 17:53 & 18:10 & 18:27 & 18:44 & 19:04 \\
Cooper St/Rutgers & 15:25 & 15:42 & 15:59 & 16:16 & 16:36 & 16:56 & 17:13 & 17:34 & 17:55 & 18:12 & 18:29 & 18:46 & 19:06 \\
Walter Rand TC & 15:29 & 15:46 & 16:03 & 16:20 & 16:40 & 17:00 & 17:17 & 17:38 & 17:59 & 18:16 & 18:33 & 18:50 & 19:10 \\
36th Street & 15:39 & 15:51 & 16:08 & 16:25 & 16:45 & 17:05 & 17:22 & 17:43 & 18:09 & 18:21 & 18:38 & 18:55 & 19:15 \\
Pennsauken TC & 15:41 & 15:53 & 16:10 & 16:27 & 16:47 & 17:07 & 17:24 & 17:45 & 18:11 & 18:23 & 18:40 & 18:57 & 19:17 \\
Route 73/Pennsauken & 15:44 & 15:56 & 16:13 & 16:30 & 16:50 & 17:10 & 17:27 & 17:48 & 18:14 & 18:26 & 18:43 & 19:00 & 19:20 \\
Palmyra & 15:47 & 15:59 & 16:16 & 16:33 & 16:53 & 17:13 & 17:30 & 17:51 & 18:17 & 18:29 & 18:46 & 19:03 & 19:23 \\
Riverton & 15:49 & 16:01 & 16:18 & 16:35 & 16:55 & 17:15 & 17:32 & 17:53 & 18:19 & 18:31 & 18:48 & 19:05 & 19:25 \\
Cinnaminson & 15:51 & 16:03 & 16:20 & 16:37 & 16:57 & 17:17 & 17:34 & 17:55 & 18:21 & 18:33 & 18:50 & 19:07 & 19:27 \\
Riverside & 15:56 & 16:08 & 16:28 & 16:42 & 17:02 & 17:22 & 17:39 & 18:00 & 18:26 & 18:38 & 18:58 & 19:12 & 19:32 \\
Delanco & 15:58 & 16:10 & 16:30 & 16:44 & 17:04 & 17:24 & 17:41 & 18:02 & 18:28 & 18:40 & 19:00 & 19:14 & 19:34 \\
Beverly/Edgewater Park & 16:02 & 16:14 & 16:34 & 16:48 & 17:08 & 17:28 & 17:45 & 18:06 & 18:32 & 18:44 & 19:04 & 19:18 & 19:38 \\
Burlington South & 16:06 & 16:18 & 16:38 & 16:52 & 17:13 & 17:34 & 17:49 & 18:10 & 18:36 & 18:48 & 19:08 & 19:22 & 19:43 \\
Burlington Towne Ctr & 16:08 & 16:20 & 16:40 & 16:54 & 17:15 & 17:36 & 17:51 & 18:12 & 18:38 & 18:50 & 19:10 & 19:24 & 19:45 \\
Florence & 16:14 & 16:26 & 16:46 & 17:00 & 17:21 & 17:42 & 17:57 & 18:18 & 18:44 & 18:56 & 19:16 & 19:30 & 19:51 \\
Roebling & 16:17 & 16:29 & 16:49 & 17:03 & 17:24 & 17:45 & 18:00 & 18:21 & 18:47 & 18:59 & 19:19 & 19:33 & 19:54 \\
Bordentown & 16:22 & 16:34 & 16:54 & 17:08 & 17:29 & 17:50 & 18:05 & 18:26 & 18:52 & 19:04 & 19:24 & 19:38 & 19:59 \\
Cass Street & 16:28 & 16:40 & 17:00 & 17:14 & 17:35 & 17:56 & 18:11 & 18:32 & 18:58 & 19:10 & 19:30 & 19:44 & 20:05 \\
Hamilton Avenue & 16:30 & 16:42 & 17:02 & 17:16 & 17:37 & 17:58 & 18:13 & 18:34 & 19:00 & 19:12 & 19:32 & 19:46 & 20:07 \\
Trenton TC & 16:33 & 16:45 & 17:05 & 17:19 & 17:40 & 18:01 & 18:16 & 18:37 & 19:03 & 19:15 & 19:35 & 19:49 & 20:10 \\
\bottomrule\end{tabular}\end{center}
\subsection*{To Trenton (read down) --- trips 40--43 of 43}
\begin{center}\footnotesize\setlength{\tabcolsep}{3pt}\begin{tabular}{@{}lrrrr@{}}\toprule
Station & 40 & 41 & 42 & 43 \\\midrule
Entertainment Center & 19:22 & 19:39 & 20:00 & 20:38 \\
Aquarium & 19:24 & 19:41 & 20:02 & 20:40 \\
Cooper St/Rutgers & 19:26 & 19:43 & 20:04 & 20:42 \\
Walter Rand TC & 19:30 & 19:47 & 20:08 & 20:46 \\
36th Street & 19:35 & 19:52 & 20:13 & 20:51 \\
Pennsauken TC & 19:37 & 19:54 & 20:15 & 20:53 \\
Route 73/Pennsauken & 19:40 & 19:57 & 20:18 & 20:56 \\
Palmyra & 19:43 & 20:00 & 20:21 & 20:59 \\
Riverton & 19:45 & 20:02 & 20:23 & 21:01 \\
Cinnaminson & 19:47 & 20:04 & 20:25 & 21:03 \\
Riverside & 19:52 & 20:09 & 20:30 & 21:08 \\
Delanco & 19:54 & 20:11 & 20:32 & 21:10 \\
Beverly/Edgewater Park & 19:58 & 20:15 & 20:36 & 21:14 \\
Burlington South & 20:04 & 20:19 & 20:40 & 21:18 \\
Burlington Towne Ctr & 20:06 & 20:21 & 20:42 & 21:20 \\
Florence & 20:12 & 20:27 & 20:48 & 21:26 \\
Roebling & 20:15 & 20:30 & 20:51 & 21:29 \\
Bordentown & 20:20 & 20:35 & 20:56 & 21:34 \\
Cass Street & 20:26 & 20:41 & 21:02 & 21:40 \\
Hamilton Avenue & 20:28 & 20:43 & 21:04 & 21:42 \\
Trenton TC & 20:31 & 20:46 & 21:07 & 21:45 \\
\bottomrule\end{tabular}\end{center}
\subsection*{To Camden (read down) --- trips 1--13 of 45}
\begin{center}\footnotesize\setlength{\tabcolsep}{3pt}\begin{tabular}{@{}lrrrrrrrrrrrrr@{}}\toprule
Station & 1 & 2 & 3 & 4 & 5 & 6 & 7 & 8 & 9 & 10 & 11 & 12 & 13 \\\midrule
Trenton TC & 5:56 & 6:17 & 6:38 & 6:52 & 7:10 & 7:25 & 7:46 & 8:07 & 8:26 & 8:47 & 9:08 & 9:22 & 9:40 \\
Hamilton Avenue & 5:58 & 6:19 & 6:40 & 6:54 & 7:12 & 7:27 & 7:48 & 8:09 & 8:28 & 8:49 & 9:10 & 9:24 & 9:42 \\
Cass Street & 6:00 & 6:21 & 6:42 & 6:56 & 7:14 & 7:29 & 7:50 & 8:11 & 8:30 & 8:51 & 9:12 & 9:26 & 9:44 \\
Bordentown & 6:07 & 6:29 & 6:49 & 7:03 & 7:21 & 7:36 & 7:57 & 8:21 & 8:37 & 8:59 & 9:19 & 9:33 & 9:51 \\
Roebling & 6:12 & 6:34 & 6:55 & 7:08 & 7:29 & 7:41 & 8:05 & 8:26 & 8:42 & 9:04 & 9:25 & 9:38 & 9:59 \\
Florence & 6:15 & 6:37 & 6:58 & 7:11 & 7:32 & 7:44 & 8:08 & 8:29 & 8:45 & 9:07 & 9:28 & 9:41 & 10:02 \\
Burlington Towne Ctr & 6:20 & 6:42 & 7:03 & 7:16 & 7:37 & 7:49 & 8:13 & 8:34 & 8:50 & 9:12 & 9:33 & 9:46 & 10:07 \\
Burlington South & 6:22 & 6:44 & 7:05 & 7:18 & 7:39 & 7:51 & 8:15 & 8:36 & 8:52 & 9:14 & 9:35 & 9:48 & 10:09 \\
Beverly/Edgewater Park & 6:26 & 6:48 & 7:09 & 7:22 & 7:43 & 7:55 & 8:19 & 8:40 & 8:56 & 9:18 & 9:39 & 9:52 & 10:13 \\
Delanco & 6:30 & 6:52 & 7:13 & 7:26 & 7:47 & 7:59 & 8:23 & 8:44 & 9:00 & 9:22 & 9:43 & 9:56 & 10:17 \\
Riverside & 6:32 & 6:54 & 7:15 & 7:28 & 7:49 & 8:01 & 8:25 & 8:46 & 9:02 & 9:24 & 9:45 & 9:58 & 10:19 \\
Cinnaminson & 6:37 & 6:59 & 7:20 & 7:33 & 7:54 & 8:06 & 8:30 & 8:51 & 9:07 & 9:29 & 9:50 & 10:03 & 10:24 \\
Riverton & 6:39 & 7:01 & 7:22 & 7:35 & 7:56 & 8:08 & 8:32 & 8:53 & 9:09 & 9:31 & 9:52 & 10:05 & 10:26 \\
Palmyra & 6:41 & 7:03 & 7:24 & 7:37 & 7:58 & 8:10 & 8:34 & 8:55 & 9:11 & 9:33 & 9:54 & 10:07 & 10:28 \\
Route 73/Pennsauken & 6:44 & 7:06 & 7:27 & 7:40 & 8:01 & 8:13 & 8:37 & 8:58 & 9:14 & 9:36 & 9:57 & 10:10 & 10:31 \\
Pennsauken TC & 6:46 & 7:08 & 7:29 & 7:42 & 8:03 & 8:15 & 8:39 & 9:00 & 9:16 & 9:38 & 9:59 & 10:12 & 10:33 \\
36th Street & 6:50 & 7:12 & 7:33 & 7:48 & 8:07 & 8:19 & 8:43 & 9:04 & 9:20 & 9:42 & 10:03 & 10:18 & 10:37 \\
Walter Rand TC & 6:55 & 7:17 & 7:38 & 7:53 & 8:12 & 8:24 & 8:48 & 9:09 & 9:25 & 9:47 & 10:08 & 10:23 & 10:42 \\
Cooper St/Rutgers & 6:58 & 7:20 & 7:41 & 7:56 & 8:15 & 8:27 & 8:51 & 9:12 & 9:28 & 9:50 & 10:11 & 10:26 & 10:45 \\
Aquarium & 7:00 & 7:22 & 7:43 & 7:58 & 8:17 & 8:29 & 8:53 & 9:14 & 9:30 & 9:52 & 10:13 & 10:28 & 10:47 \\
Entertainment Center & 7:03 & 7:25 & 7:46 & 8:01 & 8:20 & 8:32 & 8:56 & 9:17 & 9:33 & 9:55 & 10:16 & 10:31 & 10:50 \\
\bottomrule\end{tabular}\end{center}
\subsection*{To Camden (read down) --- trips 14--26 of 45}
\begin{center}\footnotesize\setlength{\tabcolsep}{3pt}\begin{tabular}{@{}lrrrrrrrrrrrrr@{}}\toprule
Station & 14 & 15 & 16 & 17 & 18 & 19 & 20 & 21 & 22 & 23 & 24 & 25 & 26 \\\midrule
Trenton TC & 9:55 & 10:16 & 10:37 & 10:56 & 11:17 & 11:52 & 12:25 & 12:46 & 13:07 & 13:26 & 13:47 & 14:22 & 14:55 \\
Hamilton Avenue & 9:57 & 10:18 & 10:39 & 10:58 & 11:19 & 11:54 & 12:27 & 12:48 & 13:09 & 13:28 & 13:49 & 14:24 & 14:57 \\
Cass Street & 9:59 & 10:20 & 10:41 & 11:00 & 11:21 & 11:56 & 12:29 & 12:50 & 13:11 & 13:30 & 13:51 & 14:26 & 14:59 \\
Bordentown & 10:06 & 10:27 & 10:51 & 11:07 & 11:29 & 12:03 & 12:36 & 12:57 & 13:21 & 13:37 & 13:59 & 14:33 & 15:06 \\
Roebling & 10:11 & 10:35 & 10:56 & 11:12 & 11:34 & 12:08 & 12:41 & 13:05 & 13:26 & 13:42 & 14:04 & 14:38 & 15:11 \\
Florence & 10:14 & 10:38 & 10:59 & 11:15 & 11:37 & 12:11 & 12:44 & 13:08 & 13:29 & 13:45 & 14:07 & 14:41 & 15:14 \\
Burlington Towne Ctr & 10:19 & 10:43 & 11:04 & 11:20 & 11:42 & 12:16 & 12:49 & 13:13 & 13:34 & 13:50 & 14:12 & 14:46 & 15:19 \\
Burlington South & 10:21 & 10:45 & 11:06 & 11:22 & 11:44 & 12:18 & 12:51 & 13:15 & 13:36 & 13:52 & 14:14 & 14:48 & 15:21 \\
Beverly/Edgewater Park & 10:25 & 10:49 & 11:10 & 11:26 & 11:48 & 12:22 & 12:55 & 13:19 & 13:40 & 13:56 & 14:18 & 14:52 & 15:25 \\
Delanco & 10:29 & 10:53 & 11:14 & 11:30 & 11:52 & 12:26 & 12:59 & 13:23 & 13:44 & 14:00 & 14:22 & 14:56 & 15:29 \\
Riverside & 10:31 & 10:55 & 11:16 & 11:32 & 11:54 & 12:28 & 13:01 & 13:25 & 13:46 & 14:02 & 14:24 & 14:58 & 15:31 \\
Cinnaminson & 10:36 & 11:00 & 11:21 & 11:37 & 11:59 & 12:33 & 13:06 & 13:30 & 13:51 & 14:07 & 14:29 & 15:03 & 15:36 \\
Riverton & 10:38 & 11:02 & 11:23 & 11:39 & 12:01 & 12:35 & 13:08 & 13:32 & 13:53 & 14:09 & 14:31 & 15:05 & 15:38 \\
Palmyra & 10:40 & 11:04 & 11:25 & 11:41 & 12:03 & 12:37 & 13:10 & 13:34 & 13:55 & 14:11 & 14:33 & 15:07 & 15:40 \\
Route 73/Pennsauken & 10:43 & 11:07 & 11:28 & 11:44 & 12:06 & 12:40 & 13:13 & 13:37 & 13:58 & 14:14 & 14:36 & 15:10 & 15:43 \\
Pennsauken TC & 10:45 & 11:09 & 11:30 & 11:46 & 12:08 & 12:42 & 13:15 & 13:39 & 14:00 & 14:16 & 14:38 & 15:12 & 15:45 \\
36th Street & 10:49 & 11:13 & 11:34 & 11:50 & 12:12 & 12:48 & 13:19 & 13:43 & 14:04 & 14:20 & 14:42 & 15:18 & 15:49 \\
Walter Rand TC & 10:54 & 11:18 & 11:39 & 11:55 & 12:17 & 12:53 & 13:24 & 13:48 & 14:09 & 14:25 & 14:47 & 15:23 & 15:54 \\
Cooper St/Rutgers & 10:57 & 11:21 & 11:42 & 11:58 & 12:20 & 12:56 & 13:27 & 13:51 & 14:12 & 14:28 & 14:50 & 15:26 & 15:57 \\
Aquarium & 10:59 & 11:23 & 11:44 & 12:00 & 12:22 & 12:58 & 13:29 & 13:53 & 14:14 & 14:30 & 14:52 & 15:28 & 15:59 \\
Entertainment Center & 11:02 & 11:26 & 11:47 & 12:03 & 12:25 & 13:01 & 13:32 & 13:56 & 14:17 & 14:33 & 14:55 & 15:31 & 16:02 \\
\bottomrule\end{tabular}\end{center}
\subsection*{To Camden (read down) --- trips 27--39 of 45}
\begin{center}\footnotesize\setlength{\tabcolsep}{3pt}\begin{tabular}{@{}lrrrrrrrrrrrrr@{}}\toprule
Station & 27 & 28 & 29 & 30 & 31 & 32 & 33 & 34 & 35 & 36 & 37 & 38 & 39 \\\midrule
Trenton TC & 15:16 & 15:37 & 15:56 & 16:17 & 16:38 & 16:52 & 17:10 & 17:25 & 17:46 & 18:07 & 18:26 & 18:47 & 19:08 \\
Hamilton Avenue & 15:18 & 15:39 & 15:58 & 16:19 & 16:40 & 16:54 & 17:12 & 17:27 & 17:48 & 18:09 & 18:28 & 18:49 & 19:10 \\
Cass Street & 15:20 & 15:41 & 16:00 & 16:21 & 16:42 & 16:56 & 17:14 & 17:29 & 17:50 & 18:11 & 18:30 & 18:51 & 19:12 \\
Bordentown & 15:27 & 15:51 & 16:07 & 16:29 & 16:49 & 17:03 & 17:21 & 17:36 & 17:57 & 18:21 & 18:37 & 18:59 & 19:19 \\
Roebling & 15:35 & 15:56 & 16:12 & 16:34 & 16:55 & 17:08 & 17:29 & 17:41 & 18:05 & 18:26 & 18:42 & 19:04 & 19:25 \\
Florence & 15:38 & 15:59 & 16:15 & 16:37 & 16:58 & 17:11 & 17:32 & 17:44 & 18:08 & 18:29 & 18:45 & 19:07 & 19:28 \\
Burlington Towne Ctr & 15:43 & 16:04 & 16:20 & 16:42 & 17:03 & 17:16 & 17:37 & 17:49 & 18:13 & 18:34 & 18:50 & 19:12 & 19:33 \\
Burlington South & 15:45 & 16:06 & 16:22 & 16:44 & 17:05 & 17:18 & 17:39 & 17:51 & 18:15 & 18:36 & 18:52 & 19:14 & 19:35 \\
Beverly/Edgewater Park & 15:49 & 16:10 & 16:26 & 16:48 & 17:09 & 17:22 & 17:43 & 17:55 & 18:19 & 18:40 & 18:56 & 19:18 & 19:39 \\
Delanco & 15:53 & 16:14 & 16:30 & 16:52 & 17:13 & 17:26 & 17:47 & 17:59 & 18:23 & 18:44 & 19:00 & 19:22 & 19:43 \\
Riverside & 15:55 & 16:16 & 16:32 & 16:54 & 17:15 & 17:28 & 17:49 & 18:01 & 18:25 & 18:46 & 19:02 & 19:24 & 19:45 \\
Cinnaminson & 16:00 & 16:21 & 16:37 & 16:59 & 17:20 & 17:33 & 17:54 & 18:06 & 18:30 & 18:51 & 19:07 & 19:29 & 19:50 \\
Riverton & 16:02 & 16:23 & 16:39 & 17:01 & 17:22 & 17:35 & 17:56 & 18:08 & 18:32 & 18:53 & 19:09 & 19:31 & 19:52 \\
Palmyra & 16:04 & 16:25 & 16:41 & 17:03 & 17:24 & 17:37 & 17:58 & 18:10 & 18:34 & 18:55 & 19:11 & 19:33 & 19:54 \\
Route 73/Pennsauken & 16:07 & 16:28 & 16:44 & 17:06 & 17:27 & 17:40 & 18:01 & 18:13 & 18:37 & 18:58 & 19:14 & 19:36 & 19:57 \\
Pennsauken TC & 16:09 & 16:30 & 16:46 & 17:08 & 17:29 & 17:42 & 18:03 & 18:15 & 18:39 & 19:00 & 19:16 & 19:38 & 19:59 \\
36th Street & 16:13 & 16:34 & 16:50 & 17:12 & 17:33 & 17:48 & 18:07 & 18:19 & 18:43 & 19:04 & 19:20 & 19:42 & 20:03 \\
Walter Rand TC & 16:18 & 16:39 & 16:55 & 17:17 & 17:38 & 17:53 & 18:12 & 18:24 & 18:48 & 19:09 & 19:25 & 19:47 & 20:08 \\
Cooper St/Rutgers & 16:21 & 16:42 & 16:58 & 17:20 & 17:41 & 17:56 & 18:15 & 18:27 & 18:51 & 19:12 & 19:28 & 19:50 & 20:11 \\
Aquarium & 16:23 & 16:44 & 17:00 & 17:22 & 17:43 & 17:58 & 18:17 & 18:29 & 18:53 & 19:14 & 19:30 & 19:52 & 20:13 \\
Entertainment Center & 16:26 & 16:47 & 17:03 & 17:25 & 17:46 & 18:01 & 18:20 & 18:32 & 18:56 & 19:17 & 19:33 & 19:55 & 20:16 \\
\bottomrule\end{tabular}\end{center}
\subsection*{To Camden (read down) --- trips 40--45 of 45}
\begin{center}\footnotesize\setlength{\tabcolsep}{3pt}\begin{tabular}{@{}lrrrrrr@{}}\toprule
Station & 40 & 41 & 42 & 43 & 44 & 45 \\\midrule
Trenton TC & 19:22 & 19:40 & 19:55 & 20:16 & 20:37 & 20:56 \\
Hamilton Avenue & 19:24 & 19:42 & 19:57 & 20:18 & 20:39 & 20:58 \\
Cass Street & 19:26 & 19:44 & 19:59 & 20:20 & 20:41 & 21:00 \\
Bordentown & 19:33 & 19:51 & 20:06 & 20:27 & 20:51 & 21:07 \\
Roebling & 19:38 & 19:59 & 20:11 & 20:35 & 20:56 & 21:12 \\
Florence & 19:41 & 20:02 & 20:14 & 20:38 & 20:59 & 21:15 \\
Burlington Towne Ctr & 19:46 & 20:07 & 20:19 & 20:43 & 21:04 & 21:20 \\
Burlington South & 19:48 & 20:09 & 20:21 & 20:45 & 21:06 & 21:22 \\
Beverly/Edgewater Park & 19:52 & 20:13 & 20:25 & 20:49 & 21:10 & 21:26 \\
Delanco & 19:56 & 20:17 & 20:29 & 20:53 & 21:14 & 21:30 \\
Riverside & 19:58 & 20:19 & 20:31 & 20:55 & 21:16 & 21:32 \\
Cinnaminson & 20:03 & 20:24 & 20:36 & 21:00 & 21:21 & 21:37 \\
Riverton & 20:05 & 20:26 & 20:38 & 21:02 & 21:23 & 21:39 \\
Palmyra & 20:07 & 20:28 & 20:40 & 21:04 & 21:25 & 21:41 \\
Route 73/Pennsauken & 20:10 & 20:31 & 20:43 & 21:07 & 21:28 & 21:44 \\
Pennsauken TC & 20:12 & 20:33 & 20:45 & 21:09 & 21:30 & 21:46 \\
36th Street & 20:18 & 20:37 & 20:49 & 21:13 & 21:34 & 21:50 \\
Walter Rand TC & 20:23 & 20:42 & 20:54 & 21:18 & 21:39 & 21:55 \\
Cooper St/Rutgers & 20:26 & 20:45 & 20:57 & 21:21 & 21:42 & 21:58 \\
Aquarium & 20:28 & 20:47 & 20:59 & 21:23 & 21:44 & 22:00 \\
Entertainment Center & 20:31 & 20:50 & 21:02 & 21:26 & 21:47 & 22:03 \\
\bottomrule\end{tabular}\end{center}
\clearpage
\section{Proposed T7 (7 peak, 6 midday)}\label{tt:proposed_T7}
\subsection*{To Trenton (read down) --- trips 1--13 of 40}
\begin{center}\footnotesize\setlength{\tabcolsep}{3pt}\begin{tabular}{@{}lrrrrrrrrrrrrr@{}}\toprule
Station & 1 & 2 & 3 & 4 & 5 & 6 & 7 & 8 & 9 & 10 & 11 & 12 & 13 \\\midrule
Entertainment Center & 5:42 & 6:08 & 6:27 & 6:46 & 7:09 & 7:30 & 7:51 & 8:12 & 8:38 & 8:57 & 9:16 & 9:39 & 10:00 \\
Aquarium & 5:44 & 6:10 & 6:29 & 6:48 & 7:11 & 7:32 & 7:53 & 8:14 & 8:40 & 8:59 & 9:18 & 9:41 & 10:02 \\
Cooper St/Rutgers & 5:46 & 6:12 & 6:31 & 6:50 & 7:13 & 7:34 & 7:55 & 8:16 & 8:42 & 9:01 & 9:20 & 9:43 & 10:04 \\
Walter Rand TC & 5:50 & 6:16 & 6:35 & 6:54 & 7:17 & 7:38 & 7:59 & 8:20 & 8:46 & 9:05 & 9:24 & 9:47 & 10:08 \\
36th Street & 5:55 & 6:21 & 6:40 & 6:59 & 7:23 & 7:43 & 8:09 & 8:25 & 8:51 & 9:10 & 9:29 & 9:53 & 10:13 \\
Pennsauken TC & 5:57 & 6:23 & 6:42 & 7:01 & 7:25 & 7:45 & 8:11 & 8:27 & 8:53 & 9:12 & 9:31 & 9:55 & 10:15 \\
Route 73/Pennsauken & 6:00 & 6:26 & 6:45 & 7:04 & 7:28 & 7:48 & 8:14 & 8:30 & 8:56 & 9:15 & 9:34 & 9:58 & 10:18 \\
Palmyra & 6:03 & 6:29 & 6:48 & 7:07 & 7:31 & 7:51 & 8:17 & 8:33 & 8:59 & 9:18 & 9:37 & 10:01 & 10:21 \\
Riverton & 6:05 & 6:31 & 6:50 & 7:09 & 7:33 & 7:53 & 8:19 & 8:35 & 9:01 & 9:20 & 9:39 & 10:03 & 10:23 \\
Cinnaminson & 6:07 & 6:33 & 6:52 & 7:11 & 7:35 & 7:55 & 8:21 & 8:37 & 9:03 & 9:22 & 9:41 & 10:05 & 10:25 \\
Riverside & 6:12 & 6:38 & 6:57 & 7:16 & 7:40 & 8:00 & 8:26 & 8:42 & 9:08 & 9:27 & 9:46 & 10:10 & 10:30 \\
Delanco & 6:14 & 6:40 & 6:59 & 7:18 & 7:42 & 8:02 & 8:28 & 8:44 & 9:10 & 9:29 & 9:48 & 10:12 & 10:32 \\
Beverly/Edgewater Park & 6:18 & 6:44 & 7:03 & 7:22 & 7:46 & 8:06 & 8:32 & 8:48 & 9:14 & 9:33 & 9:52 & 10:16 & 10:36 \\
Burlington South & 6:22 & 6:50 & 7:07 & 7:26 & 7:50 & 8:10 & 8:36 & 8:52 & 9:20 & 9:37 & 9:56 & 10:20 & 10:40 \\
Burlington Towne Ctr & 6:24 & 6:52 & 7:09 & 7:28 & 7:52 & 8:12 & 8:38 & 8:54 & 9:22 & 9:39 & 9:58 & 10:22 & 10:42 \\
Florence & 6:30 & 6:58 & 7:15 & 7:34 & 7:58 & 8:18 & 8:44 & 9:00 & 9:28 & 9:45 & 10:04 & 10:28 & 10:48 \\
Roebling & 6:33 & 7:01 & 7:18 & 7:37 & 8:01 & 8:22 & 8:47 & 9:03 & 9:31 & 9:48 & 10:07 & 10:31 & 10:52 \\
Bordentown & 6:38 & 7:06 & 7:23 & 7:42 & 8:06 & 8:27 & 8:52 & 9:08 & 9:36 & 9:53 & 10:12 & 10:36 & 10:57 \\
Cass Street & 6:44 & 7:12 & 7:29 & 7:48 & 8:12 & 8:33 & 8:59 & 9:14 & 9:42 & 9:59 & 10:18 & 10:42 & 11:03 \\
Hamilton Avenue & 6:46 & 7:14 & 7:31 & 7:50 & 8:14 & 8:35 & 9:01 & 9:16 & 9:44 & 10:01 & 10:20 & 10:44 & 11:05 \\
Trenton TC & 6:49 & 7:17 & 7:34 & 7:53 & 8:17 & 8:38 & 9:04 & 9:19 & 9:47 & 10:04 & 10:23 & 10:47 & 11:08 \\
\bottomrule\end{tabular}\end{center}
\subsection*{To Trenton (read down) --- trips 14--26 of 40}
\begin{center}\footnotesize\setlength{\tabcolsep}{3pt}\begin{tabular}{@{}lrrrrrrrrrrrrr@{}}\toprule
Station & 14 & 15 & 16 & 17 & 18 & 19 & 20 & 21 & 22 & 23 & 24 & 25 & 26 \\\midrule
Entertainment Center & 10:21 & 10:42 & 11:08 & 11:46 & 12:09 & 12:30 & 12:51 & 13:12 & 13:38 & 14:16 & 14:39 & 15:00 & 15:21 \\
Aquarium & 10:23 & 10:44 & 11:10 & 11:48 & 12:11 & 12:32 & 12:53 & 13:14 & 13:40 & 14:18 & 14:41 & 15:02 & 15:23 \\
Cooper St/Rutgers & 10:25 & 10:46 & 11:12 & 11:50 & 12:13 & 12:34 & 12:55 & 13:16 & 13:42 & 14:20 & 14:43 & 15:04 & 15:25 \\
Walter Rand TC & 10:29 & 10:50 & 11:16 & 11:54 & 12:17 & 12:38 & 12:59 & 13:20 & 13:46 & 14:24 & 14:47 & 15:08 & 15:29 \\
36th Street & 10:39 & 10:55 & 11:21 & 11:59 & 12:23 & 12:43 & 13:09 & 13:25 & 13:51 & 14:29 & 14:53 & 15:13 & 15:39 \\
Pennsauken TC & 10:41 & 10:57 & 11:23 & 12:01 & 12:25 & 12:45 & 13:11 & 13:27 & 13:53 & 14:31 & 14:55 & 15:15 & 15:41 \\
Route 73/Pennsauken & 10:44 & 11:00 & 11:26 & 12:04 & 12:28 & 12:48 & 13:14 & 13:30 & 13:56 & 14:34 & 14:58 & 15:18 & 15:44 \\
Palmyra & 10:47 & 11:03 & 11:29 & 12:07 & 12:31 & 12:51 & 13:17 & 13:33 & 13:59 & 14:37 & 15:01 & 15:21 & 15:47 \\
Riverton & 10:49 & 11:05 & 11:31 & 12:09 & 12:33 & 12:53 & 13:19 & 13:35 & 14:01 & 14:39 & 15:03 & 15:23 & 15:49 \\
Cinnaminson & 10:51 & 11:07 & 11:33 & 12:11 & 12:35 & 12:55 & 13:21 & 13:37 & 14:03 & 14:41 & 15:05 & 15:25 & 15:51 \\
Riverside & 10:56 & 11:12 & 11:38 & 12:16 & 12:40 & 13:00 & 13:26 & 13:42 & 14:08 & 14:46 & 15:10 & 15:30 & 15:56 \\
Delanco & 10:58 & 11:14 & 11:40 & 12:18 & 12:42 & 13:02 & 13:28 & 13:44 & 14:10 & 14:48 & 15:12 & 15:32 & 15:58 \\
Beverly/Edgewater Park & 11:02 & 11:18 & 11:44 & 12:22 & 12:46 & 13:06 & 13:32 & 13:48 & 14:14 & 14:52 & 15:16 & 15:36 & 16:02 \\
Burlington South & 11:06 & 11:22 & 11:50 & 12:26 & 12:50 & 13:10 & 13:36 & 13:52 & 14:20 & 14:56 & 15:20 & 15:40 & 16:06 \\
Burlington Towne Ctr & 11:08 & 11:24 & 11:52 & 12:28 & 12:52 & 13:12 & 13:38 & 13:54 & 14:22 & 14:58 & 15:22 & 15:42 & 16:08 \\
Florence & 11:14 & 11:30 & 11:58 & 12:34 & 12:58 & 13:18 & 13:44 & 14:00 & 14:28 & 15:04 & 15:28 & 15:48 & 16:14 \\
Roebling & 11:17 & 11:33 & 12:01 & 12:37 & 13:01 & 13:22 & 13:47 & 14:03 & 14:31 & 15:07 & 15:31 & 15:52 & 16:17 \\
Bordentown & 11:22 & 11:38 & 12:06 & 12:42 & 13:06 & 13:27 & 13:52 & 14:08 & 14:36 & 15:12 & 15:36 & 15:57 & 16:22 \\
Cass Street & 11:29 & 11:44 & 12:12 & 12:48 & 13:12 & 13:33 & 13:59 & 14:14 & 14:42 & 15:18 & 15:42 & 16:03 & 16:29 \\
Hamilton Avenue & 11:31 & 11:46 & 12:14 & 12:50 & 13:14 & 13:35 & 14:01 & 14:16 & 14:44 & 15:20 & 15:44 & 16:05 & 16:31 \\
Trenton TC & 11:34 & 11:49 & 12:17 & 12:53 & 13:17 & 13:38 & 14:04 & 14:19 & 14:47 & 15:23 & 15:47 & 16:08 & 16:34 \\
\bottomrule\end{tabular}\end{center}
\subsection*{To Trenton (read down) --- trips 27--39 of 40}
\begin{center}\footnotesize\setlength{\tabcolsep}{3pt}\begin{tabular}{@{}lrrrrrrrrrrrrr@{}}\toprule
Station & 27 & 28 & 29 & 30 & 31 & 32 & 33 & 34 & 35 & 36 & 37 & 38 & 39 \\\midrule
Entertainment Center & 15:42 & 16:08 & 16:27 & 16:46 & 17:09 & 17:30 & 17:51 & 18:12 & 18:38 & 19:16 & 19:39 & 20:00 & 20:21 \\
Aquarium & 15:44 & 16:10 & 16:29 & 16:48 & 17:11 & 17:32 & 17:53 & 18:14 & 18:40 & 19:18 & 19:41 & 20:02 & 20:23 \\
Cooper St/Rutgers & 15:46 & 16:12 & 16:31 & 16:50 & 17:13 & 17:34 & 17:55 & 18:16 & 18:42 & 19:20 & 19:43 & 20:04 & 20:25 \\
Walter Rand TC & 15:50 & 16:16 & 16:35 & 16:54 & 17:17 & 17:38 & 17:59 & 18:20 & 18:46 & 19:24 & 19:47 & 20:08 & 20:29 \\
36th Street & 15:55 & 16:21 & 16:40 & 16:59 & 17:23 & 17:43 & 18:09 & 18:25 & 18:51 & 19:29 & 19:53 & 20:13 & 20:39 \\
Pennsauken TC & 15:57 & 16:23 & 16:42 & 17:01 & 17:25 & 17:45 & 18:11 & 18:27 & 18:53 & 19:31 & 19:55 & 20:15 & 20:41 \\
Route 73/Pennsauken & 16:00 & 16:26 & 16:45 & 17:04 & 17:28 & 17:48 & 18:14 & 18:30 & 18:56 & 19:34 & 19:58 & 20:18 & 20:44 \\
Palmyra & 16:03 & 16:29 & 16:48 & 17:07 & 17:31 & 17:51 & 18:17 & 18:33 & 18:59 & 19:37 & 20:01 & 20:21 & 20:47 \\
Riverton & 16:05 & 16:31 & 16:50 & 17:09 & 17:33 & 17:53 & 18:19 & 18:35 & 19:01 & 19:39 & 20:03 & 20:23 & 20:49 \\
Cinnaminson & 16:07 & 16:33 & 16:52 & 17:11 & 17:35 & 17:55 & 18:21 & 18:37 & 19:03 & 19:41 & 20:05 & 20:25 & 20:51 \\
Riverside & 16:12 & 16:38 & 16:57 & 17:16 & 17:40 & 18:00 & 18:26 & 18:42 & 19:08 & 19:46 & 20:10 & 20:30 & 20:56 \\
Delanco & 16:14 & 16:40 & 16:59 & 17:18 & 17:42 & 18:02 & 18:28 & 18:44 & 19:10 & 19:48 & 20:12 & 20:32 & 20:58 \\
Beverly/Edgewater Park & 16:18 & 16:44 & 17:03 & 17:22 & 17:46 & 18:06 & 18:32 & 18:48 & 19:14 & 19:52 & 20:16 & 20:36 & 21:02 \\
Burlington South & 16:22 & 16:50 & 17:07 & 17:26 & 17:50 & 18:10 & 18:36 & 18:52 & 19:20 & 19:56 & 20:20 & 20:40 & 21:06 \\
Burlington Towne Ctr & 16:24 & 16:52 & 17:09 & 17:28 & 17:52 & 18:12 & 18:38 & 18:54 & 19:22 & 19:58 & 20:22 & 20:42 & 21:08 \\
Florence & 16:30 & 16:58 & 17:15 & 17:34 & 17:58 & 18:18 & 18:44 & 19:00 & 19:28 & 20:04 & 20:28 & 20:48 & 21:14 \\
Roebling & 16:33 & 17:01 & 17:18 & 17:37 & 18:01 & 18:22 & 18:47 & 19:03 & 19:31 & 20:07 & 20:31 & 20:52 & 21:17 \\
Bordentown & 16:38 & 17:06 & 17:23 & 17:42 & 18:06 & 18:27 & 18:52 & 19:08 & 19:36 & 20:12 & 20:36 & 20:57 & 21:22 \\
Cass Street & 16:44 & 17:12 & 17:29 & 17:48 & 18:12 & 18:33 & 18:59 & 19:14 & 19:42 & 20:18 & 20:42 & 21:03 & 21:29 \\
Hamilton Avenue & 16:46 & 17:14 & 17:31 & 17:50 & 18:14 & 18:35 & 19:01 & 19:16 & 19:44 & 20:20 & 20:44 & 21:05 & 21:31 \\
Trenton TC & 16:49 & 17:17 & 17:34 & 17:53 & 18:17 & 18:38 & 19:04 & 19:19 & 19:47 & 20:23 & 20:47 & 21:08 & 21:34 \\
\bottomrule\end{tabular}\end{center}
\subsection*{To Trenton (read down) --- trips 40--40 of 40}
\begin{center}\footnotesize\setlength{\tabcolsep}{3pt}\begin{tabular}{@{}lr@{}}\toprule
Station & 40 \\\midrule
Entertainment Center & 20:42 \\
Aquarium & 20:44 \\
Cooper St/Rutgers & 20:46 \\
Walter Rand TC & 20:50 \\
36th Street & 20:55 \\
Pennsauken TC & 20:57 \\
Route 73/Pennsauken & 21:00 \\
Palmyra & 21:03 \\
Riverton & 21:05 \\
Cinnaminson & 21:07 \\
Riverside & 21:12 \\
Delanco & 21:14 \\
Beverly/Edgewater Park & 21:18 \\
Burlington South & 21:22 \\
Burlington Towne Ctr & 21:24 \\
Florence & 21:30 \\
Roebling & 21:33 \\
Bordentown & 21:38 \\
Cass Street & 21:44 \\
Hamilton Avenue & 21:46 \\
Trenton TC & 21:49 \\
\bottomrule\end{tabular}\end{center}
\subsection*{To Camden (read down) --- trips 1--13 of 40}
\begin{center}\footnotesize\setlength{\tabcolsep}{3pt}\begin{tabular}{@{}lrrrrrrrrrrrrr@{}}\toprule
Station & 1 & 2 & 3 & 4 & 5 & 6 & 7 & 8 & 9 & 10 & 11 & 12 & 13 \\\midrule
Trenton TC & 5:52 & 6:13 & 6:39 & 6:56 & 7:22 & 7:40 & 8:01 & 8:22 & 8:43 & 9:09 & 9:26 & 9:52 & 10:10 \\
Hamilton Avenue & 5:54 & 6:15 & 6:41 & 6:58 & 7:24 & 7:42 & 8:03 & 8:24 & 8:45 & 9:11 & 9:28 & 9:54 & 10:12 \\
Cass Street & 5:56 & 6:17 & 6:43 & 7:00 & 7:26 & 7:44 & 8:05 & 8:26 & 8:47 & 9:13 & 9:30 & 9:56 & 10:14 \\
Bordentown & 6:04 & 6:24 & 6:50 & 7:07 & 7:33 & 7:51 & 8:13 & 8:34 & 8:54 & 9:20 & 9:37 & 10:03 & 10:21 \\
Roebling & 6:10 & 6:29 & 6:55 & 7:12 & 7:42 & 7:56 & 8:18 & 8:40 & 8:59 & 9:25 & 9:42 & 10:12 & 10:26 \\
Florence & 6:13 & 6:32 & 6:58 & 7:15 & 7:45 & 7:59 & 8:21 & 8:43 & 9:02 & 9:28 & 9:45 & 10:15 & 10:29 \\
Burlington Towne Ctr & 6:18 & 6:37 & 7:03 & 7:20 & 7:50 & 8:04 & 8:26 & 8:48 & 9:07 & 9:33 & 9:50 & 10:20 & 10:34 \\
Burlington South & 6:20 & 6:39 & 7:05 & 7:22 & 7:52 & 8:06 & 8:28 & 8:50 & 9:09 & 9:35 & 9:52 & 10:22 & 10:36 \\
Beverly/Edgewater Park & 6:24 & 6:43 & 7:09 & 7:26 & 7:56 & 8:10 & 8:33 & 8:54 & 9:13 & 9:39 & 9:56 & 10:26 & 10:40 \\
Delanco & 6:28 & 6:47 & 7:13 & 7:30 & 8:00 & 8:14 & 8:37 & 8:58 & 9:17 & 9:43 & 10:00 & 10:30 & 10:44 \\
Riverside & 6:30 & 6:49 & 7:15 & 7:32 & 8:02 & 8:16 & 8:39 & 9:00 & 9:19 & 9:45 & 10:02 & 10:32 & 10:46 \\
Cinnaminson & 6:35 & 6:54 & 7:20 & 7:37 & 8:07 & 8:21 & 8:45 & 9:05 & 9:24 & 9:50 & 10:07 & 10:37 & 10:51 \\
Riverton & 6:37 & 6:56 & 7:22 & 7:39 & 8:09 & 8:23 & 8:47 & 9:07 & 9:26 & 9:52 & 10:09 & 10:39 & 10:53 \\
Palmyra & 6:39 & 6:58 & 7:24 & 7:41 & 8:11 & 8:25 & 8:49 & 9:09 & 9:28 & 9:54 & 10:11 & 10:41 & 10:55 \\
Route 73/Pennsauken & 6:42 & 7:01 & 7:27 & 7:44 & 8:14 & 8:28 & 8:52 & 9:12 & 9:31 & 9:57 & 10:14 & 10:44 & 10:58 \\
Pennsauken TC & 6:44 & 7:03 & 7:29 & 7:46 & 8:16 & 8:30 & 8:54 & 9:14 & 9:33 & 9:59 & 10:16 & 10:46 & 11:00 \\
36th Street & 6:48 & 7:07 & 7:33 & 7:50 & 8:20 & 8:34 & 8:58 & 9:18 & 9:37 & 10:03 & 10:20 & 10:50 & 11:04 \\
Walter Rand TC & 6:53 & 7:12 & 7:38 & 7:55 & 8:25 & 8:39 & 9:03 & 9:23 & 9:42 & 10:08 & 10:25 & 10:55 & 11:09 \\
Cooper St/Rutgers & 6:56 & 7:15 & 7:41 & 7:58 & 8:28 & 8:42 & 9:06 & 9:26 & 9:45 & 10:11 & 10:28 & 10:58 & 11:12 \\
Aquarium & 6:58 & 7:17 & 7:43 & 8:00 & 8:30 & 8:44 & 9:08 & 9:28 & 9:47 & 10:13 & 10:30 & 11:00 & 11:14 \\
Entertainment Center & 7:01 & 7:20 & 7:46 & 8:03 & 8:33 & 8:47 & 9:11 & 9:31 & 9:50 & 10:16 & 10:33 & 11:03 & 11:17 \\
\bottomrule\end{tabular}\end{center}
\subsection*{To Camden (read down) --- trips 14--26 of 40}
\begin{center}\footnotesize\setlength{\tabcolsep}{3pt}\begin{tabular}{@{}lrrrrrrrrrrrrr@{}}\toprule
Station & 14 & 15 & 16 & 17 & 18 & 19 & 20 & 21 & 22 & 23 & 24 & 25 & 26 \\\midrule
Trenton TC & 10:31 & 10:52 & 11:13 & 11:39 & 11:56 & 12:22 & 13:01 & 13:22 & 13:43 & 14:09 & 14:26 & 14:52 & 15:31 \\
Hamilton Avenue & 10:33 & 10:54 & 11:15 & 11:41 & 11:58 & 12:24 & 13:03 & 13:24 & 13:45 & 14:11 & 14:28 & 14:54 & 15:33 \\
Cass Street & 10:35 & 10:56 & 11:17 & 11:43 & 12:00 & 12:26 & 13:05 & 13:26 & 13:47 & 14:13 & 14:30 & 14:56 & 15:35 \\
Bordentown & 10:43 & 11:04 & 11:24 & 11:50 & 12:07 & 12:33 & 13:13 & 13:34 & 13:54 & 14:20 & 14:37 & 15:03 & 15:43 \\
Roebling & 10:48 & 11:10 & 11:29 & 11:55 & 12:12 & 12:42 & 13:18 & 13:40 & 13:59 & 14:25 & 14:42 & 15:12 & 15:48 \\
Florence & 10:51 & 11:13 & 11:32 & 11:58 & 12:15 & 12:45 & 13:21 & 13:43 & 14:02 & 14:28 & 14:45 & 15:15 & 15:51 \\
Burlington Towne Ctr & 10:56 & 11:18 & 11:37 & 12:03 & 12:20 & 12:50 & 13:26 & 13:48 & 14:07 & 14:33 & 14:50 & 15:20 & 15:56 \\
Burlington South & 10:58 & 11:20 & 11:39 & 12:05 & 12:22 & 12:52 & 13:28 & 13:50 & 14:09 & 14:35 & 14:52 & 15:22 & 15:58 \\
Beverly/Edgewater Park & 11:03 & 11:24 & 11:43 & 12:09 & 12:26 & 12:56 & 13:33 & 13:54 & 14:13 & 14:39 & 14:56 & 15:26 & 16:03 \\
Delanco & 11:07 & 11:28 & 11:47 & 12:13 & 12:30 & 13:00 & 13:37 & 13:58 & 14:17 & 14:43 & 15:00 & 15:30 & 16:07 \\
Riverside & 11:09 & 11:30 & 11:49 & 12:15 & 12:32 & 13:02 & 13:39 & 14:00 & 14:19 & 14:45 & 15:02 & 15:32 & 16:09 \\
Cinnaminson & 11:15 & 11:35 & 11:54 & 12:20 & 12:37 & 13:07 & 13:45 & 14:05 & 14:24 & 14:50 & 15:07 & 15:37 & 16:15 \\
Riverton & 11:17 & 11:37 & 11:56 & 12:22 & 12:39 & 13:09 & 13:47 & 14:07 & 14:26 & 14:52 & 15:09 & 15:39 & 16:17 \\
Palmyra & 11:19 & 11:39 & 11:58 & 12:24 & 12:41 & 13:11 & 13:49 & 14:09 & 14:28 & 14:54 & 15:11 & 15:41 & 16:19 \\
Route 73/Pennsauken & 11:22 & 11:42 & 12:01 & 12:27 & 12:44 & 13:14 & 13:52 & 14:12 & 14:31 & 14:57 & 15:14 & 15:44 & 16:22 \\
Pennsauken TC & 11:24 & 11:44 & 12:03 & 12:29 & 12:46 & 13:16 & 13:54 & 14:14 & 14:33 & 14:59 & 15:16 & 15:46 & 16:24 \\
36th Street & 11:28 & 11:48 & 12:07 & 12:33 & 12:50 & 13:20 & 13:58 & 14:18 & 14:37 & 15:03 & 15:20 & 15:50 & 16:28 \\
Walter Rand TC & 11:33 & 11:53 & 12:12 & 12:38 & 12:55 & 13:25 & 14:03 & 14:23 & 14:42 & 15:08 & 15:25 & 15:55 & 16:33 \\
Cooper St/Rutgers & 11:36 & 11:56 & 12:15 & 12:41 & 12:58 & 13:28 & 14:06 & 14:26 & 14:45 & 15:11 & 15:28 & 15:58 & 16:36 \\
Aquarium & 11:38 & 11:58 & 12:17 & 12:43 & 13:00 & 13:30 & 14:08 & 14:28 & 14:47 & 15:13 & 15:30 & 16:00 & 16:38 \\
Entertainment Center & 11:41 & 12:01 & 12:20 & 12:46 & 13:03 & 13:33 & 14:11 & 14:31 & 14:50 & 15:16 & 15:33 & 16:03 & 16:41 \\
\bottomrule\end{tabular}\end{center}
\subsection*{To Camden (read down) --- trips 27--39 of 40}
\begin{center}\footnotesize\setlength{\tabcolsep}{3pt}\begin{tabular}{@{}lrrrrrrrrrrrrr@{}}\toprule
Station & 27 & 28 & 29 & 30 & 31 & 32 & 33 & 34 & 35 & 36 & 37 & 38 & 39 \\\midrule
Trenton TC & 15:52 & 16:13 & 16:39 & 16:56 & 17:22 & 17:40 & 18:01 & 18:22 & 18:43 & 19:09 & 19:26 & 19:52 & 20:31 \\
Hamilton Avenue & 15:54 & 16:15 & 16:41 & 16:58 & 17:24 & 17:42 & 18:03 & 18:24 & 18:45 & 19:11 & 19:28 & 19:54 & 20:33 \\
Cass Street & 15:56 & 16:17 & 16:43 & 17:00 & 17:26 & 17:44 & 18:05 & 18:26 & 18:47 & 19:13 & 19:30 & 19:56 & 20:35 \\
Bordentown & 16:04 & 16:24 & 16:50 & 17:07 & 17:33 & 17:51 & 18:13 & 18:34 & 18:54 & 19:20 & 19:37 & 20:03 & 20:43 \\
Roebling & 16:10 & 16:29 & 16:55 & 17:12 & 17:42 & 17:56 & 18:18 & 18:40 & 18:59 & 19:25 & 19:42 & 20:12 & 20:48 \\
Florence & 16:13 & 16:32 & 16:58 & 17:15 & 17:45 & 17:59 & 18:21 & 18:43 & 19:02 & 19:28 & 19:45 & 20:15 & 20:51 \\
Burlington Towne Ctr & 16:18 & 16:37 & 17:03 & 17:20 & 17:50 & 18:04 & 18:26 & 18:48 & 19:07 & 19:33 & 19:50 & 20:20 & 20:56 \\
Burlington South & 16:20 & 16:39 & 17:05 & 17:22 & 17:52 & 18:06 & 18:28 & 18:50 & 19:09 & 19:35 & 19:52 & 20:22 & 20:58 \\
Beverly/Edgewater Park & 16:24 & 16:43 & 17:09 & 17:26 & 17:56 & 18:10 & 18:33 & 18:54 & 19:13 & 19:39 & 19:56 & 20:26 & 21:03 \\
Delanco & 16:28 & 16:47 & 17:13 & 17:30 & 18:00 & 18:14 & 18:37 & 18:58 & 19:17 & 19:43 & 20:00 & 20:30 & 21:07 \\
Riverside & 16:30 & 16:49 & 17:15 & 17:32 & 18:02 & 18:16 & 18:39 & 19:00 & 19:19 & 19:45 & 20:02 & 20:32 & 21:09 \\
Cinnaminson & 16:35 & 16:54 & 17:20 & 17:37 & 18:07 & 18:21 & 18:45 & 19:05 & 19:24 & 19:50 & 20:07 & 20:37 & 21:15 \\
Riverton & 16:37 & 16:56 & 17:22 & 17:39 & 18:09 & 18:23 & 18:47 & 19:07 & 19:26 & 19:52 & 20:09 & 20:39 & 21:17 \\
Palmyra & 16:39 & 16:58 & 17:24 & 17:41 & 18:11 & 18:25 & 18:49 & 19:09 & 19:28 & 19:54 & 20:11 & 20:41 & 21:19 \\
Route 73/Pennsauken & 16:42 & 17:01 & 17:27 & 17:44 & 18:14 & 18:28 & 18:52 & 19:12 & 19:31 & 19:57 & 20:14 & 20:44 & 21:22 \\
Pennsauken TC & 16:44 & 17:03 & 17:29 & 17:46 & 18:16 & 18:30 & 18:54 & 19:14 & 19:33 & 19:59 & 20:16 & 20:46 & 21:24 \\
36th Street & 16:48 & 17:07 & 17:33 & 17:50 & 18:20 & 18:34 & 18:58 & 19:18 & 19:37 & 20:03 & 20:20 & 20:50 & 21:28 \\
Walter Rand TC & 16:53 & 17:12 & 17:38 & 17:55 & 18:25 & 18:39 & 19:03 & 19:23 & 19:42 & 20:08 & 20:25 & 20:55 & 21:33 \\
Cooper St/Rutgers & 16:56 & 17:15 & 17:41 & 17:58 & 18:28 & 18:42 & 19:06 & 19:26 & 19:45 & 20:11 & 20:28 & 20:58 & 21:36 \\
Aquarium & 16:58 & 17:17 & 17:43 & 18:00 & 18:30 & 18:44 & 19:08 & 19:28 & 19:47 & 20:13 & 20:30 & 21:00 & 21:38 \\
Entertainment Center & 17:01 & 17:20 & 17:46 & 18:03 & 18:33 & 18:47 & 19:11 & 19:31 & 19:50 & 20:16 & 20:33 & 21:03 & 21:41 \\
\bottomrule\end{tabular}\end{center}
\subsection*{To Camden (read down) --- trips 40--40 of 40}
\begin{center}\footnotesize\setlength{\tabcolsep}{3pt}\begin{tabular}{@{}lr@{}}\toprule
Station & 40 \\\midrule
Trenton TC & 20:52 \\
Hamilton Avenue & 20:54 \\
Cass Street & 20:56 \\
Bordentown & 21:04 \\
Roebling & 21:10 \\
Florence & 21:13 \\
Burlington Towne Ctr & 21:18 \\
Burlington South & 21:20 \\
Beverly/Edgewater Park & 21:24 \\
Delanco & 21:28 \\
Riverside & 21:30 \\
Cinnaminson & 21:35 \\
Riverton & 21:37 \\
Palmyra & 21:39 \\
Route 73/Pennsauken & 21:42 \\
Pennsauken TC & 21:44 \\
36th Street & 21:48 \\
Walter Rand TC & 21:53 \\
Cooper St/Rutgers & 21:56 \\
Aquarium & 21:58 \\
Entertainment Center & 22:01 \\
\bottomrule\end{tabular}\end{center}
\clearpage
\section{Proposed T6 (6 peak, 5 midday)}\label{tt:proposed_T6}
\subsection*{To Trenton (read down) --- trips 1--13 of 34}
\begin{center}\footnotesize\setlength{\tabcolsep}{3pt}\begin{tabular}{@{}lrrrrrrrrrrrrr@{}}\toprule
Station & 1 & 2 & 3 & 4 & 5 & 6 & 7 & 8 & 9 & 10 & 11 & 12 & 13 \\\midrule
Entertainment Center & 5:47 & 6:14 & 6:39 & 7:05 & 7:30 & 7:52 & 8:17 & 8:44 & 9:09 & 9:35 & 10:22 & 10:47 & 11:14 \\
Aquarium & 5:49 & 6:16 & 6:41 & 7:07 & 7:32 & 7:54 & 8:19 & 8:46 & 9:11 & 9:37 & 10:24 & 10:49 & 11:16 \\
Cooper St/Rutgers & 5:51 & 6:18 & 6:43 & 7:09 & 7:34 & 7:56 & 8:21 & 8:48 & 9:13 & 9:39 & 10:26 & 10:51 & 11:18 \\
Walter Rand TC & 5:55 & 6:22 & 6:47 & 7:13 & 7:38 & 8:00 & 8:25 & 8:52 & 9:17 & 9:43 & 10:30 & 10:55 & 11:22 \\
36th Street & 6:00 & 6:27 & 6:52 & 7:18 & 7:43 & 8:05 & 8:30 & 8:57 & 9:22 & 9:48 & 10:35 & 11:00 & 11:27 \\
Pennsauken TC & 6:02 & 6:29 & 6:54 & 7:20 & 7:45 & 8:07 & 8:32 & 8:59 & 9:24 & 9:50 & 10:37 & 11:02 & 11:29 \\
Route 73/Pennsauken & 6:05 & 6:32 & 6:57 & 7:23 & 7:51 & 8:10 & 8:35 & 9:02 & 9:27 & 9:53 & 10:40 & 11:05 & 11:32 \\
Palmyra & 6:08 & 6:35 & 7:00 & 7:26 & 7:54 & 8:13 & 8:38 & 9:05 & 9:30 & 9:56 & 10:43 & 11:08 & 11:35 \\
Riverton & 6:10 & 6:37 & 7:02 & 7:28 & 7:56 & 8:15 & 8:40 & 9:07 & 9:32 & 9:58 & 10:45 & 11:10 & 11:37 \\
Cinnaminson & 6:12 & 6:39 & 7:04 & 7:30 & 7:58 & 8:17 & 8:42 & 9:09 & 9:34 & 10:00 & 10:47 & 11:12 & 11:39 \\
Riverside & 6:17 & 6:45 & 7:11 & 7:35 & 8:03 & 8:22 & 8:47 & 9:15 & 9:41 & 10:05 & 10:52 & 11:17 & 11:45 \\
Delanco & 6:19 & 6:47 & 7:13 & 7:37 & 8:05 & 8:24 & 8:49 & 9:17 & 9:43 & 10:07 & 10:54 & 11:19 & 11:47 \\
Beverly/Edgewater Park & 6:23 & 6:51 & 7:17 & 7:41 & 8:09 & 8:28 & 8:53 & 9:21 & 9:47 & 10:11 & 10:58 & 11:23 & 11:51 \\
Burlington South & 6:27 & 6:55 & 7:21 & 7:47 & 8:13 & 8:32 & 8:57 & 9:25 & 9:51 & 10:17 & 11:02 & 11:27 & 11:55 \\
Burlington Towne Ctr & 6:29 & 6:57 & 7:23 & 7:49 & 8:15 & 8:34 & 8:59 & 9:27 & 9:53 & 10:19 & 11:04 & 11:29 & 11:57 \\
Florence & 6:35 & 7:03 & 7:29 & 7:55 & 8:21 & 8:40 & 9:05 & 9:33 & 9:59 & 10:25 & 11:10 & 11:35 & 12:03 \\
Roebling & 6:38 & 7:06 & 7:32 & 7:58 & 8:24 & 8:43 & 9:08 & 9:36 & 10:02 & 10:28 & 11:13 & 11:38 & 12:06 \\
Bordentown & 6:43 & 7:11 & 7:37 & 8:03 & 8:29 & 8:48 & 9:13 & 9:41 & 10:07 & 10:33 & 11:18 & 11:43 & 12:11 \\
Cass Street & 6:50 & 7:17 & 7:43 & 8:09 & 8:35 & 8:54 & 9:20 & 9:47 & 10:13 & 10:39 & 11:24 & 11:50 & 12:17 \\
Hamilton Avenue & 6:52 & 7:19 & 7:45 & 8:11 & 8:37 & 8:56 & 9:22 & 9:49 & 10:15 & 10:41 & 11:26 & 11:52 & 12:19 \\
Trenton TC & 6:55 & 7:22 & 7:48 & 8:14 & 8:40 & 8:59 & 9:25 & 9:52 & 10:18 & 10:44 & 11:29 & 11:55 & 12:22 \\
\bottomrule\end{tabular}\end{center}
\subsection*{To Trenton (read down) --- trips 14--26 of 34}
\begin{center}\footnotesize\setlength{\tabcolsep}{3pt}\begin{tabular}{@{}lrrrrrrrrrrrrr@{}}\toprule
Station & 14 & 15 & 16 & 17 & 18 & 19 & 20 & 21 & 22 & 23 & 24 & 25 & 26 \\\midrule
Entertainment Center & 11:39 & 12:05 & 12:52 & 13:17 & 13:44 & 14:09 & 14:35 & 15:00 & 15:22 & 15:47 & 16:14 & 16:39 & 17:05 \\
Aquarium & 11:41 & 12:07 & 12:54 & 13:19 & 13:46 & 14:11 & 14:37 & 15:02 & 15:24 & 15:49 & 16:16 & 16:41 & 17:07 \\
Cooper St/Rutgers & 11:43 & 12:09 & 12:56 & 13:21 & 13:48 & 14:13 & 14:39 & 15:04 & 15:26 & 15:51 & 16:18 & 16:43 & 17:09 \\
Walter Rand TC & 11:47 & 12:13 & 13:00 & 13:25 & 13:52 & 14:17 & 14:43 & 15:08 & 15:30 & 15:55 & 16:22 & 16:47 & 17:13 \\
36th Street & 11:52 & 12:18 & 13:05 & 13:30 & 13:57 & 14:22 & 14:48 & 15:13 & 15:35 & 16:00 & 16:27 & 16:52 & 17:18 \\
Pennsauken TC & 11:54 & 12:20 & 13:07 & 13:32 & 13:59 & 14:24 & 14:50 & 15:15 & 15:37 & 16:02 & 16:29 & 16:54 & 17:20 \\
Route 73/Pennsauken & 11:57 & 12:23 & 13:10 & 13:35 & 14:02 & 14:27 & 14:53 & 15:21 & 15:40 & 16:05 & 16:32 & 16:57 & 17:23 \\
Palmyra & 12:00 & 12:26 & 13:13 & 13:38 & 14:05 & 14:30 & 14:56 & 15:24 & 15:43 & 16:08 & 16:35 & 17:00 & 17:26 \\
Riverton & 12:02 & 12:28 & 13:15 & 13:40 & 14:07 & 14:32 & 14:58 & 15:26 & 15:45 & 16:10 & 16:37 & 17:02 & 17:28 \\
Cinnaminson & 12:04 & 12:30 & 13:17 & 13:42 & 14:09 & 14:34 & 15:00 & 15:28 & 15:47 & 16:12 & 16:39 & 17:04 & 17:30 \\
Riverside & 12:11 & 12:35 & 13:22 & 13:47 & 14:15 & 14:41 & 15:05 & 15:33 & 15:52 & 16:17 & 16:45 & 17:11 & 17:35 \\
Delanco & 12:13 & 12:37 & 13:24 & 13:49 & 14:17 & 14:43 & 15:07 & 15:35 & 15:54 & 16:19 & 16:47 & 17:13 & 17:37 \\
Beverly/Edgewater Park & 12:17 & 12:41 & 13:28 & 13:53 & 14:21 & 14:47 & 15:11 & 15:39 & 15:58 & 16:23 & 16:51 & 17:17 & 17:41 \\
Burlington South & 12:21 & 12:47 & 13:32 & 13:57 & 14:25 & 14:51 & 15:17 & 15:43 & 16:02 & 16:27 & 16:55 & 17:21 & 17:47 \\
Burlington Towne Ctr & 12:23 & 12:49 & 13:34 & 13:59 & 14:27 & 14:53 & 15:19 & 15:45 & 16:04 & 16:29 & 16:57 & 17:23 & 17:49 \\
Florence & 12:29 & 12:55 & 13:40 & 14:05 & 14:33 & 14:59 & 15:25 & 15:51 & 16:10 & 16:35 & 17:03 & 17:29 & 17:55 \\
Roebling & 12:32 & 12:58 & 13:43 & 14:08 & 14:36 & 15:02 & 15:28 & 15:54 & 16:13 & 16:38 & 17:06 & 17:32 & 17:58 \\
Bordentown & 12:37 & 13:03 & 13:48 & 14:13 & 14:41 & 15:07 & 15:33 & 15:59 & 16:18 & 16:43 & 17:11 & 17:37 & 18:03 \\
Cass Street & 12:43 & 13:09 & 13:54 & 14:20 & 14:47 & 15:13 & 15:39 & 16:05 & 16:24 & 16:50 & 17:17 & 17:43 & 18:09 \\
Hamilton Avenue & 12:45 & 13:11 & 13:56 & 14:22 & 14:49 & 15:15 & 15:41 & 16:07 & 16:26 & 16:52 & 17:19 & 17:45 & 18:11 \\
Trenton TC & 12:48 & 13:14 & 13:59 & 14:25 & 14:52 & 15:18 & 15:44 & 16:10 & 16:29 & 16:55 & 17:22 & 17:48 & 18:14 \\
\bottomrule\end{tabular}\end{center}
\subsection*{To Trenton (read down) --- trips 27--34 of 34}
\begin{center}\footnotesize\setlength{\tabcolsep}{3pt}\begin{tabular}{@{}lrrrrrrrr@{}}\toprule
Station & 27 & 28 & 29 & 30 & 31 & 32 & 33 & 34 \\\midrule
Entertainment Center & 17:30 & 17:52 & 18:17 & 18:44 & 19:09 & 19:35 & 20:22 & 20:47 \\
Aquarium & 17:32 & 17:54 & 18:19 & 18:46 & 19:11 & 19:37 & 20:24 & 20:49 \\
Cooper St/Rutgers & 17:34 & 17:56 & 18:21 & 18:48 & 19:13 & 19:39 & 20:26 & 20:51 \\
Walter Rand TC & 17:38 & 18:00 & 18:25 & 18:52 & 19:17 & 19:43 & 20:30 & 20:55 \\
36th Street & 17:43 & 18:05 & 18:30 & 18:57 & 19:22 & 19:48 & 20:35 & 21:00 \\
Pennsauken TC & 17:45 & 18:07 & 18:32 & 18:59 & 19:24 & 19:50 & 20:37 & 21:02 \\
Route 73/Pennsauken & 17:51 & 18:10 & 18:35 & 19:02 & 19:27 & 19:53 & 20:40 & 21:05 \\
Palmyra & 17:54 & 18:13 & 18:38 & 19:05 & 19:30 & 19:56 & 20:43 & 21:08 \\
Riverton & 17:56 & 18:15 & 18:40 & 19:07 & 19:32 & 19:58 & 20:45 & 21:10 \\
Cinnaminson & 17:58 & 18:17 & 18:42 & 19:09 & 19:34 & 20:00 & 20:47 & 21:12 \\
Riverside & 18:03 & 18:22 & 18:47 & 19:15 & 19:41 & 20:05 & 20:52 & 21:17 \\
Delanco & 18:05 & 18:24 & 18:49 & 19:17 & 19:43 & 20:07 & 20:54 & 21:19 \\
Beverly/Edgewater Park & 18:09 & 18:28 & 18:53 & 19:21 & 19:47 & 20:11 & 20:58 & 21:23 \\
Burlington South & 18:13 & 18:32 & 18:57 & 19:25 & 19:51 & 20:17 & 21:02 & 21:27 \\
Burlington Towne Ctr & 18:15 & 18:34 & 18:59 & 19:27 & 19:53 & 20:19 & 21:04 & 21:29 \\
Florence & 18:21 & 18:40 & 19:05 & 19:33 & 19:59 & 20:25 & 21:10 & 21:35 \\
Roebling & 18:24 & 18:43 & 19:08 & 19:36 & 20:02 & 20:28 & 21:13 & 21:38 \\
Bordentown & 18:29 & 18:48 & 19:13 & 19:41 & 20:07 & 20:33 & 21:18 & 21:43 \\
Cass Street & 18:35 & 18:54 & 19:20 & 19:47 & 20:13 & 20:39 & 21:24 & 21:50 \\
Hamilton Avenue & 18:37 & 18:56 & 19:22 & 19:49 & 20:15 & 20:41 & 21:26 & 21:52 \\
Trenton TC & 18:40 & 18:59 & 19:25 & 19:52 & 20:18 & 20:44 & 21:29 & 21:55 \\
\bottomrule\end{tabular}\end{center}
\subsection*{To Camden (read down) --- trips 1--13 of 35}
\begin{center}\footnotesize\setlength{\tabcolsep}{3pt}\begin{tabular}{@{}lrrrrrrrrrrrrr@{}}\toprule
Station & 1 & 2 & 3 & 4 & 5 & 6 & 7 & 8 & 9 & 10 & 11 & 12 & 13 \\\midrule
Trenton TC & 5:49 & 6:15 & 6:34 & 7:00 & 7:27 & 7:53 & 8:19 & 8:45 & 9:04 & 9:30 & 9:57 & 10:23 & 10:49 \\
Hamilton Avenue & 5:51 & 6:17 & 6:36 & 7:02 & 7:29 & 7:55 & 8:21 & 8:47 & 9:06 & 9:32 & 9:59 & 10:25 & 10:51 \\
Cass Street & 5:53 & 6:19 & 6:38 & 7:04 & 7:31 & 7:57 & 8:23 & 8:49 & 9:08 & 9:34 & 10:01 & 10:27 & 10:53 \\
Bordentown & 6:00 & 6:26 & 6:45 & 7:11 & 7:38 & 8:04 & 8:30 & 8:56 & 9:15 & 9:41 & 10:08 & 10:34 & 11:00 \\
Roebling & 6:05 & 6:31 & 6:50 & 7:16 & 7:43 & 8:09 & 8:35 & 9:01 & 9:20 & 9:46 & 10:13 & 10:39 & 11:05 \\
Florence & 6:08 & 6:34 & 6:53 & 7:19 & 7:46 & 8:12 & 8:38 & 9:04 & 9:23 & 9:49 & 10:16 & 10:42 & 11:08 \\
Burlington Towne Ctr & 6:13 & 6:39 & 6:58 & 7:24 & 7:51 & 8:17 & 8:43 & 9:09 & 9:28 & 9:54 & 10:21 & 10:47 & 11:13 \\
Burlington South & 6:15 & 6:41 & 7:00 & 7:26 & 7:53 & 8:19 & 8:45 & 9:11 & 9:30 & 9:56 & 10:23 & 10:49 & 11:15 \\
Beverly/Edgewater Park & 6:23 & 6:45 & 7:04 & 7:30 & 7:57 & 8:23 & 8:53 & 9:15 & 9:34 & 10:00 & 10:27 & 10:53 & 11:23 \\
Delanco & 6:27 & 6:49 & 7:08 & 7:34 & 8:01 & 8:27 & 8:57 & 9:19 & 9:38 & 10:04 & 10:31 & 10:57 & 11:27 \\
Riverside & 6:29 & 6:51 & 7:10 & 7:36 & 8:03 & 8:29 & 8:59 & 9:21 & 9:40 & 10:06 & 10:33 & 10:59 & 11:29 \\
Cinnaminson & 6:34 & 6:56 & 7:15 & 7:41 & 8:08 & 8:38 & 9:04 & 9:26 & 9:45 & 10:11 & 10:38 & 11:08 & 11:34 \\
Riverton & 6:36 & 6:58 & 7:17 & 7:43 & 8:10 & 8:40 & 9:06 & 9:28 & 9:47 & 10:13 & 10:40 & 11:10 & 11:36 \\
Palmyra & 6:38 & 7:03 & 7:19 & 7:45 & 8:17 & 8:42 & 9:08 & 9:33 & 9:49 & 10:15 & 10:47 & 11:12 & 11:38 \\
Route 73/Pennsauken & 6:41 & 7:06 & 7:22 & 7:48 & 8:20 & 8:45 & 9:11 & 9:36 & 9:52 & 10:18 & 10:50 & 11:15 & 11:41 \\
Pennsauken TC & 6:43 & 7:08 & 7:24 & 7:50 & 8:22 & 8:47 & 9:13 & 9:38 & 9:54 & 10:20 & 10:52 & 11:17 & 11:43 \\
36th Street & 6:47 & 7:12 & 7:28 & 7:54 & 8:26 & 8:51 & 9:17 & 9:42 & 9:58 & 10:24 & 10:56 & 11:21 & 11:47 \\
Walter Rand TC & 6:52 & 7:17 & 7:33 & 7:59 & 8:31 & 8:56 & 9:22 & 9:47 & 10:03 & 10:29 & 11:01 & 11:26 & 11:52 \\
Cooper St/Rutgers & 6:55 & 7:20 & 7:36 & 8:02 & 8:34 & 8:59 & 9:25 & 9:50 & 10:06 & 10:32 & 11:04 & 11:29 & 11:55 \\
Aquarium & 6:57 & 7:22 & 7:38 & 8:04 & 8:36 & 9:01 & 9:27 & 9:52 & 10:08 & 10:34 & 11:06 & 11:31 & 11:57 \\
Entertainment Center & 7:00 & 7:25 & 7:41 & 8:07 & 8:39 & 9:04 & 9:30 & 9:55 & 10:11 & 10:37 & 11:09 & 11:34 & 12:00 \\
\bottomrule\end{tabular}\end{center}
\subsection*{To Camden (read down) --- trips 14--26 of 35}
\begin{center}\footnotesize\setlength{\tabcolsep}{3pt}\begin{tabular}{@{}lrrrrrrrrrrrrr@{}}\toprule
Station & 14 & 15 & 16 & 17 & 18 & 19 & 20 & 21 & 22 & 23 & 24 & 25 & 26 \\\midrule
Trenton TC & 11:34 & 12:00 & 12:27 & 12:53 & 13:19 & 14:04 & 14:30 & 14:57 & 15:23 & 15:49 & 16:15 & 16:34 & 17:00 \\
Hamilton Avenue & 11:36 & 12:02 & 12:29 & 12:55 & 13:21 & 14:06 & 14:32 & 14:59 & 15:25 & 15:51 & 16:17 & 16:36 & 17:02 \\
Cass Street & 11:38 & 12:04 & 12:31 & 12:57 & 13:23 & 14:08 & 14:34 & 15:01 & 15:27 & 15:53 & 16:19 & 16:38 & 17:04 \\
Bordentown & 11:45 & 12:11 & 12:38 & 13:04 & 13:30 & 14:15 & 14:41 & 15:08 & 15:34 & 16:00 & 16:26 & 16:45 & 17:11 \\
Roebling & 11:50 & 12:16 & 12:43 & 13:09 & 13:35 & 14:20 & 14:46 & 15:13 & 15:39 & 16:05 & 16:31 & 16:50 & 17:16 \\
Florence & 11:53 & 12:19 & 12:46 & 13:12 & 13:38 & 14:23 & 14:49 & 15:16 & 15:42 & 16:08 & 16:34 & 16:53 & 17:19 \\
Burlington Towne Ctr & 11:58 & 12:24 & 12:51 & 13:17 & 13:43 & 14:28 & 14:54 & 15:21 & 15:47 & 16:13 & 16:39 & 16:58 & 17:24 \\
Burlington South & 12:00 & 12:26 & 12:53 & 13:19 & 13:45 & 14:30 & 14:56 & 15:23 & 15:49 & 16:15 & 16:41 & 17:00 & 17:26 \\
Beverly/Edgewater Park & 12:04 & 12:30 & 12:57 & 13:23 & 13:53 & 14:34 & 15:00 & 15:27 & 15:53 & 16:23 & 16:45 & 17:04 & 17:30 \\
Delanco & 12:08 & 12:34 & 13:01 & 13:27 & 13:57 & 14:38 & 15:04 & 15:31 & 15:57 & 16:27 & 16:49 & 17:08 & 17:34 \\
Riverside & 12:10 & 12:36 & 13:03 & 13:29 & 13:59 & 14:40 & 15:06 & 15:33 & 15:59 & 16:29 & 16:51 & 17:10 & 17:36 \\
Cinnaminson & 12:15 & 12:41 & 13:08 & 13:38 & 14:04 & 14:45 & 15:11 & 15:38 & 16:08 & 16:34 & 16:56 & 17:15 & 17:41 \\
Riverton & 12:17 & 12:43 & 13:10 & 13:40 & 14:06 & 14:47 & 15:13 & 15:40 & 16:10 & 16:36 & 16:58 & 17:17 & 17:43 \\
Palmyra & 12:19 & 12:45 & 13:17 & 13:42 & 14:08 & 14:49 & 15:15 & 15:47 & 16:12 & 16:38 & 17:03 & 17:19 & 17:45 \\
Route 73/Pennsauken & 12:22 & 12:48 & 13:20 & 13:45 & 14:11 & 14:52 & 15:18 & 15:50 & 16:15 & 16:41 & 17:06 & 17:22 & 17:48 \\
Pennsauken TC & 12:24 & 12:50 & 13:22 & 13:47 & 14:13 & 14:54 & 15:20 & 15:52 & 16:17 & 16:43 & 17:08 & 17:24 & 17:50 \\
36th Street & 12:28 & 12:54 & 13:26 & 13:51 & 14:17 & 14:58 & 15:24 & 15:56 & 16:21 & 16:47 & 17:12 & 17:28 & 17:54 \\
Walter Rand TC & 12:33 & 12:59 & 13:31 & 13:56 & 14:22 & 15:03 & 15:29 & 16:01 & 16:26 & 16:52 & 17:17 & 17:33 & 17:59 \\
Cooper St/Rutgers & 12:36 & 13:02 & 13:34 & 13:59 & 14:25 & 15:06 & 15:32 & 16:04 & 16:29 & 16:55 & 17:20 & 17:36 & 18:02 \\
Aquarium & 12:38 & 13:04 & 13:36 & 14:01 & 14:27 & 15:08 & 15:34 & 16:06 & 16:31 & 16:57 & 17:22 & 17:38 & 18:04 \\
Entertainment Center & 12:41 & 13:07 & 13:39 & 14:04 & 14:30 & 15:11 & 15:37 & 16:09 & 16:34 & 17:00 & 17:25 & 17:41 & 18:07 \\
\bottomrule\end{tabular}\end{center}
\subsection*{To Camden (read down) --- trips 27--35 of 35}
\begin{center}\footnotesize\setlength{\tabcolsep}{3pt}\begin{tabular}{@{}lrrrrrrrrr@{}}\toprule
Station & 27 & 28 & 29 & 30 & 31 & 32 & 33 & 34 & 35 \\\midrule
Trenton TC & 17:27 & 17:53 & 18:19 & 18:45 & 19:04 & 19:30 & 19:57 & 20:23 & 20:49 \\
Hamilton Avenue & 17:29 & 17:55 & 18:21 & 18:47 & 19:06 & 19:32 & 19:59 & 20:25 & 20:51 \\
Cass Street & 17:31 & 17:57 & 18:23 & 18:49 & 19:08 & 19:34 & 20:01 & 20:27 & 20:53 \\
Bordentown & 17:38 & 18:04 & 18:30 & 18:56 & 19:15 & 19:41 & 20:08 & 20:34 & 21:00 \\
Roebling & 17:43 & 18:09 & 18:35 & 19:01 & 19:20 & 19:46 & 20:13 & 20:39 & 21:05 \\
Florence & 17:46 & 18:12 & 18:38 & 19:04 & 19:23 & 19:49 & 20:16 & 20:42 & 21:08 \\
Burlington Towne Ctr & 17:51 & 18:17 & 18:43 & 19:09 & 19:28 & 19:54 & 20:21 & 20:47 & 21:13 \\
Burlington South & 17:53 & 18:19 & 18:45 & 19:11 & 19:30 & 19:56 & 20:23 & 20:49 & 21:15 \\
Beverly/Edgewater Park & 17:57 & 18:23 & 18:53 & 19:15 & 19:34 & 20:00 & 20:27 & 20:53 & 21:23 \\
Delanco & 18:01 & 18:27 & 18:57 & 19:19 & 19:38 & 20:04 & 20:31 & 20:57 & 21:27 \\
Riverside & 18:03 & 18:29 & 18:59 & 19:21 & 19:40 & 20:06 & 20:33 & 20:59 & 21:29 \\
Cinnaminson & 18:08 & 18:38 & 19:04 & 19:26 & 19:45 & 20:11 & 20:38 & 21:08 & 21:34 \\
Riverton & 18:10 & 18:40 & 19:06 & 19:28 & 19:47 & 20:13 & 20:40 & 21:10 & 21:36 \\
Palmyra & 18:17 & 18:42 & 19:08 & 19:33 & 19:49 & 20:15 & 20:47 & 21:12 & 21:38 \\
Route 73/Pennsauken & 18:20 & 18:45 & 19:11 & 19:36 & 19:52 & 20:18 & 20:50 & 21:15 & 21:41 \\
Pennsauken TC & 18:22 & 18:47 & 19:13 & 19:38 & 19:54 & 20:20 & 20:52 & 21:17 & 21:43 \\
36th Street & 18:26 & 18:51 & 19:17 & 19:42 & 19:58 & 20:24 & 20:56 & 21:21 & 21:47 \\
Walter Rand TC & 18:31 & 18:56 & 19:22 & 19:47 & 20:03 & 20:29 & 21:01 & 21:26 & 21:52 \\
Cooper St/Rutgers & 18:34 & 18:59 & 19:25 & 19:50 & 20:06 & 20:32 & 21:04 & 21:29 & 21:55 \\
Aquarium & 18:36 & 19:01 & 19:27 & 19:52 & 20:08 & 20:34 & 21:06 & 21:31 & 21:57 \\
Entertainment Center & 18:39 & 19:04 & 19:30 & 19:55 & 20:11 & 20:37 & 21:09 & 21:34 & 22:00 \\
\bottomrule\end{tabular}\end{center}
\clearpage
\section{Proposed T5 (5 all day)}\label{tt:proposed_T5}
\subsection*{To Trenton (read down) --- trips 1--13 of 31}
\begin{center}\footnotesize\setlength{\tabcolsep}{3pt}\begin{tabular}{@{}lrrrrrrrrrrrrr@{}}\toprule
Station & 1 & 2 & 3 & 4 & 5 & 6 & 7 & 8 & 9 & 10 & 11 & 12 & 13 \\\midrule
Entertainment Center & 6:01 & 6:30 & 7:01 & 7:30 & 8:00 & 8:31 & 9:00 & 9:31 & 10:00 & 10:30 & 11:01 & 11:30 & 12:01 \\
Aquarium & 6:03 & 6:32 & 7:03 & 7:32 & 8:02 & 8:33 & 9:02 & 9:33 & 10:02 & 10:32 & 11:03 & 11:32 & 12:03 \\
Cooper St/Rutgers & 6:05 & 6:34 & 7:05 & 7:34 & 8:04 & 8:35 & 9:04 & 9:35 & 10:04 & 10:34 & 11:05 & 11:34 & 12:05 \\
Walter Rand TC & 6:09 & 6:38 & 7:09 & 7:38 & 8:08 & 8:39 & 9:08 & 9:39 & 10:08 & 10:38 & 11:09 & 11:38 & 12:09 \\
36th Street & 6:14 & 6:43 & 7:14 & 7:43 & 8:13 & 8:44 & 9:13 & 9:44 & 10:13 & 10:43 & 11:14 & 11:43 & 12:14 \\
Pennsauken TC & 6:16 & 6:45 & 7:16 & 7:45 & 8:15 & 8:46 & 9:15 & 9:46 & 10:15 & 10:45 & 11:16 & 11:45 & 12:16 \\
Route 73/Pennsauken & 6:19 & 6:48 & 7:19 & 7:48 & 8:18 & 8:49 & 9:18 & 9:49 & 10:18 & 10:48 & 11:19 & 11:48 & 12:19 \\
Palmyra & 6:22 & 6:51 & 7:22 & 7:51 & 8:21 & 8:52 & 9:21 & 9:52 & 10:21 & 10:51 & 11:22 & 11:51 & 12:22 \\
Riverton & 6:24 & 6:53 & 7:24 & 7:53 & 8:23 & 8:54 & 9:23 & 9:54 & 10:23 & 10:53 & 11:24 & 11:53 & 12:24 \\
Cinnaminson & 6:26 & 6:55 & 7:26 & 7:55 & 8:25 & 8:56 & 9:25 & 9:56 & 10:25 & 10:55 & 11:26 & 11:55 & 12:26 \\
Riverside & 6:31 & 7:00 & 7:31 & 8:00 & 8:30 & 9:01 & 9:30 & 10:01 & 10:30 & 11:00 & 11:31 & 12:00 & 12:31 \\
Delanco & 6:33 & 7:02 & 7:33 & 8:02 & 8:32 & 9:03 & 9:32 & 10:03 & 10:32 & 11:02 & 11:33 & 12:02 & 12:33 \\
Beverly/Edgewater Park & 6:37 & 7:06 & 7:37 & 8:06 & 8:36 & 9:07 & 9:36 & 10:07 & 10:36 & 11:06 & 11:37 & 12:06 & 12:37 \\
Burlington South & 6:41 & 7:10 & 7:41 & 8:10 & 8:40 & 9:11 & 9:40 & 10:11 & 10:40 & 11:10 & 11:41 & 12:10 & 12:41 \\
Burlington Towne Ctr & 6:43 & 7:12 & 7:43 & 8:12 & 8:42 & 9:13 & 9:42 & 10:13 & 10:42 & 11:12 & 11:43 & 12:12 & 12:43 \\
Florence & 6:49 & 7:18 & 7:49 & 8:18 & 8:48 & 9:19 & 9:48 & 10:19 & 10:48 & 11:18 & 11:49 & 12:18 & 12:49 \\
Roebling & 6:52 & 7:21 & 7:52 & 8:21 & 8:51 & 9:22 & 9:51 & 10:22 & 10:51 & 11:21 & 11:52 & 12:21 & 12:52 \\
Bordentown & 6:57 & 7:26 & 7:57 & 8:26 & 8:56 & 9:27 & 9:56 & 10:27 & 10:56 & 11:26 & 11:57 & 12:26 & 12:57 \\
Cass Street & 7:03 & 7:32 & 8:03 & 8:32 & 9:02 & 9:33 & 10:02 & 10:33 & 11:02 & 11:32 & 12:03 & 12:32 & 13:03 \\
Hamilton Avenue & 7:05 & 7:34 & 8:05 & 8:34 & 9:04 & 9:35 & 10:04 & 10:35 & 11:04 & 11:34 & 12:05 & 12:34 & 13:05 \\
Trenton TC & 7:08 & 7:37 & 8:08 & 8:37 & 9:07 & 9:38 & 10:07 & 10:38 & 11:07 & 11:37 & 12:08 & 12:37 & 13:08 \\
\bottomrule\end{tabular}\end{center}
\subsection*{To Trenton (read down) --- trips 14--26 of 31}
\begin{center}\footnotesize\setlength{\tabcolsep}{3pt}\begin{tabular}{@{}lrrrrrrrrrrrrr@{}}\toprule
Station & 14 & 15 & 16 & 17 & 18 & 19 & 20 & 21 & 22 & 23 & 24 & 25 & 26 \\\midrule
Entertainment Center & 12:30 & 13:00 & 13:31 & 14:00 & 14:31 & 15:00 & 15:30 & 16:01 & 16:30 & 17:01 & 17:30 & 18:00 & 18:31 \\
Aquarium & 12:32 & 13:02 & 13:33 & 14:02 & 14:33 & 15:02 & 15:32 & 16:03 & 16:32 & 17:03 & 17:32 & 18:02 & 18:33 \\
Cooper St/Rutgers & 12:34 & 13:04 & 13:35 & 14:04 & 14:35 & 15:04 & 15:34 & 16:05 & 16:34 & 17:05 & 17:34 & 18:04 & 18:35 \\
Walter Rand TC & 12:38 & 13:08 & 13:39 & 14:08 & 14:39 & 15:08 & 15:38 & 16:09 & 16:38 & 17:09 & 17:38 & 18:08 & 18:39 \\
36th Street & 12:43 & 13:13 & 13:44 & 14:13 & 14:44 & 15:13 & 15:43 & 16:14 & 16:43 & 17:14 & 17:43 & 18:13 & 18:44 \\
Pennsauken TC & 12:45 & 13:15 & 13:46 & 14:15 & 14:46 & 15:15 & 15:45 & 16:16 & 16:45 & 17:16 & 17:45 & 18:15 & 18:46 \\
Route 73/Pennsauken & 12:48 & 13:18 & 13:49 & 14:18 & 14:49 & 15:18 & 15:48 & 16:19 & 16:48 & 17:19 & 17:48 & 18:18 & 18:49 \\
Palmyra & 12:51 & 13:21 & 13:52 & 14:21 & 14:52 & 15:21 & 15:51 & 16:22 & 16:51 & 17:22 & 17:51 & 18:21 & 18:52 \\
Riverton & 12:53 & 13:23 & 13:54 & 14:23 & 14:54 & 15:23 & 15:53 & 16:24 & 16:53 & 17:24 & 17:53 & 18:23 & 18:54 \\
Cinnaminson & 12:55 & 13:25 & 13:56 & 14:25 & 14:56 & 15:25 & 15:55 & 16:26 & 16:55 & 17:26 & 17:55 & 18:25 & 18:56 \\
Riverside & 13:00 & 13:30 & 14:01 & 14:30 & 15:01 & 15:30 & 16:00 & 16:31 & 17:00 & 17:31 & 18:00 & 18:30 & 19:01 \\
Delanco & 13:02 & 13:32 & 14:03 & 14:32 & 15:03 & 15:32 & 16:02 & 16:33 & 17:02 & 17:33 & 18:02 & 18:32 & 19:03 \\
Beverly/Edgewater Park & 13:06 & 13:36 & 14:07 & 14:36 & 15:07 & 15:36 & 16:06 & 16:37 & 17:06 & 17:37 & 18:06 & 18:36 & 19:07 \\
Burlington South & 13:10 & 13:40 & 14:11 & 14:40 & 15:11 & 15:40 & 16:10 & 16:41 & 17:10 & 17:41 & 18:10 & 18:40 & 19:11 \\
Burlington Towne Ctr & 13:12 & 13:42 & 14:13 & 14:42 & 15:13 & 15:42 & 16:12 & 16:43 & 17:12 & 17:43 & 18:12 & 18:42 & 19:13 \\
Florence & 13:18 & 13:48 & 14:19 & 14:48 & 15:19 & 15:48 & 16:18 & 16:49 & 17:18 & 17:49 & 18:18 & 18:48 & 19:19 \\
Roebling & 13:21 & 13:51 & 14:22 & 14:51 & 15:22 & 15:51 & 16:21 & 16:52 & 17:21 & 17:52 & 18:21 & 18:51 & 19:22 \\
Bordentown & 13:26 & 13:56 & 14:27 & 14:56 & 15:27 & 15:56 & 16:26 & 16:57 & 17:26 & 17:57 & 18:26 & 18:56 & 19:27 \\
Cass Street & 13:32 & 14:02 & 14:33 & 15:02 & 15:33 & 16:02 & 16:32 & 17:03 & 17:32 & 18:03 & 18:32 & 19:02 & 19:33 \\
Hamilton Avenue & 13:34 & 14:04 & 14:35 & 15:04 & 15:35 & 16:04 & 16:34 & 17:05 & 17:34 & 18:05 & 18:34 & 19:04 & 19:35 \\
Trenton TC & 13:37 & 14:07 & 14:38 & 15:07 & 15:38 & 16:07 & 16:37 & 17:08 & 17:37 & 18:08 & 18:37 & 19:07 & 19:38 \\
\bottomrule\end{tabular}\end{center}
\subsection*{To Trenton (read down) --- trips 27--31 of 31}
\begin{center}\footnotesize\setlength{\tabcolsep}{3pt}\begin{tabular}{@{}lrrrrr@{}}\toprule
Station & 27 & 28 & 29 & 30 & 31 \\\midrule
Entertainment Center & 19:00 & 19:31 & 20:00 & 20:30 & 21:01 \\
Aquarium & 19:02 & 19:33 & 20:02 & 20:32 & 21:03 \\
Cooper St/Rutgers & 19:04 & 19:35 & 20:04 & 20:34 & 21:05 \\
Walter Rand TC & 19:08 & 19:39 & 20:08 & 20:38 & 21:09 \\
36th Street & 19:13 & 19:44 & 20:13 & 20:43 & 21:14 \\
Pennsauken TC & 19:15 & 19:46 & 20:15 & 20:45 & 21:16 \\
Route 73/Pennsauken & 19:18 & 19:49 & 20:18 & 20:48 & 21:19 \\
Palmyra & 19:21 & 19:52 & 20:21 & 20:51 & 21:22 \\
Riverton & 19:23 & 19:54 & 20:23 & 20:53 & 21:24 \\
Cinnaminson & 19:25 & 19:56 & 20:25 & 20:55 & 21:26 \\
Riverside & 19:30 & 20:01 & 20:30 & 21:00 & 21:31 \\
Delanco & 19:32 & 20:03 & 20:32 & 21:02 & 21:33 \\
Beverly/Edgewater Park & 19:36 & 20:07 & 20:36 & 21:06 & 21:37 \\
Burlington South & 19:40 & 20:11 & 20:40 & 21:10 & 21:41 \\
Burlington Towne Ctr & 19:42 & 20:13 & 20:42 & 21:12 & 21:43 \\
Florence & 19:48 & 20:19 & 20:48 & 21:18 & 21:49 \\
Roebling & 19:51 & 20:22 & 20:51 & 21:21 & 21:52 \\
Bordentown & 19:56 & 20:27 & 20:56 & 21:26 & 21:57 \\
Cass Street & 20:02 & 20:33 & 21:02 & 21:32 & 22:03 \\
Hamilton Avenue & 20:04 & 20:35 & 21:04 & 21:34 & 22:05 \\
Trenton TC & 20:07 & 20:38 & 21:07 & 21:37 & 22:08 \\
\bottomrule\end{tabular}\end{center}
\subsection*{To Camden (read down) --- trips 1--13 of 31}
\begin{center}\footnotesize\setlength{\tabcolsep}{3pt}\begin{tabular}{@{}lrrrrrrrrrrrrr@{}}\toprule
Station & 1 & 2 & 3 & 4 & 5 & 6 & 7 & 8 & 9 & 10 & 11 & 12 & 13 \\\midrule
Trenton TC & 5:43 & 6:13 & 6:42 & 7:13 & 7:44 & 8:13 & 8:43 & 9:12 & 9:43 & 10:14 & 10:43 & 11:13 & 11:42 \\
Hamilton Avenue & 5:45 & 6:15 & 6:44 & 7:15 & 7:46 & 8:15 & 8:45 & 9:14 & 9:45 & 10:16 & 10:45 & 11:15 & 11:44 \\
Cass Street & 5:47 & 6:17 & 6:46 & 7:17 & 7:48 & 8:17 & 8:47 & 9:16 & 9:47 & 10:18 & 10:47 & 11:17 & 11:46 \\
Bordentown & 5:54 & 6:24 & 6:53 & 7:24 & 7:55 & 8:24 & 8:54 & 9:23 & 9:54 & 10:25 & 10:54 & 11:24 & 11:53 \\
Roebling & 5:59 & 6:29 & 6:58 & 7:29 & 8:00 & 8:29 & 8:59 & 9:28 & 9:59 & 10:30 & 10:59 & 11:29 & 11:58 \\
Florence & 6:02 & 6:32 & 7:01 & 7:32 & 8:03 & 8:32 & 9:02 & 9:31 & 10:02 & 10:33 & 11:02 & 11:32 & 12:01 \\
Burlington Towne Ctr & 6:07 & 6:37 & 7:06 & 7:37 & 8:08 & 8:37 & 9:07 & 9:36 & 10:07 & 10:38 & 11:07 & 11:37 & 12:06 \\
Burlington South & 6:09 & 6:39 & 7:08 & 7:39 & 8:10 & 8:39 & 9:09 & 9:38 & 10:09 & 10:40 & 11:09 & 11:39 & 12:08 \\
Beverly/Edgewater Park & 6:13 & 6:43 & 7:12 & 7:43 & 8:14 & 8:43 & 9:13 & 9:42 & 10:13 & 10:44 & 11:13 & 11:43 & 12:12 \\
Delanco & 6:17 & 6:47 & 7:16 & 7:47 & 8:18 & 8:47 & 9:17 & 9:46 & 10:17 & 10:48 & 11:17 & 11:47 & 12:16 \\
Riverside & 6:19 & 6:49 & 7:18 & 7:49 & 8:20 & 8:49 & 9:19 & 9:48 & 10:19 & 10:50 & 11:19 & 11:49 & 12:18 \\
Cinnaminson & 6:24 & 6:54 & 7:23 & 7:54 & 8:25 & 8:54 & 9:24 & 9:53 & 10:24 & 10:55 & 11:24 & 11:54 & 12:23 \\
Riverton & 6:26 & 6:56 & 7:25 & 7:56 & 8:27 & 8:56 & 9:26 & 9:55 & 10:26 & 10:57 & 11:26 & 11:56 & 12:25 \\
Palmyra & 6:28 & 6:58 & 7:27 & 7:58 & 8:29 & 8:58 & 9:28 & 9:57 & 10:28 & 10:59 & 11:28 & 11:58 & 12:27 \\
Route 73/Pennsauken & 6:31 & 7:01 & 7:30 & 8:01 & 8:32 & 9:01 & 9:31 & 10:00 & 10:31 & 11:02 & 11:31 & 12:01 & 12:30 \\
Pennsauken TC & 6:33 & 7:03 & 7:32 & 8:03 & 8:34 & 9:03 & 9:33 & 10:02 & 10:33 & 11:04 & 11:33 & 12:03 & 12:32 \\
36th Street & 6:37 & 7:07 & 7:36 & 8:07 & 8:38 & 9:07 & 9:37 & 10:06 & 10:37 & 11:08 & 11:37 & 12:07 & 12:36 \\
Walter Rand TC & 6:42 & 7:12 & 7:41 & 8:12 & 8:43 & 9:12 & 9:42 & 10:11 & 10:42 & 11:13 & 11:42 & 12:12 & 12:41 \\
Cooper St/Rutgers & 6:45 & 7:15 & 7:44 & 8:15 & 8:46 & 9:15 & 9:45 & 10:14 & 10:45 & 11:16 & 11:45 & 12:15 & 12:44 \\
Aquarium & 6:47 & 7:17 & 7:46 & 8:17 & 8:48 & 9:17 & 9:47 & 10:16 & 10:47 & 11:18 & 11:47 & 12:17 & 12:46 \\
Entertainment Center & 6:50 & 7:20 & 7:49 & 8:20 & 8:51 & 9:20 & 9:50 & 10:19 & 10:50 & 11:21 & 11:50 & 12:20 & 12:49 \\
\bottomrule\end{tabular}\end{center}
\subsection*{To Camden (read down) --- trips 14--26 of 31}
\begin{center}\footnotesize\setlength{\tabcolsep}{3pt}\begin{tabular}{@{}lrrrrrrrrrrrrr@{}}\toprule
Station & 14 & 15 & 16 & 17 & 18 & 19 & 20 & 21 & 22 & 23 & 24 & 25 & 26 \\\midrule
Trenton TC & 12:13 & 12:44 & 13:13 & 13:43 & 14:12 & 14:43 & 15:14 & 15:43 & 16:13 & 16:42 & 17:13 & 17:44 & 18:13 \\
Hamilton Avenue & 12:15 & 12:46 & 13:15 & 13:45 & 14:14 & 14:45 & 15:16 & 15:45 & 16:15 & 16:44 & 17:15 & 17:46 & 18:15 \\
Cass Street & 12:17 & 12:48 & 13:17 & 13:47 & 14:16 & 14:47 & 15:18 & 15:47 & 16:17 & 16:46 & 17:17 & 17:48 & 18:17 \\
Bordentown & 12:24 & 12:55 & 13:24 & 13:54 & 14:23 & 14:54 & 15:25 & 15:54 & 16:24 & 16:53 & 17:24 & 17:55 & 18:24 \\
Roebling & 12:29 & 13:00 & 13:29 & 13:59 & 14:28 & 14:59 & 15:30 & 15:59 & 16:29 & 16:58 & 17:29 & 18:00 & 18:29 \\
Florence & 12:32 & 13:03 & 13:32 & 14:02 & 14:31 & 15:02 & 15:33 & 16:02 & 16:32 & 17:01 & 17:32 & 18:03 & 18:32 \\
Burlington Towne Ctr & 12:37 & 13:08 & 13:37 & 14:07 & 14:36 & 15:07 & 15:38 & 16:07 & 16:37 & 17:06 & 17:37 & 18:08 & 18:37 \\
Burlington South & 12:39 & 13:10 & 13:39 & 14:09 & 14:38 & 15:09 & 15:40 & 16:09 & 16:39 & 17:08 & 17:39 & 18:10 & 18:39 \\
Beverly/Edgewater Park & 12:43 & 13:14 & 13:43 & 14:13 & 14:42 & 15:13 & 15:44 & 16:13 & 16:43 & 17:12 & 17:43 & 18:14 & 18:43 \\
Delanco & 12:47 & 13:18 & 13:47 & 14:17 & 14:46 & 15:17 & 15:48 & 16:17 & 16:47 & 17:16 & 17:47 & 18:18 & 18:47 \\
Riverside & 12:49 & 13:20 & 13:49 & 14:19 & 14:48 & 15:19 & 15:50 & 16:19 & 16:49 & 17:18 & 17:49 & 18:20 & 18:49 \\
Cinnaminson & 12:54 & 13:25 & 13:54 & 14:24 & 14:53 & 15:24 & 15:55 & 16:24 & 16:54 & 17:23 & 17:54 & 18:25 & 18:54 \\
Riverton & 12:56 & 13:27 & 13:56 & 14:26 & 14:55 & 15:26 & 15:57 & 16:26 & 16:56 & 17:25 & 17:56 & 18:27 & 18:56 \\
Palmyra & 12:58 & 13:29 & 13:58 & 14:28 & 14:57 & 15:28 & 15:59 & 16:28 & 16:58 & 17:27 & 17:58 & 18:29 & 18:58 \\
Route 73/Pennsauken & 13:01 & 13:32 & 14:01 & 14:31 & 15:00 & 15:31 & 16:02 & 16:31 & 17:01 & 17:30 & 18:01 & 18:32 & 19:01 \\
Pennsauken TC & 13:03 & 13:34 & 14:03 & 14:33 & 15:02 & 15:33 & 16:04 & 16:33 & 17:03 & 17:32 & 18:03 & 18:34 & 19:03 \\
36th Street & 13:07 & 13:38 & 14:07 & 14:37 & 15:06 & 15:37 & 16:08 & 16:37 & 17:07 & 17:36 & 18:07 & 18:38 & 19:07 \\
Walter Rand TC & 13:12 & 13:43 & 14:12 & 14:42 & 15:11 & 15:42 & 16:13 & 16:42 & 17:12 & 17:41 & 18:12 & 18:43 & 19:12 \\
Cooper St/Rutgers & 13:15 & 13:46 & 14:15 & 14:45 & 15:14 & 15:45 & 16:16 & 16:45 & 17:15 & 17:44 & 18:15 & 18:46 & 19:15 \\
Aquarium & 13:17 & 13:48 & 14:17 & 14:47 & 15:16 & 15:47 & 16:18 & 16:47 & 17:17 & 17:46 & 18:17 & 18:48 & 19:17 \\
Entertainment Center & 13:20 & 13:51 & 14:20 & 14:50 & 15:19 & 15:50 & 16:21 & 16:50 & 17:20 & 17:49 & 18:20 & 18:51 & 19:20 \\
\bottomrule\end{tabular}\end{center}
\subsection*{To Camden (read down) --- trips 27--31 of 31}
\begin{center}\footnotesize\setlength{\tabcolsep}{3pt}\begin{tabular}{@{}lrrrrr@{}}\toprule
Station & 27 & 28 & 29 & 30 & 31 \\\midrule
Trenton TC & 18:43 & 19:12 & 19:43 & 20:14 & 20:43 \\
Hamilton Avenue & 18:45 & 19:14 & 19:45 & 20:16 & 20:45 \\
Cass Street & 18:47 & 19:16 & 19:47 & 20:18 & 20:47 \\
Bordentown & 18:54 & 19:23 & 19:54 & 20:25 & 20:54 \\
Roebling & 18:59 & 19:28 & 19:59 & 20:30 & 20:59 \\
Florence & 19:02 & 19:31 & 20:02 & 20:33 & 21:02 \\
Burlington Towne Ctr & 19:07 & 19:36 & 20:07 & 20:38 & 21:07 \\
Burlington South & 19:09 & 19:38 & 20:09 & 20:40 & 21:09 \\
Beverly/Edgewater Park & 19:13 & 19:42 & 20:13 & 20:44 & 21:13 \\
Delanco & 19:17 & 19:46 & 20:17 & 20:48 & 21:17 \\
Riverside & 19:19 & 19:48 & 20:19 & 20:50 & 21:19 \\
Cinnaminson & 19:24 & 19:53 & 20:24 & 20:55 & 21:24 \\
Riverton & 19:26 & 19:55 & 20:26 & 20:57 & 21:26 \\
Palmyra & 19:28 & 19:57 & 20:28 & 20:59 & 21:28 \\
Route 73/Pennsauken & 19:31 & 20:00 & 20:31 & 21:02 & 21:31 \\
Pennsauken TC & 19:33 & 20:02 & 20:33 & 21:04 & 21:33 \\
36th Street & 19:37 & 20:06 & 20:37 & 21:08 & 21:37 \\
Walter Rand TC & 19:42 & 20:11 & 20:42 & 21:13 & 21:42 \\
Cooper St/Rutgers & 19:45 & 20:14 & 20:45 & 21:16 & 21:45 \\
Aquarium & 19:47 & 20:16 & 20:47 & 21:18 & 21:47 \\
Entertainment Center & 19:50 & 20:19 & 20:50 & 21:21 & 21:50 \\
\bottomrule\end{tabular}\end{center}
\clearpage
\section{Proposed T4 (4 all day)}\label{tt:proposed_T4}
\subsection*{To Trenton (read down) --- trips 1--13 of 24}
\begin{center}\footnotesize\setlength{\tabcolsep}{3pt}\begin{tabular}{@{}lrrrrrrrrrrrrr@{}}\toprule
Station & 1 & 2 & 3 & 4 & 5 & 6 & 7 & 8 & 9 & 10 & 11 & 12 & 13 \\\midrule
Entertainment Center & 6:15 & 6:52 & 7:30 & 8:07 & 8:45 & 9:22 & 10:00 & 10:37 & 11:15 & 11:52 & 12:30 & 13:07 & 13:45 \\
Aquarium & 6:17 & 6:54 & 7:32 & 8:09 & 8:47 & 9:24 & 10:02 & 10:39 & 11:17 & 11:54 & 12:32 & 13:09 & 13:47 \\
Cooper St/Rutgers & 6:19 & 6:56 & 7:34 & 8:11 & 8:49 & 9:26 & 10:04 & 10:41 & 11:19 & 11:56 & 12:34 & 13:11 & 13:49 \\
Walter Rand TC & 6:23 & 7:00 & 7:38 & 8:15 & 8:53 & 9:30 & 10:08 & 10:45 & 11:23 & 12:00 & 12:38 & 13:15 & 13:53 \\
36th Street & 6:28 & 7:05 & 7:43 & 8:20 & 8:58 & 9:35 & 10:13 & 10:50 & 11:28 & 12:05 & 12:43 & 13:20 & 13:58 \\
Pennsauken TC & 6:30 & 7:07 & 7:45 & 8:22 & 9:00 & 9:37 & 10:15 & 10:52 & 11:30 & 12:07 & 12:45 & 13:22 & 14:00 \\
Route 73/Pennsauken & 6:33 & 7:10 & 7:48 & 8:25 & 9:03 & 9:40 & 10:18 & 10:55 & 11:33 & 12:10 & 12:48 & 13:25 & 14:03 \\
Palmyra & 6:36 & 7:13 & 7:51 & 8:28 & 9:06 & 9:43 & 10:21 & 10:58 & 11:36 & 12:13 & 12:51 & 13:28 & 14:06 \\
Riverton & 6:38 & 7:15 & 7:53 & 8:30 & 9:08 & 9:45 & 10:23 & 11:00 & 11:38 & 12:15 & 12:53 & 13:30 & 14:08 \\
Cinnaminson & 6:40 & 7:17 & 7:55 & 8:32 & 9:10 & 9:47 & 10:25 & 11:02 & 11:40 & 12:17 & 12:55 & 13:32 & 14:10 \\
Riverside & 6:45 & 7:22 & 8:00 & 8:37 & 9:15 & 9:52 & 10:30 & 11:07 & 11:45 & 12:22 & 13:00 & 13:37 & 14:15 \\
Delanco & 6:47 & 7:24 & 8:02 & 8:39 & 9:17 & 9:54 & 10:32 & 11:09 & 11:47 & 12:24 & 13:02 & 13:39 & 14:17 \\
Beverly/Edgewater Park & 6:51 & 7:28 & 8:06 & 8:43 & 9:21 & 9:58 & 10:36 & 11:13 & 11:51 & 12:28 & 13:06 & 13:43 & 14:21 \\
Burlington South & 6:55 & 7:32 & 8:10 & 8:47 & 9:25 & 10:02 & 10:40 & 11:17 & 11:55 & 12:32 & 13:10 & 13:47 & 14:25 \\
Burlington Towne Ctr & 6:57 & 7:34 & 8:12 & 8:49 & 9:27 & 10:04 & 10:42 & 11:19 & 11:57 & 12:34 & 13:12 & 13:49 & 14:27 \\
Florence & 7:03 & 7:40 & 8:18 & 8:55 & 9:33 & 10:10 & 10:48 & 11:25 & 12:03 & 12:40 & 13:18 & 13:55 & 14:33 \\
Roebling & 7:06 & 7:43 & 8:21 & 8:58 & 9:36 & 10:13 & 10:51 & 11:28 & 12:06 & 12:43 & 13:21 & 13:58 & 14:36 \\
Bordentown & 7:11 & 7:48 & 8:26 & 9:03 & 9:41 & 10:18 & 10:56 & 11:33 & 12:11 & 12:48 & 13:26 & 14:03 & 14:41 \\
Cass Street & 7:17 & 7:54 & 8:32 & 9:09 & 9:47 & 10:24 & 11:02 & 11:39 & 12:17 & 12:54 & 13:32 & 14:09 & 14:47 \\
Hamilton Avenue & 7:19 & 7:56 & 8:34 & 9:11 & 9:49 & 10:26 & 11:04 & 11:41 & 12:19 & 12:56 & 13:34 & 14:11 & 14:49 \\
Trenton TC & 7:22 & 7:59 & 8:37 & 9:14 & 9:52 & 10:29 & 11:07 & 11:44 & 12:22 & 12:59 & 13:37 & 14:14 & 14:52 \\
\bottomrule\end{tabular}\end{center}
\subsection*{To Trenton (read down) --- trips 14--24 of 24}
\begin{center}\footnotesize\setlength{\tabcolsep}{3pt}\begin{tabular}{@{}lrrrrrrrrrrr@{}}\toprule
Station & 14 & 15 & 16 & 17 & 18 & 19 & 20 & 21 & 22 & 23 & 24 \\\midrule
Entertainment Center & 14:22 & 15:00 & 15:37 & 16:15 & 16:52 & 17:30 & 18:07 & 18:45 & 19:22 & 20:00 & 20:37 \\
Aquarium & 14:24 & 15:02 & 15:39 & 16:17 & 16:54 & 17:32 & 18:09 & 18:47 & 19:24 & 20:02 & 20:39 \\
Cooper St/Rutgers & 14:26 & 15:04 & 15:41 & 16:19 & 16:56 & 17:34 & 18:11 & 18:49 & 19:26 & 20:04 & 20:41 \\
Walter Rand TC & 14:30 & 15:08 & 15:45 & 16:23 & 17:00 & 17:38 & 18:15 & 18:53 & 19:30 & 20:08 & 20:45 \\
36th Street & 14:35 & 15:13 & 15:50 & 16:28 & 17:05 & 17:43 & 18:20 & 18:58 & 19:35 & 20:13 & 20:50 \\
Pennsauken TC & 14:37 & 15:15 & 15:52 & 16:30 & 17:07 & 17:45 & 18:22 & 19:00 & 19:37 & 20:15 & 20:52 \\
Route 73/Pennsauken & 14:40 & 15:18 & 15:55 & 16:33 & 17:10 & 17:48 & 18:25 & 19:03 & 19:40 & 20:18 & 20:55 \\
Palmyra & 14:43 & 15:21 & 15:58 & 16:36 & 17:13 & 17:51 & 18:28 & 19:06 & 19:43 & 20:21 & 20:58 \\
Riverton & 14:45 & 15:23 & 16:00 & 16:38 & 17:15 & 17:53 & 18:30 & 19:08 & 19:45 & 20:23 & 21:00 \\
Cinnaminson & 14:47 & 15:25 & 16:02 & 16:40 & 17:17 & 17:55 & 18:32 & 19:10 & 19:47 & 20:25 & 21:02 \\
Riverside & 14:52 & 15:30 & 16:07 & 16:45 & 17:22 & 18:00 & 18:37 & 19:15 & 19:52 & 20:30 & 21:07 \\
Delanco & 14:54 & 15:32 & 16:09 & 16:47 & 17:24 & 18:02 & 18:39 & 19:17 & 19:54 & 20:32 & 21:09 \\
Beverly/Edgewater Park & 14:58 & 15:36 & 16:13 & 16:51 & 17:28 & 18:06 & 18:43 & 19:21 & 19:58 & 20:36 & 21:13 \\
Burlington South & 15:02 & 15:40 & 16:17 & 16:55 & 17:32 & 18:10 & 18:47 & 19:25 & 20:02 & 20:40 & 21:17 \\
Burlington Towne Ctr & 15:04 & 15:42 & 16:19 & 16:57 & 17:34 & 18:12 & 18:49 & 19:27 & 20:04 & 20:42 & 21:19 \\
Florence & 15:10 & 15:48 & 16:25 & 17:03 & 17:40 & 18:18 & 18:55 & 19:33 & 20:10 & 20:48 & 21:25 \\
Roebling & 15:13 & 15:51 & 16:28 & 17:06 & 17:43 & 18:21 & 18:58 & 19:36 & 20:13 & 20:51 & 21:28 \\
Bordentown & 15:18 & 15:56 & 16:33 & 17:11 & 17:48 & 18:26 & 19:03 & 19:41 & 20:18 & 20:56 & 21:33 \\
Cass Street & 15:24 & 16:02 & 16:39 & 17:17 & 17:54 & 18:32 & 19:09 & 19:47 & 20:24 & 21:02 & 21:39 \\
Hamilton Avenue & 15:26 & 16:04 & 16:41 & 17:19 & 17:56 & 18:34 & 19:11 & 19:49 & 20:26 & 21:04 & 21:41 \\
Trenton TC & 15:29 & 16:07 & 16:44 & 17:22 & 17:59 & 18:37 & 19:14 & 19:52 & 20:29 & 21:07 & 21:44 \\
\bottomrule\end{tabular}\end{center}
\subsection*{To Camden (read down) --- trips 1--13 of 24}
\begin{center}\footnotesize\setlength{\tabcolsep}{3pt}\begin{tabular}{@{}lrrrrrrrrrrrrr@{}}\toprule
Station & 1 & 2 & 3 & 4 & 5 & 6 & 7 & 8 & 9 & 10 & 11 & 12 & 13 \\\midrule
Trenton TC & 6:17 & 6:55 & 7:32 & 8:10 & 8:47 & 9:25 & 10:02 & 10:40 & 11:17 & 11:55 & 12:32 & 13:10 & 13:47 \\
Hamilton Avenue & 6:19 & 6:57 & 7:34 & 8:12 & 8:49 & 9:27 & 10:04 & 10:42 & 11:19 & 11:57 & 12:34 & 13:12 & 13:49 \\
Cass Street & 6:21 & 6:59 & 7:36 & 8:14 & 8:51 & 9:29 & 10:06 & 10:44 & 11:21 & 11:59 & 12:36 & 13:14 & 13:51 \\
Bordentown & 6:28 & 7:06 & 7:43 & 8:21 & 8:58 & 9:36 & 10:13 & 10:51 & 11:28 & 12:06 & 12:43 & 13:21 & 13:58 \\
Roebling & 6:33 & 7:11 & 7:48 & 8:26 & 9:03 & 9:41 & 10:18 & 10:56 & 11:33 & 12:11 & 12:48 & 13:26 & 14:03 \\
Florence & 6:36 & 7:14 & 7:51 & 8:29 & 9:06 & 9:44 & 10:21 & 10:59 & 11:36 & 12:14 & 12:51 & 13:29 & 14:06 \\
Burlington Towne Ctr & 6:41 & 7:19 & 7:56 & 8:34 & 9:11 & 9:49 & 10:26 & 11:04 & 11:41 & 12:19 & 12:56 & 13:34 & 14:11 \\
Burlington South & 6:43 & 7:21 & 7:58 & 8:36 & 9:13 & 9:51 & 10:28 & 11:06 & 11:43 & 12:21 & 12:58 & 13:36 & 14:13 \\
Beverly/Edgewater Park & 6:47 & 7:25 & 8:02 & 8:40 & 9:17 & 9:55 & 10:32 & 11:10 & 11:47 & 12:25 & 13:02 & 13:40 & 14:17 \\
Delanco & 6:51 & 7:29 & 8:06 & 8:44 & 9:21 & 9:59 & 10:36 & 11:14 & 11:51 & 12:29 & 13:06 & 13:44 & 14:21 \\
Riverside & 6:53 & 7:31 & 8:08 & 8:46 & 9:23 & 10:01 & 10:38 & 11:16 & 11:53 & 12:31 & 13:08 & 13:46 & 14:23 \\
Cinnaminson & 6:58 & 7:36 & 8:13 & 8:51 & 9:28 & 10:06 & 10:43 & 11:21 & 11:58 & 12:36 & 13:13 & 13:51 & 14:28 \\
Riverton & 7:00 & 7:38 & 8:15 & 8:53 & 9:30 & 10:08 & 10:45 & 11:23 & 12:00 & 12:38 & 13:15 & 13:53 & 14:30 \\
Palmyra & 7:02 & 7:40 & 8:17 & 8:55 & 9:32 & 10:10 & 10:47 & 11:25 & 12:02 & 12:40 & 13:17 & 13:55 & 14:32 \\
Route 73/Pennsauken & 7:05 & 7:43 & 8:20 & 8:58 & 9:35 & 10:13 & 10:50 & 11:28 & 12:05 & 12:43 & 13:20 & 13:58 & 14:35 \\
Pennsauken TC & 7:07 & 7:45 & 8:22 & 9:00 & 9:37 & 10:15 & 10:52 & 11:30 & 12:07 & 12:45 & 13:22 & 14:00 & 14:37 \\
36th Street & 7:11 & 7:49 & 8:26 & 9:04 & 9:41 & 10:19 & 10:56 & 11:34 & 12:11 & 12:49 & 13:26 & 14:04 & 14:41 \\
Walter Rand TC & 7:16 & 7:54 & 8:31 & 9:09 & 9:46 & 10:24 & 11:01 & 11:39 & 12:16 & 12:54 & 13:31 & 14:09 & 14:46 \\
Cooper St/Rutgers & 7:19 & 7:57 & 8:34 & 9:12 & 9:49 & 10:27 & 11:04 & 11:42 & 12:19 & 12:57 & 13:34 & 14:12 & 14:49 \\
Aquarium & 7:21 & 7:59 & 8:36 & 9:14 & 9:51 & 10:29 & 11:06 & 11:44 & 12:21 & 12:59 & 13:36 & 14:14 & 14:51 \\
Entertainment Center & 7:24 & 8:02 & 8:39 & 9:17 & 9:54 & 10:32 & 11:09 & 11:47 & 12:24 & 13:02 & 13:39 & 14:17 & 14:54 \\
\bottomrule\end{tabular}\end{center}
\subsection*{To Camden (read down) --- trips 14--24 of 24}
\begin{center}\footnotesize\setlength{\tabcolsep}{3pt}\begin{tabular}{@{}lrrrrrrrrrrr@{}}\toprule
Station & 14 & 15 & 16 & 17 & 18 & 19 & 20 & 21 & 22 & 23 & 24 \\\midrule
Trenton TC & 14:25 & 15:02 & 15:40 & 16:17 & 16:55 & 17:32 & 18:10 & 18:47 & 19:25 & 20:02 & 20:40 \\
Hamilton Avenue & 14:27 & 15:04 & 15:42 & 16:19 & 16:57 & 17:34 & 18:12 & 18:49 & 19:27 & 20:04 & 20:42 \\
Cass Street & 14:29 & 15:06 & 15:44 & 16:21 & 16:59 & 17:36 & 18:14 & 18:51 & 19:29 & 20:06 & 20:44 \\
Bordentown & 14:36 & 15:13 & 15:51 & 16:28 & 17:06 & 17:43 & 18:21 & 18:58 & 19:36 & 20:13 & 20:51 \\
Roebling & 14:41 & 15:18 & 15:56 & 16:33 & 17:11 & 17:48 & 18:26 & 19:03 & 19:41 & 20:18 & 20:56 \\
Florence & 14:44 & 15:21 & 15:59 & 16:36 & 17:14 & 17:51 & 18:29 & 19:06 & 19:44 & 20:21 & 20:59 \\
Burlington Towne Ctr & 14:49 & 15:26 & 16:04 & 16:41 & 17:19 & 17:56 & 18:34 & 19:11 & 19:49 & 20:26 & 21:04 \\
Burlington South & 14:51 & 15:28 & 16:06 & 16:43 & 17:21 & 17:58 & 18:36 & 19:13 & 19:51 & 20:28 & 21:06 \\
Beverly/Edgewater Park & 14:55 & 15:32 & 16:10 & 16:47 & 17:25 & 18:02 & 18:40 & 19:17 & 19:55 & 20:32 & 21:10 \\
Delanco & 14:59 & 15:36 & 16:14 & 16:51 & 17:29 & 18:06 & 18:44 & 19:21 & 19:59 & 20:36 & 21:14 \\
Riverside & 15:01 & 15:38 & 16:16 & 16:53 & 17:31 & 18:08 & 18:46 & 19:23 & 20:01 & 20:38 & 21:16 \\
Cinnaminson & 15:06 & 15:43 & 16:21 & 16:58 & 17:36 & 18:13 & 18:51 & 19:28 & 20:06 & 20:43 & 21:21 \\
Riverton & 15:08 & 15:45 & 16:23 & 17:00 & 17:38 & 18:15 & 18:53 & 19:30 & 20:08 & 20:45 & 21:23 \\
Palmyra & 15:10 & 15:47 & 16:25 & 17:02 & 17:40 & 18:17 & 18:55 & 19:32 & 20:10 & 20:47 & 21:25 \\
Route 73/Pennsauken & 15:13 & 15:50 & 16:28 & 17:05 & 17:43 & 18:20 & 18:58 & 19:35 & 20:13 & 20:50 & 21:28 \\
Pennsauken TC & 15:15 & 15:52 & 16:30 & 17:07 & 17:45 & 18:22 & 19:00 & 19:37 & 20:15 & 20:52 & 21:30 \\
36th Street & 15:19 & 15:56 & 16:34 & 17:11 & 17:49 & 18:26 & 19:04 & 19:41 & 20:19 & 20:56 & 21:34 \\
Walter Rand TC & 15:24 & 16:01 & 16:39 & 17:16 & 17:54 & 18:31 & 19:09 & 19:46 & 20:24 & 21:01 & 21:39 \\
Cooper St/Rutgers & 15:27 & 16:04 & 16:42 & 17:19 & 17:57 & 18:34 & 19:12 & 19:49 & 20:27 & 21:04 & 21:42 \\
Aquarium & 15:29 & 16:06 & 16:44 & 17:21 & 17:59 & 18:36 & 19:14 & 19:51 & 20:29 & 21:06 & 21:44 \\
Entertainment Center & 15:32 & 16:09 & 16:47 & 17:24 & 18:02 & 18:39 & 19:17 & 19:54 & 20:32 & 21:09 & 21:47 \\
\bottomrule\end{tabular}\end{center}
\clearpage
\section{Proposed T3 (3 all day)}\label{tt:proposed_T3}
\subsection*{To Trenton (read down) --- trips 1--13 of 19}
\begin{center}\footnotesize\setlength{\tabcolsep}{3pt}\begin{tabular}{@{}lrrrrrrrrrrrrr@{}}\toprule
Station & 1 & 2 & 3 & 4 & 5 & 6 & 7 & 8 & 9 & 10 & 11 & 12 & 13 \\\midrule
Entertainment Center & 5:51 & 6:39 & 7:30 & 8:21 & 9:09 & 10:00 & 10:51 & 11:39 & 12:30 & 13:21 & 14:09 & 15:00 & 15:51 \\
Aquarium & 5:53 & 6:41 & 7:32 & 8:23 & 9:11 & 10:02 & 10:53 & 11:41 & 12:32 & 13:23 & 14:11 & 15:02 & 15:53 \\
Cooper St/Rutgers & 5:55 & 6:43 & 7:34 & 8:25 & 9:13 & 10:04 & 10:55 & 11:43 & 12:34 & 13:25 & 14:13 & 15:04 & 15:55 \\
Walter Rand TC & 5:59 & 6:47 & 7:38 & 8:29 & 9:17 & 10:08 & 10:59 & 11:47 & 12:38 & 13:29 & 14:17 & 15:08 & 15:59 \\
36th Street & 6:07 & 6:52 & 7:43 & 8:37 & 9:22 & 10:13 & 11:07 & 11:52 & 12:43 & 13:37 & 14:22 & 15:13 & 16:07 \\
Pennsauken TC & 6:09 & 6:54 & 7:45 & 8:39 & 9:24 & 10:15 & 11:09 & 11:54 & 12:45 & 13:39 & 14:24 & 15:15 & 16:09 \\
Route 73/Pennsauken & 6:12 & 6:57 & 7:48 & 8:42 & 9:27 & 10:18 & 11:12 & 11:57 & 12:48 & 13:42 & 14:27 & 15:18 & 16:12 \\
Palmyra & 6:15 & 7:00 & 7:51 & 8:45 & 9:30 & 10:21 & 11:15 & 12:00 & 12:51 & 13:45 & 14:30 & 15:21 & 16:15 \\
Riverton & 6:17 & 7:02 & 7:53 & 8:47 & 9:32 & 10:23 & 11:17 & 12:02 & 12:53 & 13:47 & 14:32 & 15:23 & 16:17 \\
Cinnaminson & 6:19 & 7:04 & 7:55 & 8:49 & 9:34 & 10:25 & 11:19 & 12:04 & 12:55 & 13:49 & 14:34 & 15:25 & 16:19 \\
Riverside & 6:24 & 7:09 & 8:00 & 8:54 & 9:39 & 10:30 & 11:24 & 12:09 & 13:00 & 13:54 & 14:39 & 15:30 & 16:24 \\
Delanco & 6:26 & 7:11 & 8:02 & 8:56 & 9:41 & 10:32 & 11:26 & 12:11 & 13:02 & 13:56 & 14:41 & 15:32 & 16:26 \\
Beverly/Edgewater Park & 6:30 & 7:15 & 8:06 & 9:00 & 9:45 & 10:36 & 11:30 & 12:15 & 13:06 & 14:00 & 14:45 & 15:36 & 16:30 \\
Burlington South & 6:34 & 7:19 & 8:10 & 9:04 & 9:49 & 10:40 & 11:34 & 12:19 & 13:10 & 14:04 & 14:49 & 15:40 & 16:34 \\
Burlington Towne Ctr & 6:36 & 7:21 & 8:12 & 9:06 & 9:51 & 10:42 & 11:36 & 12:21 & 13:12 & 14:06 & 14:51 & 15:42 & 16:36 \\
Florence & 6:42 & 7:27 & 8:19 & 9:12 & 9:57 & 10:49 & 11:42 & 12:27 & 13:19 & 14:12 & 14:57 & 15:49 & 16:42 \\
Roebling & 6:45 & 7:30 & 8:22 & 9:15 & 10:00 & 10:52 & 11:45 & 12:30 & 13:22 & 14:15 & 15:00 & 15:52 & 16:45 \\
Bordentown & 6:50 & 7:35 & 8:27 & 9:20 & 10:05 & 10:57 & 11:50 & 12:35 & 13:27 & 14:20 & 15:05 & 15:57 & 16:50 \\
Cass Street & 6:56 & 7:41 & 8:33 & 9:26 & 10:11 & 11:03 & 11:56 & 12:41 & 13:33 & 14:26 & 15:11 & 16:03 & 16:56 \\
Hamilton Avenue & 6:58 & 7:43 & 8:35 & 9:28 & 10:13 & 11:05 & 11:58 & 12:43 & 13:35 & 14:28 & 15:13 & 16:05 & 16:58 \\
Trenton TC & 7:01 & 7:46 & 8:38 & 9:31 & 10:16 & 11:08 & 12:01 & 12:46 & 13:38 & 14:31 & 15:16 & 16:08 & 17:01 \\
\bottomrule\end{tabular}\end{center}
\subsection*{To Trenton (read down) --- trips 14--19 of 19}
\begin{center}\footnotesize\setlength{\tabcolsep}{3pt}\begin{tabular}{@{}lrrrrrr@{}}\toprule
Station & 14 & 15 & 16 & 17 & 18 & 19 \\\midrule
Entertainment Center & 16:39 & 17:30 & 18:21 & 19:09 & 20:00 & 20:51 \\
Aquarium & 16:41 & 17:32 & 18:23 & 19:11 & 20:02 & 20:53 \\
Cooper St/Rutgers & 16:43 & 17:34 & 18:25 & 19:13 & 20:04 & 20:55 \\
Walter Rand TC & 16:47 & 17:38 & 18:29 & 19:17 & 20:08 & 20:59 \\
36th Street & 16:52 & 17:43 & 18:37 & 19:22 & 20:13 & 21:07 \\
Pennsauken TC & 16:54 & 17:45 & 18:39 & 19:24 & 20:15 & 21:09 \\
Route 73/Pennsauken & 16:57 & 17:48 & 18:42 & 19:27 & 20:18 & 21:12 \\
Palmyra & 17:00 & 17:51 & 18:45 & 19:30 & 20:21 & 21:15 \\
Riverton & 17:02 & 17:53 & 18:47 & 19:32 & 20:23 & 21:17 \\
Cinnaminson & 17:04 & 17:55 & 18:49 & 19:34 & 20:25 & 21:19 \\
Riverside & 17:09 & 18:00 & 18:54 & 19:39 & 20:30 & 21:24 \\
Delanco & 17:11 & 18:02 & 18:56 & 19:41 & 20:32 & 21:26 \\
Beverly/Edgewater Park & 17:15 & 18:06 & 19:00 & 19:45 & 20:36 & 21:30 \\
Burlington South & 17:19 & 18:10 & 19:04 & 19:49 & 20:40 & 21:34 \\
Burlington Towne Ctr & 17:21 & 18:12 & 19:06 & 19:51 & 20:42 & 21:36 \\
Florence & 17:27 & 18:19 & 19:12 & 19:57 & 20:49 & 21:42 \\
Roebling & 17:30 & 18:22 & 19:15 & 20:00 & 20:52 & 21:45 \\
Bordentown & 17:35 & 18:27 & 19:20 & 20:05 & 20:57 & 21:50 \\
Cass Street & 17:41 & 18:33 & 19:26 & 20:11 & 21:03 & 21:56 \\
Hamilton Avenue & 17:43 & 18:35 & 19:28 & 20:13 & 21:05 & 21:58 \\
Trenton TC & 17:46 & 18:38 & 19:31 & 20:16 & 21:08 & 22:01 \\
\bottomrule\end{tabular}\end{center}
\subsection*{To Camden (read down) --- trips 1--13 of 18}
\begin{center}\footnotesize\setlength{\tabcolsep}{3pt}\begin{tabular}{@{}lrrrrrrrrrrrrr@{}}\toprule
Station & 1 & 2 & 3 & 4 & 5 & 6 & 7 & 8 & 9 & 10 & 11 & 12 & 13 \\\midrule
Trenton TC & 6:13 & 7:06 & 7:51 & 8:43 & 9:36 & 10:21 & 11:13 & 12:06 & 12:51 & 13:43 & 14:36 & 15:21 & 16:13 \\
Hamilton Avenue & 6:15 & 7:08 & 7:53 & 8:45 & 9:38 & 10:23 & 11:15 & 12:08 & 12:53 & 13:45 & 14:38 & 15:23 & 16:15 \\
Cass Street & 6:17 & 7:10 & 7:55 & 8:47 & 9:40 & 10:25 & 11:17 & 12:10 & 12:55 & 13:47 & 14:40 & 15:25 & 16:17 \\
Bordentown & 6:24 & 7:17 & 8:02 & 8:54 & 9:47 & 10:32 & 11:24 & 12:17 & 13:02 & 13:54 & 14:47 & 15:32 & 16:24 \\
Roebling & 6:29 & 7:22 & 8:07 & 8:59 & 9:52 & 10:37 & 11:29 & 12:22 & 13:07 & 13:59 & 14:52 & 15:37 & 16:29 \\
Florence & 6:32 & 7:25 & 8:10 & 9:02 & 9:55 & 10:40 & 11:32 & 12:25 & 13:10 & 14:02 & 14:55 & 15:40 & 16:32 \\
Burlington Towne Ctr & 6:37 & 7:30 & 8:15 & 9:07 & 10:00 & 10:45 & 11:37 & 12:30 & 13:15 & 14:07 & 15:00 & 15:45 & 16:37 \\
Burlington South & 6:39 & 7:32 & 8:17 & 9:09 & 10:02 & 10:47 & 11:39 & 12:32 & 13:17 & 14:09 & 15:02 & 15:47 & 16:39 \\
Beverly/Edgewater Park & 6:43 & 7:36 & 8:21 & 9:13 & 10:06 & 10:51 & 11:43 & 12:36 & 13:21 & 14:13 & 15:06 & 15:51 & 16:43 \\
Delanco & 6:47 & 7:40 & 8:25 & 9:17 & 10:10 & 10:55 & 11:47 & 12:40 & 13:25 & 14:17 & 15:10 & 15:55 & 16:47 \\
Riverside & 6:49 & 7:42 & 8:27 & 9:19 & 10:12 & 10:57 & 11:49 & 12:42 & 13:27 & 14:19 & 15:12 & 15:57 & 16:49 \\
Cinnaminson & 6:54 & 7:50 & 8:32 & 9:24 & 10:20 & 11:02 & 11:54 & 12:50 & 13:32 & 14:24 & 15:20 & 16:02 & 16:54 \\
Riverton & 6:56 & 7:52 & 8:34 & 9:26 & 10:22 & 11:04 & 11:56 & 12:52 & 13:34 & 14:26 & 15:22 & 16:04 & 16:56 \\
Palmyra & 7:03 & 7:54 & 8:36 & 9:33 & 10:24 & 11:06 & 12:03 & 12:54 & 13:36 & 14:33 & 15:24 & 16:06 & 17:03 \\
Route 73/Pennsauken & 7:06 & 7:57 & 8:39 & 9:36 & 10:27 & 11:09 & 12:06 & 12:57 & 13:39 & 14:36 & 15:27 & 16:09 & 17:06 \\
Pennsauken TC & 7:08 & 7:59 & 8:41 & 9:38 & 10:29 & 11:11 & 12:08 & 12:59 & 13:41 & 14:38 & 15:29 & 16:11 & 17:08 \\
36th Street & 7:12 & 8:03 & 8:45 & 9:42 & 10:33 & 11:15 & 12:12 & 13:03 & 13:45 & 14:42 & 15:33 & 16:15 & 17:12 \\
Walter Rand TC & 7:17 & 8:08 & 8:50 & 9:47 & 10:38 & 11:20 & 12:17 & 13:08 & 13:50 & 14:47 & 15:38 & 16:20 & 17:17 \\
Cooper St/Rutgers & 7:20 & 8:11 & 8:53 & 9:50 & 10:41 & 11:23 & 12:20 & 13:11 & 13:53 & 14:50 & 15:41 & 16:23 & 17:20 \\
Aquarium & 7:22 & 8:13 & 8:55 & 9:52 & 10:43 & 11:25 & 12:22 & 13:13 & 13:55 & 14:52 & 15:43 & 16:25 & 17:22 \\
Entertainment Center & 7:25 & 8:16 & 8:58 & 9:55 & 10:46 & 11:28 & 12:25 & 13:16 & 13:58 & 14:55 & 15:46 & 16:28 & 17:25 \\
\bottomrule\end{tabular}\end{center}
\subsection*{To Camden (read down) --- trips 14--18 of 18}
\begin{center}\footnotesize\setlength{\tabcolsep}{3pt}\begin{tabular}{@{}lrrrrr@{}}\toprule
Station & 14 & 15 & 16 & 17 & 18 \\\midrule
Trenton TC & 17:06 & 17:51 & 18:43 & 19:36 & 20:21 \\
Hamilton Avenue & 17:08 & 17:53 & 18:45 & 19:38 & 20:23 \\
Cass Street & 17:10 & 17:55 & 18:47 & 19:40 & 20:25 \\
Bordentown & 17:17 & 18:02 & 18:54 & 19:47 & 20:32 \\
Roebling & 17:22 & 18:07 & 18:59 & 19:52 & 20:37 \\
Florence & 17:25 & 18:10 & 19:02 & 19:55 & 20:40 \\
Burlington Towne Ctr & 17:30 & 18:15 & 19:07 & 20:00 & 20:45 \\
Burlington South & 17:32 & 18:17 & 19:09 & 20:02 & 20:47 \\
Beverly/Edgewater Park & 17:36 & 18:21 & 19:13 & 20:06 & 20:51 \\
Delanco & 17:40 & 18:25 & 19:17 & 20:10 & 20:55 \\
Riverside & 17:42 & 18:27 & 19:19 & 20:12 & 20:57 \\
Cinnaminson & 17:50 & 18:32 & 19:24 & 20:20 & 21:02 \\
Riverton & 17:52 & 18:34 & 19:26 & 20:22 & 21:04 \\
Palmyra & 17:54 & 18:36 & 19:33 & 20:24 & 21:06 \\
Route 73/Pennsauken & 17:57 & 18:39 & 19:36 & 20:27 & 21:09 \\
Pennsauken TC & 17:59 & 18:41 & 19:38 & 20:29 & 21:11 \\
36th Street & 18:03 & 18:45 & 19:42 & 20:33 & 21:15 \\
Walter Rand TC & 18:08 & 18:50 & 19:47 & 20:38 & 21:20 \\
Cooper St/Rutgers & 18:11 & 18:53 & 19:50 & 20:41 & 21:23 \\
Aquarium & 18:13 & 18:55 & 19:52 & 20:43 & 21:25 \\
Entertainment Center & 18:16 & 18:58 & 19:55 & 20:46 & 21:28 \\
\bottomrule\end{tabular}\end{center}
\clearpage
\section{Proposal B8: 8 trains, 160-min cycle, 2.5-min meet buffers, 9-min layovers}\label{tt:proposal_B8}
\subsection*{To Trenton (read down) --- trips 1--13 of 41}
\begin{center}\footnotesize\setlength{\tabcolsep}{3pt}\begin{tabular}{@{}lrrrrrrrrrrrrr@{}}\toprule
Station & 1 & 2 & 3 & 4 & 5 & 6 & 7 & 8 & 9 & 10 & 11 & 12 & 13 \\\midrule
Entertainment Center & 5:45 & 6:08 & 6:33 & 6:56 & 7:12 & 7:28 & 7:43 & 8:00 & 8:25 & 8:48 & 9:13 & 9:36 & 10:08 \\
Aquarium & 5:47 & 6:10 & 6:35 & 6:58 & 7:14 & 7:30 & 7:45 & 8:02 & 8:27 & 8:50 & 9:15 & 9:38 & 10:10 \\
Cooper St/Rutgers & 5:49 & 6:12 & 6:37 & 7:00 & 7:16 & 7:32 & 7:47 & 8:04 & 8:29 & 8:52 & 9:17 & 9:40 & 10:12 \\
Walter Rand TC & 5:53 & 6:16 & 6:41 & 7:04 & 7:20 & 7:36 & 7:51 & 8:08 & 8:33 & 8:56 & 9:21 & 9:44 & 10:16 \\
36th Street & 6:02 & 6:21 & 6:47 & 7:09 & 7:25 & 7:41 & 7:56 & 8:13 & 8:42 & 9:01 & 9:27 & 9:49 & 10:21 \\
Pennsauken TC & 6:04 & 6:23 & 6:49 & 7:11 & 7:27 & 7:43 & 7:58 & 8:15 & 8:44 & 9:03 & 9:29 & 9:51 & 10:23 \\
Route 73/Pennsauken & 6:07 & 6:26 & 6:55 & 7:14 & 7:30 & 7:46 & 8:01 & 8:18 & 8:47 & 9:06 & 9:35 & 9:54 & 10:26 \\
Palmyra & 6:10 & 6:29 & 6:58 & 7:17 & 7:33 & 7:49 & 8:04 & 8:21 & 8:50 & 9:09 & 9:38 & 9:57 & 10:29 \\
Riverton & 6:12 & 6:31 & 7:00 & 7:19 & 7:35 & 7:51 & 8:06 & 8:23 & 8:52 & 9:11 & 9:40 & 9:59 & 10:31 \\
Cinnaminson & 6:14 & 6:33 & 7:02 & 7:21 & 7:37 & 7:53 & 8:08 & 8:25 & 8:54 & 9:13 & 9:42 & 10:01 & 10:33 \\
Riverside & 6:19 & 6:38 & 7:08 & 7:26 & 7:42 & 7:58 & 8:13 & 8:30 & 8:59 & 9:18 & 9:48 & 10:06 & 10:38 \\
Delanco & 6:21 & 6:40 & 7:10 & 7:28 & 7:44 & 8:00 & 8:15 & 8:32 & 9:01 & 9:20 & 9:50 & 10:08 & 10:40 \\
Beverly/Edgewater Park & 6:25 & 6:44 & 7:14 & 7:32 & 7:48 & 8:04 & 8:19 & 8:36 & 9:05 & 9:24 & 9:54 & 10:12 & 10:44 \\
Burlington South & 6:32 & 6:48 & 7:18 & 7:36 & 7:52 & 8:08 & 8:23 & 8:40 & 9:12 & 9:28 & 9:58 & 10:16 & 10:48 \\
Burlington Towne Ctr & 6:34 & 6:50 & 7:20 & 7:38 & 7:54 & 8:10 & 8:25 & 8:42 & 9:14 & 9:30 & 10:00 & 10:18 & 10:50 \\
Florence & 6:40 & 6:56 & 7:28 & 7:44 & 8:00 & 8:16 & 8:31 & 8:48 & 9:20 & 9:36 & 10:08 & 10:24 & 10:56 \\
Roebling & 6:44 & 6:59 & 7:31 & 7:47 & 8:03 & 8:19 & 8:35 & 8:52 & 9:24 & 9:39 & 10:11 & 10:27 & 10:59 \\
Bordentown & 6:49 & 7:04 & 7:36 & 7:52 & 8:08 & 8:24 & 8:40 & 8:57 & 9:29 & 9:44 & 10:16 & 10:32 & 11:04 \\
Cass Street & 6:55 & 7:10 & 7:42 & 7:58 & 8:14 & 8:31 & 8:46 & 9:03 & 9:35 & 9:50 & 10:22 & 10:38 & 11:11 \\
Hamilton Avenue & 6:57 & 7:12 & 7:44 & 8:00 & 8:16 & 8:33 & 8:48 & 9:05 & 9:37 & 9:52 & 10:24 & 10:40 & 11:13 \\
Trenton TC & 7:00 & 7:15 & 7:47 & 8:03 & 8:19 & 8:36 & 8:51 & 9:08 & 9:40 & 9:55 & 10:27 & 10:43 & 11:16 \\
\bottomrule\end{tabular}\end{center}
\subsection*{To Trenton (read down) --- trips 14--26 of 41}
\begin{center}\footnotesize\setlength{\tabcolsep}{3pt}\begin{tabular}{@{}lrrrrrrrrrrrrr@{}}\toprule
Station & 14 & 15 & 16 & 17 & 18 & 19 & 20 & 21 & 22 & 23 & 24 & 25 & 26 \\\midrule
Entertainment Center & 10:40 & 11:05 & 11:28 & 11:53 & 12:16 & 12:48 & 13:20 & 13:45 & 14:08 & 14:33 & 14:56 & 15:12 & 15:28 \\
Aquarium & 10:42 & 11:07 & 11:30 & 11:55 & 12:18 & 12:50 & 13:22 & 13:47 & 14:10 & 14:35 & 14:58 & 15:14 & 15:30 \\
Cooper St/Rutgers & 10:44 & 11:09 & 11:32 & 11:57 & 12:20 & 12:52 & 13:24 & 13:49 & 14:12 & 14:37 & 15:00 & 15:16 & 15:32 \\
Walter Rand TC & 10:48 & 11:13 & 11:36 & 12:01 & 12:24 & 12:56 & 13:28 & 13:53 & 14:16 & 14:41 & 15:04 & 15:20 & 15:36 \\
36th Street & 10:53 & 11:22 & 11:41 & 12:07 & 12:29 & 13:01 & 13:33 & 14:02 & 14:21 & 14:47 & 15:09 & 15:25 & 15:41 \\
Pennsauken TC & 10:55 & 11:24 & 11:43 & 12:09 & 12:31 & 13:03 & 13:35 & 14:04 & 14:23 & 14:49 & 15:11 & 15:27 & 15:43 \\
Route 73/Pennsauken & 10:58 & 11:27 & 11:46 & 12:15 & 12:34 & 13:06 & 13:38 & 14:07 & 14:26 & 14:55 & 15:14 & 15:30 & 15:46 \\
Palmyra & 11:01 & 11:30 & 11:49 & 12:18 & 12:37 & 13:09 & 13:41 & 14:10 & 14:29 & 14:58 & 15:17 & 15:33 & 15:49 \\
Riverton & 11:03 & 11:32 & 11:51 & 12:20 & 12:39 & 13:11 & 13:43 & 14:12 & 14:31 & 15:00 & 15:19 & 15:35 & 15:51 \\
Cinnaminson & 11:05 & 11:34 & 11:53 & 12:22 & 12:41 & 13:13 & 13:45 & 14:14 & 14:33 & 15:02 & 15:21 & 15:37 & 15:53 \\
Riverside & 11:10 & 11:39 & 11:58 & 12:28 & 12:46 & 13:18 & 13:50 & 14:19 & 14:38 & 15:08 & 15:26 & 15:42 & 15:58 \\
Delanco & 11:12 & 11:41 & 12:00 & 12:30 & 12:48 & 13:20 & 13:52 & 14:21 & 14:40 & 15:10 & 15:28 & 15:44 & 16:00 \\
Beverly/Edgewater Park & 11:16 & 11:45 & 12:04 & 12:34 & 12:52 & 13:24 & 13:56 & 14:25 & 14:44 & 15:14 & 15:32 & 15:48 & 16:04 \\
Burlington South & 11:20 & 11:52 & 12:08 & 12:38 & 12:56 & 13:28 & 14:00 & 14:32 & 14:48 & 15:18 & 15:36 & 15:52 & 16:08 \\
Burlington Towne Ctr & 11:22 & 11:54 & 12:10 & 12:40 & 12:58 & 13:30 & 14:02 & 14:34 & 14:50 & 15:20 & 15:38 & 15:54 & 16:10 \\
Florence & 11:28 & 12:00 & 12:16 & 12:48 & 13:04 & 13:36 & 14:08 & 14:40 & 14:56 & 15:28 & 15:44 & 16:00 & 16:16 \\
Roebling & 11:32 & 12:04 & 12:19 & 12:51 & 13:07 & 13:39 & 14:12 & 14:44 & 14:59 & 15:31 & 15:47 & 16:03 & 16:19 \\
Bordentown & 11:37 & 12:09 & 12:24 & 12:56 & 13:12 & 13:44 & 14:17 & 14:49 & 15:04 & 15:36 & 15:52 & 16:08 & 16:24 \\
Cass Street & 11:43 & 12:15 & 12:30 & 13:02 & 13:18 & 13:51 & 14:23 & 14:55 & 15:10 & 15:42 & 15:58 & 16:14 & 16:31 \\
Hamilton Avenue & 11:45 & 12:17 & 12:32 & 13:04 & 13:20 & 13:53 & 14:25 & 14:57 & 15:12 & 15:44 & 16:00 & 16:16 & 16:33 \\
Trenton TC & 11:48 & 12:20 & 12:35 & 13:07 & 13:23 & 13:56 & 14:28 & 15:00 & 15:15 & 15:47 & 16:03 & 16:19 & 16:36 \\
\bottomrule\end{tabular}\end{center}
\subsection*{To Trenton (read down) --- trips 27--39 of 41}
\begin{center}\footnotesize\setlength{\tabcolsep}{3pt}\begin{tabular}{@{}lrrrrrrrrrrrrr@{}}\toprule
Station & 27 & 28 & 29 & 30 & 31 & 32 & 33 & 34 & 35 & 36 & 37 & 38 & 39 \\\midrule
Entertainment Center & 15:43 & 16:00 & 16:25 & 16:48 & 17:13 & 17:36 & 17:52 & 18:08 & 18:23 & 18:40 & 19:05 & 19:28 & 19:53 \\
Aquarium & 15:45 & 16:02 & 16:27 & 16:50 & 17:15 & 17:38 & 17:54 & 18:10 & 18:25 & 18:42 & 19:07 & 19:30 & 19:55 \\
Cooper St/Rutgers & 15:47 & 16:04 & 16:29 & 16:52 & 17:17 & 17:40 & 17:56 & 18:12 & 18:27 & 18:44 & 19:09 & 19:32 & 19:57 \\
Walter Rand TC & 15:51 & 16:08 & 16:33 & 16:56 & 17:21 & 17:44 & 18:00 & 18:16 & 18:31 & 18:48 & 19:13 & 19:36 & 20:01 \\
36th Street & 15:56 & 16:13 & 16:42 & 17:01 & 17:27 & 17:49 & 18:05 & 18:21 & 18:36 & 18:53 & 19:22 & 19:41 & 20:07 \\
Pennsauken TC & 15:58 & 16:15 & 16:44 & 17:03 & 17:29 & 17:51 & 18:07 & 18:23 & 18:38 & 18:55 & 19:24 & 19:43 & 20:09 \\
Route 73/Pennsauken & 16:01 & 16:18 & 16:47 & 17:06 & 17:35 & 17:54 & 18:10 & 18:26 & 18:41 & 18:58 & 19:27 & 19:46 & 20:15 \\
Palmyra & 16:04 & 16:21 & 16:50 & 17:09 & 17:38 & 17:57 & 18:13 & 18:29 & 18:44 & 19:01 & 19:30 & 19:49 & 20:18 \\
Riverton & 16:06 & 16:23 & 16:52 & 17:11 & 17:40 & 17:59 & 18:15 & 18:31 & 18:46 & 19:03 & 19:32 & 19:51 & 20:20 \\
Cinnaminson & 16:08 & 16:25 & 16:54 & 17:13 & 17:42 & 18:01 & 18:17 & 18:33 & 18:48 & 19:05 & 19:34 & 19:53 & 20:22 \\
Riverside & 16:13 & 16:30 & 16:59 & 17:18 & 17:48 & 18:06 & 18:22 & 18:38 & 18:53 & 19:10 & 19:39 & 19:58 & 20:28 \\
Delanco & 16:15 & 16:32 & 17:01 & 17:20 & 17:50 & 18:08 & 18:24 & 18:40 & 18:55 & 19:12 & 19:41 & 20:00 & 20:30 \\
Beverly/Edgewater Park & 16:19 & 16:36 & 17:05 & 17:24 & 17:54 & 18:12 & 18:28 & 18:44 & 18:59 & 19:16 & 19:45 & 20:04 & 20:34 \\
Burlington South & 16:23 & 16:40 & 17:12 & 17:28 & 17:58 & 18:16 & 18:32 & 18:48 & 19:03 & 19:20 & 19:52 & 20:08 & 20:38 \\
Burlington Towne Ctr & 16:25 & 16:42 & 17:14 & 17:30 & 18:00 & 18:18 & 18:34 & 18:50 & 19:05 & 19:22 & 19:54 & 20:10 & 20:40 \\
Florence & 16:31 & 16:48 & 17:20 & 17:36 & 18:08 & 18:24 & 18:40 & 18:56 & 19:11 & 19:28 & 20:00 & 20:16 & 20:48 \\
Roebling & 16:35 & 16:52 & 17:24 & 17:39 & 18:11 & 18:27 & 18:43 & 18:59 & 19:15 & 19:32 & 20:04 & 20:19 & 20:51 \\
Bordentown & 16:40 & 16:57 & 17:29 & 17:44 & 18:16 & 18:32 & 18:48 & 19:04 & 19:20 & 19:37 & 20:09 & 20:24 & 20:56 \\
Cass Street & 16:46 & 17:03 & 17:35 & 17:50 & 18:22 & 18:38 & 18:54 & 19:11 & 19:26 & 19:43 & 20:15 & 20:30 & 21:02 \\
Hamilton Avenue & 16:48 & 17:05 & 17:37 & 17:52 & 18:24 & 18:40 & 18:56 & 19:13 & 19:28 & 19:45 & 20:17 & 20:32 & 21:04 \\
Trenton TC & 16:51 & 17:08 & 17:40 & 17:55 & 18:27 & 18:43 & 18:59 & 19:16 & 19:31 & 19:48 & 20:20 & 20:35 & 21:07 \\
\bottomrule\end{tabular}\end{center}
\subsection*{To Trenton (read down) --- trips 40--41 of 41}
\begin{center}\footnotesize\setlength{\tabcolsep}{3pt}\begin{tabular}{@{}lrr@{}}\toprule
Station & 40 & 41 \\\midrule
Entertainment Center & 20:16 & 20:48 \\
Aquarium & 20:18 & 20:50 \\
Cooper St/Rutgers & 20:20 & 20:52 \\
Walter Rand TC & 20:24 & 20:56 \\
36th Street & 20:29 & 21:01 \\
Pennsauken TC & 20:31 & 21:03 \\
Route 73/Pennsauken & 20:34 & 21:06 \\
Palmyra & 20:37 & 21:09 \\
Riverton & 20:39 & 21:11 \\
Cinnaminson & 20:41 & 21:13 \\
Riverside & 20:46 & 21:18 \\
Delanco & 20:48 & 21:20 \\
Beverly/Edgewater Park & 20:52 & 21:24 \\
Burlington South & 20:56 & 21:28 \\
Burlington Towne Ctr & 20:58 & 21:30 \\
Florence & 21:04 & 21:36 \\
Roebling & 21:07 & 21:39 \\
Bordentown & 21:12 & 21:44 \\
Cass Street & 21:18 & 21:51 \\
Hamilton Avenue & 21:20 & 21:53 \\
Trenton TC & 21:23 & 21:56 \\
\bottomrule\end{tabular}\end{center}
\subsection*{To Camden (read down) --- trips 1--13 of 42}
\begin{center}\footnotesize\setlength{\tabcolsep}{3pt}\begin{tabular}{@{}lrrrrrrrrrrrrr@{}}\toprule
Station & 1 & 2 & 3 & 4 & 5 & 6 & 7 & 8 & 9 & 10 & 11 & 12 & 13 \\\midrule
Trenton TC & 5:48 & 6:05 & 6:22 & 6:37 & 7:09 & 7:24 & 7:56 & 8:13 & 8:28 & 8:45 & 9:02 & 9:17 & 9:49 \\
Hamilton Avenue & 5:50 & 6:07 & 6:24 & 6:39 & 7:11 & 7:26 & 7:58 & 8:15 & 8:30 & 8:47 & 9:04 & 9:19 & 9:51 \\
Cass Street & 5:52 & 6:09 & 6:26 & 6:41 & 7:13 & 7:28 & 8:00 & 8:17 & 8:32 & 8:49 & 9:06 & 9:21 & 9:53 \\
Bordentown & 5:59 & 6:16 & 6:33 & 6:48 & 7:20 & 7:35 & 8:07 & 8:24 & 8:39 & 8:56 & 9:13 & 9:28 & 10:00 \\
Roebling & 6:05 & 6:21 & 6:38 & 6:53 & 7:25 & 7:40 & 8:12 & 8:29 & 8:45 & 9:01 & 9:18 & 9:33 & 10:05 \\
Florence & 6:08 & 6:24 & 6:41 & 6:56 & 7:28 & 7:43 & 8:15 & 8:32 & 8:48 & 9:04 & 9:21 & 9:36 & 10:08 \\
Burlington Towne Ctr & 6:13 & 6:29 & 6:46 & 7:01 & 7:33 & 7:49 & 8:21 & 8:37 & 8:53 & 9:09 & 9:26 & 9:41 & 10:13 \\
Burlington South & 6:15 & 6:31 & 6:48 & 7:03 & 7:35 & 7:51 & 8:23 & 8:39 & 8:55 & 9:11 & 9:28 & 9:43 & 10:15 \\
Beverly/Edgewater Park & 6:19 & 6:35 & 6:52 & 7:07 & 7:39 & 7:55 & 8:27 & 8:43 & 8:59 & 9:15 & 9:32 & 9:47 & 10:19 \\
Delanco & 6:23 & 6:39 & 6:56 & 7:11 & 7:43 & 7:59 & 8:31 & 8:47 & 9:03 & 9:19 & 9:36 & 9:51 & 10:23 \\
Riverside & 6:25 & 6:41 & 6:58 & 7:13 & 7:45 & 8:01 & 8:33 & 8:49 & 9:05 & 9:21 & 9:38 & 9:53 & 10:25 \\
Cinnaminson & 6:30 & 6:46 & 7:03 & 7:18 & 7:50 & 8:06 & 8:38 & 8:54 & 9:10 & 9:26 & 9:43 & 9:58 & 10:30 \\
Riverton & 6:32 & 6:48 & 7:05 & 7:20 & 7:52 & 8:08 & 8:40 & 8:56 & 9:12 & 9:28 & 9:45 & 10:00 & 10:32 \\
Palmyra & 6:34 & 6:50 & 7:07 & 7:23 & 7:54 & 8:11 & 8:42 & 8:59 & 9:14 & 9:30 & 9:47 & 10:03 & 10:34 \\
Route 73/Pennsauken & 6:37 & 6:53 & 7:10 & 7:26 & 7:57 & 8:14 & 8:45 & 9:02 & 9:17 & 9:33 & 9:50 & 10:06 & 10:37 \\
Pennsauken TC & 6:39 & 6:55 & 7:12 & 7:28 & 7:59 & 8:16 & 8:47 & 9:04 & 9:19 & 9:35 & 9:52 & 10:08 & 10:39 \\
36th Street & 6:43 & 6:59 & 7:16 & 7:32 & 8:03 & 8:20 & 8:51 & 9:08 & 9:23 & 9:39 & 9:56 & 10:12 & 10:43 \\
Walter Rand TC & 6:48 & 7:04 & 7:21 & 7:37 & 8:08 & 8:25 & 8:56 & 9:13 & 9:28 & 9:44 & 10:01 & 10:17 & 10:48 \\
Cooper St/Rutgers & 6:51 & 7:07 & 7:24 & 7:40 & 8:11 & 8:28 & 8:59 & 9:16 & 9:31 & 9:47 & 10:04 & 10:20 & 10:51 \\
Aquarium & 6:53 & 7:09 & 7:26 & 7:42 & 8:13 & 8:30 & 9:01 & 9:18 & 9:33 & 9:49 & 10:06 & 10:22 & 10:53 \\
Entertainment Center & 6:56 & 7:12 & 7:29 & 7:45 & 8:16 & 8:33 & 9:04 & 9:21 & 9:36 & 9:52 & 10:09 & 10:25 & 10:56 \\
\bottomrule\end{tabular}\end{center}
\subsection*{To Camden (read down) --- trips 14--26 of 42}
\begin{center}\footnotesize\setlength{\tabcolsep}{3pt}\begin{tabular}{@{}lrrrrrrrrrrrrr@{}}\toprule
Station & 14 & 15 & 16 & 17 & 18 & 19 & 20 & 21 & 22 & 23 & 24 & 25 & 26 \\\midrule
Trenton TC & 10:04 & 10:36 & 10:53 & 11:25 & 11:57 & 12:29 & 12:44 & 13:16 & 13:33 & 14:05 & 14:37 & 15:09 & 15:24 \\
Hamilton Avenue & 10:06 & 10:38 & 10:55 & 11:27 & 11:59 & 12:31 & 12:46 & 13:18 & 13:35 & 14:07 & 14:39 & 15:11 & 15:26 \\
Cass Street & 10:08 & 10:40 & 10:57 & 11:29 & 12:01 & 12:33 & 12:48 & 13:20 & 13:37 & 14:09 & 14:41 & 15:13 & 15:28 \\
Bordentown & 10:15 & 10:47 & 11:04 & 11:36 & 12:08 & 12:40 & 12:55 & 13:27 & 13:44 & 14:16 & 14:48 & 15:20 & 15:35 \\
Roebling & 10:20 & 10:52 & 11:09 & 11:41 & 12:13 & 12:45 & 13:00 & 13:32 & 13:49 & 14:21 & 14:53 & 15:25 & 15:40 \\
Florence & 10:23 & 10:55 & 11:12 & 11:44 & 12:16 & 12:48 & 13:03 & 13:35 & 13:52 & 14:24 & 14:56 & 15:28 & 15:43 \\
Burlington Towne Ctr & 10:29 & 11:01 & 11:17 & 11:49 & 12:21 & 12:53 & 13:09 & 13:41 & 13:57 & 14:29 & 15:01 & 15:33 & 15:49 \\
Burlington South & 10:31 & 11:03 & 11:19 & 11:51 & 12:23 & 12:55 & 13:11 & 13:43 & 13:59 & 14:31 & 15:03 & 15:35 & 15:51 \\
Beverly/Edgewater Park & 10:35 & 11:07 & 11:23 & 11:55 & 12:27 & 12:59 & 13:15 & 13:47 & 14:03 & 14:35 & 15:07 & 15:39 & 15:55 \\
Delanco & 10:39 & 11:11 & 11:27 & 11:59 & 12:31 & 13:03 & 13:19 & 13:51 & 14:07 & 14:39 & 15:11 & 15:43 & 15:59 \\
Riverside & 10:41 & 11:13 & 11:29 & 12:01 & 12:33 & 13:05 & 13:21 & 13:53 & 14:09 & 14:41 & 15:13 & 15:45 & 16:01 \\
Cinnaminson & 10:46 & 11:18 & 11:34 & 12:06 & 12:38 & 13:10 & 13:26 & 13:58 & 14:14 & 14:46 & 15:18 & 15:50 & 16:06 \\
Riverton & 10:48 & 11:20 & 11:36 & 12:08 & 12:40 & 13:12 & 13:28 & 14:00 & 14:16 & 14:48 & 15:20 & 15:52 & 16:08 \\
Palmyra & 10:51 & 11:22 & 11:39 & 12:10 & 12:43 & 13:14 & 13:31 & 14:02 & 14:19 & 14:50 & 15:23 & 15:54 & 16:11 \\
Route 73/Pennsauken & 10:54 & 11:25 & 11:42 & 12:13 & 12:46 & 13:17 & 13:34 & 14:05 & 14:22 & 14:53 & 15:26 & 15:57 & 16:14 \\
Pennsauken TC & 10:56 & 11:27 & 11:44 & 12:15 & 12:48 & 13:19 & 13:36 & 14:07 & 14:24 & 14:55 & 15:28 & 15:59 & 16:16 \\
36th Street & 11:00 & 11:31 & 11:48 & 12:19 & 12:52 & 13:23 & 13:40 & 14:11 & 14:28 & 14:59 & 15:32 & 16:03 & 16:20 \\
Walter Rand TC & 11:05 & 11:36 & 11:53 & 12:24 & 12:57 & 13:28 & 13:45 & 14:16 & 14:33 & 15:04 & 15:37 & 16:08 & 16:25 \\
Cooper St/Rutgers & 11:08 & 11:39 & 11:56 & 12:27 & 13:00 & 13:31 & 13:48 & 14:19 & 14:36 & 15:07 & 15:40 & 16:11 & 16:28 \\
Aquarium & 11:10 & 11:41 & 11:58 & 12:29 & 13:02 & 13:33 & 13:50 & 14:21 & 14:38 & 15:09 & 15:42 & 16:13 & 16:30 \\
Entertainment Center & 11:13 & 11:44 & 12:01 & 12:32 & 13:05 & 13:36 & 13:53 & 14:24 & 14:41 & 15:12 & 15:45 & 16:16 & 16:33 \\
\bottomrule\end{tabular}\end{center}
\subsection*{To Camden (read down) --- trips 27--39 of 42}
\begin{center}\footnotesize\setlength{\tabcolsep}{3pt}\begin{tabular}{@{}lrrrrrrrrrrrrr@{}}\toprule
Station & 27 & 28 & 29 & 30 & 31 & 32 & 33 & 34 & 35 & 36 & 37 & 38 & 39 \\\midrule
Trenton TC & 15:56 & 16:13 & 16:28 & 16:45 & 17:02 & 17:17 & 17:49 & 18:04 & 18:36 & 18:53 & 19:08 & 19:25 & 19:42 \\
Hamilton Avenue & 15:58 & 16:15 & 16:30 & 16:47 & 17:04 & 17:19 & 17:51 & 18:06 & 18:38 & 18:55 & 19:10 & 19:27 & 19:44 \\
Cass Street & 16:00 & 16:17 & 16:32 & 16:49 & 17:06 & 17:21 & 17:53 & 18:08 & 18:40 & 18:57 & 19:12 & 19:29 & 19:46 \\
Bordentown & 16:07 & 16:24 & 16:39 & 16:56 & 17:13 & 17:28 & 18:00 & 18:15 & 18:47 & 19:04 & 19:19 & 19:36 & 19:53 \\
Roebling & 16:12 & 16:29 & 16:45 & 17:01 & 17:18 & 17:33 & 18:05 & 18:20 & 18:52 & 19:09 & 19:25 & 19:41 & 19:58 \\
Florence & 16:15 & 16:32 & 16:48 & 17:04 & 17:21 & 17:36 & 18:08 & 18:23 & 18:55 & 19:12 & 19:28 & 19:44 & 20:01 \\
Burlington Towne Ctr & 16:21 & 16:37 & 16:53 & 17:09 & 17:26 & 17:41 & 18:13 & 18:29 & 19:01 & 19:17 & 19:33 & 19:49 & 20:06 \\
Burlington South & 16:23 & 16:39 & 16:55 & 17:11 & 17:28 & 17:43 & 18:15 & 18:31 & 19:03 & 19:19 & 19:35 & 19:51 & 20:08 \\
Beverly/Edgewater Park & 16:27 & 16:43 & 16:59 & 17:15 & 17:32 & 17:47 & 18:19 & 18:35 & 19:07 & 19:23 & 19:39 & 19:55 & 20:12 \\
Delanco & 16:31 & 16:47 & 17:03 & 17:19 & 17:36 & 17:51 & 18:23 & 18:39 & 19:11 & 19:27 & 19:43 & 19:59 & 20:16 \\
Riverside & 16:33 & 16:49 & 17:05 & 17:21 & 17:38 & 17:53 & 18:25 & 18:41 & 19:13 & 19:29 & 19:45 & 20:01 & 20:18 \\
Cinnaminson & 16:38 & 16:54 & 17:10 & 17:26 & 17:43 & 17:58 & 18:30 & 18:46 & 19:18 & 19:34 & 19:50 & 20:06 & 20:23 \\
Riverton & 16:40 & 16:56 & 17:12 & 17:28 & 17:45 & 18:00 & 18:32 & 18:48 & 19:20 & 19:36 & 19:52 & 20:08 & 20:25 \\
Palmyra & 16:42 & 16:59 & 17:14 & 17:30 & 17:47 & 18:03 & 18:34 & 18:51 & 19:22 & 19:39 & 19:54 & 20:10 & 20:27 \\
Route 73/Pennsauken & 16:45 & 17:02 & 17:17 & 17:33 & 17:50 & 18:06 & 18:37 & 18:54 & 19:25 & 19:42 & 19:57 & 20:13 & 20:30 \\
Pennsauken TC & 16:47 & 17:04 & 17:19 & 17:35 & 17:52 & 18:08 & 18:39 & 18:56 & 19:27 & 19:44 & 19:59 & 20:15 & 20:32 \\
36th Street & 16:51 & 17:08 & 17:23 & 17:39 & 17:56 & 18:12 & 18:43 & 19:00 & 19:31 & 19:48 & 20:03 & 20:19 & 20:36 \\
Walter Rand TC & 16:56 & 17:13 & 17:28 & 17:44 & 18:01 & 18:17 & 18:48 & 19:05 & 19:36 & 19:53 & 20:08 & 20:24 & 20:41 \\
Cooper St/Rutgers & 16:59 & 17:16 & 17:31 & 17:47 & 18:04 & 18:20 & 18:51 & 19:08 & 19:39 & 19:56 & 20:11 & 20:27 & 20:44 \\
Aquarium & 17:01 & 17:18 & 17:33 & 17:49 & 18:06 & 18:22 & 18:53 & 19:10 & 19:41 & 19:58 & 20:13 & 20:29 & 20:46 \\
Entertainment Center & 17:04 & 17:21 & 17:36 & 17:52 & 18:09 & 18:25 & 18:56 & 19:13 & 19:44 & 20:01 & 20:16 & 20:32 & 20:49 \\
\bottomrule\end{tabular}\end{center}
\subsection*{To Camden (read down) --- trips 40--42 of 42}
\begin{center}\footnotesize\setlength{\tabcolsep}{3pt}\begin{tabular}{@{}lrrr@{}}\toprule
Station & 40 & 41 & 42 \\\midrule
Trenton TC & 19:57 & 20:29 & 20:44 \\
Hamilton Avenue & 19:59 & 20:31 & 20:46 \\
Cass Street & 20:01 & 20:33 & 20:48 \\
Bordentown & 20:08 & 20:40 & 20:55 \\
Roebling & 20:13 & 20:45 & 21:00 \\
Florence & 20:16 & 20:48 & 21:03 \\
Burlington Towne Ctr & 20:21 & 20:53 & 21:09 \\
Burlington South & 20:23 & 20:55 & 21:11 \\
Beverly/Edgewater Park & 20:27 & 20:59 & 21:15 \\
Delanco & 20:31 & 21:03 & 21:19 \\
Riverside & 20:33 & 21:05 & 21:21 \\
Cinnaminson & 20:38 & 21:10 & 21:26 \\
Riverton & 20:40 & 21:12 & 21:28 \\
Palmyra & 20:43 & 21:14 & 21:31 \\
Route 73/Pennsauken & 20:46 & 21:17 & 21:34 \\
Pennsauken TC & 20:48 & 21:19 & 21:36 \\
36th Street & 20:52 & 21:23 & 21:40 \\
Walter Rand TC & 20:57 & 21:28 & 21:45 \\
Cooper St/Rutgers & 21:00 & 21:31 & 21:48 \\
Aquarium & 21:02 & 21:33 & 21:50 \\
Entertainment Center & 21:05 & 21:36 & 21:53 \\
\bottomrule\end{tabular}\end{center}
\clearpage
\section{Proposal B7: 7 trains, 160-min cycle, 2.5-min meet buffers, 9-min layovers}\label{tt:proposal_B7}
\subsection*{To Trenton (read down) --- trips 1--13 of 38}
\begin{center}\footnotesize\setlength{\tabcolsep}{3pt}\begin{tabular}{@{}lrrrrrrrrrrrrr@{}}\toprule
Station & 1 & 2 & 3 & 4 & 5 & 6 & 7 & 8 & 9 & 10 & 11 & 12 & 13 \\\midrule
Entertainment Center & 5:46 & 6:10 & 6:35 & 6:58 & 7:25 & 7:43 & 8:00 & 8:26 & 8:50 & 9:15 & 9:38 & 10:05 & 10:40 \\
Aquarium & 5:48 & 6:12 & 6:37 & 7:00 & 7:27 & 7:45 & 8:02 & 8:28 & 8:52 & 9:17 & 9:40 & 10:07 & 10:42 \\
Cooper St/Rutgers & 5:50 & 6:14 & 6:39 & 7:02 & 7:29 & 7:47 & 8:04 & 8:30 & 8:54 & 9:19 & 9:42 & 10:09 & 10:44 \\
Walter Rand TC & 5:54 & 6:18 & 6:43 & 7:06 & 7:33 & 7:51 & 8:08 & 8:34 & 8:58 & 9:23 & 9:46 & 10:13 & 10:48 \\
36th Street & 6:06 & 6:23 & 6:52 & 7:11 & 7:42 & 7:56 & 8:14 & 8:46 & 9:03 & 9:32 & 9:51 & 10:22 & 10:54 \\
Pennsauken TC & 6:08 & 6:25 & 6:54 & 7:13 & 7:44 & 7:58 & 8:16 & 8:48 & 9:05 & 9:34 & 9:53 & 10:24 & 10:56 \\
Route 73/Pennsauken & 6:11 & 6:28 & 6:57 & 7:16 & 7:47 & 8:01 & 8:19 & 8:51 & 9:08 & 9:37 & 9:56 & 10:27 & 10:59 \\
Palmyra & 6:14 & 6:31 & 7:00 & 7:19 & 7:50 & 8:04 & 8:22 & 8:54 & 9:11 & 9:40 & 9:59 & 10:30 & 11:02 \\
Riverton & 6:16 & 6:33 & 7:02 & 7:21 & 7:52 & 8:06 & 8:24 & 8:56 & 9:13 & 9:42 & 10:01 & 10:32 & 11:04 \\
Cinnaminson & 6:18 & 6:35 & 7:04 & 7:23 & 7:54 & 8:08 & 8:26 & 8:58 & 9:15 & 9:44 & 10:03 & 10:34 & 11:06 \\
Riverside & 6:23 & 6:40 & 7:09 & 7:28 & 7:59 & 8:13 & 8:31 & 9:03 & 9:20 & 9:49 & 10:08 & 10:39 & 11:11 \\
Delanco & 6:25 & 6:42 & 7:11 & 7:30 & 8:01 & 8:15 & 8:33 & 9:05 & 9:22 & 9:51 & 10:10 & 10:41 & 11:13 \\
Beverly/Edgewater Park & 6:29 & 6:46 & 7:15 & 7:34 & 8:05 & 8:19 & 8:37 & 9:09 & 9:26 & 9:55 & 10:14 & 10:45 & 11:17 \\
Burlington South & 6:34 & 6:50 & 7:21 & 7:38 & 8:10 & 8:23 & 8:41 & 9:14 & 9:30 & 10:01 & 10:18 & 10:50 & 11:21 \\
Burlington Towne Ctr & 6:36 & 6:52 & 7:23 & 7:40 & 8:12 & 8:25 & 8:43 & 9:16 & 9:32 & 10:03 & 10:20 & 10:52 & 11:23 \\
Florence & 6:42 & 6:58 & 7:29 & 7:46 & 8:18 & 8:33 & 8:50 & 9:22 & 9:38 & 10:09 & 10:26 & 10:58 & 11:30 \\
Roebling & 6:45 & 7:01 & 7:34 & 7:49 & 8:21 & 8:36 & 8:53 & 9:25 & 9:41 & 10:14 & 10:29 & 11:01 & 11:33 \\
Bordentown & 6:50 & 7:06 & 7:39 & 7:54 & 8:26 & 8:41 & 8:58 & 9:30 & 9:46 & 10:19 & 10:34 & 11:06 & 11:38 \\
Cass Street & 6:56 & 7:12 & 7:45 & 8:00 & 8:32 & 8:47 & 9:04 & 9:36 & 9:52 & 10:25 & 10:40 & 11:12 & 11:44 \\
Hamilton Avenue & 6:58 & 7:14 & 7:47 & 8:02 & 8:34 & 8:49 & 9:06 & 9:38 & 9:54 & 10:27 & 10:42 & 11:14 & 11:46 \\
Trenton TC & 7:01 & 7:17 & 7:50 & 8:05 & 8:37 & 8:52 & 9:09 & 9:41 & 9:57 & 10:30 & 10:45 & 11:17 & 11:49 \\
\bottomrule\end{tabular}\end{center}
\subsection*{To Trenton (read down) --- trips 14--26 of 38}
\begin{center}\footnotesize\setlength{\tabcolsep}{3pt}\begin{tabular}{@{}lrrrrrrrrrrrrr@{}}\toprule
Station & 14 & 15 & 16 & 17 & 18 & 19 & 20 & 21 & 22 & 23 & 24 & 25 & 26 \\\midrule
Entertainment Center & 11:06 & 11:30 & 11:55 & 12:18 & 12:45 & 13:20 & 13:46 & 14:10 & 14:35 & 14:58 & 15:25 & 15:43 & 16:00 \\
Aquarium & 11:08 & 11:32 & 11:57 & 12:20 & 12:47 & 13:22 & 13:48 & 14:12 & 14:37 & 15:00 & 15:27 & 15:45 & 16:02 \\
Cooper St/Rutgers & 11:10 & 11:34 & 11:59 & 12:22 & 12:49 & 13:24 & 13:50 & 14:14 & 14:39 & 15:02 & 15:29 & 15:47 & 16:04 \\
Walter Rand TC & 11:14 & 11:38 & 12:03 & 12:26 & 12:53 & 13:28 & 13:54 & 14:18 & 14:43 & 15:06 & 15:33 & 15:51 & 16:08 \\
36th Street & 11:26 & 11:43 & 12:12 & 12:31 & 13:02 & 13:34 & 14:06 & 14:23 & 14:52 & 15:11 & 15:42 & 15:56 & 16:14 \\
Pennsauken TC & 11:28 & 11:45 & 12:14 & 12:33 & 13:04 & 13:36 & 14:08 & 14:25 & 14:54 & 15:13 & 15:44 & 15:58 & 16:16 \\
Route 73/Pennsauken & 11:31 & 11:48 & 12:17 & 12:36 & 13:07 & 13:39 & 14:11 & 14:28 & 14:57 & 15:16 & 15:47 & 16:01 & 16:19 \\
Palmyra & 11:34 & 11:51 & 12:20 & 12:39 & 13:10 & 13:42 & 14:14 & 14:31 & 15:00 & 15:19 & 15:50 & 16:04 & 16:22 \\
Riverton & 11:36 & 11:53 & 12:22 & 12:41 & 13:12 & 13:44 & 14:16 & 14:33 & 15:02 & 15:21 & 15:52 & 16:06 & 16:24 \\
Cinnaminson & 11:38 & 11:55 & 12:24 & 12:43 & 13:14 & 13:46 & 14:18 & 14:35 & 15:04 & 15:23 & 15:54 & 16:08 & 16:26 \\
Riverside & 11:43 & 12:00 & 12:29 & 12:48 & 13:19 & 13:51 & 14:23 & 14:40 & 15:09 & 15:28 & 15:59 & 16:13 & 16:31 \\
Delanco & 11:45 & 12:02 & 12:31 & 12:50 & 13:21 & 13:53 & 14:25 & 14:42 & 15:11 & 15:30 & 16:01 & 16:15 & 16:33 \\
Beverly/Edgewater Park & 11:49 & 12:06 & 12:35 & 12:54 & 13:25 & 13:57 & 14:29 & 14:46 & 15:15 & 15:34 & 16:05 & 16:19 & 16:37 \\
Burlington South & 11:54 & 12:10 & 12:41 & 12:58 & 13:30 & 14:01 & 14:34 & 14:50 & 15:21 & 15:38 & 16:10 & 16:23 & 16:41 \\
Burlington Towne Ctr & 11:56 & 12:12 & 12:43 & 13:00 & 13:32 & 14:03 & 14:36 & 14:52 & 15:23 & 15:40 & 16:12 & 16:25 & 16:43 \\
Florence & 12:02 & 12:18 & 12:49 & 13:06 & 13:38 & 14:10 & 14:42 & 14:58 & 15:29 & 15:46 & 16:18 & 16:33 & 16:50 \\
Roebling & 12:05 & 12:21 & 12:54 & 13:09 & 13:41 & 14:13 & 14:45 & 15:01 & 15:34 & 15:49 & 16:21 & 16:36 & 16:53 \\
Bordentown & 12:10 & 12:26 & 12:59 & 13:14 & 13:46 & 14:18 & 14:50 & 15:06 & 15:39 & 15:54 & 16:26 & 16:41 & 16:58 \\
Cass Street & 12:16 & 12:32 & 13:05 & 13:20 & 13:52 & 14:24 & 14:56 & 15:12 & 15:45 & 16:00 & 16:32 & 16:47 & 17:04 \\
Hamilton Avenue & 12:18 & 12:34 & 13:07 & 13:22 & 13:54 & 14:26 & 14:58 & 15:14 & 15:47 & 16:02 & 16:34 & 16:49 & 17:06 \\
Trenton TC & 12:21 & 12:37 & 13:10 & 13:25 & 13:57 & 14:29 & 15:01 & 15:17 & 15:50 & 16:05 & 16:37 & 16:52 & 17:09 \\
\bottomrule\end{tabular}\end{center}
\subsection*{To Trenton (read down) --- trips 27--38 of 38}
\begin{center}\footnotesize\setlength{\tabcolsep}{3pt}\begin{tabular}{@{}lrrrrrrrrrrrr@{}}\toprule
Station & 27 & 28 & 29 & 30 & 31 & 32 & 33 & 34 & 35 & 36 & 37 & 38 \\\midrule
Entertainment Center & 16:26 & 16:50 & 17:15 & 17:38 & 18:05 & 18:23 & 18:40 & 19:06 & 19:30 & 19:55 & 20:18 & 20:45 \\
Aquarium & 16:28 & 16:52 & 17:17 & 17:40 & 18:07 & 18:25 & 18:42 & 19:08 & 19:32 & 19:57 & 20:20 & 20:47 \\
Cooper St/Rutgers & 16:30 & 16:54 & 17:19 & 17:42 & 18:09 & 18:27 & 18:44 & 19:10 & 19:34 & 19:59 & 20:22 & 20:49 \\
Walter Rand TC & 16:34 & 16:58 & 17:23 & 17:46 & 18:13 & 18:31 & 18:48 & 19:14 & 19:38 & 20:03 & 20:26 & 20:53 \\
36th Street & 16:46 & 17:03 & 17:32 & 17:51 & 18:22 & 18:36 & 18:54 & 19:26 & 19:43 & 20:12 & 20:31 & 21:02 \\
Pennsauken TC & 16:48 & 17:05 & 17:34 & 17:53 & 18:24 & 18:38 & 18:56 & 19:28 & 19:45 & 20:14 & 20:33 & 21:04 \\
Route 73/Pennsauken & 16:51 & 17:08 & 17:37 & 17:56 & 18:27 & 18:41 & 18:59 & 19:31 & 19:48 & 20:17 & 20:36 & 21:07 \\
Palmyra & 16:54 & 17:11 & 17:40 & 17:59 & 18:30 & 18:44 & 19:02 & 19:34 & 19:51 & 20:20 & 20:39 & 21:10 \\
Riverton & 16:56 & 17:13 & 17:42 & 18:01 & 18:32 & 18:46 & 19:04 & 19:36 & 19:53 & 20:22 & 20:41 & 21:12 \\
Cinnaminson & 16:58 & 17:15 & 17:44 & 18:03 & 18:34 & 18:48 & 19:06 & 19:38 & 19:55 & 20:24 & 20:43 & 21:14 \\
Riverside & 17:03 & 17:20 & 17:49 & 18:08 & 18:39 & 18:53 & 19:11 & 19:43 & 20:00 & 20:29 & 20:48 & 21:19 \\
Delanco & 17:05 & 17:22 & 17:51 & 18:10 & 18:41 & 18:55 & 19:13 & 19:45 & 20:02 & 20:31 & 20:50 & 21:21 \\
Beverly/Edgewater Park & 17:09 & 17:26 & 17:55 & 18:14 & 18:45 & 18:59 & 19:17 & 19:49 & 20:06 & 20:35 & 20:54 & 21:25 \\
Burlington South & 17:14 & 17:30 & 18:01 & 18:18 & 18:50 & 19:03 & 19:21 & 19:54 & 20:10 & 20:41 & 20:58 & 21:30 \\
Burlington Towne Ctr & 17:16 & 17:32 & 18:03 & 18:20 & 18:52 & 19:05 & 19:23 & 19:56 & 20:12 & 20:43 & 21:00 & 21:32 \\
Florence & 17:22 & 17:38 & 18:09 & 18:26 & 18:58 & 19:13 & 19:30 & 20:02 & 20:18 & 20:49 & 21:06 & 21:38 \\
Roebling & 17:25 & 17:41 & 18:14 & 18:29 & 19:01 & 19:16 & 19:33 & 20:05 & 20:21 & 20:54 & 21:09 & 21:41 \\
Bordentown & 17:30 & 17:46 & 18:19 & 18:34 & 19:06 & 19:21 & 19:38 & 20:10 & 20:26 & 20:59 & 21:14 & 21:46 \\
Cass Street & 17:36 & 17:52 & 18:25 & 18:40 & 19:12 & 19:27 & 19:44 & 20:16 & 20:32 & 21:05 & 21:20 & 21:52 \\
Hamilton Avenue & 17:38 & 17:54 & 18:27 & 18:42 & 19:14 & 19:29 & 19:46 & 20:18 & 20:34 & 21:07 & 21:22 & 21:54 \\
Trenton TC & 17:41 & 17:57 & 18:30 & 18:45 & 19:17 & 19:32 & 19:49 & 20:21 & 20:37 & 21:10 & 21:25 & 21:57 \\
\bottomrule\end{tabular}\end{center}
\subsection*{To Camden (read down) --- trips 1--13 of 38}
\begin{center}\footnotesize\setlength{\tabcolsep}{3pt}\begin{tabular}{@{}lrrrrrrrrrrrrr@{}}\toprule
Station & 1 & 2 & 3 & 4 & 5 & 6 & 7 & 8 & 9 & 10 & 11 & 12 & 13 \\\midrule
Trenton TC & 6:06 & 6:23 & 6:38 & 7:10 & 7:27 & 7:59 & 8:14 & 8:46 & 9:03 & 9:18 & 9:50 & 10:07 & 10:39 \\
Hamilton Avenue & 6:08 & 6:25 & 6:40 & 7:12 & 7:29 & 8:01 & 8:16 & 8:48 & 9:05 & 9:20 & 9:52 & 10:09 & 10:41 \\
Cass Street & 6:10 & 6:27 & 6:42 & 7:14 & 7:31 & 8:03 & 8:18 & 8:50 & 9:07 & 9:22 & 9:54 & 10:11 & 10:43 \\
Bordentown & 6:17 & 6:34 & 6:49 & 7:21 & 7:38 & 8:10 & 8:25 & 8:57 & 9:14 & 9:29 & 10:01 & 10:18 & 10:50 \\
Roebling & 6:22 & 6:39 & 6:55 & 7:26 & 7:43 & 8:15 & 8:30 & 9:02 & 9:19 & 9:35 & 10:06 & 10:23 & 10:55 \\
Florence & 6:25 & 6:42 & 6:58 & 7:29 & 7:46 & 8:18 & 8:33 & 9:05 & 9:22 & 9:38 & 10:09 & 10:26 & 10:58 \\
Burlington Towne Ctr & 6:30 & 6:47 & 7:03 & 7:34 & 7:51 & 8:23 & 8:38 & 9:10 & 9:27 & 9:43 & 10:14 & 10:31 & 11:03 \\
Burlington South & 6:32 & 6:49 & 7:05 & 7:36 & 7:53 & 8:25 & 8:40 & 9:12 & 9:29 & 9:45 & 10:16 & 10:33 & 11:05 \\
Beverly/Edgewater Park & 6:36 & 6:53 & 7:09 & 7:40 & 7:57 & 8:29 & 8:44 & 9:16 & 9:33 & 9:49 & 10:20 & 10:37 & 11:09 \\
Delanco & 6:40 & 6:57 & 7:13 & 7:44 & 8:01 & 8:33 & 8:48 & 9:20 & 9:37 & 9:53 & 10:24 & 10:41 & 11:13 \\
Riverside & 6:42 & 6:59 & 7:15 & 7:46 & 8:03 & 8:35 & 8:50 & 9:22 & 9:39 & 9:55 & 10:26 & 10:43 & 11:15 \\
Cinnaminson & 6:47 & 7:04 & 7:20 & 7:51 & 8:08 & 8:40 & 8:55 & 9:27 & 9:44 & 10:00 & 10:31 & 10:48 & 11:20 \\
Riverton & 6:49 & 7:06 & 7:22 & 7:53 & 8:10 & 8:42 & 8:57 & 9:29 & 9:46 & 10:02 & 10:33 & 10:50 & 11:22 \\
Palmyra & 6:51 & 7:08 & 7:24 & 7:55 & 8:12 & 8:44 & 8:59 & 9:31 & 9:48 & 10:04 & 10:35 & 10:52 & 11:24 \\
Route 73/Pennsauken & 6:54 & 7:11 & 7:27 & 7:58 & 8:15 & 8:47 & 9:02 & 9:34 & 9:51 & 10:07 & 10:38 & 10:55 & 11:27 \\
Pennsauken TC & 6:56 & 7:13 & 7:29 & 8:00 & 8:17 & 8:49 & 9:04 & 9:36 & 9:53 & 10:09 & 10:40 & 10:57 & 11:29 \\
36th Street & 7:00 & 7:18 & 7:33 & 8:04 & 8:21 & 8:53 & 9:10 & 9:40 & 9:58 & 10:13 & 10:44 & 11:01 & 11:33 \\
Walter Rand TC & 7:05 & 7:23 & 7:38 & 8:09 & 8:26 & 8:58 & 9:15 & 9:45 & 10:03 & 10:18 & 10:49 & 11:06 & 11:38 \\
Cooper St/Rutgers & 7:08 & 7:26 & 7:41 & 8:12 & 8:29 & 9:01 & 9:18 & 9:48 & 10:06 & 10:21 & 10:52 & 11:09 & 11:41 \\
Aquarium & 7:10 & 7:28 & 7:43 & 8:14 & 8:31 & 9:03 & 9:20 & 9:50 & 10:08 & 10:23 & 10:54 & 11:11 & 11:43 \\
Entertainment Center & 7:13 & 7:31 & 7:46 & 8:17 & 8:34 & 9:06 & 9:23 & 9:53 & 10:11 & 10:26 & 10:57 & 11:14 & 11:46 \\
\bottomrule\end{tabular}\end{center}
\subsection*{To Camden (read down) --- trips 14--26 of 38}
\begin{center}\footnotesize\setlength{\tabcolsep}{3pt}\begin{tabular}{@{}lrrrrrrrrrrrrr@{}}\toprule
Station & 14 & 15 & 16 & 17 & 18 & 19 & 20 & 21 & 22 & 23 & 24 & 25 & 26 \\\midrule
Trenton TC & 10:54 & 11:26 & 11:58 & 12:30 & 12:47 & 13:19 & 13:34 & 14:06 & 14:38 & 15:10 & 15:27 & 15:59 & 16:14 \\
Hamilton Avenue & 10:56 & 11:28 & 12:00 & 12:32 & 12:49 & 13:21 & 13:36 & 14:08 & 14:40 & 15:12 & 15:29 & 16:01 & 16:16 \\
Cass Street & 10:58 & 11:30 & 12:02 & 12:34 & 12:51 & 13:23 & 13:38 & 14:10 & 14:42 & 15:14 & 15:31 & 16:03 & 16:18 \\
Bordentown & 11:05 & 11:37 & 12:09 & 12:41 & 12:58 & 13:30 & 13:45 & 14:17 & 14:49 & 15:21 & 15:38 & 16:10 & 16:25 \\
Roebling & 11:10 & 11:42 & 12:15 & 12:46 & 13:03 & 13:35 & 13:50 & 14:22 & 14:55 & 15:26 & 15:43 & 16:15 & 16:30 \\
Florence & 11:13 & 11:45 & 12:18 & 12:49 & 13:06 & 13:38 & 13:53 & 14:25 & 14:58 & 15:29 & 15:46 & 16:18 & 16:33 \\
Burlington Towne Ctr & 11:18 & 11:50 & 12:23 & 12:54 & 13:11 & 13:43 & 13:58 & 14:30 & 15:03 & 15:34 & 15:51 & 16:23 & 16:38 \\
Burlington South & 11:20 & 11:52 & 12:25 & 12:56 & 13:13 & 13:45 & 14:00 & 14:32 & 15:05 & 15:36 & 15:53 & 16:25 & 16:40 \\
Beverly/Edgewater Park & 11:24 & 11:56 & 12:29 & 13:00 & 13:17 & 13:49 & 14:04 & 14:36 & 15:09 & 15:40 & 15:57 & 16:29 & 16:44 \\
Delanco & 11:28 & 12:00 & 12:33 & 13:04 & 13:21 & 13:53 & 14:08 & 14:40 & 15:13 & 15:44 & 16:01 & 16:33 & 16:48 \\
Riverside & 11:30 & 12:02 & 12:35 & 13:06 & 13:23 & 13:55 & 14:10 & 14:42 & 15:15 & 15:46 & 16:03 & 16:35 & 16:50 \\
Cinnaminson & 11:35 & 12:07 & 12:40 & 13:11 & 13:28 & 14:00 & 14:15 & 14:47 & 15:20 & 15:51 & 16:08 & 16:40 & 16:55 \\
Riverton & 11:37 & 12:09 & 12:42 & 13:13 & 13:30 & 14:02 & 14:17 & 14:49 & 15:22 & 15:53 & 16:10 & 16:42 & 16:57 \\
Palmyra & 11:39 & 12:11 & 12:44 & 13:15 & 13:32 & 14:04 & 14:19 & 14:51 & 15:24 & 15:55 & 16:12 & 16:44 & 16:59 \\
Route 73/Pennsauken & 11:42 & 12:14 & 12:47 & 13:18 & 13:35 & 14:07 & 14:22 & 14:54 & 15:27 & 15:58 & 16:15 & 16:47 & 17:02 \\
Pennsauken TC & 11:44 & 12:16 & 12:49 & 13:20 & 13:37 & 14:09 & 14:24 & 14:56 & 15:29 & 16:00 & 16:17 & 16:49 & 17:04 \\
36th Street & 11:50 & 12:20 & 12:53 & 13:24 & 13:41 & 14:13 & 14:30 & 15:00 & 15:33 & 16:04 & 16:21 & 16:53 & 17:10 \\
Walter Rand TC & 11:55 & 12:25 & 12:58 & 13:29 & 13:46 & 14:18 & 14:35 & 15:05 & 15:38 & 16:09 & 16:26 & 16:58 & 17:15 \\
Cooper St/Rutgers & 11:58 & 12:28 & 13:01 & 13:32 & 13:49 & 14:21 & 14:38 & 15:08 & 15:41 & 16:12 & 16:29 & 17:01 & 17:18 \\
Aquarium & 12:00 & 12:30 & 13:03 & 13:34 & 13:51 & 14:23 & 14:40 & 15:10 & 15:43 & 16:14 & 16:31 & 17:03 & 17:20 \\
Entertainment Center & 12:03 & 12:33 & 13:06 & 13:37 & 13:54 & 14:26 & 14:43 & 15:13 & 15:46 & 16:17 & 16:34 & 17:06 & 17:23 \\
\bottomrule\end{tabular}\end{center}
\subsection*{To Camden (read down) --- trips 27--38 of 38}
\begin{center}\footnotesize\setlength{\tabcolsep}{3pt}\begin{tabular}{@{}lrrrrrrrrrrrr@{}}\toprule
Station & 27 & 28 & 29 & 30 & 31 & 32 & 33 & 34 & 35 & 36 & 37 & 38 \\\midrule
Trenton TC & 16:46 & 17:03 & 17:18 & 17:50 & 18:07 & 18:39 & 18:54 & 19:26 & 19:43 & 19:58 & 20:30 & 20:47 \\
Hamilton Avenue & 16:48 & 17:05 & 17:20 & 17:52 & 18:09 & 18:41 & 18:56 & 19:28 & 19:45 & 20:00 & 20:32 & 20:49 \\
Cass Street & 16:50 & 17:07 & 17:22 & 17:54 & 18:11 & 18:43 & 18:58 & 19:30 & 19:47 & 20:02 & 20:34 & 20:51 \\
Bordentown & 16:57 & 17:14 & 17:29 & 18:01 & 18:18 & 18:50 & 19:05 & 19:37 & 19:54 & 20:09 & 20:41 & 20:58 \\
Roebling & 17:02 & 17:19 & 17:35 & 18:06 & 18:23 & 18:55 & 19:10 & 19:42 & 19:59 & 20:15 & 20:46 & 21:03 \\
Florence & 17:05 & 17:22 & 17:38 & 18:09 & 18:26 & 18:58 & 19:13 & 19:45 & 20:02 & 20:18 & 20:49 & 21:06 \\
Burlington Towne Ctr & 17:10 & 17:27 & 17:43 & 18:14 & 18:31 & 19:03 & 19:18 & 19:50 & 20:07 & 20:23 & 20:54 & 21:11 \\
Burlington South & 17:12 & 17:29 & 17:45 & 18:16 & 18:33 & 19:05 & 19:20 & 19:52 & 20:09 & 20:25 & 20:56 & 21:13 \\
Beverly/Edgewater Park & 17:16 & 17:33 & 17:49 & 18:20 & 18:37 & 19:09 & 19:24 & 19:56 & 20:13 & 20:29 & 21:00 & 21:17 \\
Delanco & 17:20 & 17:37 & 17:53 & 18:24 & 18:41 & 19:13 & 19:28 & 20:00 & 20:17 & 20:33 & 21:04 & 21:21 \\
Riverside & 17:22 & 17:39 & 17:55 & 18:26 & 18:43 & 19:15 & 19:30 & 20:02 & 20:19 & 20:35 & 21:06 & 21:23 \\
Cinnaminson & 17:27 & 17:44 & 18:00 & 18:31 & 18:48 & 19:20 & 19:35 & 20:07 & 20:24 & 20:40 & 21:11 & 21:28 \\
Riverton & 17:29 & 17:46 & 18:02 & 18:33 & 18:50 & 19:22 & 19:37 & 20:09 & 20:26 & 20:42 & 21:13 & 21:30 \\
Palmyra & 17:31 & 17:48 & 18:04 & 18:35 & 18:52 & 19:24 & 19:39 & 20:11 & 20:28 & 20:44 & 21:15 & 21:32 \\
Route 73/Pennsauken & 17:34 & 17:51 & 18:07 & 18:38 & 18:55 & 19:27 & 19:42 & 20:14 & 20:31 & 20:47 & 21:18 & 21:35 \\
Pennsauken TC & 17:36 & 17:53 & 18:09 & 18:40 & 18:57 & 19:29 & 19:44 & 20:16 & 20:33 & 20:49 & 21:20 & 21:37 \\
36th Street & 17:40 & 17:58 & 18:13 & 18:44 & 19:01 & 19:33 & 19:50 & 20:20 & 20:38 & 20:53 & 21:24 & 21:41 \\
Walter Rand TC & 17:45 & 18:03 & 18:18 & 18:49 & 19:06 & 19:38 & 19:55 & 20:25 & 20:43 & 20:58 & 21:29 & 21:46 \\
Cooper St/Rutgers & 17:48 & 18:06 & 18:21 & 18:52 & 19:09 & 19:41 & 19:58 & 20:28 & 20:46 & 21:01 & 21:32 & 21:49 \\
Aquarium & 17:50 & 18:08 & 18:23 & 18:54 & 19:11 & 19:43 & 20:00 & 20:30 & 20:48 & 21:03 & 21:34 & 21:51 \\
Entertainment Center & 17:53 & 18:11 & 18:26 & 18:57 & 19:14 & 19:46 & 20:03 & 20:33 & 20:51 & 21:06 & 21:37 & 21:54 \\
\bottomrule\end{tabular}\end{center}
\clearpage
\section{Proposal E5: 8 peak, 5 midday (frees a car)}\label{tt:proposal_E5}
\subsection*{To Trenton (read down) --- trips 1--13 of 40}
\begin{center}\footnotesize\setlength{\tabcolsep}{3pt}\begin{tabular}{@{}lrrrrrrrrrrrrr@{}}\toprule
Station & 1 & 2 & 3 & 4 & 5 & 6 & 7 & 8 & 9 & 10 & 11 & 12 & 13 \\\midrule
Entertainment Center & 5:42 & 5:59 & 6:24 & 6:41 & 6:57 & 7:13 & 7:30 & 7:55 & 8:12 & 8:29 & 8:54 & 9:27 & 10:00 \\
Aquarium & 5:44 & 6:01 & 6:26 & 6:43 & 6:59 & 7:15 & 7:32 & 7:57 & 8:14 & 8:31 & 8:56 & 9:29 & 10:02 \\
Cooper St/Rutgers & 5:46 & 6:03 & 6:28 & 6:45 & 7:01 & 7:17 & 7:34 & 7:59 & 8:16 & 8:33 & 8:58 & 9:31 & 10:04 \\
Walter Rand TC & 5:50 & 6:07 & 6:32 & 6:49 & 7:05 & 7:21 & 7:38 & 8:03 & 8:20 & 8:37 & 9:02 & 9:35 & 10:08 \\
36th Street & 5:55 & 6:12 & 6:37 & 6:56 & 7:10 & 7:26 & 7:43 & 8:08 & 8:25 & 8:42 & 9:07 & 9:40 & 10:13 \\
Pennsauken TC & 5:57 & 6:14 & 6:39 & 6:58 & 7:12 & 7:28 & 7:45 & 8:10 & 8:27 & 8:44 & 9:09 & 9:42 & 10:15 \\
Route 73/Pennsauken & 6:00 & 6:17 & 6:44 & 7:01 & 7:15 & 7:31 & 7:48 & 8:13 & 8:30 & 8:47 & 9:14 & 9:45 & 10:18 \\
Palmyra & 6:03 & 6:20 & 6:47 & 7:04 & 7:18 & 7:34 & 7:51 & 8:16 & 8:33 & 8:50 & 9:17 & 9:48 & 10:21 \\
Riverton & 6:05 & 6:22 & 6:49 & 7:06 & 7:20 & 7:36 & 7:53 & 8:18 & 8:35 & 8:52 & 9:19 & 9:50 & 10:23 \\
Cinnaminson & 6:07 & 6:24 & 6:51 & 7:08 & 7:22 & 7:38 & 7:55 & 8:20 & 8:37 & 8:54 & 9:21 & 9:52 & 10:25 \\
Riverside & 6:12 & 6:29 & 6:56 & 7:13 & 7:27 & 7:43 & 8:00 & 8:25 & 8:42 & 8:59 & 9:26 & 9:57 & 10:30 \\
Delanco & 6:14 & 6:31 & 6:58 & 7:15 & 7:29 & 7:45 & 8:02 & 8:27 & 8:44 & 9:01 & 9:28 & 9:59 & 10:32 \\
Beverly/Edgewater Park & 6:18 & 6:35 & 7:02 & 7:19 & 7:33 & 7:49 & 8:06 & 8:31 & 8:48 & 9:05 & 9:32 & 10:03 & 10:36 \\
Burlington South & 6:22 & 6:39 & 7:06 & 7:23 & 7:37 & 7:53 & 8:10 & 8:35 & 8:52 & 9:09 & 9:36 & 10:07 & 10:40 \\
Burlington Towne Ctr & 6:24 & 6:41 & 7:08 & 7:25 & 7:39 & 7:55 & 8:12 & 8:37 & 8:54 & 9:11 & 9:38 & 10:09 & 10:42 \\
Florence & 6:30 & 6:47 & 7:14 & 7:31 & 7:45 & 8:01 & 8:18 & 8:43 & 9:00 & 9:17 & 9:44 & 10:15 & 10:48 \\
Roebling & 6:33 & 6:50 & 7:17 & 7:34 & 7:48 & 8:04 & 8:21 & 8:46 & 9:03 & 9:20 & 9:47 & 10:18 & 10:51 \\
Bordentown & 6:38 & 6:55 & 7:22 & 7:39 & 7:53 & 8:09 & 8:26 & 8:51 & 9:08 & 9:25 & 9:52 & 10:23 & 10:56 \\
Cass Street & 6:44 & 7:01 & 7:28 & 7:45 & 7:59 & 8:15 & 8:32 & 8:58 & 9:14 & 9:31 & 9:58 & 10:29 & 11:02 \\
Hamilton Avenue & 6:46 & 7:03 & 7:30 & 7:47 & 8:01 & 8:17 & 8:34 & 9:00 & 9:16 & 9:33 & 10:00 & 10:31 & 11:04 \\
Trenton TC & 6:49 & 7:06 & 7:33 & 7:50 & 8:04 & 8:20 & 8:37 & 9:03 & 9:19 & 9:36 & 10:03 & 10:34 & 11:07 \\
\bottomrule\end{tabular}\end{center}
\subsection*{To Trenton (read down) --- trips 14--26 of 40}
\begin{center}\footnotesize\setlength{\tabcolsep}{3pt}\begin{tabular}{@{}lrrrrrrrrrrrrr@{}}\toprule
Station & 14 & 15 & 16 & 17 & 18 & 19 & 20 & 21 & 22 & 23 & 24 & 25 & 26 \\\midrule
Entertainment Center & 10:25 & 10:59 & 11:24 & 11:57 & 12:30 & 12:55 & 13:29 & 13:54 & 14:27 & 14:43 & 15:00 & 15:25 & 15:42 \\
Aquarium & 10:27 & 11:01 & 11:26 & 11:59 & 12:32 & 12:57 & 13:31 & 13:56 & 14:29 & 14:45 & 15:02 & 15:27 & 15:44 \\
Cooper St/Rutgers & 10:29 & 11:03 & 11:28 & 12:01 & 12:34 & 12:59 & 13:33 & 13:58 & 14:31 & 14:47 & 15:04 & 15:29 & 15:46 \\
Walter Rand TC & 10:33 & 11:07 & 11:32 & 12:05 & 12:38 & 13:03 & 13:37 & 14:02 & 14:35 & 14:51 & 15:08 & 15:33 & 15:50 \\
36th Street & 10:38 & 11:12 & 11:37 & 12:10 & 12:43 & 13:08 & 13:42 & 14:07 & 14:40 & 14:56 & 15:13 & 15:38 & 15:55 \\
Pennsauken TC & 10:40 & 11:14 & 11:39 & 12:12 & 12:45 & 13:10 & 13:44 & 14:09 & 14:42 & 14:58 & 15:15 & 15:40 & 15:57 \\
Route 73/Pennsauken & 10:43 & 11:17 & 11:44 & 12:15 & 12:48 & 13:13 & 13:47 & 14:14 & 14:45 & 15:01 & 15:18 & 15:43 & 16:00 \\
Palmyra & 10:46 & 11:20 & 11:47 & 12:18 & 12:51 & 13:16 & 13:50 & 14:17 & 14:48 & 15:04 & 15:21 & 15:46 & 16:03 \\
Riverton & 10:48 & 11:22 & 11:49 & 12:20 & 12:53 & 13:18 & 13:52 & 14:19 & 14:50 & 15:06 & 15:23 & 15:48 & 16:05 \\
Cinnaminson & 10:50 & 11:24 & 11:51 & 12:22 & 12:55 & 13:20 & 13:54 & 14:21 & 14:52 & 15:08 & 15:25 & 15:50 & 16:07 \\
Riverside & 10:55 & 11:29 & 11:56 & 12:27 & 13:00 & 13:25 & 13:59 & 14:26 & 14:57 & 15:13 & 15:30 & 15:55 & 16:12 \\
Delanco & 10:57 & 11:31 & 11:58 & 12:29 & 13:02 & 13:27 & 14:01 & 14:28 & 14:59 & 15:15 & 15:32 & 15:57 & 16:14 \\
Beverly/Edgewater Park & 11:01 & 11:35 & 12:02 & 12:33 & 13:06 & 13:31 & 14:05 & 14:32 & 15:03 & 15:19 & 15:36 & 16:01 & 16:18 \\
Burlington South & 11:05 & 11:39 & 12:06 & 12:37 & 13:10 & 13:35 & 14:09 & 14:36 & 15:07 & 15:23 & 15:40 & 16:05 & 16:22 \\
Burlington Towne Ctr & 11:07 & 11:41 & 12:08 & 12:39 & 13:12 & 13:37 & 14:11 & 14:38 & 15:09 & 15:25 & 15:42 & 16:07 & 16:24 \\
Florence & 11:13 & 11:47 & 12:14 & 12:45 & 13:18 & 13:43 & 14:17 & 14:44 & 15:15 & 15:31 & 15:48 & 16:13 & 16:30 \\
Roebling & 11:16 & 11:50 & 12:17 & 12:48 & 13:21 & 13:46 & 14:20 & 14:47 & 15:18 & 15:34 & 15:51 & 16:16 & 16:33 \\
Bordentown & 11:21 & 11:55 & 12:22 & 12:53 & 13:26 & 13:51 & 14:25 & 14:52 & 15:23 & 15:39 & 15:56 & 16:21 & 16:38 \\
Cass Street & 11:28 & 12:01 & 12:28 & 12:59 & 13:32 & 13:58 & 14:31 & 14:58 & 15:29 & 15:45 & 16:02 & 16:28 & 16:44 \\
Hamilton Avenue & 11:30 & 12:03 & 12:30 & 13:01 & 13:34 & 14:00 & 14:33 & 15:00 & 15:31 & 15:47 & 16:04 & 16:30 & 16:46 \\
Trenton TC & 11:33 & 12:06 & 12:33 & 13:04 & 13:37 & 14:03 & 14:36 & 15:03 & 15:34 & 15:50 & 16:07 & 16:33 & 16:49 \\
\bottomrule\end{tabular}\end{center}
\subsection*{To Trenton (read down) --- trips 27--39 of 40}
\begin{center}\footnotesize\setlength{\tabcolsep}{3pt}\begin{tabular}{@{}lrrrrrrrrrrrrr@{}}\toprule
Station & 27 & 28 & 29 & 30 & 31 & 32 & 33 & 34 & 35 & 36 & 37 & 38 & 39 \\\midrule
Entertainment Center & 15:59 & 16:24 & 16:41 & 16:57 & 17:13 & 17:30 & 17:55 & 18:12 & 18:29 & 18:54 & 19:27 & 20:00 & 20:25 \\
Aquarium & 16:01 & 16:26 & 16:43 & 16:59 & 17:15 & 17:32 & 17:57 & 18:14 & 18:31 & 18:56 & 19:29 & 20:02 & 20:27 \\
Cooper St/Rutgers & 16:03 & 16:28 & 16:45 & 17:01 & 17:17 & 17:34 & 17:59 & 18:16 & 18:33 & 18:58 & 19:31 & 20:04 & 20:29 \\
Walter Rand TC & 16:07 & 16:32 & 16:49 & 17:05 & 17:21 & 17:38 & 18:03 & 18:20 & 18:37 & 19:02 & 19:35 & 20:08 & 20:33 \\
36th Street & 16:12 & 16:37 & 16:56 & 17:10 & 17:26 & 17:43 & 18:08 & 18:25 & 18:42 & 19:07 & 19:40 & 20:13 & 20:38 \\
Pennsauken TC & 16:14 & 16:39 & 16:58 & 17:12 & 17:28 & 17:45 & 18:10 & 18:27 & 18:44 & 19:09 & 19:42 & 20:15 & 20:40 \\
Route 73/Pennsauken & 16:17 & 16:44 & 17:01 & 17:15 & 17:31 & 17:48 & 18:13 & 18:30 & 18:47 & 19:14 & 19:45 & 20:18 & 20:43 \\
Palmyra & 16:20 & 16:47 & 17:04 & 17:18 & 17:34 & 17:51 & 18:16 & 18:33 & 18:50 & 19:17 & 19:48 & 20:21 & 20:46 \\
Riverton & 16:22 & 16:49 & 17:06 & 17:20 & 17:36 & 17:53 & 18:18 & 18:35 & 18:52 & 19:19 & 19:50 & 20:23 & 20:48 \\
Cinnaminson & 16:24 & 16:51 & 17:08 & 17:22 & 17:38 & 17:55 & 18:20 & 18:37 & 18:54 & 19:21 & 19:52 & 20:25 & 20:50 \\
Riverside & 16:29 & 16:56 & 17:13 & 17:27 & 17:43 & 18:00 & 18:25 & 18:42 & 18:59 & 19:26 & 19:57 & 20:30 & 20:55 \\
Delanco & 16:31 & 16:58 & 17:15 & 17:29 & 17:45 & 18:02 & 18:27 & 18:44 & 19:01 & 19:28 & 19:59 & 20:32 & 20:57 \\
Beverly/Edgewater Park & 16:35 & 17:02 & 17:19 & 17:33 & 17:49 & 18:06 & 18:31 & 18:48 & 19:05 & 19:32 & 20:03 & 20:36 & 21:01 \\
Burlington South & 16:39 & 17:06 & 17:23 & 17:37 & 17:53 & 18:10 & 18:35 & 18:52 & 19:09 & 19:36 & 20:07 & 20:40 & 21:05 \\
Burlington Towne Ctr & 16:41 & 17:08 & 17:25 & 17:39 & 17:55 & 18:12 & 18:37 & 18:54 & 19:11 & 19:38 & 20:09 & 20:42 & 21:07 \\
Florence & 16:47 & 17:14 & 17:31 & 17:45 & 18:01 & 18:18 & 18:43 & 19:00 & 19:17 & 19:44 & 20:15 & 20:48 & 21:13 \\
Roebling & 16:50 & 17:17 & 17:34 & 17:48 & 18:04 & 18:21 & 18:46 & 19:03 & 19:20 & 19:47 & 20:18 & 20:51 & 21:16 \\
Bordentown & 16:55 & 17:22 & 17:39 & 17:53 & 18:09 & 18:26 & 18:51 & 19:08 & 19:25 & 19:52 & 20:23 & 20:56 & 21:21 \\
Cass Street & 17:01 & 17:28 & 17:45 & 17:59 & 18:15 & 18:32 & 18:58 & 19:14 & 19:31 & 19:58 & 20:29 & 21:02 & 21:28 \\
Hamilton Avenue & 17:03 & 17:30 & 17:47 & 18:01 & 18:17 & 18:34 & 19:00 & 19:16 & 19:33 & 20:00 & 20:31 & 21:04 & 21:30 \\
Trenton TC & 17:06 & 17:33 & 17:50 & 18:04 & 18:20 & 18:37 & 19:03 & 19:19 & 19:36 & 20:03 & 20:34 & 21:07 & 21:33 \\
\bottomrule\end{tabular}\end{center}
\subsection*{To Trenton (read down) --- trips 40--40 of 40}
\begin{center}\footnotesize\setlength{\tabcolsep}{3pt}\begin{tabular}{@{}lr@{}}\toprule
Station & 40 \\\midrule
Entertainment Center & 20:59 \\
Aquarium & 21:01 \\
Cooper St/Rutgers & 21:03 \\
Walter Rand TC & 21:07 \\
36th Street & 21:12 \\
Pennsauken TC & 21:14 \\
Route 73/Pennsauken & 21:17 \\
Palmyra & 21:20 \\
Riverton & 21:22 \\
Cinnaminson & 21:24 \\
Riverside & 21:29 \\
Delanco & 21:31 \\
Beverly/Edgewater Park & 21:35 \\
Burlington South & 21:39 \\
Burlington Towne Ctr & 21:41 \\
Florence & 21:47 \\
Roebling & 21:50 \\
Bordentown & 21:55 \\
Cass Street & 22:01 \\
Hamilton Avenue & 22:03 \\
Trenton TC & 22:06 \\
\bottomrule\end{tabular}\end{center}
\subsection*{To Camden (read down) --- trips 1--13 of 40}
\begin{center}\footnotesize\setlength{\tabcolsep}{3pt}\begin{tabular}{@{}lrrrrrrrrrrrrr@{}}\toprule
Station & 1 & 2 & 3 & 4 & 5 & 6 & 7 & 8 & 9 & 10 & 11 & 12 & 13 \\\midrule
Trenton TC & 5:55 & 6:12 & 6:38 & 6:55 & 7:12 & 7:38 & 7:55 & 8:10 & 8:25 & 8:42 & 9:08 & 9:25 & 9:42 \\
Hamilton Avenue & 5:57 & 6:14 & 6:40 & 6:57 & 7:14 & 7:40 & 7:57 & 8:12 & 8:27 & 8:44 & 9:10 & 9:27 & 9:44 \\
Cass Street & 5:59 & 6:16 & 6:42 & 6:59 & 7:16 & 7:42 & 7:59 & 8:14 & 8:29 & 8:46 & 9:12 & 9:29 & 9:46 \\
Bordentown & 6:06 & 6:23 & 6:49 & 7:06 & 7:23 & 7:49 & 8:06 & 8:21 & 8:36 & 8:53 & 9:19 & 9:36 & 9:53 \\
Roebling & 6:11 & 6:28 & 6:55 & 7:11 & 7:28 & 7:54 & 8:11 & 8:26 & 8:41 & 8:58 & 9:25 & 9:41 & 9:58 \\
Florence & 6:14 & 6:31 & 6:58 & 7:14 & 7:31 & 7:57 & 8:14 & 8:29 & 8:44 & 9:01 & 9:28 & 9:44 & 10:01 \\
Burlington Towne Ctr & 6:19 & 6:36 & 7:03 & 7:19 & 7:36 & 8:04 & 8:21 & 8:34 & 8:49 & 9:06 & 9:33 & 9:49 & 10:06 \\
Burlington South & 6:21 & 6:38 & 7:05 & 7:21 & 7:38 & 8:06 & 8:23 & 8:36 & 8:51 & 9:08 & 9:35 & 9:51 & 10:08 \\
Beverly/Edgewater Park & 6:25 & 6:42 & 7:09 & 7:25 & 7:42 & 8:10 & 8:27 & 8:40 & 8:55 & 9:12 & 9:39 & 9:55 & 10:12 \\
Delanco & 6:29 & 6:46 & 7:13 & 7:29 & 7:46 & 8:14 & 8:31 & 8:44 & 8:59 & 9:16 & 9:43 & 9:59 & 10:16 \\
Riverside & 6:31 & 6:48 & 7:15 & 7:31 & 7:48 & 8:16 & 8:33 & 8:46 & 9:01 & 9:18 & 9:45 & 10:01 & 10:18 \\
Cinnaminson & 6:36 & 6:53 & 7:20 & 7:36 & 7:53 & 8:21 & 8:38 & 8:51 & 9:06 & 9:23 & 9:50 & 10:06 & 10:23 \\
Riverton & 6:38 & 6:55 & 7:22 & 7:38 & 7:55 & 8:23 & 8:40 & 8:53 & 9:08 & 9:25 & 9:52 & 10:08 & 10:25 \\
Palmyra & 6:40 & 6:57 & 7:24 & 7:40 & 7:57 & 8:25 & 8:42 & 8:55 & 9:10 & 9:27 & 9:54 & 10:10 & 10:27 \\
Route 73/Pennsauken & 6:43 & 7:00 & 7:27 & 7:43 & 8:00 & 8:28 & 8:45 & 8:58 & 9:13 & 9:30 & 9:57 & 10:13 & 10:30 \\
Pennsauken TC & 6:45 & 7:02 & 7:29 & 7:45 & 8:02 & 8:30 & 8:47 & 9:00 & 9:15 & 9:32 & 9:59 & 10:15 & 10:32 \\
36th Street & 6:49 & 7:06 & 7:33 & 7:49 & 8:06 & 8:34 & 8:51 & 9:04 & 9:19 & 9:36 & 10:03 & 10:19 & 10:36 \\
Walter Rand TC & 6:54 & 7:11 & 7:38 & 7:54 & 8:11 & 8:39 & 8:56 & 9:09 & 9:24 & 9:41 & 10:08 & 10:24 & 10:41 \\
Cooper St/Rutgers & 6:57 & 7:14 & 7:41 & 7:57 & 8:14 & 8:42 & 8:59 & 9:12 & 9:27 & 9:44 & 10:11 & 10:27 & 10:44 \\
Aquarium & 6:59 & 7:16 & 7:43 & 7:59 & 8:16 & 8:44 & 9:01 & 9:14 & 9:29 & 9:46 & 10:13 & 10:29 & 10:46 \\
Entertainment Center & 7:02 & 7:19 & 7:46 & 8:02 & 8:19 & 8:47 & 9:04 & 9:17 & 9:32 & 9:49 & 10:16 & 10:32 & 10:49 \\
\bottomrule\end{tabular}\end{center}
\subsection*{To Camden (read down) --- trips 14--26 of 40}
\begin{center}\footnotesize\setlength{\tabcolsep}{3pt}\begin{tabular}{@{}lrrrrrrrrrrrrr@{}}\toprule
Station & 14 & 15 & 16 & 17 & 18 & 19 & 20 & 21 & 22 & 23 & 24 & 25 & 26 \\\midrule
Trenton TC & 10:08 & 10:40 & 11:12 & 11:38 & 12:12 & 12:38 & 13:10 & 13:42 & 14:08 & 14:42 & 15:08 & 15:40 & 15:55 \\
Hamilton Avenue & 10:10 & 10:42 & 11:14 & 11:40 & 12:14 & 12:40 & 13:12 & 13:44 & 14:10 & 14:44 & 15:10 & 15:42 & 15:57 \\
Cass Street & 10:12 & 10:44 & 11:16 & 11:42 & 12:16 & 12:42 & 13:14 & 13:46 & 14:12 & 14:46 & 15:12 & 15:44 & 15:59 \\
Bordentown & 10:19 & 10:51 & 11:23 & 11:49 & 12:23 & 12:49 & 13:21 & 13:53 & 14:19 & 14:53 & 15:19 & 15:51 & 16:06 \\
Roebling & 10:24 & 10:56 & 11:28 & 11:55 & 12:28 & 12:54 & 13:26 & 13:58 & 14:25 & 14:58 & 15:24 & 15:56 & 16:11 \\
Florence & 10:27 & 10:59 & 11:31 & 11:58 & 12:31 & 12:57 & 13:29 & 14:01 & 14:28 & 15:01 & 15:27 & 15:59 & 16:14 \\
Burlington Towne Ctr & 10:34 & 11:04 & 11:36 & 12:03 & 12:36 & 13:04 & 13:34 & 14:06 & 14:33 & 15:06 & 15:34 & 16:04 & 16:19 \\
Burlington South & 10:36 & 11:06 & 11:38 & 12:05 & 12:38 & 13:06 & 13:36 & 14:08 & 14:35 & 15:08 & 15:36 & 16:06 & 16:21 \\
Beverly/Edgewater Park & 10:40 & 11:10 & 11:42 & 12:09 & 12:42 & 13:10 & 13:40 & 14:12 & 14:39 & 15:12 & 15:40 & 16:10 & 16:25 \\
Delanco & 10:44 & 11:14 & 11:46 & 12:13 & 12:46 & 13:14 & 13:44 & 14:16 & 14:43 & 15:16 & 15:44 & 16:14 & 16:29 \\
Riverside & 10:46 & 11:16 & 11:48 & 12:15 & 12:48 & 13:16 & 13:46 & 14:18 & 14:45 & 15:18 & 15:46 & 16:16 & 16:31 \\
Cinnaminson & 10:51 & 11:21 & 11:53 & 12:20 & 12:53 & 13:21 & 13:51 & 14:23 & 14:50 & 15:23 & 15:51 & 16:21 & 16:36 \\
Riverton & 10:53 & 11:23 & 11:55 & 12:22 & 12:55 & 13:23 & 13:53 & 14:25 & 14:52 & 15:25 & 15:53 & 16:23 & 16:38 \\
Palmyra & 10:55 & 11:25 & 11:57 & 12:24 & 12:57 & 13:25 & 13:55 & 14:27 & 14:54 & 15:27 & 15:55 & 16:25 & 16:40 \\
Route 73/Pennsauken & 10:58 & 11:28 & 12:00 & 12:27 & 13:00 & 13:28 & 13:58 & 14:30 & 14:57 & 15:30 & 15:58 & 16:28 & 16:43 \\
Pennsauken TC & 11:00 & 11:30 & 12:02 & 12:29 & 13:02 & 13:30 & 14:00 & 14:32 & 14:59 & 15:32 & 16:00 & 16:30 & 16:45 \\
36th Street & 11:04 & 11:34 & 12:06 & 12:33 & 13:06 & 13:34 & 14:04 & 14:36 & 15:03 & 15:36 & 16:04 & 16:34 & 16:49 \\
Walter Rand TC & 11:09 & 11:39 & 12:11 & 12:38 & 13:11 & 13:39 & 14:09 & 14:41 & 15:08 & 15:41 & 16:09 & 16:39 & 16:54 \\
Cooper St/Rutgers & 11:12 & 11:42 & 12:14 & 12:41 & 13:14 & 13:42 & 14:12 & 14:44 & 15:11 & 15:44 & 16:12 & 16:42 & 16:57 \\
Aquarium & 11:14 & 11:44 & 12:16 & 12:43 & 13:16 & 13:44 & 14:14 & 14:46 & 15:13 & 15:46 & 16:14 & 16:44 & 16:59 \\
Entertainment Center & 11:17 & 11:47 & 12:19 & 12:46 & 13:19 & 13:47 & 14:17 & 14:49 & 15:16 & 15:49 & 16:17 & 16:47 & 17:02 \\
\bottomrule\end{tabular}\end{center}
\subsection*{To Camden (read down) --- trips 27--39 of 40}
\begin{center}\footnotesize\setlength{\tabcolsep}{3pt}\begin{tabular}{@{}lrrrrrrrrrrrrr@{}}\toprule
Station & 27 & 28 & 29 & 30 & 31 & 32 & 33 & 34 & 35 & 36 & 37 & 38 & 39 \\\midrule
Trenton TC & 16:12 & 16:38 & 16:55 & 17:12 & 17:38 & 17:55 & 18:10 & 18:25 & 18:42 & 19:08 & 19:25 & 19:42 & 20:08 \\
Hamilton Avenue & 16:14 & 16:40 & 16:57 & 17:14 & 17:40 & 17:57 & 18:12 & 18:27 & 18:44 & 19:10 & 19:27 & 19:44 & 20:10 \\
Cass Street & 16:16 & 16:42 & 16:59 & 17:16 & 17:42 & 17:59 & 18:14 & 18:29 & 18:46 & 19:12 & 19:29 & 19:46 & 20:12 \\
Bordentown & 16:23 & 16:49 & 17:06 & 17:23 & 17:49 & 18:06 & 18:21 & 18:36 & 18:53 & 19:19 & 19:36 & 19:53 & 20:19 \\
Roebling & 16:28 & 16:55 & 17:11 & 17:28 & 17:54 & 18:11 & 18:26 & 18:41 & 18:58 & 19:25 & 19:41 & 19:58 & 20:24 \\
Florence & 16:31 & 16:58 & 17:14 & 17:31 & 17:57 & 18:14 & 18:29 & 18:44 & 19:01 & 19:28 & 19:44 & 20:01 & 20:27 \\
Burlington Towne Ctr & 16:36 & 17:03 & 17:19 & 17:36 & 18:04 & 18:21 & 18:34 & 18:49 & 19:06 & 19:33 & 19:49 & 20:06 & 20:34 \\
Burlington South & 16:38 & 17:05 & 17:21 & 17:38 & 18:06 & 18:23 & 18:36 & 18:51 & 19:08 & 19:35 & 19:51 & 20:08 & 20:36 \\
Beverly/Edgewater Park & 16:42 & 17:09 & 17:25 & 17:42 & 18:10 & 18:27 & 18:40 & 18:55 & 19:12 & 19:39 & 19:55 & 20:12 & 20:40 \\
Delanco & 16:46 & 17:13 & 17:29 & 17:46 & 18:14 & 18:31 & 18:44 & 18:59 & 19:16 & 19:43 & 19:59 & 20:16 & 20:44 \\
Riverside & 16:48 & 17:15 & 17:31 & 17:48 & 18:16 & 18:33 & 18:46 & 19:01 & 19:18 & 19:45 & 20:01 & 20:18 & 20:46 \\
Cinnaminson & 16:53 & 17:20 & 17:36 & 17:53 & 18:21 & 18:38 & 18:51 & 19:06 & 19:23 & 19:50 & 20:06 & 20:23 & 20:51 \\
Riverton & 16:55 & 17:22 & 17:38 & 17:55 & 18:23 & 18:40 & 18:53 & 19:08 & 19:25 & 19:52 & 20:08 & 20:25 & 20:53 \\
Palmyra & 16:57 & 17:24 & 17:40 & 17:57 & 18:25 & 18:42 & 18:55 & 19:10 & 19:27 & 19:54 & 20:10 & 20:27 & 20:55 \\
Route 73/Pennsauken & 17:00 & 17:27 & 17:43 & 18:00 & 18:28 & 18:45 & 18:58 & 19:13 & 19:30 & 19:57 & 20:13 & 20:30 & 20:58 \\
Pennsauken TC & 17:02 & 17:29 & 17:45 & 18:02 & 18:30 & 18:47 & 19:00 & 19:15 & 19:32 & 19:59 & 20:15 & 20:32 & 21:00 \\
36th Street & 17:06 & 17:33 & 17:49 & 18:06 & 18:34 & 18:51 & 19:04 & 19:19 & 19:36 & 20:03 & 20:19 & 20:36 & 21:04 \\
Walter Rand TC & 17:11 & 17:38 & 17:54 & 18:11 & 18:39 & 18:56 & 19:09 & 19:24 & 19:41 & 20:08 & 20:24 & 20:41 & 21:09 \\
Cooper St/Rutgers & 17:14 & 17:41 & 17:57 & 18:14 & 18:42 & 18:59 & 19:12 & 19:27 & 19:44 & 20:11 & 20:27 & 20:44 & 21:12 \\
Aquarium & 17:16 & 17:43 & 17:59 & 18:16 & 18:44 & 19:01 & 19:14 & 19:29 & 19:46 & 20:13 & 20:29 & 20:46 & 21:14 \\
Entertainment Center & 17:19 & 17:46 & 18:02 & 18:19 & 18:47 & 19:04 & 19:17 & 19:32 & 19:49 & 20:16 & 20:32 & 20:49 & 21:17 \\
\bottomrule\end{tabular}\end{center}
\subsection*{To Camden (read down) --- trips 40--40 of 40}
\begin{center}\footnotesize\setlength{\tabcolsep}{3pt}\begin{tabular}{@{}lr@{}}\toprule
Station & 40 \\\midrule
Trenton TC & 20:40 \\
Hamilton Avenue & 20:42 \\
Cass Street & 20:44 \\
Bordentown & 20:51 \\
Roebling & 20:56 \\
Florence & 20:59 \\
Burlington Towne Ctr & 21:04 \\
Burlington South & 21:06 \\
Beverly/Edgewater Park & 21:10 \\
Delanco & 21:14 \\
Riverside & 21:16 \\
Cinnaminson & 21:21 \\
Riverton & 21:23 \\
Palmyra & 21:25 \\
Route 73/Pennsauken & 21:28 \\
Pennsauken TC & 21:30 \\
36th Street & 21:34 \\
Walter Rand TC & 21:39 \\
Cooper St/Rutgers & 21:42 \\
Aquarium & 21:44 \\
Entertainment Center & 21:47 \\
\bottomrule\end{tabular}\end{center}
\clearpage
\section{Proposal E6: 8 peak, 6 midday, even midday}\label{tt:proposal_E6}
\subsection*{To Trenton (read down) --- trips 1--13 of 43}
\begin{center}\footnotesize\setlength{\tabcolsep}{3pt}\begin{tabular}{@{}lrrrrrrrrrrrrr@{}}\toprule
Station & 1 & 2 & 3 & 4 & 5 & 6 & 7 & 8 & 9 & 10 & 11 & 12 & 13 \\\midrule
Entertainment Center & 5:46 & 6:07 & 6:30 & 6:45 & 7:00 & 7:15 & 7:30 & 7:53 & 8:16 & 8:37 & 9:00 & 9:30 & 10:00 \\
Aquarium & 5:48 & 6:09 & 6:32 & 6:47 & 7:02 & 7:17 & 7:32 & 7:55 & 8:18 & 8:39 & 9:02 & 9:32 & 10:02 \\
Cooper St/Rutgers & 5:50 & 6:11 & 6:34 & 6:49 & 7:04 & 7:19 & 7:34 & 7:57 & 8:20 & 8:41 & 9:04 & 9:34 & 10:04 \\
Walter Rand TC & 5:54 & 6:15 & 6:38 & 6:53 & 7:08 & 7:23 & 7:38 & 8:01 & 8:24 & 8:45 & 9:08 & 9:38 & 10:08 \\
36th Street & 5:59 & 6:20 & 6:43 & 6:58 & 7:13 & 7:28 & 7:43 & 8:06 & 8:29 & 8:50 & 9:13 & 9:43 & 10:13 \\
Pennsauken TC & 6:01 & 6:22 & 6:45 & 7:00 & 7:15 & 7:30 & 7:45 & 8:08 & 8:31 & 8:52 & 9:15 & 9:45 & 10:15 \\
Route 73/Pennsauken & 6:04 & 6:25 & 6:48 & 7:03 & 7:18 & 7:33 & 7:48 & 8:11 & 8:34 & 8:55 & 9:18 & 9:48 & 10:18 \\
Palmyra & 6:07 & 6:28 & 6:51 & 7:06 & 7:21 & 7:36 & 7:51 & 8:14 & 8:37 & 8:58 & 9:21 & 9:51 & 10:21 \\
Riverton & 6:09 & 6:30 & 6:53 & 7:08 & 7:23 & 7:38 & 7:53 & 8:16 & 8:39 & 9:00 & 9:23 & 9:53 & 10:23 \\
Cinnaminson & 6:11 & 6:32 & 6:55 & 7:10 & 7:25 & 7:40 & 7:55 & 8:18 & 8:41 & 9:02 & 9:25 & 9:55 & 10:25 \\
Riverside & 6:16 & 6:37 & 7:00 & 7:16 & 7:30 & 7:45 & 8:00 & 8:23 & 8:46 & 9:07 & 9:30 & 10:00 & 10:30 \\
Delanco & 6:18 & 6:39 & 7:02 & 7:18 & 7:32 & 7:47 & 8:02 & 8:25 & 8:48 & 9:09 & 9:32 & 10:02 & 10:32 \\
Beverly/Edgewater Park & 6:22 & 6:43 & 7:06 & 7:22 & 7:36 & 7:51 & 8:06 & 8:29 & 8:52 & 9:13 & 9:36 & 10:06 & 10:36 \\
Burlington South & 6:26 & 6:47 & 7:12 & 7:26 & 7:40 & 7:55 & 8:10 & 8:35 & 8:56 & 9:17 & 9:42 & 10:10 & 10:40 \\
Burlington Towne Ctr & 6:28 & 6:49 & 7:14 & 7:28 & 7:42 & 7:57 & 8:12 & 8:37 & 8:58 & 9:19 & 9:44 & 10:12 & 10:42 \\
Florence & 6:34 & 6:56 & 7:20 & 7:37 & 7:48 & 8:03 & 8:18 & 8:43 & 9:04 & 9:26 & 9:50 & 10:18 & 10:48 \\
Roebling & 6:37 & 6:59 & 7:25 & 7:40 & 7:51 & 8:06 & 8:21 & 8:47 & 9:07 & 9:29 & 9:55 & 10:21 & 10:51 \\
Bordentown & 6:42 & 7:04 & 7:30 & 7:45 & 7:56 & 8:11 & 8:26 & 8:52 & 9:12 & 9:34 & 10:00 & 10:26 & 10:56 \\
Cass Street & 6:48 & 7:10 & 7:36 & 7:51 & 8:02 & 8:17 & 8:32 & 8:58 & 9:18 & 9:40 & 10:06 & 10:32 & 11:02 \\
Hamilton Avenue & 6:50 & 7:12 & 7:38 & 7:53 & 8:04 & 8:19 & 8:34 & 9:00 & 9:20 & 9:42 & 10:08 & 10:34 & 11:04 \\
Trenton TC & 6:53 & 7:15 & 7:41 & 7:56 & 8:07 & 8:22 & 8:37 & 9:03 & 9:23 & 9:45 & 10:11 & 10:37 & 11:07 \\
\bottomrule\end{tabular}\end{center}
\subsection*{To Trenton (read down) --- trips 14--26 of 43}
\begin{center}\footnotesize\setlength{\tabcolsep}{3pt}\begin{tabular}{@{}lrrrrrrrrrrrrr@{}}\toprule
Station & 14 & 15 & 16 & 17 & 18 & 19 & 20 & 21 & 22 & 23 & 24 & 25 & 26 \\\midrule
Entertainment Center & 10:23 & 10:46 & 11:07 & 11:30 & 12:00 & 12:30 & 12:53 & 13:16 & 13:37 & 14:00 & 14:15 & 14:30 & 14:45 \\
Aquarium & 10:25 & 10:48 & 11:09 & 11:32 & 12:02 & 12:32 & 12:55 & 13:18 & 13:39 & 14:02 & 14:17 & 14:32 & 14:47 \\
Cooper St/Rutgers & 10:27 & 10:50 & 11:11 & 11:34 & 12:04 & 12:34 & 12:57 & 13:20 & 13:41 & 14:04 & 14:19 & 14:34 & 14:49 \\
Walter Rand TC & 10:31 & 10:54 & 11:15 & 11:38 & 12:08 & 12:38 & 13:01 & 13:24 & 13:45 & 14:08 & 14:23 & 14:38 & 14:53 \\
36th Street & 10:36 & 10:59 & 11:20 & 11:43 & 12:13 & 12:43 & 13:06 & 13:29 & 13:50 & 14:13 & 14:28 & 14:43 & 14:58 \\
Pennsauken TC & 10:38 & 11:01 & 11:22 & 11:45 & 12:15 & 12:45 & 13:08 & 13:31 & 13:52 & 14:15 & 14:30 & 14:45 & 15:00 \\
Route 73/Pennsauken & 10:41 & 11:04 & 11:25 & 11:48 & 12:18 & 12:48 & 13:11 & 13:34 & 13:55 & 14:18 & 14:33 & 14:48 & 15:03 \\
Palmyra & 10:44 & 11:07 & 11:28 & 11:51 & 12:21 & 12:51 & 13:14 & 13:37 & 13:58 & 14:21 & 14:36 & 14:51 & 15:06 \\
Riverton & 10:46 & 11:09 & 11:30 & 11:53 & 12:23 & 12:53 & 13:16 & 13:39 & 14:00 & 14:23 & 14:38 & 14:53 & 15:08 \\
Cinnaminson & 10:48 & 11:11 & 11:32 & 11:55 & 12:25 & 12:55 & 13:18 & 13:41 & 14:02 & 14:25 & 14:40 & 14:55 & 15:10 \\
Riverside & 10:53 & 11:16 & 11:37 & 12:00 & 12:30 & 13:00 & 13:23 & 13:46 & 14:07 & 14:30 & 14:46 & 15:00 & 15:15 \\
Delanco & 10:55 & 11:18 & 11:39 & 12:02 & 12:32 & 13:02 & 13:25 & 13:48 & 14:09 & 14:32 & 14:48 & 15:02 & 15:17 \\
Beverly/Edgewater Park & 10:59 & 11:22 & 11:43 & 12:06 & 12:36 & 13:06 & 13:29 & 13:52 & 14:13 & 14:36 & 14:52 & 15:06 & 15:21 \\
Burlington South & 11:05 & 11:26 & 11:47 & 12:12 & 12:40 & 13:10 & 13:35 & 13:56 & 14:17 & 14:42 & 14:56 & 15:10 & 15:25 \\
Burlington Towne Ctr & 11:07 & 11:28 & 11:49 & 12:14 & 12:42 & 13:12 & 13:37 & 13:58 & 14:19 & 14:44 & 14:58 & 15:12 & 15:27 \\
Florence & 11:13 & 11:34 & 11:56 & 12:20 & 12:48 & 13:18 & 13:43 & 14:04 & 14:26 & 14:50 & 15:07 & 15:18 & 15:33 \\
Roebling & 11:17 & 11:37 & 11:59 & 12:25 & 12:51 & 13:21 & 13:47 & 14:07 & 14:29 & 14:55 & 15:10 & 15:21 & 15:36 \\
Bordentown & 11:22 & 11:42 & 12:04 & 12:30 & 12:56 & 13:26 & 13:52 & 14:12 & 14:34 & 15:00 & 15:15 & 15:26 & 15:41 \\
Cass Street & 11:28 & 11:48 & 12:10 & 12:36 & 13:02 & 13:32 & 13:58 & 14:18 & 14:40 & 15:06 & 15:21 & 15:32 & 15:47 \\
Hamilton Avenue & 11:30 & 11:50 & 12:12 & 12:38 & 13:04 & 13:34 & 14:00 & 14:20 & 14:42 & 15:08 & 15:23 & 15:34 & 15:49 \\
Trenton TC & 11:33 & 11:53 & 12:15 & 12:41 & 13:07 & 13:37 & 14:03 & 14:23 & 14:45 & 15:11 & 15:26 & 15:37 & 15:52 \\
\bottomrule\end{tabular}\end{center}
\subsection*{To Trenton (read down) --- trips 27--39 of 43}
\begin{center}\footnotesize\setlength{\tabcolsep}{3pt}\begin{tabular}{@{}lrrrrrrrrrrrrr@{}}\toprule
Station & 27 & 28 & 29 & 30 & 31 & 32 & 33 & 34 & 35 & 36 & 37 & 38 & 39 \\\midrule
Entertainment Center & 15:00 & 15:23 & 15:46 & 16:07 & 16:30 & 16:45 & 17:00 & 17:15 & 17:30 & 17:53 & 18:16 & 18:37 & 19:00 \\
Aquarium & 15:02 & 15:25 & 15:48 & 16:09 & 16:32 & 16:47 & 17:02 & 17:17 & 17:32 & 17:55 & 18:18 & 18:39 & 19:02 \\
Cooper St/Rutgers & 15:04 & 15:27 & 15:50 & 16:11 & 16:34 & 16:49 & 17:04 & 17:19 & 17:34 & 17:57 & 18:20 & 18:41 & 19:04 \\
Walter Rand TC & 15:08 & 15:31 & 15:54 & 16:15 & 16:38 & 16:53 & 17:08 & 17:23 & 17:38 & 18:01 & 18:24 & 18:45 & 19:08 \\
36th Street & 15:13 & 15:36 & 15:59 & 16:20 & 16:43 & 16:58 & 17:13 & 17:28 & 17:43 & 18:06 & 18:29 & 18:50 & 19:13 \\
Pennsauken TC & 15:15 & 15:38 & 16:01 & 16:22 & 16:45 & 17:00 & 17:15 & 17:30 & 17:45 & 18:08 & 18:31 & 18:52 & 19:15 \\
Route 73/Pennsauken & 15:18 & 15:41 & 16:04 & 16:25 & 16:48 & 17:03 & 17:18 & 17:33 & 17:48 & 18:11 & 18:34 & 18:55 & 19:18 \\
Palmyra & 15:21 & 15:44 & 16:07 & 16:28 & 16:51 & 17:06 & 17:21 & 17:36 & 17:51 & 18:14 & 18:37 & 18:58 & 19:21 \\
Riverton & 15:23 & 15:46 & 16:09 & 16:30 & 16:53 & 17:08 & 17:23 & 17:38 & 17:53 & 18:16 & 18:39 & 19:00 & 19:23 \\
Cinnaminson & 15:25 & 15:48 & 16:11 & 16:32 & 16:55 & 17:10 & 17:25 & 17:40 & 17:55 & 18:18 & 18:41 & 19:02 & 19:25 \\
Riverside & 15:30 & 15:53 & 16:16 & 16:37 & 17:00 & 17:16 & 17:30 & 17:45 & 18:00 & 18:23 & 18:46 & 19:07 & 19:30 \\
Delanco & 15:32 & 15:55 & 16:18 & 16:39 & 17:02 & 17:18 & 17:32 & 17:47 & 18:02 & 18:25 & 18:48 & 19:09 & 19:32 \\
Beverly/Edgewater Park & 15:36 & 15:59 & 16:22 & 16:43 & 17:06 & 17:22 & 17:36 & 17:51 & 18:06 & 18:29 & 18:52 & 19:13 & 19:36 \\
Burlington South & 15:40 & 16:05 & 16:26 & 16:47 & 17:12 & 17:26 & 17:40 & 17:55 & 18:10 & 18:35 & 18:56 & 19:17 & 19:42 \\
Burlington Towne Ctr & 15:42 & 16:07 & 16:28 & 16:49 & 17:14 & 17:28 & 17:42 & 17:57 & 18:12 & 18:37 & 18:58 & 19:19 & 19:44 \\
Florence & 15:48 & 16:13 & 16:34 & 16:56 & 17:20 & 17:37 & 17:48 & 18:03 & 18:18 & 18:43 & 19:04 & 19:26 & 19:50 \\
Roebling & 15:51 & 16:17 & 16:37 & 16:59 & 17:25 & 17:40 & 17:51 & 18:06 & 18:21 & 18:47 & 19:07 & 19:29 & 19:55 \\
Bordentown & 15:56 & 16:22 & 16:42 & 17:04 & 17:30 & 17:45 & 17:56 & 18:11 & 18:26 & 18:52 & 19:12 & 19:34 & 20:00 \\
Cass Street & 16:02 & 16:28 & 16:48 & 17:10 & 17:36 & 17:51 & 18:02 & 18:17 & 18:32 & 18:58 & 19:18 & 19:40 & 20:06 \\
Hamilton Avenue & 16:04 & 16:30 & 16:50 & 17:12 & 17:38 & 17:53 & 18:04 & 18:19 & 18:34 & 19:00 & 19:20 & 19:42 & 20:08 \\
Trenton TC & 16:07 & 16:33 & 16:53 & 17:15 & 17:41 & 17:56 & 18:07 & 18:22 & 18:37 & 19:03 & 19:23 & 19:45 & 20:11 \\
\bottomrule\end{tabular}\end{center}
\subsection*{To Trenton (read down) --- trips 40--43 of 43}
\begin{center}\footnotesize\setlength{\tabcolsep}{3pt}\begin{tabular}{@{}lrrrr@{}}\toprule
Station & 40 & 41 & 42 & 43 \\\midrule
Entertainment Center & 19:30 & 20:00 & 20:23 & 20:46 \\
Aquarium & 19:32 & 20:02 & 20:25 & 20:48 \\
Cooper St/Rutgers & 19:34 & 20:04 & 20:27 & 20:50 \\
Walter Rand TC & 19:38 & 20:08 & 20:31 & 20:54 \\
36th Street & 19:43 & 20:13 & 20:36 & 20:59 \\
Pennsauken TC & 19:45 & 20:15 & 20:38 & 21:01 \\
Route 73/Pennsauken & 19:48 & 20:18 & 20:41 & 21:04 \\
Palmyra & 19:51 & 20:21 & 20:44 & 21:07 \\
Riverton & 19:53 & 20:23 & 20:46 & 21:09 \\
Cinnaminson & 19:55 & 20:25 & 20:48 & 21:11 \\
Riverside & 20:00 & 20:30 & 20:53 & 21:16 \\
Delanco & 20:02 & 20:32 & 20:55 & 21:18 \\
Beverly/Edgewater Park & 20:06 & 20:36 & 20:59 & 21:22 \\
Burlington South & 20:10 & 20:40 & 21:05 & 21:26 \\
Burlington Towne Ctr & 20:12 & 20:42 & 21:07 & 21:28 \\
Florence & 20:18 & 20:48 & 21:13 & 21:34 \\
Roebling & 20:21 & 20:51 & 21:17 & 21:37 \\
Bordentown & 20:26 & 20:56 & 21:22 & 21:42 \\
Cass Street & 20:32 & 21:02 & 21:28 & 21:48 \\
Hamilton Avenue & 20:34 & 21:04 & 21:30 & 21:50 \\
Trenton TC & 20:37 & 21:07 & 21:33 & 21:53 \\
\bottomrule\end{tabular}\end{center}
\subsection*{To Camden (read down) --- trips 1--13 of 44}
\begin{center}\footnotesize\setlength{\tabcolsep}{3pt}\begin{tabular}{@{}lrrrrrrrrrrrrr@{}}\toprule
Station & 1 & 2 & 3 & 4 & 5 & 6 & 7 & 8 & 9 & 10 & 11 & 12 & 13 \\\midrule
Trenton TC & 5:44 & 5:57 & 6:12 & 6:38 & 6:59 & 7:20 & 7:46 & 8:01 & 8:14 & 8:27 & 8:42 & 9:08 & 9:29 \\
Hamilton Avenue & 5:46 & 5:59 & 6:14 & 6:40 & 7:01 & 7:22 & 7:48 & 8:03 & 8:16 & 8:29 & 8:44 & 9:10 & 9:31 \\
Cass Street & 5:48 & 6:01 & 6:16 & 6:42 & 7:03 & 7:24 & 7:50 & 8:05 & 8:18 & 8:31 & 8:46 & 9:12 & 9:33 \\
Bordentown & 5:55 & 6:08 & 6:23 & 6:49 & 7:11 & 7:31 & 7:57 & 8:12 & 8:25 & 8:38 & 8:53 & 9:19 & 9:41 \\
Roebling & 6:00 & 6:13 & 6:30 & 6:54 & 7:16 & 7:36 & 8:02 & 8:17 & 8:30 & 8:43 & 9:00 & 9:24 & 9:46 \\
Florence & 6:03 & 6:16 & 6:33 & 6:57 & 7:19 & 7:39 & 8:05 & 8:20 & 8:33 & 8:46 & 9:03 & 9:27 & 9:49 \\
Burlington Towne Ctr & 6:08 & 6:21 & 6:38 & 7:02 & 7:24 & 7:44 & 8:10 & 8:25 & 8:38 & 8:51 & 9:08 & 9:32 & 9:54 \\
Burlington South & 6:10 & 6:23 & 6:40 & 7:04 & 7:26 & 7:46 & 8:12 & 8:27 & 8:40 & 8:53 & 9:10 & 9:34 & 9:56 \\
Beverly/Edgewater Park & 6:17 & 6:29 & 6:44 & 7:08 & 7:30 & 7:50 & 8:16 & 8:31 & 8:47 & 8:59 & 9:14 & 9:38 & 10:00 \\
Delanco & 6:21 & 6:33 & 6:48 & 7:12 & 7:34 & 7:54 & 8:20 & 8:35 & 8:51 & 9:03 & 9:18 & 9:42 & 10:04 \\
Riverside & 6:23 & 6:35 & 6:50 & 7:14 & 7:36 & 7:56 & 8:22 & 8:37 & 8:53 & 9:05 & 9:20 & 9:44 & 10:06 \\
Cinnaminson & 6:28 & 6:40 & 6:55 & 7:19 & 7:41 & 8:03 & 8:27 & 8:42 & 8:58 & 9:10 & 9:25 & 9:49 & 10:11 \\
Riverton & 6:30 & 6:42 & 6:57 & 7:21 & 7:43 & 8:05 & 8:29 & 8:44 & 9:00 & 9:12 & 9:27 & 9:51 & 10:13 \\
Palmyra & 6:32 & 6:44 & 6:59 & 7:25 & 7:45 & 8:07 & 8:31 & 8:46 & 9:02 & 9:14 & 9:29 & 9:55 & 10:15 \\
Route 73/Pennsauken & 6:35 & 6:47 & 7:02 & 7:28 & 7:48 & 8:10 & 8:34 & 8:49 & 9:05 & 9:17 & 9:32 & 9:58 & 10:18 \\
Pennsauken TC & 6:37 & 6:49 & 7:04 & 7:30 & 7:50 & 8:12 & 8:36 & 8:51 & 9:07 & 9:19 & 9:34 & 10:00 & 10:20 \\
36th Street & 6:41 & 6:53 & 7:08 & 7:34 & 7:54 & 8:16 & 8:40 & 8:55 & 9:11 & 9:23 & 9:38 & 10:04 & 10:24 \\
Walter Rand TC & 6:46 & 6:58 & 7:13 & 7:39 & 7:59 & 8:21 & 8:45 & 9:00 & 9:16 & 9:28 & 9:43 & 10:09 & 10:29 \\
Cooper St/Rutgers & 6:49 & 7:01 & 7:16 & 7:42 & 8:02 & 8:24 & 8:48 & 9:03 & 9:19 & 9:31 & 9:46 & 10:12 & 10:32 \\
Aquarium & 6:51 & 7:03 & 7:18 & 7:44 & 8:04 & 8:26 & 8:50 & 9:05 & 9:21 & 9:33 & 9:48 & 10:14 & 10:34 \\
Entertainment Center & 6:54 & 7:06 & 7:21 & 7:47 & 8:07 & 8:29 & 8:53 & 9:08 & 9:24 & 9:36 & 9:51 & 10:17 & 10:37 \\
\bottomrule\end{tabular}\end{center}
\subsection*{To Camden (read down) --- trips 14--26 of 44}
\begin{center}\footnotesize\setlength{\tabcolsep}{3pt}\begin{tabular}{@{}lrrrrrrrrrrrrr@{}}\toprule
Station & 14 & 15 & 16 & 17 & 18 & 19 & 20 & 21 & 22 & 23 & 24 & 25 & 26 \\\midrule
Trenton TC & 9:50 & 10:16 & 10:44 & 11:12 & 11:38 & 11:59 & 12:20 & 12:46 & 13:14 & 13:42 & 14:08 & 14:29 & 14:50 \\
Hamilton Avenue & 9:52 & 10:18 & 10:46 & 11:14 & 11:40 & 12:01 & 12:22 & 12:48 & 13:16 & 13:44 & 14:10 & 14:31 & 14:52 \\
Cass Street & 9:54 & 10:20 & 10:48 & 11:16 & 11:42 & 12:03 & 12:24 & 12:50 & 13:18 & 13:46 & 14:12 & 14:33 & 14:54 \\
Bordentown & 10:01 & 10:27 & 10:55 & 11:23 & 11:49 & 12:11 & 12:31 & 12:57 & 13:25 & 13:53 & 14:19 & 14:41 & 15:01 \\
Roebling & 10:06 & 10:32 & 11:00 & 11:30 & 11:54 & 12:16 & 12:36 & 13:02 & 13:30 & 14:00 & 14:24 & 14:46 & 15:06 \\
Florence & 10:09 & 10:35 & 11:03 & 11:33 & 11:57 & 12:19 & 12:39 & 13:05 & 13:33 & 14:03 & 14:27 & 14:49 & 15:09 \\
Burlington Towne Ctr & 10:14 & 10:40 & 11:08 & 11:38 & 12:02 & 12:24 & 12:44 & 13:10 & 13:38 & 14:08 & 14:32 & 14:54 & 15:14 \\
Burlington South & 10:16 & 10:42 & 11:10 & 11:40 & 12:04 & 12:26 & 12:46 & 13:12 & 13:40 & 14:10 & 14:34 & 14:56 & 15:16 \\
Beverly/Edgewater Park & 10:20 & 10:46 & 11:17 & 11:44 & 12:08 & 12:30 & 12:50 & 13:16 & 13:47 & 14:14 & 14:38 & 15:00 & 15:20 \\
Delanco & 10:24 & 10:50 & 11:21 & 11:48 & 12:12 & 12:34 & 12:54 & 13:20 & 13:51 & 14:18 & 14:42 & 15:04 & 15:24 \\
Riverside & 10:26 & 10:52 & 11:23 & 11:50 & 12:14 & 12:36 & 12:56 & 13:22 & 13:53 & 14:20 & 14:44 & 15:06 & 15:26 \\
Cinnaminson & 10:33 & 10:57 & 11:28 & 11:55 & 12:19 & 12:41 & 13:03 & 13:27 & 13:58 & 14:25 & 14:49 & 15:11 & 15:33 \\
Riverton & 10:35 & 10:59 & 11:30 & 11:57 & 12:21 & 12:43 & 13:05 & 13:29 & 14:00 & 14:27 & 14:51 & 15:13 & 15:35 \\
Palmyra & 10:37 & 11:01 & 11:32 & 11:59 & 12:25 & 12:45 & 13:07 & 13:31 & 14:02 & 14:29 & 14:55 & 15:15 & 15:37 \\
Route 73/Pennsauken & 10:40 & 11:04 & 11:35 & 12:02 & 12:28 & 12:48 & 13:10 & 13:34 & 14:05 & 14:32 & 14:58 & 15:18 & 15:40 \\
Pennsauken TC & 10:42 & 11:06 & 11:37 & 12:04 & 12:30 & 12:50 & 13:12 & 13:36 & 14:07 & 14:34 & 15:00 & 15:20 & 15:42 \\
36th Street & 10:46 & 11:10 & 11:41 & 12:08 & 12:34 & 12:54 & 13:16 & 13:40 & 14:11 & 14:38 & 15:04 & 15:24 & 15:46 \\
Walter Rand TC & 10:51 & 11:15 & 11:46 & 12:13 & 12:39 & 12:59 & 13:21 & 13:45 & 14:16 & 14:43 & 15:09 & 15:29 & 15:51 \\
Cooper St/Rutgers & 10:54 & 11:18 & 11:49 & 12:16 & 12:42 & 13:02 & 13:24 & 13:48 & 14:19 & 14:46 & 15:12 & 15:32 & 15:54 \\
Aquarium & 10:56 & 11:20 & 11:51 & 12:18 & 12:44 & 13:04 & 13:26 & 13:50 & 14:21 & 14:48 & 15:14 & 15:34 & 15:56 \\
Entertainment Center & 10:59 & 11:23 & 11:54 & 12:21 & 12:47 & 13:07 & 13:29 & 13:53 & 14:24 & 14:51 & 15:17 & 15:37 & 15:59 \\
\bottomrule\end{tabular}\end{center}
\subsection*{To Camden (read down) --- trips 27--39 of 44}
\begin{center}\footnotesize\setlength{\tabcolsep}{3pt}\begin{tabular}{@{}lrrrrrrrrrrrrr@{}}\toprule
Station & 27 & 28 & 29 & 30 & 31 & 32 & 33 & 34 & 35 & 36 & 37 & 38 & 39 \\\midrule
Trenton TC & 15:16 & 15:31 & 15:44 & 15:57 & 16:12 & 16:38 & 16:59 & 17:20 & 17:46 & 18:01 & 18:14 & 18:27 & 18:42 \\
Hamilton Avenue & 15:18 & 15:33 & 15:46 & 15:59 & 16:14 & 16:40 & 17:01 & 17:22 & 17:48 & 18:03 & 18:16 & 18:29 & 18:44 \\
Cass Street & 15:20 & 15:35 & 15:48 & 16:01 & 16:16 & 16:42 & 17:03 & 17:24 & 17:50 & 18:05 & 18:18 & 18:31 & 18:46 \\
Bordentown & 15:27 & 15:42 & 15:55 & 16:08 & 16:23 & 16:49 & 17:11 & 17:31 & 17:57 & 18:12 & 18:25 & 18:38 & 18:53 \\
Roebling & 15:32 & 15:47 & 16:00 & 16:13 & 16:30 & 16:54 & 17:16 & 17:36 & 18:02 & 18:17 & 18:30 & 18:43 & 19:00 \\
Florence & 15:35 & 15:50 & 16:03 & 16:16 & 16:33 & 16:57 & 17:19 & 17:39 & 18:05 & 18:20 & 18:33 & 18:46 & 19:03 \\
Burlington Towne Ctr & 15:40 & 15:55 & 16:08 & 16:21 & 16:38 & 17:02 & 17:24 & 17:44 & 18:10 & 18:25 & 18:38 & 18:51 & 19:08 \\
Burlington South & 15:42 & 15:57 & 16:10 & 16:23 & 16:40 & 17:04 & 17:26 & 17:46 & 18:12 & 18:27 & 18:40 & 18:53 & 19:10 \\
Beverly/Edgewater Park & 15:46 & 16:01 & 16:17 & 16:29 & 16:44 & 17:08 & 17:30 & 17:50 & 18:16 & 18:31 & 18:47 & 18:59 & 19:14 \\
Delanco & 15:50 & 16:05 & 16:21 & 16:33 & 16:48 & 17:12 & 17:34 & 17:54 & 18:20 & 18:35 & 18:51 & 19:03 & 19:18 \\
Riverside & 15:52 & 16:07 & 16:23 & 16:35 & 16:50 & 17:14 & 17:36 & 17:56 & 18:22 & 18:37 & 18:53 & 19:05 & 19:20 \\
Cinnaminson & 15:57 & 16:12 & 16:28 & 16:40 & 16:55 & 17:19 & 17:41 & 18:03 & 18:27 & 18:42 & 18:58 & 19:10 & 19:25 \\
Riverton & 15:59 & 16:14 & 16:30 & 16:42 & 16:57 & 17:21 & 17:43 & 18:05 & 18:29 & 18:44 & 19:00 & 19:12 & 19:27 \\
Palmyra & 16:01 & 16:16 & 16:32 & 16:44 & 16:59 & 17:25 & 17:45 & 18:07 & 18:31 & 18:46 & 19:02 & 19:14 & 19:29 \\
Route 73/Pennsauken & 16:04 & 16:19 & 16:35 & 16:47 & 17:02 & 17:28 & 17:48 & 18:10 & 18:34 & 18:49 & 19:05 & 19:17 & 19:32 \\
Pennsauken TC & 16:06 & 16:21 & 16:37 & 16:49 & 17:04 & 17:30 & 17:50 & 18:12 & 18:36 & 18:51 & 19:07 & 19:19 & 19:34 \\
36th Street & 16:10 & 16:25 & 16:41 & 16:53 & 17:08 & 17:34 & 17:54 & 18:16 & 18:40 & 18:55 & 19:11 & 19:23 & 19:38 \\
Walter Rand TC & 16:15 & 16:30 & 16:46 & 16:58 & 17:13 & 17:39 & 17:59 & 18:21 & 18:45 & 19:00 & 19:16 & 19:28 & 19:43 \\
Cooper St/Rutgers & 16:18 & 16:33 & 16:49 & 17:01 & 17:16 & 17:42 & 18:02 & 18:24 & 18:48 & 19:03 & 19:19 & 19:31 & 19:46 \\
Aquarium & 16:20 & 16:35 & 16:51 & 17:03 & 17:18 & 17:44 & 18:04 & 18:26 & 18:50 & 19:05 & 19:21 & 19:33 & 19:48 \\
Entertainment Center & 16:23 & 16:38 & 16:54 & 17:06 & 17:21 & 17:47 & 18:07 & 18:29 & 18:53 & 19:08 & 19:24 & 19:36 & 19:51 \\
\bottomrule\end{tabular}\end{center}
\subsection*{To Camden (read down) --- trips 40--44 of 44}
\begin{center}\footnotesize\setlength{\tabcolsep}{3pt}\begin{tabular}{@{}lrrrrr@{}}\toprule
Station & 40 & 41 & 42 & 43 & 44 \\\midrule
Trenton TC & 19:08 & 19:29 & 19:50 & 20:16 & 20:44 \\
Hamilton Avenue & 19:10 & 19:31 & 19:52 & 20:18 & 20:46 \\
Cass Street & 19:12 & 19:33 & 19:54 & 20:20 & 20:48 \\
Bordentown & 19:19 & 19:41 & 20:01 & 20:27 & 20:55 \\
Roebling & 19:24 & 19:46 & 20:06 & 20:32 & 21:00 \\
Florence & 19:27 & 19:49 & 20:09 & 20:35 & 21:03 \\
Burlington Towne Ctr & 19:32 & 19:54 & 20:14 & 20:40 & 21:08 \\
Burlington South & 19:34 & 19:56 & 20:16 & 20:42 & 21:10 \\
Beverly/Edgewater Park & 19:38 & 20:00 & 20:20 & 20:46 & 21:17 \\
Delanco & 19:42 & 20:04 & 20:24 & 20:50 & 21:21 \\
Riverside & 19:44 & 20:06 & 20:26 & 20:52 & 21:23 \\
Cinnaminson & 19:49 & 20:11 & 20:33 & 20:57 & 21:28 \\
Riverton & 19:51 & 20:13 & 20:35 & 20:59 & 21:30 \\
Palmyra & 19:55 & 20:15 & 20:37 & 21:01 & 21:32 \\
Route 73/Pennsauken & 19:58 & 20:18 & 20:40 & 21:04 & 21:35 \\
Pennsauken TC & 20:00 & 20:20 & 20:42 & 21:06 & 21:37 \\
36th Street & 20:04 & 20:24 & 20:46 & 21:10 & 21:41 \\
Walter Rand TC & 20:09 & 20:29 & 20:51 & 21:15 & 21:46 \\
Cooper St/Rutgers & 20:12 & 20:32 & 20:54 & 21:18 & 21:49 \\
Aquarium & 20:14 & 20:34 & 20:56 & 21:20 & 21:51 \\
Entertainment Center & 20:17 & 20:37 & 20:59 & 21:23 & 21:54 \\
\bottomrule\end{tabular}\end{center}
\clearpage
\section{Proposal G8: 8 trains, 2-min meet buffers}\label{tt:proposal_G8}
\subsection*{To Trenton (read down) --- trips 1--13 of 44}
\begin{center}\footnotesize\setlength{\tabcolsep}{3pt}\begin{tabular}{@{}lrrrrrrrrrrrrr@{}}\toprule
Station & 1 & 2 & 3 & 4 & 5 & 6 & 7 & 8 & 9 & 10 & 11 & 12 & 13 \\\midrule
Entertainment Center & 5:45 & 6:05 & 6:20 & 6:34 & 7:01 & 7:16 & 7:30 & 7:50 & 8:15 & 8:35 & 8:50 & 9:04 & 9:31 \\
Aquarium & 5:47 & 6:07 & 6:22 & 6:36 & 7:03 & 7:18 & 7:32 & 7:52 & 8:17 & 8:37 & 8:52 & 9:06 & 9:33 \\
Cooper St/Rutgers & 5:49 & 6:09 & 6:24 & 6:38 & 7:05 & 7:20 & 7:34 & 7:54 & 8:19 & 8:39 & 8:54 & 9:08 & 9:35 \\
Walter Rand TC & 5:53 & 6:13 & 6:28 & 6:42 & 7:09 & 7:24 & 7:38 & 7:58 & 8:23 & 8:43 & 8:58 & 9:12 & 9:39 \\
36th Street & 6:02 & 6:18 & 6:34 & 6:47 & 7:19 & 7:31 & 7:48 & 8:03 & 8:32 & 8:48 & 9:04 & 9:17 & 9:49 \\
Pennsauken TC & 6:04 & 6:20 & 6:36 & 6:49 & 7:21 & 7:33 & 7:50 & 8:05 & 8:34 & 8:50 & 9:06 & 9:19 & 9:51 \\
Route 73/Pennsauken & 6:07 & 6:23 & 6:39 & 6:52 & 7:24 & 7:36 & 7:53 & 8:08 & 8:37 & 8:53 & 9:09 & 9:22 & 9:54 \\
Palmyra & 6:10 & 6:26 & 6:42 & 6:55 & 7:27 & 7:39 & 7:56 & 8:11 & 8:40 & 8:56 & 9:12 & 9:25 & 9:57 \\
Riverton & 6:12 & 6:28 & 6:44 & 6:57 & 7:29 & 7:41 & 7:58 & 8:13 & 8:42 & 8:58 & 9:14 & 9:27 & 9:59 \\
Cinnaminson & 6:14 & 6:30 & 6:46 & 6:59 & 7:31 & 7:43 & 8:00 & 8:15 & 8:44 & 9:00 & 9:16 & 9:29 & 10:01 \\
Riverside & 6:19 & 6:35 & 6:51 & 7:04 & 7:36 & 7:48 & 8:05 & 8:20 & 8:49 & 9:05 & 9:21 & 9:34 & 10:06 \\
Delanco & 6:21 & 6:37 & 6:53 & 7:06 & 7:38 & 7:50 & 8:07 & 8:22 & 8:51 & 9:07 & 9:23 & 9:36 & 10:08 \\
Beverly/Edgewater Park & 6:25 & 6:41 & 6:57 & 7:10 & 7:42 & 7:54 & 8:11 & 8:26 & 8:55 & 9:11 & 9:27 & 9:40 & 10:12 \\
Burlington South & 6:29 & 6:45 & 7:01 & 7:14 & 7:46 & 7:58 & 8:15 & 8:30 & 8:59 & 9:15 & 9:31 & 9:44 & 10:16 \\
Burlington Towne Ctr & 6:31 & 6:47 & 7:03 & 7:16 & 7:48 & 8:00 & 8:17 & 8:32 & 9:01 & 9:17 & 9:33 & 9:46 & 10:18 \\
Florence & 6:37 & 6:53 & 7:09 & 7:22 & 7:54 & 8:06 & 8:23 & 8:38 & 9:07 & 9:23 & 9:39 & 9:52 & 10:24 \\
Roebling & 6:40 & 6:56 & 7:12 & 7:25 & 7:57 & 8:09 & 8:26 & 8:41 & 9:10 & 9:26 & 9:42 & 9:55 & 10:27 \\
Bordentown & 6:45 & 7:01 & 7:17 & 7:30 & 8:02 & 8:14 & 8:31 & 8:46 & 9:15 & 9:31 & 9:47 & 10:00 & 10:32 \\
Cass Street & 6:51 & 7:07 & 7:23 & 7:36 & 8:08 & 8:20 & 8:37 & 8:52 & 9:21 & 9:37 & 9:53 & 10:06 & 10:38 \\
Hamilton Avenue & 6:53 & 7:09 & 7:25 & 7:38 & 8:10 & 8:22 & 8:39 & 8:54 & 9:23 & 9:39 & 9:55 & 10:08 & 10:40 \\
Trenton TC & 6:56 & 7:12 & 7:28 & 7:41 & 8:13 & 8:25 & 8:42 & 8:57 & 9:26 & 9:42 & 9:58 & 10:11 & 10:43 \\
\bottomrule\end{tabular}\end{center}
\subsection*{To Trenton (read down) --- trips 14--26 of 44}
\begin{center}\footnotesize\setlength{\tabcolsep}{3pt}\begin{tabular}{@{}lrrrrrrrrrrrrr@{}}\toprule
Station & 14 & 15 & 16 & 17 & 18 & 19 & 20 & 21 & 22 & 23 & 24 & 25 & 26 \\\midrule
Entertainment Center & 10:00 & 10:20 & 10:45 & 11:05 & 11:34 & 12:01 & 12:30 & 12:50 & 13:15 & 13:35 & 14:04 & 14:31 & 14:46 \\
Aquarium & 10:02 & 10:22 & 10:47 & 11:07 & 11:36 & 12:03 & 12:32 & 12:52 & 13:17 & 13:37 & 14:06 & 14:33 & 14:48 \\
Cooper St/Rutgers & 10:04 & 10:24 & 10:49 & 11:09 & 11:38 & 12:05 & 12:34 & 12:54 & 13:19 & 13:39 & 14:08 & 14:35 & 14:50 \\
Walter Rand TC & 10:08 & 10:28 & 10:53 & 11:13 & 11:42 & 12:09 & 12:38 & 12:58 & 13:23 & 13:43 & 14:12 & 14:39 & 14:54 \\
36th Street & 10:18 & 10:33 & 11:02 & 11:18 & 11:47 & 12:19 & 12:48 & 13:03 & 13:32 & 13:48 & 14:17 & 14:49 & 15:01 \\
Pennsauken TC & 10:20 & 10:35 & 11:04 & 11:20 & 11:49 & 12:21 & 12:50 & 13:05 & 13:34 & 13:50 & 14:19 & 14:51 & 15:03 \\
Route 73/Pennsauken & 10:23 & 10:38 & 11:07 & 11:23 & 11:52 & 12:24 & 12:53 & 13:08 & 13:37 & 13:53 & 14:22 & 14:54 & 15:06 \\
Palmyra & 10:26 & 10:41 & 11:10 & 11:26 & 11:55 & 12:27 & 12:56 & 13:11 & 13:40 & 13:56 & 14:25 & 14:57 & 15:09 \\
Riverton & 10:28 & 10:43 & 11:12 & 11:28 & 11:57 & 12:29 & 12:58 & 13:13 & 13:42 & 13:58 & 14:27 & 14:59 & 15:11 \\
Cinnaminson & 10:30 & 10:45 & 11:14 & 11:30 & 11:59 & 12:31 & 13:00 & 13:15 & 13:44 & 14:00 & 14:29 & 15:01 & 15:13 \\
Riverside & 10:35 & 10:50 & 11:19 & 11:35 & 12:04 & 12:36 & 13:05 & 13:20 & 13:49 & 14:05 & 14:34 & 15:06 & 15:18 \\
Delanco & 10:37 & 10:52 & 11:21 & 11:37 & 12:06 & 12:38 & 13:07 & 13:22 & 13:51 & 14:07 & 14:36 & 15:08 & 15:20 \\
Beverly/Edgewater Park & 10:41 & 10:56 & 11:25 & 11:41 & 12:10 & 12:42 & 13:11 & 13:26 & 13:55 & 14:11 & 14:40 & 15:12 & 15:24 \\
Burlington South & 10:45 & 11:00 & 11:29 & 11:45 & 12:14 & 12:46 & 13:15 & 13:30 & 13:59 & 14:15 & 14:44 & 15:16 & 15:28 \\
Burlington Towne Ctr & 10:47 & 11:02 & 11:31 & 11:47 & 12:16 & 12:48 & 13:17 & 13:32 & 14:01 & 14:17 & 14:46 & 15:18 & 15:30 \\
Florence & 10:53 & 11:08 & 11:37 & 11:53 & 12:22 & 12:54 & 13:23 & 13:38 & 14:07 & 14:23 & 14:52 & 15:24 & 15:36 \\
Roebling & 10:56 & 11:11 & 11:40 & 11:56 & 12:25 & 12:57 & 13:26 & 13:41 & 14:10 & 14:26 & 14:55 & 15:27 & 15:39 \\
Bordentown & 11:01 & 11:16 & 11:45 & 12:01 & 12:30 & 13:02 & 13:31 & 13:46 & 14:15 & 14:31 & 15:00 & 15:32 & 15:44 \\
Cass Street & 11:07 & 11:22 & 11:51 & 12:07 & 12:36 & 13:08 & 13:37 & 13:52 & 14:21 & 14:37 & 15:06 & 15:38 & 15:50 \\
Hamilton Avenue & 11:09 & 11:24 & 11:53 & 12:09 & 12:38 & 13:10 & 13:39 & 13:54 & 14:23 & 14:39 & 15:08 & 15:40 & 15:52 \\
Trenton TC & 11:12 & 11:27 & 11:56 & 12:12 & 12:41 & 13:13 & 13:42 & 13:57 & 14:26 & 14:42 & 15:11 & 15:43 & 15:55 \\
\bottomrule\end{tabular}\end{center}
\subsection*{To Trenton (read down) --- trips 27--39 of 44}
\begin{center}\footnotesize\setlength{\tabcolsep}{3pt}\begin{tabular}{@{}lrrrrrrrrrrrrr@{}}\toprule
Station & 27 & 28 & 29 & 30 & 31 & 32 & 33 & 34 & 35 & 36 & 37 & 38 & 39 \\\midrule
Entertainment Center & 15:00 & 15:20 & 15:45 & 16:05 & 16:20 & 16:34 & 17:01 & 17:16 & 17:30 & 17:50 & 18:15 & 18:35 & 19:04 \\
Aquarium & 15:02 & 15:22 & 15:47 & 16:07 & 16:22 & 16:36 & 17:03 & 17:18 & 17:32 & 17:52 & 18:17 & 18:37 & 19:06 \\
Cooper St/Rutgers & 15:04 & 15:24 & 15:49 & 16:09 & 16:24 & 16:38 & 17:05 & 17:20 & 17:34 & 17:54 & 18:19 & 18:39 & 19:08 \\
Walter Rand TC & 15:08 & 15:28 & 15:53 & 16:13 & 16:28 & 16:42 & 17:09 & 17:24 & 17:38 & 17:58 & 18:23 & 18:43 & 19:12 \\
36th Street & 15:18 & 15:33 & 16:02 & 16:18 & 16:34 & 16:47 & 17:19 & 17:31 & 17:48 & 18:03 & 18:32 & 18:48 & 19:17 \\
Pennsauken TC & 15:20 & 15:35 & 16:04 & 16:20 & 16:36 & 16:49 & 17:21 & 17:33 & 17:50 & 18:05 & 18:34 & 18:50 & 19:19 \\
Route 73/Pennsauken & 15:23 & 15:38 & 16:07 & 16:23 & 16:39 & 16:52 & 17:24 & 17:36 & 17:53 & 18:08 & 18:37 & 18:53 & 19:22 \\
Palmyra & 15:26 & 15:41 & 16:10 & 16:26 & 16:42 & 16:55 & 17:27 & 17:39 & 17:56 & 18:11 & 18:40 & 18:56 & 19:25 \\
Riverton & 15:28 & 15:43 & 16:12 & 16:28 & 16:44 & 16:57 & 17:29 & 17:41 & 17:58 & 18:13 & 18:42 & 18:58 & 19:27 \\
Cinnaminson & 15:30 & 15:45 & 16:14 & 16:30 & 16:46 & 16:59 & 17:31 & 17:43 & 18:00 & 18:15 & 18:44 & 19:00 & 19:29 \\
Riverside & 15:35 & 15:50 & 16:19 & 16:35 & 16:51 & 17:04 & 17:36 & 17:48 & 18:05 & 18:20 & 18:49 & 19:05 & 19:34 \\
Delanco & 15:37 & 15:52 & 16:21 & 16:37 & 16:53 & 17:06 & 17:38 & 17:50 & 18:07 & 18:22 & 18:51 & 19:07 & 19:36 \\
Beverly/Edgewater Park & 15:41 & 15:56 & 16:25 & 16:41 & 16:57 & 17:10 & 17:42 & 17:54 & 18:11 & 18:26 & 18:55 & 19:11 & 19:40 \\
Burlington South & 15:45 & 16:00 & 16:29 & 16:45 & 17:01 & 17:14 & 17:46 & 17:58 & 18:15 & 18:30 & 18:59 & 19:15 & 19:44 \\
Burlington Towne Ctr & 15:47 & 16:02 & 16:31 & 16:47 & 17:03 & 17:16 & 17:48 & 18:00 & 18:17 & 18:32 & 19:01 & 19:17 & 19:46 \\
Florence & 15:53 & 16:08 & 16:37 & 16:53 & 17:09 & 17:22 & 17:54 & 18:06 & 18:23 & 18:38 & 19:07 & 19:23 & 19:52 \\
Roebling & 15:56 & 16:11 & 16:40 & 16:56 & 17:12 & 17:25 & 17:57 & 18:09 & 18:26 & 18:41 & 19:10 & 19:26 & 19:55 \\
Bordentown & 16:01 & 16:16 & 16:45 & 17:01 & 17:17 & 17:30 & 18:02 & 18:14 & 18:31 & 18:46 & 19:15 & 19:31 & 20:00 \\
Cass Street & 16:07 & 16:22 & 16:51 & 17:07 & 17:23 & 17:36 & 18:08 & 18:20 & 18:37 & 18:52 & 19:21 & 19:37 & 20:06 \\
Hamilton Avenue & 16:09 & 16:24 & 16:53 & 17:09 & 17:25 & 17:38 & 18:10 & 18:22 & 18:39 & 18:54 & 19:23 & 19:39 & 20:08 \\
Trenton TC & 16:12 & 16:27 & 16:56 & 17:12 & 17:28 & 17:41 & 18:13 & 18:25 & 18:42 & 18:57 & 19:26 & 19:42 & 20:11 \\
\bottomrule\end{tabular}\end{center}
\subsection*{To Trenton (read down) --- trips 40--44 of 44}
\begin{center}\footnotesize\setlength{\tabcolsep}{3pt}\begin{tabular}{@{}lrrrrr@{}}\toprule
Station & 40 & 41 & 42 & 43 & 44 \\\midrule
Entertainment Center & 19:31 & 20:00 & 20:20 & 20:45 & 21:05 \\
Aquarium & 19:33 & 20:02 & 20:22 & 20:47 & 21:07 \\
Cooper St/Rutgers & 19:35 & 20:04 & 20:24 & 20:49 & 21:09 \\
Walter Rand TC & 19:39 & 20:08 & 20:28 & 20:53 & 21:13 \\
36th Street & 19:49 & 20:18 & 20:33 & 21:02 & 21:18 \\
Pennsauken TC & 19:51 & 20:20 & 20:35 & 21:04 & 21:20 \\
Route 73/Pennsauken & 19:54 & 20:23 & 20:38 & 21:07 & 21:23 \\
Palmyra & 19:57 & 20:26 & 20:41 & 21:10 & 21:26 \\
Riverton & 19:59 & 20:28 & 20:43 & 21:12 & 21:28 \\
Cinnaminson & 20:01 & 20:30 & 20:45 & 21:14 & 21:30 \\
Riverside & 20:06 & 20:35 & 20:50 & 21:19 & 21:35 \\
Delanco & 20:08 & 20:37 & 20:52 & 21:21 & 21:37 \\
Beverly/Edgewater Park & 20:12 & 20:41 & 20:56 & 21:25 & 21:41 \\
Burlington South & 20:16 & 20:45 & 21:00 & 21:29 & 21:45 \\
Burlington Towne Ctr & 20:18 & 20:47 & 21:02 & 21:31 & 21:47 \\
Florence & 20:24 & 20:53 & 21:08 & 21:37 & 21:53 \\
Roebling & 20:27 & 20:56 & 21:11 & 21:40 & 21:56 \\
Bordentown & 20:32 & 21:01 & 21:16 & 21:45 & 22:01 \\
Cass Street & 20:38 & 21:07 & 21:22 & 21:51 & 22:07 \\
Hamilton Avenue & 20:40 & 21:09 & 21:24 & 21:53 & 22:09 \\
Trenton TC & 20:43 & 21:12 & 21:27 & 21:56 & 22:12 \\
\bottomrule\end{tabular}\end{center}
\subsection*{To Camden (read down) --- trips 1--13 of 44}
\begin{center}\footnotesize\setlength{\tabcolsep}{3pt}\begin{tabular}{@{}lrrrrrrrrrrrrr@{}}\toprule
Station & 1 & 2 & 3 & 4 & 5 & 6 & 7 & 8 & 9 & 10 & 11 & 12 & 13 \\\midrule
Trenton TC & 5:48 & 6:02 & 6:17 & 6:34 & 7:01 & 7:18 & 7:33 & 7:47 & 8:18 & 8:32 & 8:47 & 9:04 & 9:31 \\
Hamilton Avenue & 5:50 & 6:04 & 6:19 & 6:36 & 7:03 & 7:20 & 7:35 & 7:49 & 8:20 & 8:34 & 8:49 & 9:06 & 9:33 \\
Cass Street & 5:52 & 6:06 & 6:21 & 6:38 & 7:05 & 7:22 & 7:37 & 7:51 & 8:22 & 8:36 & 8:51 & 9:08 & 9:35 \\
Bordentown & 5:59 & 6:13 & 6:29 & 6:45 & 7:12 & 7:29 & 7:45 & 7:58 & 8:29 & 8:43 & 8:59 & 9:15 & 9:42 \\
Roebling & 6:05 & 6:18 & 6:34 & 6:50 & 7:19 & 7:34 & 7:50 & 8:03 & 8:35 & 8:48 & 9:04 & 9:20 & 9:49 \\
Florence & 6:08 & 6:21 & 6:37 & 6:53 & 7:22 & 7:37 & 7:53 & 8:06 & 8:38 & 8:51 & 9:07 & 9:23 & 9:52 \\
Burlington Towne Ctr & 6:13 & 6:26 & 6:42 & 6:58 & 7:27 & 7:42 & 7:58 & 8:11 & 8:43 & 8:56 & 9:12 & 9:28 & 9:57 \\
Burlington South & 6:15 & 6:28 & 6:44 & 7:00 & 7:29 & 7:44 & 8:00 & 8:13 & 8:45 & 8:58 & 9:14 & 9:30 & 9:59 \\
Beverly/Edgewater Park & 6:19 & 6:32 & 6:48 & 7:04 & 7:33 & 7:48 & 8:04 & 8:17 & 8:49 & 9:02 & 9:18 & 9:34 & 10:03 \\
Delanco & 6:23 & 6:36 & 6:52 & 7:08 & 7:37 & 7:52 & 8:08 & 8:21 & 8:53 & 9:06 & 9:22 & 9:38 & 10:07 \\
Riverside & 6:25 & 6:38 & 6:54 & 7:10 & 7:39 & 7:54 & 8:10 & 8:23 & 8:55 & 9:08 & 9:24 & 9:40 & 10:09 \\
Cinnaminson & 6:30 & 6:43 & 6:59 & 7:15 & 7:44 & 7:59 & 8:15 & 8:28 & 9:00 & 9:13 & 9:29 & 9:45 & 10:14 \\
Riverton & 6:32 & 6:45 & 7:01 & 7:17 & 7:46 & 8:01 & 8:17 & 8:30 & 9:02 & 9:15 & 9:31 & 9:47 & 10:16 \\
Palmyra & 6:34 & 6:47 & 7:03 & 7:19 & 7:48 & 8:03 & 8:19 & 8:32 & 9:04 & 9:17 & 9:33 & 9:49 & 10:18 \\
Route 73/Pennsauken & 6:37 & 6:50 & 7:06 & 7:22 & 7:51 & 8:06 & 8:22 & 8:35 & 9:07 & 9:20 & 9:36 & 9:52 & 10:21 \\
Pennsauken TC & 6:39 & 6:52 & 7:08 & 7:24 & 7:53 & 8:08 & 8:24 & 8:37 & 9:09 & 9:22 & 9:38 & 9:54 & 10:23 \\
36th Street & 6:43 & 6:56 & 7:12 & 7:28 & 7:57 & 8:12 & 8:28 & 8:41 & 9:13 & 9:26 & 9:42 & 9:58 & 10:27 \\
Walter Rand TC & 6:48 & 7:01 & 7:17 & 7:33 & 8:02 & 8:17 & 8:33 & 8:46 & 9:18 & 9:31 & 9:47 & 10:03 & 10:32 \\
Cooper St/Rutgers & 6:51 & 7:04 & 7:20 & 7:36 & 8:05 & 8:20 & 8:36 & 8:49 & 9:21 & 9:34 & 9:50 & 10:06 & 10:35 \\
Aquarium & 6:53 & 7:06 & 7:22 & 7:38 & 8:07 & 8:22 & 8:38 & 8:51 & 9:23 & 9:36 & 9:52 & 10:08 & 10:37 \\
Entertainment Center & 6:56 & 7:09 & 7:25 & 7:41 & 8:10 & 8:25 & 8:41 & 8:54 & 9:26 & 9:39 & 9:55 & 10:11 & 10:40 \\
\bottomrule\end{tabular}\end{center}
\subsection*{To Camden (read down) --- trips 14--26 of 44}
\begin{center}\footnotesize\setlength{\tabcolsep}{3pt}\begin{tabular}{@{}lrrrrrrrrrrrrr@{}}\toprule
Station & 14 & 15 & 16 & 17 & 18 & 19 & 20 & 21 & 22 & 23 & 24 & 25 & 26 \\\midrule
Trenton TC & 9:48 & 10:03 & 10:17 & 10:48 & 11:17 & 11:34 & 12:01 & 12:18 & 12:47 & 13:18 & 13:47 & 14:04 & 14:31 \\
Hamilton Avenue & 9:50 & 10:05 & 10:19 & 10:50 & 11:19 & 11:36 & 12:03 & 12:20 & 12:49 & 13:20 & 13:49 & 14:06 & 14:33 \\
Cass Street & 9:52 & 10:07 & 10:21 & 10:52 & 11:21 & 11:38 & 12:05 & 12:22 & 12:51 & 13:22 & 13:51 & 14:08 & 14:35 \\
Bordentown & 9:59 & 10:15 & 10:28 & 10:59 & 11:29 & 11:45 & 12:12 & 12:29 & 12:58 & 13:29 & 13:59 & 14:15 & 14:42 \\
Roebling & 10:04 & 10:20 & 10:33 & 11:05 & 11:34 & 11:50 & 12:19 & 12:34 & 13:03 & 13:35 & 14:04 & 14:20 & 14:49 \\
Florence & 10:07 & 10:23 & 10:36 & 11:08 & 11:37 & 11:53 & 12:22 & 12:37 & 13:06 & 13:38 & 14:07 & 14:23 & 14:52 \\
Burlington Towne Ctr & 10:12 & 10:28 & 10:41 & 11:13 & 11:42 & 11:58 & 12:27 & 12:42 & 13:11 & 13:43 & 14:12 & 14:28 & 14:57 \\
Burlington South & 10:14 & 10:30 & 10:43 & 11:15 & 11:44 & 12:00 & 12:29 & 12:44 & 13:13 & 13:45 & 14:14 & 14:30 & 14:59 \\
Beverly/Edgewater Park & 10:18 & 10:34 & 10:47 & 11:19 & 11:48 & 12:04 & 12:33 & 12:48 & 13:17 & 13:49 & 14:18 & 14:34 & 15:03 \\
Delanco & 10:22 & 10:38 & 10:51 & 11:23 & 11:52 & 12:08 & 12:37 & 12:52 & 13:21 & 13:53 & 14:22 & 14:38 & 15:07 \\
Riverside & 10:24 & 10:40 & 10:53 & 11:25 & 11:54 & 12:10 & 12:39 & 12:54 & 13:23 & 13:55 & 14:24 & 14:40 & 15:09 \\
Cinnaminson & 10:29 & 10:45 & 10:58 & 11:30 & 11:59 & 12:15 & 12:44 & 12:59 & 13:28 & 14:00 & 14:29 & 14:45 & 15:14 \\
Riverton & 10:31 & 10:47 & 11:00 & 11:32 & 12:01 & 12:17 & 12:46 & 13:01 & 13:30 & 14:02 & 14:31 & 14:47 & 15:16 \\
Palmyra & 10:33 & 10:49 & 11:02 & 11:34 & 12:03 & 12:19 & 12:48 & 13:03 & 13:32 & 14:04 & 14:33 & 14:49 & 15:18 \\
Route 73/Pennsauken & 10:36 & 10:52 & 11:05 & 11:37 & 12:06 & 12:22 & 12:51 & 13:06 & 13:35 & 14:07 & 14:36 & 14:52 & 15:21 \\
Pennsauken TC & 10:38 & 10:54 & 11:07 & 11:39 & 12:08 & 12:24 & 12:53 & 13:08 & 13:37 & 14:09 & 14:38 & 14:54 & 15:23 \\
36th Street & 10:42 & 10:58 & 11:11 & 11:43 & 12:12 & 12:28 & 12:57 & 13:12 & 13:41 & 14:13 & 14:42 & 14:58 & 15:27 \\
Walter Rand TC & 10:47 & 11:03 & 11:16 & 11:48 & 12:17 & 12:33 & 13:02 & 13:17 & 13:46 & 14:18 & 14:47 & 15:03 & 15:32 \\
Cooper St/Rutgers & 10:50 & 11:06 & 11:19 & 11:51 & 12:20 & 12:36 & 13:05 & 13:20 & 13:49 & 14:21 & 14:50 & 15:06 & 15:35 \\
Aquarium & 10:52 & 11:08 & 11:21 & 11:53 & 12:22 & 12:38 & 13:07 & 13:22 & 13:51 & 14:23 & 14:52 & 15:08 & 15:37 \\
Entertainment Center & 10:55 & 11:11 & 11:24 & 11:56 & 12:25 & 12:41 & 13:10 & 13:25 & 13:54 & 14:26 & 14:55 & 15:11 & 15:40 \\
\bottomrule\end{tabular}\end{center}
\subsection*{To Camden (read down) --- trips 27--39 of 44}
\begin{center}\footnotesize\setlength{\tabcolsep}{3pt}\begin{tabular}{@{}lrrrrrrrrrrrrr@{}}\toprule
Station & 27 & 28 & 29 & 30 & 31 & 32 & 33 & 34 & 35 & 36 & 37 & 38 & 39 \\\midrule
Trenton TC & 14:48 & 15:17 & 15:48 & 16:02 & 16:17 & 16:34 & 17:01 & 17:18 & 17:33 & 17:47 & 18:18 & 18:32 & 18:47 \\
Hamilton Avenue & 14:50 & 15:19 & 15:50 & 16:04 & 16:19 & 16:36 & 17:03 & 17:20 & 17:35 & 17:49 & 18:20 & 18:34 & 18:49 \\
Cass Street & 14:52 & 15:21 & 15:52 & 16:06 & 16:21 & 16:38 & 17:05 & 17:22 & 17:37 & 17:51 & 18:22 & 18:36 & 18:51 \\
Bordentown & 14:59 & 15:28 & 15:59 & 16:13 & 16:29 & 16:45 & 17:12 & 17:29 & 17:45 & 17:58 & 18:29 & 18:43 & 18:59 \\
Roebling & 15:04 & 15:33 & 16:05 & 16:18 & 16:34 & 16:50 & 17:19 & 17:34 & 17:50 & 18:03 & 18:35 & 18:48 & 19:04 \\
Florence & 15:07 & 15:36 & 16:08 & 16:21 & 16:37 & 16:53 & 17:22 & 17:37 & 17:53 & 18:06 & 18:38 & 18:51 & 19:07 \\
Burlington Towne Ctr & 15:12 & 15:41 & 16:13 & 16:26 & 16:42 & 16:58 & 17:27 & 17:42 & 17:58 & 18:11 & 18:43 & 18:56 & 19:12 \\
Burlington South & 15:14 & 15:43 & 16:15 & 16:28 & 16:44 & 17:00 & 17:29 & 17:44 & 18:00 & 18:13 & 18:45 & 18:58 & 19:14 \\
Beverly/Edgewater Park & 15:18 & 15:47 & 16:19 & 16:32 & 16:48 & 17:04 & 17:33 & 17:48 & 18:04 & 18:17 & 18:49 & 19:02 & 19:18 \\
Delanco & 15:22 & 15:51 & 16:23 & 16:36 & 16:52 & 17:08 & 17:37 & 17:52 & 18:08 & 18:21 & 18:53 & 19:06 & 19:22 \\
Riverside & 15:24 & 15:53 & 16:25 & 16:38 & 16:54 & 17:10 & 17:39 & 17:54 & 18:10 & 18:23 & 18:55 & 19:08 & 19:24 \\
Cinnaminson & 15:29 & 15:58 & 16:30 & 16:43 & 16:59 & 17:15 & 17:44 & 17:59 & 18:15 & 18:28 & 19:00 & 19:13 & 19:29 \\
Riverton & 15:31 & 16:00 & 16:32 & 16:45 & 17:01 & 17:17 & 17:46 & 18:01 & 18:17 & 18:30 & 19:02 & 19:15 & 19:31 \\
Palmyra & 15:33 & 16:02 & 16:34 & 16:47 & 17:03 & 17:19 & 17:48 & 18:03 & 18:19 & 18:32 & 19:04 & 19:17 & 19:33 \\
Route 73/Pennsauken & 15:36 & 16:05 & 16:37 & 16:50 & 17:06 & 17:22 & 17:51 & 18:06 & 18:22 & 18:35 & 19:07 & 19:20 & 19:36 \\
Pennsauken TC & 15:38 & 16:07 & 16:39 & 16:52 & 17:08 & 17:24 & 17:53 & 18:08 & 18:24 & 18:37 & 19:09 & 19:22 & 19:38 \\
36th Street & 15:42 & 16:11 & 16:43 & 16:56 & 17:12 & 17:28 & 17:57 & 18:12 & 18:28 & 18:41 & 19:13 & 19:26 & 19:42 \\
Walter Rand TC & 15:47 & 16:16 & 16:48 & 17:01 & 17:17 & 17:33 & 18:02 & 18:17 & 18:33 & 18:46 & 19:18 & 19:31 & 19:47 \\
Cooper St/Rutgers & 15:50 & 16:19 & 16:51 & 17:04 & 17:20 & 17:36 & 18:05 & 18:20 & 18:36 & 18:49 & 19:21 & 19:34 & 19:50 \\
Aquarium & 15:52 & 16:21 & 16:53 & 17:06 & 17:22 & 17:38 & 18:07 & 18:22 & 18:38 & 18:51 & 19:23 & 19:36 & 19:52 \\
Entertainment Center & 15:55 & 16:24 & 16:56 & 17:09 & 17:25 & 17:41 & 18:10 & 18:25 & 18:41 & 18:54 & 19:26 & 19:39 & 19:55 \\
\bottomrule\end{tabular}\end{center}
\subsection*{To Camden (read down) --- trips 40--44 of 44}
\begin{center}\footnotesize\setlength{\tabcolsep}{3pt}\begin{tabular}{@{}lrrrrr@{}}\toprule
Station & 40 & 41 & 42 & 43 & 44 \\\midrule
Trenton TC & 19:04 & 19:31 & 19:48 & 20:17 & 20:48 \\
Hamilton Avenue & 19:06 & 19:33 & 19:50 & 20:19 & 20:50 \\
Cass Street & 19:08 & 19:35 & 19:52 & 20:21 & 20:52 \\
Bordentown & 19:15 & 19:42 & 19:59 & 20:28 & 20:59 \\
Roebling & 19:20 & 19:49 & 20:04 & 20:33 & 21:05 \\
Florence & 19:23 & 19:52 & 20:07 & 20:36 & 21:08 \\
Burlington Towne Ctr & 19:28 & 19:57 & 20:12 & 20:41 & 21:13 \\
Burlington South & 19:30 & 19:59 & 20:14 & 20:43 & 21:15 \\
Beverly/Edgewater Park & 19:34 & 20:03 & 20:18 & 20:47 & 21:19 \\
Delanco & 19:38 & 20:07 & 20:22 & 20:51 & 21:23 \\
Riverside & 19:40 & 20:09 & 20:24 & 20:53 & 21:25 \\
Cinnaminson & 19:45 & 20:14 & 20:29 & 20:58 & 21:30 \\
Riverton & 19:47 & 20:16 & 20:31 & 21:00 & 21:32 \\
Palmyra & 19:49 & 20:18 & 20:33 & 21:02 & 21:34 \\
Route 73/Pennsauken & 19:52 & 20:21 & 20:36 & 21:05 & 21:37 \\
Pennsauken TC & 19:54 & 20:23 & 20:38 & 21:07 & 21:39 \\
36th Street & 19:58 & 20:27 & 20:42 & 21:11 & 21:43 \\
Walter Rand TC & 20:03 & 20:32 & 20:47 & 21:16 & 21:48 \\
Cooper St/Rutgers & 20:06 & 20:35 & 20:50 & 21:19 & 21:51 \\
Aquarium & 20:08 & 20:37 & 20:52 & 21:21 & 21:53 \\
Entertainment Center & 20:11 & 20:40 & 20:55 & 21:24 & 21:56 \\
\bottomrule\end{tabular}\end{center}
\clearpage
\section{Proposal S8: 8 trains, slack held before single track}\label{tt:proposal_S8}
\subsection*{To Trenton (read down) --- trips 1--13 of 43}
\begin{center}\footnotesize\setlength{\tabcolsep}{3pt}\begin{tabular}{@{}lrrrrrrrrrrrrr@{}}\toprule
Station & 1 & 2 & 3 & 4 & 5 & 6 & 7 & 8 & 9 & 10 & 11 & 12 & 13 \\\midrule
Entertainment Center & 5:46 & 6:09 & 6:31 & 6:46 & 7:00 & 7:15 & 7:30 & 7:53 & 8:16 & 8:39 & 9:01 & 9:30 & 10:00 \\
Aquarium & 5:47 & 6:10 & 6:33 & 6:48 & 7:02 & 7:17 & 7:32 & 7:55 & 8:17 & 8:40 & 9:03 & 9:32 & 10:02 \\
Cooper St/Rutgers & 5:49 & 6:12 & 6:35 & 6:50 & 7:04 & 7:19 & 7:34 & 7:57 & 8:19 & 8:42 & 9:05 & 9:34 & 10:04 \\
Walter Rand TC & 5:53 & 6:16 & 6:38 & 6:53 & 7:07 & 7:23 & 7:38 & 8:01 & 8:23 & 8:46 & 9:08 & 9:37 & 10:08 \\
36th Street & 5:58 & 6:21 & 6:44 & 6:59 & 7:13 & 7:28 & 7:46 & 8:11 & 8:28 & 8:51 & 9:14 & 9:43 & 10:16 \\
Pennsauken TC & 6:00 & 6:23 & 6:45 & 7:00 & 7:14 & 7:30 & 7:47 & 8:13 & 8:30 & 8:53 & 9:15 & 9:44 & 10:17 \\
Route 73/Pennsauken & 6:04 & 6:27 & 6:49 & 7:04 & 7:18 & 7:33 & 7:51 & 8:16 & 8:34 & 8:57 & 9:19 & 9:48 & 10:21 \\
Palmyra & 6:06 & 6:29 & 6:52 & 7:07 & 7:21 & 7:36 & 7:54 & 8:19 & 8:36 & 8:59 & 9:22 & 9:51 & 10:24 \\
Riverton & 6:08 & 6:31 & 6:54 & 7:09 & 7:23 & 7:38 & 7:55 & 8:21 & 8:38 & 9:01 & 9:24 & 9:53 & 10:25 \\
Cinnaminson & 6:10 & 6:33 & 6:56 & 7:11 & 7:24 & 7:40 & 7:57 & 8:23 & 8:40 & 9:03 & 9:26 & 9:54 & 10:27 \\
Riverside & 6:16 & 6:39 & 7:01 & 7:16 & 7:30 & 7:45 & 8:03 & 8:28 & 8:46 & 9:09 & 9:31 & 10:00 & 10:33 \\
Delanco & 6:17 & 6:40 & 7:03 & 7:18 & 7:32 & 7:47 & 8:05 & 8:30 & 8:47 & 9:10 & 9:33 & 10:02 & 10:35 \\
Beverly/Edgewater Park & 6:21 & 6:44 & 7:07 & 7:22 & 7:35 & 7:51 & 8:08 & 8:34 & 8:51 & 9:14 & 9:37 & 10:05 & 10:38 \\
Burlington South & 6:26 & 6:49 & 7:11 & 7:26 & 7:40 & 7:55 & 8:13 & 8:38 & 8:56 & 9:19 & 9:41 & 10:10 & 10:43 \\
Burlington Towne Ctr & 6:27 & 6:51 & 7:13 & 7:28 & 7:42 & 7:57 & 8:14 & 8:40 & 8:57 & 9:21 & 9:43 & 10:12 & 10:44 \\
Florence & 6:34 & 6:57 & 7:19 & 7:34 & 7:48 & 8:04 & 8:21 & 8:46 & 9:04 & 9:27 & 9:49 & 10:18 & 10:51 \\
Roebling & 6:37 & 7:01 & 7:23 & 7:37 & 7:51 & 8:07 & 8:24 & 8:50 & 9:07 & 9:31 & 9:53 & 10:21 & 10:54 \\
Bordentown & 6:42 & 7:05 & 7:27 & 7:42 & 7:56 & 8:12 & 8:29 & 8:55 & 9:12 & 9:35 & 9:57 & 10:26 & 10:59 \\
Cass Street & 6:48 & 7:12 & 7:34 & 7:48 & 8:02 & 8:18 & 8:35 & 9:01 & 9:18 & 9:42 & 10:04 & 10:32 & 11:05 \\
Hamilton Avenue & 6:50 & 7:13 & 7:36 & 7:50 & 8:04 & 8:20 & 8:37 & 9:03 & 9:20 & 9:43 & 10:06 & 10:34 & 11:07 \\
Trenton TC & 6:53 & 7:16 & 7:38 & 7:53 & 8:07 & 8:23 & 8:40 & 9:05 & 9:23 & 9:46 & 10:08 & 10:37 & 11:10 \\
\bottomrule\end{tabular}\end{center}
\subsection*{To Trenton (read down) --- trips 14--26 of 43}
\begin{center}\footnotesize\setlength{\tabcolsep}{3pt}\begin{tabular}{@{}lrrrrrrrrrrrrr@{}}\toprule
Station & 14 & 15 & 16 & 17 & 18 & 19 & 20 & 21 & 22 & 23 & 24 & 25 & 26 \\\midrule
Entertainment Center & 10:23 & 10:46 & 11:09 & 11:31 & 12:00 & 12:30 & 12:53 & 13:16 & 13:39 & 14:01 & 14:16 & 14:30 & 14:45 \\
Aquarium & 10:25 & 10:47 & 11:10 & 11:33 & 12:02 & 12:32 & 12:55 & 13:17 & 13:40 & 14:03 & 14:18 & 14:32 & 14:47 \\
Cooper St/Rutgers & 10:27 & 10:49 & 11:12 & 11:35 & 12:04 & 12:34 & 12:57 & 13:19 & 13:42 & 14:05 & 14:20 & 14:34 & 14:49 \\
Walter Rand TC & 10:31 & 10:53 & 11:16 & 11:38 & 12:07 & 12:38 & 13:01 & 13:23 & 13:46 & 14:08 & 14:23 & 14:37 & 14:53 \\
36th Street & 10:41 & 10:58 & 11:21 & 11:44 & 12:13 & 12:46 & 13:11 & 13:28 & 13:51 & 14:14 & 14:29 & 14:43 & 14:58 \\
Pennsauken TC & 10:43 & 11:00 & 11:23 & 11:45 & 12:14 & 12:47 & 13:13 & 13:30 & 13:53 & 14:15 & 14:30 & 14:44 & 15:00 \\
Route 73/Pennsauken & 10:46 & 11:04 & 11:27 & 11:49 & 12:18 & 12:51 & 13:16 & 13:34 & 13:57 & 14:19 & 14:34 & 14:48 & 15:03 \\
Palmyra & 10:49 & 11:06 & 11:29 & 11:52 & 12:21 & 12:54 & 13:19 & 13:36 & 13:59 & 14:22 & 14:37 & 14:51 & 15:06 \\
Riverton & 10:51 & 11:08 & 11:31 & 11:54 & 12:23 & 12:55 & 13:21 & 13:38 & 14:01 & 14:24 & 14:39 & 14:53 & 15:08 \\
Cinnaminson & 10:53 & 11:10 & 11:33 & 11:56 & 12:24 & 12:57 & 13:23 & 13:40 & 14:03 & 14:26 & 14:41 & 14:54 & 15:10 \\
Riverside & 10:58 & 11:16 & 11:39 & 12:01 & 12:30 & 13:03 & 13:28 & 13:46 & 14:09 & 14:31 & 14:46 & 15:00 & 15:15 \\
Delanco & 11:00 & 11:17 & 11:40 & 12:03 & 12:32 & 13:05 & 13:30 & 13:47 & 14:10 & 14:33 & 14:48 & 15:02 & 15:17 \\
Beverly/Edgewater Park & 11:04 & 11:21 & 11:44 & 12:07 & 12:35 & 13:08 & 13:34 & 13:51 & 14:14 & 14:37 & 14:52 & 15:05 & 15:21 \\
Burlington South & 11:08 & 11:26 & 11:49 & 12:11 & 12:40 & 13:13 & 13:38 & 13:56 & 14:19 & 14:41 & 14:56 & 15:10 & 15:25 \\
Burlington Towne Ctr & 11:10 & 11:27 & 11:51 & 12:13 & 12:42 & 13:14 & 13:40 & 13:57 & 14:21 & 14:43 & 14:58 & 15:12 & 15:27 \\
Florence & 11:16 & 11:34 & 11:57 & 12:19 & 12:48 & 13:21 & 13:46 & 14:04 & 14:27 & 14:49 & 15:04 & 15:18 & 15:34 \\
Roebling & 11:20 & 11:37 & 12:01 & 12:23 & 12:51 & 13:24 & 13:50 & 14:07 & 14:31 & 14:53 & 15:07 & 15:21 & 15:37 \\
Bordentown & 11:25 & 11:42 & 12:05 & 12:27 & 12:56 & 13:29 & 13:55 & 14:12 & 14:35 & 14:57 & 15:12 & 15:26 & 15:42 \\
Cass Street & 11:31 & 11:48 & 12:12 & 12:34 & 13:02 & 13:35 & 14:01 & 14:18 & 14:42 & 15:04 & 15:18 & 15:32 & 15:48 \\
Hamilton Avenue & 11:33 & 11:50 & 12:13 & 12:36 & 13:04 & 13:37 & 14:03 & 14:20 & 14:43 & 15:06 & 15:20 & 15:34 & 15:50 \\
Trenton TC & 11:35 & 11:53 & 12:16 & 12:38 & 13:07 & 13:40 & 14:05 & 14:23 & 14:46 & 15:08 & 15:23 & 15:37 & 15:53 \\
\bottomrule\end{tabular}\end{center}
\subsection*{To Trenton (read down) --- trips 27--39 of 43}
\begin{center}\footnotesize\setlength{\tabcolsep}{3pt}\begin{tabular}{@{}lrrrrrrrrrrrrr@{}}\toprule
Station & 27 & 28 & 29 & 30 & 31 & 32 & 33 & 34 & 35 & 36 & 37 & 38 & 39 \\\midrule
Entertainment Center & 15:00 & 15:23 & 15:46 & 16:09 & 16:31 & 16:46 & 17:00 & 17:15 & 17:30 & 17:53 & 18:16 & 18:39 & 19:01 \\
Aquarium & 15:02 & 15:25 & 15:47 & 16:10 & 16:33 & 16:48 & 17:02 & 17:17 & 17:32 & 17:55 & 18:17 & 18:40 & 19:03 \\
Cooper St/Rutgers & 15:04 & 15:27 & 15:49 & 16:12 & 16:35 & 16:50 & 17:04 & 17:19 & 17:34 & 17:57 & 18:19 & 18:42 & 19:05 \\
Walter Rand TC & 15:08 & 15:31 & 15:53 & 16:16 & 16:38 & 16:53 & 17:07 & 17:23 & 17:38 & 18:01 & 18:23 & 18:46 & 19:08 \\
36th Street & 15:16 & 15:41 & 15:58 & 16:21 & 16:44 & 16:59 & 17:13 & 17:28 & 17:46 & 18:11 & 18:28 & 18:51 & 19:14 \\
Pennsauken TC & 15:17 & 15:43 & 16:00 & 16:23 & 16:45 & 17:00 & 17:14 & 17:30 & 17:47 & 18:13 & 18:30 & 18:53 & 19:15 \\
Route 73/Pennsauken & 15:21 & 15:46 & 16:04 & 16:27 & 16:49 & 17:04 & 17:18 & 17:33 & 17:51 & 18:16 & 18:34 & 18:57 & 19:19 \\
Palmyra & 15:24 & 15:49 & 16:06 & 16:29 & 16:52 & 17:07 & 17:21 & 17:36 & 17:54 & 18:19 & 18:36 & 18:59 & 19:22 \\
Riverton & 15:25 & 15:51 & 16:08 & 16:31 & 16:54 & 17:09 & 17:23 & 17:38 & 17:55 & 18:21 & 18:38 & 19:01 & 19:24 \\
Cinnaminson & 15:27 & 15:53 & 16:10 & 16:33 & 16:56 & 17:11 & 17:24 & 17:40 & 17:57 & 18:23 & 18:40 & 19:03 & 19:26 \\
Riverside & 15:33 & 15:58 & 16:16 & 16:39 & 17:01 & 17:16 & 17:30 & 17:45 & 18:03 & 18:28 & 18:46 & 19:09 & 19:31 \\
Delanco & 15:35 & 16:00 & 16:17 & 16:40 & 17:03 & 17:18 & 17:32 & 17:47 & 18:05 & 18:30 & 18:47 & 19:10 & 19:33 \\
Beverly/Edgewater Park & 15:38 & 16:04 & 16:21 & 16:44 & 17:07 & 17:22 & 17:35 & 17:51 & 18:08 & 18:34 & 18:51 & 19:14 & 19:37 \\
Burlington South & 15:43 & 16:08 & 16:26 & 16:49 & 17:11 & 17:26 & 17:40 & 17:55 & 18:13 & 18:38 & 18:56 & 19:19 & 19:41 \\
Burlington Towne Ctr & 15:44 & 16:10 & 16:27 & 16:51 & 17:13 & 17:28 & 17:42 & 17:57 & 18:14 & 18:40 & 18:57 & 19:21 & 19:43 \\
Florence & 15:51 & 16:16 & 16:34 & 16:57 & 17:19 & 17:34 & 17:48 & 18:04 & 18:21 & 18:46 & 19:04 & 19:27 & 19:49 \\
Roebling & 15:54 & 16:20 & 16:37 & 17:01 & 17:23 & 17:37 & 17:51 & 18:07 & 18:24 & 18:50 & 19:07 & 19:31 & 19:53 \\
Bordentown & 15:59 & 16:25 & 16:42 & 17:05 & 17:27 & 17:42 & 17:56 & 18:12 & 18:29 & 18:55 & 19:12 & 19:35 & 19:57 \\
Cass Street & 16:05 & 16:31 & 16:48 & 17:12 & 17:34 & 17:48 & 18:02 & 18:18 & 18:35 & 19:01 & 19:18 & 19:42 & 20:04 \\
Hamilton Avenue & 16:07 & 16:33 & 16:50 & 17:13 & 17:36 & 17:50 & 18:04 & 18:20 & 18:37 & 19:03 & 19:20 & 19:43 & 20:06 \\
Trenton TC & 16:10 & 16:35 & 16:53 & 17:16 & 17:38 & 17:53 & 18:07 & 18:23 & 18:40 & 19:05 & 19:23 & 19:46 & 20:08 \\
\bottomrule\end{tabular}\end{center}
\subsection*{To Trenton (read down) --- trips 40--43 of 43}
\begin{center}\footnotesize\setlength{\tabcolsep}{3pt}\begin{tabular}{@{}lrrrr@{}}\toprule
Station & 40 & 41 & 42 & 43 \\\midrule
Entertainment Center & 19:30 & 20:00 & 20:23 & 20:46 \\
Aquarium & 19:32 & 20:02 & 20:25 & 20:47 \\
Cooper St/Rutgers & 19:34 & 20:04 & 20:27 & 20:49 \\
Walter Rand TC & 19:37 & 20:08 & 20:31 & 20:53 \\
36th Street & 19:43 & 20:16 & 20:41 & 20:58 \\
Pennsauken TC & 19:44 & 20:17 & 20:43 & 21:00 \\
Route 73/Pennsauken & 19:48 & 20:21 & 20:46 & 21:04 \\
Palmyra & 19:51 & 20:24 & 20:49 & 21:06 \\
Riverton & 19:53 & 20:25 & 20:51 & 21:08 \\
Cinnaminson & 19:54 & 20:27 & 20:53 & 21:10 \\
Riverside & 20:00 & 20:33 & 20:58 & 21:16 \\
Delanco & 20:02 & 20:35 & 21:00 & 21:17 \\
Beverly/Edgewater Park & 20:05 & 20:38 & 21:04 & 21:21 \\
Burlington South & 20:10 & 20:43 & 21:08 & 21:26 \\
Burlington Towne Ctr & 20:12 & 20:44 & 21:10 & 21:27 \\
Florence & 20:18 & 20:51 & 21:16 & 21:34 \\
Roebling & 20:21 & 20:54 & 21:20 & 21:37 \\
Bordentown & 20:26 & 20:59 & 21:25 & 21:42 \\
Cass Street & 20:32 & 21:05 & 21:31 & 21:48 \\
Hamilton Avenue & 20:34 & 21:07 & 21:33 & 21:50 \\
Trenton TC & 20:37 & 21:10 & 21:35 & 21:53 \\
\bottomrule\end{tabular}\end{center}
\subsection*{To Camden (read down) --- trips 1--13 of 44}
\begin{center}\footnotesize\setlength{\tabcolsep}{3pt}\begin{tabular}{@{}lrrrrrrrrrrrrr@{}}\toprule
Station & 1 & 2 & 3 & 4 & 5 & 6 & 7 & 8 & 9 & 10 & 11 & 12 & 13 \\\midrule
Trenton TC & 5:45 & 6:00 & 6:15 & 6:40 & 7:02 & 7:24 & 7:45 & 8:00 & 8:15 & 8:30 & 8:45 & 9:10 & 9:32 \\
Hamilton Avenue & 5:46 & 6:01 & 6:17 & 6:42 & 7:04 & 7:25 & 7:46 & 8:01 & 8:16 & 8:31 & 8:47 & 9:12 & 9:34 \\
Cass Street & 5:48 & 6:03 & 6:19 & 6:44 & 7:05 & 7:27 & 7:48 & 8:03 & 8:18 & 8:33 & 8:49 & 9:14 & 9:35 \\
Bordentown & 5:56 & 6:11 & 6:26 & 6:51 & 7:13 & 7:35 & 7:55 & 8:11 & 8:26 & 8:41 & 8:56 & 9:21 & 9:43 \\
Roebling & 6:03 & 6:16 & 6:31 & 6:57 & 7:18 & 7:42 & 8:01 & 8:16 & 8:33 & 8:46 & 9:01 & 9:27 & 9:48 \\
Florence & 6:06 & 6:19 & 6:34 & 6:59 & 7:21 & 7:45 & 8:04 & 8:19 & 8:36 & 8:49 & 9:04 & 9:29 & 9:51 \\
Burlington Towne Ctr & 6:11 & 6:24 & 6:39 & 7:05 & 7:26 & 7:51 & 8:09 & 8:24 & 8:41 & 8:54 & 9:09 & 9:35 & 9:56 \\
Burlington South & 6:13 & 6:26 & 6:41 & 7:07 & 7:28 & 7:53 & 8:11 & 8:26 & 8:43 & 8:56 & 9:11 & 9:37 & 9:58 \\
Beverly/Edgewater Park & 6:18 & 6:31 & 6:45 & 7:11 & 7:32 & 7:57 & 8:15 & 8:30 & 8:48 & 9:01 & 9:15 & 9:41 & 10:02 \\
Delanco & 6:21 & 6:34 & 6:49 & 7:15 & 7:36 & 8:01 & 8:19 & 8:34 & 8:51 & 9:04 & 9:19 & 9:45 & 10:06 \\
Riverside & 6:23 & 6:36 & 6:51 & 7:17 & 7:38 & 8:03 & 8:21 & 8:36 & 8:53 & 9:06 & 9:21 & 9:47 & 10:08 \\
Cinnaminson & 6:28 & 6:42 & 6:56 & 7:22 & 7:43 & 8:08 & 8:26 & 8:41 & 8:58 & 9:12 & 9:26 & 9:52 & 10:13 \\
Riverton & 6:30 & 6:44 & 6:58 & 7:24 & 7:45 & 8:10 & 8:28 & 8:43 & 9:00 & 9:14 & 9:28 & 9:54 & 10:15 \\
Palmyra & 6:33 & 6:46 & 7:01 & 7:27 & 7:47 & 8:12 & 8:30 & 8:47 & 9:03 & 9:16 & 9:31 & 9:57 & 10:17 \\
Route 73/Pennsauken & 6:36 & 6:49 & 7:04 & 7:30 & 7:50 & 8:15 & 8:33 & 8:50 & 9:06 & 9:19 & 9:34 & 10:00 & 10:20 \\
Pennsauken TC & 6:37 & 6:51 & 7:05 & 7:31 & 7:52 & 8:17 & 8:35 & 8:52 & 9:07 & 9:21 & 9:35 & 10:01 & 10:22 \\
36th Street & 6:42 & 6:55 & 7:10 & 7:36 & 7:56 & 8:21 & 8:39 & 8:56 & 9:12 & 9:25 & 9:40 & 10:06 & 10:26 \\
Walter Rand TC & 6:47 & 7:00 & 7:14 & 7:40 & 8:01 & 8:26 & 8:44 & 9:01 & 9:17 & 9:30 & 9:44 & 10:10 & 10:31 \\
Cooper St/Rutgers & 6:49 & 7:03 & 7:17 & 7:43 & 8:04 & 8:29 & 8:47 & 9:04 & 9:19 & 9:33 & 9:47 & 10:13 & 10:34 \\
Aquarium & 6:51 & 7:04 & 7:19 & 7:45 & 8:06 & 8:31 & 8:49 & 9:06 & 9:21 & 9:34 & 9:49 & 10:15 & 10:36 \\
Entertainment Center & 6:54 & 7:07 & 7:22 & 7:48 & 8:09 & 8:34 & 8:52 & 9:08 & 9:24 & 9:37 & 9:52 & 10:18 & 10:39 \\
\bottomrule\end{tabular}\end{center}
\subsection*{To Camden (read down) --- trips 14--26 of 44}
\begin{center}\footnotesize\setlength{\tabcolsep}{3pt}\begin{tabular}{@{}lrrrrrrrrrrrrr@{}}\toprule
Station & 14 & 15 & 16 & 17 & 18 & 19 & 20 & 21 & 22 & 23 & 24 & 25 & 26 \\\midrule
Trenton TC & 9:54 & 10:15 & 10:45 & 11:15 & 11:40 & 12:02 & 12:24 & 12:45 & 13:15 & 13:45 & 14:10 & 14:32 & 14:54 \\
Hamilton Avenue & 9:55 & 10:16 & 10:46 & 11:17 & 11:42 & 12:04 & 12:25 & 12:46 & 13:16 & 13:47 & 14:12 & 14:34 & 14:55 \\
Cass Street & 9:57 & 10:18 & 10:48 & 11:19 & 11:44 & 12:05 & 12:27 & 12:48 & 13:18 & 13:49 & 14:14 & 14:35 & 14:57 \\
Bordentown & 10:05 & 10:25 & 10:56 & 11:26 & 11:51 & 12:13 & 12:35 & 12:55 & 13:26 & 13:56 & 14:21 & 14:43 & 15:05 \\
Roebling & 10:12 & 10:31 & 11:03 & 11:31 & 11:57 & 12:18 & 12:42 & 13:01 & 13:33 & 14:01 & 14:27 & 14:48 & 15:12 \\
Florence & 10:15 & 10:34 & 11:06 & 11:34 & 11:59 & 12:21 & 12:45 & 13:04 & 13:36 & 14:04 & 14:29 & 14:51 & 15:15 \\
Burlington Towne Ctr & 10:21 & 10:39 & 11:11 & 11:39 & 12:05 & 12:26 & 12:51 & 13:09 & 13:41 & 14:09 & 14:35 & 14:56 & 15:21 \\
Burlington South & 10:23 & 10:41 & 11:13 & 11:41 & 12:07 & 12:28 & 12:53 & 13:11 & 13:43 & 14:11 & 14:37 & 14:58 & 15:23 \\
Beverly/Edgewater Park & 10:27 & 10:45 & 11:18 & 11:45 & 12:11 & 12:32 & 12:57 & 13:15 & 13:48 & 14:15 & 14:41 & 15:02 & 15:27 \\
Delanco & 10:31 & 10:49 & 11:21 & 11:49 & 12:15 & 12:36 & 13:01 & 13:19 & 13:51 & 14:19 & 14:45 & 15:06 & 15:31 \\
Riverside & 10:33 & 10:51 & 11:23 & 11:51 & 12:17 & 12:38 & 13:03 & 13:21 & 13:53 & 14:21 & 14:47 & 15:08 & 15:33 \\
Cinnaminson & 10:38 & 10:56 & 11:28 & 11:56 & 12:22 & 12:43 & 13:08 & 13:26 & 13:58 & 14:26 & 14:52 & 15:13 & 15:38 \\
Riverton & 10:40 & 10:58 & 11:30 & 11:58 & 12:24 & 12:45 & 13:10 & 13:28 & 14:00 & 14:28 & 14:54 & 15:15 & 15:40 \\
Palmyra & 10:42 & 11:00 & 11:33 & 12:01 & 12:27 & 12:47 & 13:12 & 13:30 & 14:03 & 14:31 & 14:57 & 15:17 & 15:42 \\
Route 73/Pennsauken & 10:45 & 11:03 & 11:36 & 12:04 & 12:30 & 12:50 & 13:15 & 13:33 & 14:06 & 14:34 & 15:00 & 15:20 & 15:45 \\
Pennsauken TC & 10:47 & 11:05 & 11:37 & 12:05 & 12:31 & 12:52 & 13:17 & 13:35 & 14:07 & 14:35 & 15:01 & 15:22 & 15:47 \\
36th Street & 10:51 & 11:09 & 11:42 & 12:10 & 12:36 & 12:56 & 13:21 & 13:39 & 14:12 & 14:40 & 15:06 & 15:26 & 15:51 \\
Walter Rand TC & 10:56 & 11:14 & 11:47 & 12:14 & 12:40 & 13:01 & 13:26 & 13:44 & 14:17 & 14:44 & 15:10 & 15:31 & 15:56 \\
Cooper St/Rutgers & 10:59 & 11:17 & 11:49 & 12:17 & 12:43 & 13:04 & 13:29 & 13:47 & 14:19 & 14:47 & 15:13 & 15:34 & 15:59 \\
Aquarium & 11:01 & 11:19 & 11:51 & 12:19 & 12:45 & 13:06 & 13:31 & 13:49 & 14:21 & 14:49 & 15:15 & 15:36 & 16:01 \\
Entertainment Center & 11:04 & 11:22 & 11:54 & 12:22 & 12:48 & 13:09 & 13:34 & 13:52 & 14:24 & 14:52 & 15:18 & 15:39 & 16:04 \\
\bottomrule\end{tabular}\end{center}
\subsection*{To Camden (read down) --- trips 27--39 of 44}
\begin{center}\footnotesize\setlength{\tabcolsep}{3pt}\begin{tabular}{@{}lrrrrrrrrrrrrr@{}}\toprule
Station & 27 & 28 & 29 & 30 & 31 & 32 & 33 & 34 & 35 & 36 & 37 & 38 & 39 \\\midrule
Trenton TC & 15:15 & 15:30 & 15:45 & 16:00 & 16:15 & 16:40 & 17:02 & 17:24 & 17:45 & 18:00 & 18:15 & 18:30 & 18:45 \\
Hamilton Avenue & 15:16 & 15:31 & 15:46 & 16:01 & 16:17 & 16:42 & 17:04 & 17:25 & 17:46 & 18:01 & 18:16 & 18:31 & 18:47 \\
Cass Street & 15:18 & 15:33 & 15:48 & 16:03 & 16:19 & 16:44 & 17:05 & 17:27 & 17:48 & 18:03 & 18:18 & 18:33 & 18:49 \\
Bordentown & 15:25 & 15:41 & 15:56 & 16:11 & 16:26 & 16:51 & 17:13 & 17:35 & 17:55 & 18:11 & 18:26 & 18:41 & 18:56 \\
Roebling & 15:31 & 15:46 & 16:03 & 16:16 & 16:31 & 16:57 & 17:18 & 17:42 & 18:01 & 18:16 & 18:33 & 18:46 & 19:01 \\
Florence & 15:34 & 15:49 & 16:06 & 16:19 & 16:34 & 16:59 & 17:21 & 17:45 & 18:04 & 18:19 & 18:36 & 18:49 & 19:04 \\
Burlington Towne Ctr & 15:39 & 15:54 & 16:11 & 16:24 & 16:39 & 17:05 & 17:26 & 17:51 & 18:09 & 18:24 & 18:41 & 18:54 & 19:09 \\
Burlington South & 15:41 & 15:56 & 16:13 & 16:26 & 16:41 & 17:07 & 17:28 & 17:53 & 18:11 & 18:26 & 18:43 & 18:56 & 19:11 \\
Beverly/Edgewater Park & 15:45 & 16:00 & 16:18 & 16:31 & 16:45 & 17:11 & 17:32 & 17:57 & 18:15 & 18:30 & 18:48 & 19:01 & 19:15 \\
Delanco & 15:49 & 16:04 & 16:21 & 16:34 & 16:49 & 17:15 & 17:36 & 18:01 & 18:19 & 18:34 & 18:51 & 19:04 & 19:19 \\
Riverside & 15:51 & 16:06 & 16:23 & 16:36 & 16:51 & 17:17 & 17:38 & 18:03 & 18:21 & 18:36 & 18:53 & 19:06 & 19:21 \\
Cinnaminson & 15:56 & 16:11 & 16:28 & 16:42 & 16:56 & 17:22 & 17:43 & 18:08 & 18:26 & 18:41 & 18:58 & 19:12 & 19:26 \\
Riverton & 15:58 & 16:13 & 16:30 & 16:44 & 16:58 & 17:24 & 17:45 & 18:10 & 18:28 & 18:43 & 19:00 & 19:14 & 19:28 \\
Palmyra & 16:00 & 16:17 & 16:33 & 16:46 & 17:01 & 17:27 & 17:47 & 18:12 & 18:30 & 18:47 & 19:03 & 19:16 & 19:31 \\
Route 73/Pennsauken & 16:03 & 16:20 & 16:36 & 16:49 & 17:04 & 17:30 & 17:50 & 18:15 & 18:33 & 18:50 & 19:06 & 19:19 & 19:34 \\
Pennsauken TC & 16:05 & 16:22 & 16:37 & 16:51 & 17:05 & 17:31 & 17:52 & 18:17 & 18:35 & 18:52 & 19:07 & 19:21 & 19:35 \\
36th Street & 16:09 & 16:26 & 16:42 & 16:55 & 17:10 & 17:36 & 17:56 & 18:21 & 18:39 & 18:56 & 19:12 & 19:25 & 19:40 \\
Walter Rand TC & 16:14 & 16:31 & 16:47 & 17:00 & 17:14 & 17:40 & 18:01 & 18:26 & 18:44 & 19:01 & 19:17 & 19:30 & 19:44 \\
Cooper St/Rutgers & 16:17 & 16:34 & 16:49 & 17:03 & 17:17 & 17:43 & 18:04 & 18:29 & 18:47 & 19:04 & 19:19 & 19:33 & 19:47 \\
Aquarium & 16:19 & 16:36 & 16:51 & 17:04 & 17:19 & 17:45 & 18:06 & 18:31 & 18:49 & 19:06 & 19:21 & 19:34 & 19:49 \\
Entertainment Center & 16:22 & 16:38 & 16:54 & 17:07 & 17:22 & 17:48 & 18:09 & 18:34 & 18:52 & 19:08 & 19:24 & 19:37 & 19:52 \\
\bottomrule\end{tabular}\end{center}
\subsection*{To Camden (read down) --- trips 40--44 of 44}
\begin{center}\footnotesize\setlength{\tabcolsep}{3pt}\begin{tabular}{@{}lrrrrr@{}}\toprule
Station & 40 & 41 & 42 & 43 & 44 \\\midrule
Trenton TC & 19:10 & 19:32 & 19:54 & 20:15 & 20:45 \\
Hamilton Avenue & 19:12 & 19:34 & 19:55 & 20:16 & 20:46 \\
Cass Street & 19:14 & 19:35 & 19:57 & 20:18 & 20:48 \\
Bordentown & 19:21 & 19:43 & 20:05 & 20:25 & 20:56 \\
Roebling & 19:27 & 19:48 & 20:12 & 20:31 & 21:03 \\
Florence & 19:29 & 19:51 & 20:15 & 20:34 & 21:06 \\
Burlington Towne Ctr & 19:35 & 19:56 & 20:21 & 20:39 & 21:11 \\
Burlington South & 19:37 & 19:58 & 20:23 & 20:41 & 21:13 \\
Beverly/Edgewater Park & 19:41 & 20:02 & 20:27 & 20:45 & 21:18 \\
Delanco & 19:45 & 20:06 & 20:31 & 20:49 & 21:21 \\
Riverside & 19:47 & 20:08 & 20:33 & 20:51 & 21:23 \\
Cinnaminson & 19:52 & 20:13 & 20:38 & 20:56 & 21:28 \\
Riverton & 19:54 & 20:15 & 20:40 & 20:58 & 21:30 \\
Palmyra & 19:57 & 20:17 & 20:42 & 21:00 & 21:33 \\
Route 73/Pennsauken & 20:00 & 20:20 & 20:45 & 21:03 & 21:36 \\
Pennsauken TC & 20:01 & 20:22 & 20:47 & 21:05 & 21:37 \\
36th Street & 20:06 & 20:26 & 20:51 & 21:09 & 21:42 \\
Walter Rand TC & 20:10 & 20:31 & 20:56 & 21:14 & 21:47 \\
Cooper St/Rutgers & 20:13 & 20:34 & 20:59 & 21:17 & 21:49 \\
Aquarium & 20:15 & 20:36 & 21:01 & 21:19 & 21:51 \\
Entertainment Center & 20:18 & 20:39 & 21:04 & 21:22 & 21:54 \\
\bottomrule\end{tabular}\end{center}
\clearpage
\section{Proposal R8: 8 trains, 160-min cycle, 8-min layovers}\label{tt:proposal_R8}
\subsection*{To Trenton (read down) --- trips 1--13 of 41}
\begin{center}\footnotesize\setlength{\tabcolsep}{3pt}\begin{tabular}{@{}lrrrrrrrrrrrrr@{}}\toprule
Station & 1 & 2 & 3 & 4 & 5 & 6 & 7 & 8 & 9 & 10 & 11 & 12 & 13 \\\midrule
Entertainment Center & 5:45 & 6:08 & 6:31 & 6:55 & 7:10 & 7:27 & 7:44 & 8:00 & 8:25 & 8:48 & 9:11 & 9:35 & 10:07 \\
Aquarium & 5:47 & 6:10 & 6:33 & 6:57 & 7:12 & 7:29 & 7:46 & 8:02 & 8:27 & 8:50 & 9:13 & 9:37 & 10:09 \\
Cooper St/Rutgers & 5:49 & 6:12 & 6:35 & 6:59 & 7:14 & 7:31 & 7:48 & 8:04 & 8:29 & 8:52 & 9:15 & 9:39 & 10:11 \\
Walter Rand TC & 5:53 & 6:16 & 6:39 & 7:03 & 7:18 & 7:35 & 7:52 & 8:08 & 8:33 & 8:56 & 9:19 & 9:43 & 10:15 \\
36th Street & 5:58 & 6:21 & 6:44 & 7:08 & 7:23 & 7:46 & 7:58 & 8:13 & 8:38 & 9:01 & 9:24 & 9:48 & 10:26 \\
Pennsauken TC & 6:00 & 6:23 & 6:46 & 7:10 & 7:25 & 7:48 & 8:00 & 8:15 & 8:40 & 9:03 & 9:26 & 9:50 & 10:28 \\
Route 73/Pennsauken & 6:03 & 6:26 & 6:49 & 7:14 & 7:28 & 7:51 & 8:03 & 8:18 & 8:43 & 9:06 & 9:29 & 9:54 & 10:31 \\
Palmyra & 6:06 & 6:29 & 6:52 & 7:17 & 7:31 & 7:54 & 8:06 & 8:21 & 8:46 & 9:09 & 9:32 & 9:57 & 10:34 \\
Riverton & 6:08 & 6:31 & 6:54 & 7:19 & 7:33 & 7:56 & 8:08 & 8:23 & 8:48 & 9:11 & 9:34 & 9:59 & 10:36 \\
Cinnaminson & 6:10 & 6:33 & 6:56 & 7:21 & 7:35 & 7:58 & 8:10 & 8:25 & 8:50 & 9:13 & 9:36 & 10:01 & 10:38 \\
Riverside & 6:15 & 6:38 & 7:02 & 7:26 & 7:40 & 8:03 & 8:15 & 8:30 & 8:55 & 9:18 & 9:42 & 10:06 & 10:43 \\
Delanco & 6:17 & 6:40 & 7:04 & 7:28 & 7:42 & 8:05 & 8:17 & 8:32 & 8:57 & 9:20 & 9:44 & 10:08 & 10:45 \\
Beverly/Edgewater Park & 6:21 & 6:44 & 7:08 & 7:32 & 7:46 & 8:09 & 8:21 & 8:36 & 9:01 & 9:24 & 9:48 & 10:12 & 10:49 \\
Burlington South & 6:25 & 6:48 & 7:13 & 7:36 & 7:50 & 8:13 & 8:25 & 8:40 & 9:05 & 9:28 & 9:53 & 10:16 & 10:53 \\
Burlington Towne Ctr & 6:27 & 6:50 & 7:15 & 7:38 & 7:52 & 8:15 & 8:27 & 8:42 & 9:07 & 9:30 & 9:55 & 10:18 & 10:55 \\
Florence & 6:33 & 6:56 & 7:21 & 7:44 & 7:58 & 8:21 & 8:33 & 8:48 & 9:13 & 9:36 & 10:01 & 10:24 & 11:01 \\
Roebling & 6:36 & 6:59 & 7:24 & 7:47 & 8:01 & 8:24 & 8:36 & 8:51 & 9:16 & 9:39 & 10:04 & 10:27 & 11:04 \\
Bordentown & 6:41 & 7:04 & 7:29 & 7:52 & 8:06 & 8:29 & 8:41 & 8:56 & 9:21 & 9:44 & 10:09 & 10:32 & 11:09 \\
Cass Street & 6:47 & 7:10 & 7:35 & 7:58 & 8:12 & 8:35 & 8:47 & 9:04 & 9:27 & 9:50 & 10:15 & 10:38 & 11:15 \\
Hamilton Avenue & 6:49 & 7:12 & 7:37 & 8:00 & 8:14 & 8:37 & 8:49 & 9:06 & 9:29 & 9:52 & 10:17 & 10:40 & 11:17 \\
Trenton TC & 6:52 & 7:15 & 7:40 & 8:03 & 8:17 & 8:40 & 8:52 & 9:09 & 9:32 & 9:55 & 10:20 & 10:43 & 11:20 \\
\bottomrule\end{tabular}\end{center}
\subsection*{To Trenton (read down) --- trips 14--26 of 41}
\begin{center}\footnotesize\setlength{\tabcolsep}{3pt}\begin{tabular}{@{}lrrrrrrrrrrrrr@{}}\toprule
Station & 14 & 15 & 16 & 17 & 18 & 19 & 20 & 21 & 22 & 23 & 24 & 25 & 26 \\\midrule
Entertainment Center & 10:40 & 11:05 & 11:28 & 11:51 & 12:15 & 12:47 & 13:20 & 13:45 & 14:08 & 14:31 & 14:55 & 15:10 & 15:27 \\
Aquarium & 10:42 & 11:07 & 11:30 & 11:53 & 12:17 & 12:49 & 13:22 & 13:47 & 14:10 & 14:33 & 14:57 & 15:12 & 15:29 \\
Cooper St/Rutgers & 10:44 & 11:09 & 11:32 & 11:55 & 12:19 & 12:51 & 13:24 & 13:49 & 14:12 & 14:35 & 14:59 & 15:14 & 15:31 \\
Walter Rand TC & 10:48 & 11:13 & 11:36 & 11:59 & 12:23 & 12:55 & 13:28 & 13:53 & 14:16 & 14:39 & 15:03 & 15:18 & 15:35 \\
36th Street & 10:53 & 11:18 & 11:41 & 12:04 & 12:28 & 13:06 & 13:33 & 13:58 & 14:21 & 14:44 & 15:08 & 15:23 & 15:46 \\
Pennsauken TC & 10:55 & 11:20 & 11:43 & 12:06 & 12:30 & 13:08 & 13:35 & 14:00 & 14:23 & 14:46 & 15:10 & 15:25 & 15:48 \\
Route 73/Pennsauken & 10:58 & 11:23 & 11:46 & 12:09 & 12:34 & 13:11 & 13:38 & 14:03 & 14:26 & 14:49 & 15:14 & 15:28 & 15:51 \\
Palmyra & 11:01 & 11:26 & 11:49 & 12:12 & 12:37 & 13:14 & 13:41 & 14:06 & 14:29 & 14:52 & 15:17 & 15:31 & 15:54 \\
Riverton & 11:03 & 11:28 & 11:51 & 12:14 & 12:39 & 13:16 & 13:43 & 14:08 & 14:31 & 14:54 & 15:19 & 15:33 & 15:56 \\
Cinnaminson & 11:05 & 11:30 & 11:53 & 12:16 & 12:41 & 13:18 & 13:45 & 14:10 & 14:33 & 14:56 & 15:21 & 15:35 & 15:58 \\
Riverside & 11:10 & 11:35 & 11:58 & 12:22 & 12:46 & 13:23 & 13:50 & 14:15 & 14:38 & 15:02 & 15:26 & 15:40 & 16:03 \\
Delanco & 11:12 & 11:37 & 12:00 & 12:24 & 12:48 & 13:25 & 13:52 & 14:17 & 14:40 & 15:04 & 15:28 & 15:42 & 16:05 \\
Beverly/Edgewater Park & 11:16 & 11:41 & 12:04 & 12:28 & 12:52 & 13:29 & 13:56 & 14:21 & 14:44 & 15:08 & 15:32 & 15:46 & 16:09 \\
Burlington South & 11:20 & 11:45 & 12:08 & 12:33 & 12:56 & 13:33 & 14:00 & 14:25 & 14:48 & 15:13 & 15:36 & 15:50 & 16:13 \\
Burlington Towne Ctr & 11:22 & 11:47 & 12:10 & 12:35 & 12:58 & 13:35 & 14:02 & 14:27 & 14:50 & 15:15 & 15:38 & 15:52 & 16:15 \\
Florence & 11:28 & 11:53 & 12:16 & 12:41 & 13:04 & 13:41 & 14:08 & 14:33 & 14:56 & 15:21 & 15:44 & 15:58 & 16:21 \\
Roebling & 11:31 & 11:56 & 12:19 & 12:44 & 13:07 & 13:44 & 14:11 & 14:36 & 14:59 & 15:24 & 15:47 & 16:01 & 16:24 \\
Bordentown & 11:36 & 12:01 & 12:24 & 12:49 & 13:12 & 13:49 & 14:16 & 14:41 & 15:04 & 15:29 & 15:52 & 16:06 & 16:29 \\
Cass Street & 11:44 & 12:07 & 12:30 & 12:55 & 13:18 & 13:55 & 14:24 & 14:47 & 15:10 & 15:35 & 15:58 & 16:12 & 16:35 \\
Hamilton Avenue & 11:46 & 12:09 & 12:32 & 12:57 & 13:20 & 13:57 & 14:26 & 14:49 & 15:12 & 15:37 & 16:00 & 16:14 & 16:37 \\
Trenton TC & 11:49 & 12:12 & 12:35 & 13:00 & 13:23 & 14:00 & 14:29 & 14:52 & 15:15 & 15:40 & 16:03 & 16:17 & 16:40 \\
\bottomrule\end{tabular}\end{center}
\subsection*{To Trenton (read down) --- trips 27--39 of 41}
\begin{center}\footnotesize\setlength{\tabcolsep}{3pt}\begin{tabular}{@{}lrrrrrrrrrrrrr@{}}\toprule
Station & 27 & 28 & 29 & 30 & 31 & 32 & 33 & 34 & 35 & 36 & 37 & 38 & 39 \\\midrule
Entertainment Center & 15:44 & 16:00 & 16:25 & 16:48 & 17:11 & 17:35 & 17:50 & 18:07 & 18:24 & 18:40 & 19:05 & 19:28 & 19:51 \\
Aquarium & 15:46 & 16:02 & 16:27 & 16:50 & 17:13 & 17:37 & 17:52 & 18:09 & 18:26 & 18:42 & 19:07 & 19:30 & 19:53 \\
Cooper St/Rutgers & 15:48 & 16:04 & 16:29 & 16:52 & 17:15 & 17:39 & 17:54 & 18:11 & 18:28 & 18:44 & 19:09 & 19:32 & 19:55 \\
Walter Rand TC & 15:52 & 16:08 & 16:33 & 16:56 & 17:19 & 17:43 & 17:58 & 18:15 & 18:32 & 18:48 & 19:13 & 19:36 & 19:59 \\
36th Street & 15:58 & 16:13 & 16:38 & 17:01 & 17:24 & 17:48 & 18:03 & 18:26 & 18:38 & 18:53 & 19:18 & 19:41 & 20:04 \\
Pennsauken TC & 16:00 & 16:15 & 16:40 & 17:03 & 17:26 & 17:50 & 18:05 & 18:28 & 18:40 & 18:55 & 19:20 & 19:43 & 20:06 \\
Route 73/Pennsauken & 16:03 & 16:18 & 16:43 & 17:06 & 17:29 & 17:54 & 18:08 & 18:31 & 18:43 & 18:58 & 19:23 & 19:46 & 20:09 \\
Palmyra & 16:06 & 16:21 & 16:46 & 17:09 & 17:32 & 17:57 & 18:11 & 18:34 & 18:46 & 19:01 & 19:26 & 19:49 & 20:12 \\
Riverton & 16:08 & 16:23 & 16:48 & 17:11 & 17:34 & 17:59 & 18:13 & 18:36 & 18:48 & 19:03 & 19:28 & 19:51 & 20:14 \\
Cinnaminson & 16:10 & 16:25 & 16:50 & 17:13 & 17:36 & 18:01 & 18:15 & 18:38 & 18:50 & 19:05 & 19:30 & 19:53 & 20:16 \\
Riverside & 16:15 & 16:30 & 16:55 & 17:18 & 17:42 & 18:06 & 18:20 & 18:43 & 18:55 & 19:10 & 19:35 & 19:58 & 20:22 \\
Delanco & 16:17 & 16:32 & 16:57 & 17:20 & 17:44 & 18:08 & 18:22 & 18:45 & 18:57 & 19:12 & 19:37 & 20:00 & 20:24 \\
Beverly/Edgewater Park & 16:21 & 16:36 & 17:01 & 17:24 & 17:48 & 18:12 & 18:26 & 18:49 & 19:01 & 19:16 & 19:41 & 20:04 & 20:28 \\
Burlington South & 16:25 & 16:40 & 17:05 & 17:28 & 17:53 & 18:16 & 18:30 & 18:53 & 19:05 & 19:20 & 19:45 & 20:08 & 20:33 \\
Burlington Towne Ctr & 16:27 & 16:42 & 17:07 & 17:30 & 17:55 & 18:18 & 18:32 & 18:55 & 19:07 & 19:22 & 19:47 & 20:10 & 20:35 \\
Florence & 16:33 & 16:48 & 17:13 & 17:36 & 18:01 & 18:24 & 18:38 & 19:01 & 19:13 & 19:28 & 19:53 & 20:16 & 20:41 \\
Roebling & 16:36 & 16:51 & 17:16 & 17:39 & 18:04 & 18:27 & 18:41 & 19:04 & 19:16 & 19:31 & 19:56 & 20:19 & 20:44 \\
Bordentown & 16:41 & 16:56 & 17:21 & 17:44 & 18:09 & 18:32 & 18:46 & 19:09 & 19:21 & 19:36 & 20:01 & 20:24 & 20:49 \\
Cass Street & 16:47 & 17:04 & 17:27 & 17:50 & 18:15 & 18:38 & 18:52 & 19:15 & 19:27 & 19:44 & 20:07 & 20:30 & 20:55 \\
Hamilton Avenue & 16:49 & 17:06 & 17:29 & 17:52 & 18:17 & 18:40 & 18:54 & 19:17 & 19:29 & 19:46 & 20:09 & 20:32 & 20:57 \\
Trenton TC & 16:52 & 17:09 & 17:32 & 17:55 & 18:20 & 18:43 & 18:57 & 19:20 & 19:32 & 19:49 & 20:12 & 20:35 & 21:00 \\
\bottomrule\end{tabular}\end{center}
\subsection*{To Trenton (read down) --- trips 40--41 of 41}
\begin{center}\footnotesize\setlength{\tabcolsep}{3pt}\begin{tabular}{@{}lrr@{}}\toprule
Station & 40 & 41 \\\midrule
Entertainment Center & 20:15 & 20:47 \\
Aquarium & 20:17 & 20:49 \\
Cooper St/Rutgers & 20:19 & 20:51 \\
Walter Rand TC & 20:23 & 20:55 \\
36th Street & 20:28 & 21:06 \\
Pennsauken TC & 20:30 & 21:08 \\
Route 73/Pennsauken & 20:34 & 21:11 \\
Palmyra & 20:37 & 21:14 \\
Riverton & 20:39 & 21:16 \\
Cinnaminson & 20:41 & 21:18 \\
Riverside & 20:46 & 21:23 \\
Delanco & 20:48 & 21:25 \\
Beverly/Edgewater Park & 20:52 & 21:29 \\
Burlington South & 20:56 & 21:33 \\
Burlington Towne Ctr & 20:58 & 21:35 \\
Florence & 21:04 & 21:41 \\
Roebling & 21:07 & 21:44 \\
Bordentown & 21:12 & 21:49 \\
Cass Street & 21:18 & 21:55 \\
Hamilton Avenue & 21:20 & 21:57 \\
Trenton TC & 21:23 & 22:00 \\
\bottomrule\end{tabular}\end{center}
\subsection*{To Camden (read down) --- trips 1--13 of 42}
\begin{center}\footnotesize\setlength{\tabcolsep}{3pt}\begin{tabular}{@{}lrrrrrrrrrrrrr@{}}\toprule
Station & 1 & 2 & 3 & 4 & 5 & 6 & 7 & 8 & 9 & 10 & 11 & 12 & 13 \\\midrule
Trenton TC & 5:51 & 6:08 & 6:24 & 6:39 & 7:02 & 7:27 & 7:51 & 8:16 & 8:31 & 8:48 & 9:04 & 9:19 & 9:42 \\
Hamilton Avenue & 5:53 & 6:10 & 6:26 & 6:41 & 7:04 & 7:29 & 7:53 & 8:18 & 8:33 & 8:50 & 9:06 & 9:21 & 9:44 \\
Cass Street & 5:55 & 6:12 & 6:28 & 6:43 & 7:06 & 7:31 & 7:55 & 8:20 & 8:35 & 8:52 & 9:08 & 9:23 & 9:46 \\
Bordentown & 6:02 & 6:19 & 6:35 & 6:50 & 7:13 & 7:38 & 8:02 & 8:27 & 8:42 & 8:59 & 9:15 & 9:30 & 9:53 \\
Roebling & 6:07 & 6:24 & 6:41 & 6:55 & 7:18 & 7:43 & 8:07 & 8:32 & 8:47 & 9:04 & 9:21 & 9:35 & 9:58 \\
Florence & 6:10 & 6:27 & 6:44 & 6:58 & 7:21 & 7:46 & 8:10 & 8:35 & 8:50 & 9:07 & 9:24 & 9:38 & 10:01 \\
Burlington Towne Ctr & 6:15 & 6:36 & 6:49 & 7:03 & 7:26 & 7:51 & 8:15 & 8:40 & 8:55 & 9:16 & 9:29 & 9:43 & 10:06 \\
Burlington South & 6:17 & 6:38 & 6:51 & 7:05 & 7:28 & 7:53 & 8:17 & 8:42 & 8:57 & 9:18 & 9:31 & 9:45 & 10:08 \\
Beverly/Edgewater Park & 6:21 & 6:42 & 6:55 & 7:09 & 7:32 & 7:57 & 8:21 & 8:46 & 9:01 & 9:22 & 9:35 & 9:49 & 10:12 \\
Delanco & 6:25 & 6:46 & 6:59 & 7:13 & 7:36 & 8:01 & 8:25 & 8:50 & 9:05 & 9:26 & 9:39 & 9:53 & 10:16 \\
Riverside & 6:27 & 6:48 & 7:01 & 7:15 & 7:38 & 8:03 & 8:27 & 8:52 & 9:07 & 9:28 & 9:41 & 9:55 & 10:18 \\
Cinnaminson & 6:36 & 6:53 & 7:06 & 7:20 & 7:43 & 8:08 & 8:33 & 8:58 & 9:16 & 9:33 & 9:46 & 10:00 & 10:23 \\
Riverton & 6:38 & 6:55 & 7:08 & 7:22 & 7:45 & 8:10 & 8:35 & 9:00 & 9:18 & 9:35 & 9:48 & 10:02 & 10:25 \\
Palmyra & 6:40 & 6:57 & 7:10 & 7:24 & 7:47 & 8:12 & 8:37 & 9:02 & 9:20 & 9:37 & 9:50 & 10:04 & 10:27 \\
Route 73/Pennsauken & 6:43 & 7:00 & 7:13 & 7:27 & 7:50 & 8:15 & 8:40 & 9:05 & 9:23 & 9:40 & 9:53 & 10:07 & 10:30 \\
Pennsauken TC & 6:45 & 7:02 & 7:15 & 7:29 & 7:52 & 8:17 & 8:42 & 9:07 & 9:25 & 9:42 & 9:55 & 10:09 & 10:32 \\
36th Street & 6:49 & 7:06 & 7:19 & 7:33 & 7:56 & 8:21 & 8:46 & 9:11 & 9:29 & 9:46 & 9:59 & 10:13 & 10:36 \\
Walter Rand TC & 6:54 & 7:11 & 7:24 & 7:38 & 8:01 & 8:26 & 8:51 & 9:16 & 9:34 & 9:51 & 10:04 & 10:18 & 10:41 \\
Cooper St/Rutgers & 6:57 & 7:14 & 7:27 & 7:41 & 8:04 & 8:29 & 8:54 & 9:19 & 9:37 & 9:54 & 10:07 & 10:21 & 10:44 \\
Aquarium & 6:59 & 7:16 & 7:29 & 7:43 & 8:06 & 8:31 & 8:56 & 9:21 & 9:39 & 9:56 & 10:09 & 10:23 & 10:46 \\
Entertainment Center & 7:02 & 7:19 & 7:32 & 7:46 & 8:09 & 8:34 & 8:59 & 9:24 & 9:42 & 9:59 & 10:12 & 10:26 & 10:49 \\
\bottomrule\end{tabular}\end{center}
\subsection*{To Camden (read down) --- trips 14--26 of 42}
\begin{center}\footnotesize\setlength{\tabcolsep}{3pt}\begin{tabular}{@{}lrrrrrrrrrrrrr@{}}\toprule
Station & 14 & 15 & 16 & 17 & 18 & 19 & 20 & 21 & 22 & 23 & 24 & 25 & 26 \\\midrule
Trenton TC & 10:07 & 10:31 & 10:56 & 11:28 & 11:59 & 12:22 & 12:47 & 13:11 & 13:36 & 14:08 & 14:39 & 15:02 & 15:27 \\
Hamilton Avenue & 10:09 & 10:33 & 10:58 & 11:30 & 12:01 & 12:24 & 12:49 & 13:13 & 13:38 & 14:10 & 14:41 & 15:04 & 15:29 \\
Cass Street & 10:11 & 10:35 & 11:00 & 11:32 & 12:03 & 12:26 & 12:51 & 13:15 & 13:40 & 14:12 & 14:43 & 15:06 & 15:31 \\
Bordentown & 10:18 & 10:42 & 11:07 & 11:39 & 12:10 & 12:33 & 12:58 & 13:22 & 13:47 & 14:19 & 14:50 & 15:13 & 15:38 \\
Roebling & 10:23 & 10:47 & 11:12 & 11:44 & 12:15 & 12:38 & 13:03 & 13:27 & 13:52 & 14:24 & 14:55 & 15:18 & 15:43 \\
Florence & 10:26 & 10:50 & 11:15 & 11:47 & 12:18 & 12:41 & 13:06 & 13:30 & 13:55 & 14:27 & 14:58 & 15:21 & 15:46 \\
Burlington Towne Ctr & 10:31 & 10:55 & 11:20 & 11:56 & 12:23 & 12:46 & 13:11 & 13:35 & 14:00 & 14:36 & 15:03 & 15:26 & 15:51 \\
Burlington South & 10:33 & 10:57 & 11:22 & 11:58 & 12:25 & 12:48 & 13:13 & 13:37 & 14:02 & 14:38 & 15:05 & 15:28 & 15:53 \\
Beverly/Edgewater Park & 10:37 & 11:01 & 11:26 & 12:02 & 12:29 & 12:52 & 13:17 & 13:41 & 14:06 & 14:42 & 15:09 & 15:32 & 15:57 \\
Delanco & 10:41 & 11:05 & 11:30 & 12:06 & 12:33 & 12:56 & 13:21 & 13:45 & 14:10 & 14:46 & 15:13 & 15:36 & 16:01 \\
Riverside & 10:43 & 11:07 & 11:32 & 12:08 & 12:35 & 12:58 & 13:23 & 13:47 & 14:12 & 14:48 & 15:15 & 15:38 & 16:03 \\
Cinnaminson & 10:48 & 11:13 & 11:38 & 12:13 & 12:40 & 13:03 & 13:28 & 13:53 & 14:18 & 14:53 & 15:20 & 15:43 & 16:08 \\
Riverton & 10:50 & 11:15 & 11:40 & 12:15 & 12:42 & 13:05 & 13:30 & 13:55 & 14:20 & 14:55 & 15:22 & 15:45 & 16:10 \\
Palmyra & 10:52 & 11:17 & 11:42 & 12:17 & 12:44 & 13:07 & 13:32 & 13:57 & 14:22 & 14:57 & 15:24 & 15:47 & 16:12 \\
Route 73/Pennsauken & 10:55 & 11:20 & 11:45 & 12:20 & 12:47 & 13:10 & 13:35 & 14:00 & 14:25 & 15:00 & 15:27 & 15:50 & 16:15 \\
Pennsauken TC & 10:57 & 11:22 & 11:47 & 12:22 & 12:49 & 13:12 & 13:37 & 14:02 & 14:27 & 15:02 & 15:29 & 15:52 & 16:17 \\
36th Street & 11:01 & 11:26 & 11:51 & 12:26 & 12:53 & 13:16 & 13:41 & 14:06 & 14:31 & 15:06 & 15:33 & 15:56 & 16:21 \\
Walter Rand TC & 11:06 & 11:31 & 11:56 & 12:31 & 12:58 & 13:21 & 13:46 & 14:11 & 14:36 & 15:11 & 15:38 & 16:01 & 16:26 \\
Cooper St/Rutgers & 11:09 & 11:34 & 11:59 & 12:34 & 13:01 & 13:24 & 13:49 & 14:14 & 14:39 & 15:14 & 15:41 & 16:04 & 16:29 \\
Aquarium & 11:11 & 11:36 & 12:01 & 12:36 & 13:03 & 13:26 & 13:51 & 14:16 & 14:41 & 15:16 & 15:43 & 16:06 & 16:31 \\
Entertainment Center & 11:14 & 11:39 & 12:04 & 12:39 & 13:06 & 13:29 & 13:54 & 14:19 & 14:44 & 15:19 & 15:46 & 16:09 & 16:34 \\
\bottomrule\end{tabular}\end{center}
\subsection*{To Camden (read down) --- trips 27--39 of 42}
\begin{center}\footnotesize\setlength{\tabcolsep}{3pt}\begin{tabular}{@{}lrrrrrrrrrrrrr@{}}\toprule
Station & 27 & 28 & 29 & 30 & 31 & 32 & 33 & 34 & 35 & 36 & 37 & 38 & 39 \\\midrule
Trenton TC & 15:51 & 16:16 & 16:31 & 16:48 & 17:04 & 17:19 & 17:42 & 18:07 & 18:31 & 18:56 & 19:11 & 19:28 & 19:44 \\
Hamilton Avenue & 15:53 & 16:18 & 16:33 & 16:50 & 17:06 & 17:21 & 17:44 & 18:09 & 18:33 & 18:58 & 19:13 & 19:30 & 19:46 \\
Cass Street & 15:55 & 16:20 & 16:35 & 16:52 & 17:08 & 17:23 & 17:46 & 18:11 & 18:35 & 19:00 & 19:15 & 19:32 & 19:48 \\
Bordentown & 16:02 & 16:27 & 16:42 & 16:59 & 17:15 & 17:30 & 17:53 & 18:18 & 18:42 & 19:07 & 19:22 & 19:39 & 19:55 \\
Roebling & 16:07 & 16:32 & 16:47 & 17:04 & 17:21 & 17:35 & 17:58 & 18:23 & 18:47 & 19:12 & 19:27 & 19:44 & 20:01 \\
Florence & 16:10 & 16:35 & 16:50 & 17:07 & 17:24 & 17:38 & 18:01 & 18:26 & 18:50 & 19:15 & 19:30 & 19:47 & 20:04 \\
Burlington Towne Ctr & 16:15 & 16:40 & 16:55 & 17:16 & 17:29 & 17:43 & 18:06 & 18:31 & 18:55 & 19:20 & 19:35 & 19:56 & 20:09 \\
Burlington South & 16:17 & 16:42 & 16:57 & 17:18 & 17:31 & 17:45 & 18:08 & 18:33 & 18:57 & 19:22 & 19:37 & 19:58 & 20:11 \\
Beverly/Edgewater Park & 16:21 & 16:46 & 17:01 & 17:22 & 17:35 & 17:49 & 18:12 & 18:37 & 19:01 & 19:26 & 19:41 & 20:02 & 20:15 \\
Delanco & 16:25 & 16:50 & 17:05 & 17:26 & 17:39 & 17:53 & 18:16 & 18:41 & 19:05 & 19:30 & 19:45 & 20:06 & 20:19 \\
Riverside & 16:27 & 16:52 & 17:07 & 17:28 & 17:41 & 17:55 & 18:18 & 18:43 & 19:07 & 19:32 & 19:47 & 20:08 & 20:21 \\
Cinnaminson & 16:33 & 16:58 & 17:16 & 17:33 & 17:46 & 18:00 & 18:23 & 18:48 & 19:13 & 19:38 & 19:56 & 20:13 & 20:26 \\
Riverton & 16:35 & 17:00 & 17:18 & 17:35 & 17:48 & 18:02 & 18:25 & 18:50 & 19:15 & 19:40 & 19:58 & 20:15 & 20:28 \\
Palmyra & 16:37 & 17:02 & 17:20 & 17:37 & 17:50 & 18:04 & 18:27 & 18:52 & 19:17 & 19:42 & 20:00 & 20:17 & 20:30 \\
Route 73/Pennsauken & 16:40 & 17:05 & 17:23 & 17:40 & 17:53 & 18:07 & 18:30 & 18:55 & 19:20 & 19:45 & 20:03 & 20:20 & 20:33 \\
Pennsauken TC & 16:42 & 17:07 & 17:25 & 17:42 & 17:55 & 18:09 & 18:32 & 18:57 & 19:22 & 19:47 & 20:05 & 20:22 & 20:35 \\
36th Street & 16:46 & 17:11 & 17:29 & 17:46 & 17:59 & 18:13 & 18:36 & 19:01 & 19:26 & 19:51 & 20:09 & 20:26 & 20:39 \\
Walter Rand TC & 16:51 & 17:16 & 17:34 & 17:51 & 18:04 & 18:18 & 18:41 & 19:06 & 19:31 & 19:56 & 20:14 & 20:31 & 20:44 \\
Cooper St/Rutgers & 16:54 & 17:19 & 17:37 & 17:54 & 18:07 & 18:21 & 18:44 & 19:09 & 19:34 & 19:59 & 20:17 & 20:34 & 20:47 \\
Aquarium & 16:56 & 17:21 & 17:39 & 17:56 & 18:09 & 18:23 & 18:46 & 19:11 & 19:36 & 20:01 & 20:19 & 20:36 & 20:49 \\
Entertainment Center & 16:59 & 17:24 & 17:42 & 17:59 & 18:12 & 18:26 & 18:49 & 19:14 & 19:39 & 20:04 & 20:22 & 20:39 & 20:52 \\
\bottomrule\end{tabular}\end{center}
\subsection*{To Camden (read down) --- trips 40--42 of 42}
\begin{center}\footnotesize\setlength{\tabcolsep}{3pt}\begin{tabular}{@{}lrrr@{}}\toprule
Station & 40 & 41 & 42 \\\midrule
Trenton TC & 19:59 & 20:22 & 20:47 \\
Hamilton Avenue & 20:01 & 20:24 & 20:49 \\
Cass Street & 20:03 & 20:26 & 20:51 \\
Bordentown & 20:10 & 20:33 & 20:58 \\
Roebling & 20:15 & 20:38 & 21:03 \\
Florence & 20:18 & 20:41 & 21:06 \\
Burlington Towne Ctr & 20:23 & 20:46 & 21:11 \\
Burlington South & 20:25 & 20:48 & 21:13 \\
Beverly/Edgewater Park & 20:29 & 20:52 & 21:17 \\
Delanco & 20:33 & 20:56 & 21:21 \\
Riverside & 20:35 & 20:58 & 21:23 \\
Cinnaminson & 20:40 & 21:03 & 21:28 \\
Riverton & 20:42 & 21:05 & 21:30 \\
Palmyra & 20:44 & 21:07 & 21:32 \\
Route 73/Pennsauken & 20:47 & 21:10 & 21:35 \\
Pennsauken TC & 20:49 & 21:12 & 21:37 \\
36th Street & 20:53 & 21:16 & 21:41 \\
Walter Rand TC & 20:58 & 21:21 & 21:46 \\
Cooper St/Rutgers & 21:01 & 21:24 & 21:49 \\
Aquarium & 21:03 & 21:26 & 21:51 \\
Entertainment Center & 21:06 & 21:29 & 21:54 \\
\bottomrule\end{tabular}\end{center}
\clearpage
\section{Proposal W7: summer (+6\% run time), 7 trains}\label{tt:proposal_W7}
\subsection*{To Trenton (read down) --- trips 1--13 of 35}
\begin{center}\footnotesize\setlength{\tabcolsep}{3pt}\begin{tabular}{@{}lrrrrrrrrrrrrr@{}}\toprule
Station & 1 & 2 & 3 & 4 & 5 & 6 & 7 & 8 & 9 & 10 & 11 & 12 & 13 \\\midrule
Entertainment Center & 5:49 & 6:18 & 6:36 & 6:55 & 7:23 & 7:41 & 8:00 & 8:29 & 8:58 & 9:35 & 10:03 & 10:40 & 11:09 \\
Aquarium & 5:51 & 6:20 & 6:38 & 6:57 & 7:25 & 7:43 & 8:02 & 8:31 & 9:00 & 9:37 & 10:05 & 10:42 & 11:11 \\
Cooper St/Rutgers & 5:53 & 6:22 & 6:40 & 6:59 & 7:27 & 7:45 & 8:04 & 8:33 & 9:02 & 9:39 & 10:07 & 10:44 & 11:13 \\
Walter Rand TC & 5:57 & 6:26 & 6:44 & 7:03 & 7:31 & 7:49 & 8:08 & 8:37 & 9:06 & 9:43 & 10:11 & 10:48 & 11:17 \\
36th Street & 6:03 & 6:32 & 6:50 & 7:09 & 7:37 & 7:55 & 8:14 & 8:43 & 9:12 & 9:49 & 10:17 & 10:54 & 11:23 \\
Pennsauken TC & 6:05 & 6:34 & 6:52 & 7:11 & 7:39 & 7:57 & 8:16 & 8:45 & 9:14 & 9:51 & 10:19 & 10:56 & 11:25 \\
Route 73/Pennsauken & 6:08 & 6:37 & 6:55 & 7:14 & 7:42 & 8:00 & 8:19 & 8:48 & 9:17 & 9:54 & 10:22 & 10:59 & 11:28 \\
Palmyra & 6:11 & 6:40 & 6:58 & 7:17 & 7:45 & 8:03 & 8:22 & 8:51 & 9:20 & 9:57 & 10:25 & 11:02 & 11:31 \\
Riverton & 6:13 & 6:42 & 7:00 & 7:19 & 7:47 & 8:05 & 8:24 & 8:53 & 9:22 & 9:59 & 10:27 & 11:04 & 11:33 \\
Cinnaminson & 6:16 & 6:44 & 7:02 & 7:21 & 7:49 & 8:08 & 8:26 & 8:56 & 9:24 & 10:01 & 10:29 & 11:06 & 11:36 \\
Riverside & 6:21 & 6:50 & 7:08 & 7:27 & 7:55 & 8:13 & 8:32 & 9:01 & 9:30 & 10:07 & 10:35 & 11:12 & 11:41 \\
Delanco & 6:23 & 6:52 & 7:10 & 7:29 & 7:57 & 8:15 & 8:34 & 9:03 & 9:32 & 10:09 & 10:37 & 11:14 & 11:43 \\
Beverly/Edgewater Park & 6:27 & 6:56 & 7:14 & 7:33 & 8:01 & 8:19 & 8:38 & 9:07 & 9:36 & 10:13 & 10:41 & 11:18 & 11:47 \\
Burlington South & 6:31 & 7:00 & 7:18 & 7:37 & 8:05 & 8:23 & 8:42 & 9:11 & 9:40 & 10:17 & 10:45 & 11:22 & 11:51 \\
Burlington Towne Ctr & 6:34 & 7:02 & 7:20 & 7:39 & 8:08 & 8:26 & 8:45 & 9:14 & 9:42 & 10:19 & 10:48 & 11:25 & 11:54 \\
Florence & 6:40 & 7:09 & 7:27 & 7:46 & 8:14 & 8:32 & 8:51 & 9:20 & 9:49 & 10:26 & 10:54 & 11:31 & 12:00 \\
Roebling & 6:43 & 7:12 & 7:30 & 7:49 & 8:17 & 8:35 & 8:54 & 9:23 & 9:52 & 10:29 & 10:57 & 11:34 & 12:03 \\
Bordentown & 6:48 & 7:17 & 7:35 & 7:54 & 8:22 & 8:40 & 8:59 & 9:28 & 9:57 & 10:34 & 11:02 & 11:39 & 12:08 \\
Cass Street & 6:55 & 7:24 & 7:42 & 8:01 & 8:29 & 8:47 & 9:06 & 9:35 & 10:04 & 10:41 & 11:09 & 11:46 & 12:15 \\
Hamilton Avenue & 6:57 & 7:26 & 7:44 & 8:03 & 8:31 & 8:49 & 9:08 & 9:37 & 10:06 & 10:43 & 11:11 & 11:48 & 12:17 \\
Trenton TC & 7:00 & 7:29 & 7:47 & 8:06 & 8:34 & 8:52 & 9:11 & 9:40 & 10:09 & 10:46 & 11:14 & 11:51 & 12:20 \\
\bottomrule\end{tabular}\end{center}
\subsection*{To Trenton (read down) --- trips 14--26 of 35}
\begin{center}\footnotesize\setlength{\tabcolsep}{3pt}\begin{tabular}{@{}lrrrrrrrrrrrrr@{}}\toprule
Station & 14 & 15 & 16 & 17 & 18 & 19 & 20 & 21 & 22 & 23 & 24 & 25 & 26 \\\midrule
Entertainment Center & 11:38 & 12:15 & 12:43 & 13:20 & 13:49 & 14:18 & 14:36 & 14:55 & 15:23 & 15:41 & 16:00 & 16:29 & 16:58 \\
Aquarium & 11:40 & 12:17 & 12:45 & 13:22 & 13:51 & 14:20 & 14:38 & 14:57 & 15:25 & 15:43 & 16:02 & 16:31 & 17:00 \\
Cooper St/Rutgers & 11:42 & 12:19 & 12:47 & 13:24 & 13:53 & 14:22 & 14:40 & 14:59 & 15:27 & 15:45 & 16:04 & 16:33 & 17:02 \\
Walter Rand TC & 11:46 & 12:23 & 12:51 & 13:28 & 13:57 & 14:26 & 14:44 & 15:03 & 15:31 & 15:49 & 16:08 & 16:37 & 17:06 \\
36th Street & 11:52 & 12:29 & 12:57 & 13:34 & 14:03 & 14:32 & 14:50 & 15:09 & 15:37 & 15:55 & 16:14 & 16:43 & 17:12 \\
Pennsauken TC & 11:54 & 12:31 & 12:59 & 13:36 & 14:05 & 14:34 & 14:52 & 15:11 & 15:39 & 15:57 & 16:16 & 16:45 & 17:14 \\
Route 73/Pennsauken & 11:57 & 12:34 & 13:02 & 13:39 & 14:08 & 14:37 & 14:55 & 15:14 & 15:42 & 16:00 & 16:19 & 16:48 & 17:17 \\
Palmyra & 12:00 & 12:37 & 13:05 & 13:42 & 14:11 & 14:40 & 14:58 & 15:17 & 15:45 & 16:03 & 16:22 & 16:51 & 17:20 \\
Riverton & 12:02 & 12:39 & 13:07 & 13:44 & 14:13 & 14:42 & 15:00 & 15:19 & 15:47 & 16:05 & 16:24 & 16:53 & 17:22 \\
Cinnaminson & 12:04 & 12:41 & 13:09 & 13:46 & 14:16 & 14:44 & 15:02 & 15:21 & 15:49 & 16:08 & 16:26 & 16:56 & 17:24 \\
Riverside & 12:10 & 12:47 & 13:15 & 13:52 & 14:21 & 14:50 & 15:08 & 15:27 & 15:55 & 16:13 & 16:32 & 17:01 & 17:30 \\
Delanco & 12:12 & 12:49 & 13:17 & 13:54 & 14:23 & 14:52 & 15:10 & 15:29 & 15:57 & 16:15 & 16:34 & 17:03 & 17:32 \\
Beverly/Edgewater Park & 12:16 & 12:53 & 13:21 & 13:58 & 14:27 & 14:56 & 15:14 & 15:33 & 16:01 & 16:19 & 16:38 & 17:07 & 17:36 \\
Burlington South & 12:20 & 12:57 & 13:25 & 14:02 & 14:31 & 15:00 & 15:18 & 15:37 & 16:05 & 16:23 & 16:42 & 17:11 & 17:40 \\
Burlington Towne Ctr & 12:22 & 12:59 & 13:28 & 14:05 & 14:34 & 15:02 & 15:20 & 15:39 & 16:08 & 16:26 & 16:45 & 17:14 & 17:42 \\
Florence & 12:29 & 13:06 & 13:34 & 14:11 & 14:40 & 15:09 & 15:27 & 15:46 & 16:14 & 16:32 & 16:51 & 17:20 & 17:49 \\
Roebling & 12:32 & 13:09 & 13:37 & 14:14 & 14:43 & 15:12 & 15:30 & 15:49 & 16:17 & 16:35 & 16:54 & 17:23 & 17:52 \\
Bordentown & 12:37 & 13:14 & 13:42 & 14:19 & 14:48 & 15:17 & 15:35 & 15:54 & 16:22 & 16:40 & 16:59 & 17:28 & 17:57 \\
Cass Street & 12:44 & 13:21 & 13:49 & 14:26 & 14:55 & 15:24 & 15:42 & 16:01 & 16:29 & 16:47 & 17:06 & 17:35 & 18:04 \\
Hamilton Avenue & 12:46 & 13:23 & 13:51 & 14:28 & 14:57 & 15:26 & 15:44 & 16:03 & 16:31 & 16:49 & 17:08 & 17:37 & 18:06 \\
Trenton TC & 12:49 & 13:26 & 13:54 & 14:31 & 15:00 & 15:29 & 15:47 & 16:06 & 16:34 & 16:52 & 17:11 & 17:40 & 18:09 \\
\bottomrule\end{tabular}\end{center}
\subsection*{To Trenton (read down) --- trips 27--35 of 35}
\begin{center}\footnotesize\setlength{\tabcolsep}{3pt}\begin{tabular}{@{}lrrrrrrrrr@{}}\toprule
Station & 27 & 28 & 29 & 30 & 31 & 32 & 33 & 34 & 35 \\\midrule
Entertainment Center & 17:16 & 17:35 & 18:03 & 18:21 & 18:40 & 19:09 & 19:38 & 20:15 & 20:43 \\
Aquarium & 17:18 & 17:37 & 18:05 & 18:23 & 18:42 & 19:11 & 19:40 & 20:17 & 20:45 \\
Cooper St/Rutgers & 17:20 & 17:39 & 18:07 & 18:25 & 18:44 & 19:13 & 19:42 & 20:19 & 20:47 \\
Walter Rand TC & 17:24 & 17:43 & 18:11 & 18:29 & 18:48 & 19:17 & 19:46 & 20:23 & 20:51 \\
36th Street & 17:30 & 17:49 & 18:17 & 18:35 & 18:54 & 19:23 & 19:52 & 20:29 & 20:57 \\
Pennsauken TC & 17:32 & 17:51 & 18:19 & 18:37 & 18:56 & 19:25 & 19:54 & 20:31 & 20:59 \\
Route 73/Pennsauken & 17:35 & 17:54 & 18:22 & 18:40 & 18:59 & 19:28 & 19:57 & 20:34 & 21:02 \\
Palmyra & 17:38 & 17:57 & 18:25 & 18:43 & 19:02 & 19:31 & 20:00 & 20:37 & 21:05 \\
Riverton & 17:40 & 17:59 & 18:27 & 18:45 & 19:04 & 19:33 & 20:02 & 20:39 & 21:07 \\
Cinnaminson & 17:42 & 18:01 & 18:29 & 18:48 & 19:06 & 19:36 & 20:04 & 20:41 & 21:09 \\
Riverside & 17:48 & 18:07 & 18:35 & 18:53 & 19:12 & 19:41 & 20:10 & 20:47 & 21:15 \\
Delanco & 17:50 & 18:09 & 18:37 & 18:55 & 19:14 & 19:43 & 20:12 & 20:49 & 21:17 \\
Beverly/Edgewater Park & 17:54 & 18:13 & 18:41 & 18:59 & 19:18 & 19:47 & 20:16 & 20:53 & 21:21 \\
Burlington South & 17:58 & 18:17 & 18:45 & 19:03 & 19:22 & 19:51 & 20:20 & 20:57 & 21:25 \\
Burlington Towne Ctr & 18:00 & 18:19 & 18:48 & 19:06 & 19:25 & 19:54 & 20:22 & 20:59 & 21:28 \\
Florence & 18:07 & 18:26 & 18:54 & 19:12 & 19:31 & 20:00 & 20:29 & 21:06 & 21:34 \\
Roebling & 18:10 & 18:29 & 18:57 & 19:15 & 19:34 & 20:03 & 20:32 & 21:09 & 21:37 \\
Bordentown & 18:15 & 18:34 & 19:02 & 19:20 & 19:39 & 20:08 & 20:37 & 21:14 & 21:42 \\
Cass Street & 18:22 & 18:41 & 19:09 & 19:27 & 19:46 & 20:15 & 20:44 & 21:21 & 21:49 \\
Hamilton Avenue & 18:24 & 18:43 & 19:11 & 19:29 & 19:48 & 20:17 & 20:46 & 21:23 & 21:51 \\
Trenton TC & 18:27 & 18:46 & 19:14 & 19:32 & 19:51 & 20:20 & 20:49 & 21:26 & 21:54 \\
\bottomrule\end{tabular}\end{center}
\subsection*{To Camden (read down) --- trips 1--13 of 36}
\begin{center}\footnotesize\setlength{\tabcolsep}{3pt}\begin{tabular}{@{}lrrrrrrrrrrrrr@{}}\toprule
Station & 1 & 2 & 3 & 4 & 5 & 6 & 7 & 8 & 9 & 10 & 11 & 12 & 13 \\\midrule
Trenton TC & 6:01 & 6:20 & 6:38 & 7:06 & 7:35 & 7:53 & 8:12 & 8:41 & 9:00 & 9:18 & 9:46 & 10:15 & 10:52 \\
Hamilton Avenue & 6:03 & 6:22 & 6:40 & 7:08 & 7:37 & 7:55 & 8:14 & 8:43 & 9:02 & 9:20 & 9:48 & 10:17 & 10:54 \\
Cass Street & 6:05 & 6:24 & 6:42 & 7:10 & 7:39 & 7:57 & 8:16 & 8:45 & 9:04 & 9:22 & 9:50 & 10:19 & 10:56 \\
Bordentown & 6:14 & 6:31 & 6:49 & 7:17 & 7:49 & 8:04 & 8:23 & 8:54 & 9:11 & 9:29 & 9:57 & 10:29 & 11:03 \\
Roebling & 6:19 & 6:37 & 6:55 & 7:23 & 7:54 & 8:10 & 8:29 & 8:59 & 9:17 & 9:35 & 10:03 & 10:34 & 11:09 \\
Florence & 6:22 & 6:40 & 6:58 & 7:26 & 7:57 & 8:13 & 8:32 & 9:02 & 9:20 & 9:38 & 10:06 & 10:37 & 11:12 \\
Burlington Towne Ctr & 6:28 & 6:45 & 7:03 & 7:31 & 8:03 & 8:18 & 8:37 & 9:08 & 9:25 & 9:43 & 10:11 & 10:43 & 11:17 \\
Burlington South & 6:30 & 6:47 & 7:05 & 7:33 & 8:05 & 8:20 & 8:39 & 9:10 & 9:27 & 9:45 & 10:13 & 10:45 & 11:19 \\
Beverly/Edgewater Park & 6:34 & 6:52 & 7:10 & 7:38 & 8:09 & 8:25 & 8:44 & 9:14 & 9:32 & 9:50 & 10:18 & 10:49 & 11:24 \\
Delanco & 6:38 & 6:56 & 7:14 & 7:42 & 8:13 & 8:29 & 8:48 & 9:18 & 9:36 & 9:54 & 10:22 & 10:53 & 11:28 \\
Riverside & 6:40 & 6:58 & 7:16 & 7:44 & 8:15 & 8:31 & 8:50 & 9:20 & 9:38 & 9:56 & 10:24 & 10:55 & 11:30 \\
Cinnaminson & 6:46 & 7:03 & 7:21 & 7:49 & 8:22 & 8:36 & 8:55 & 9:26 & 9:43 & 10:01 & 10:29 & 11:02 & 11:35 \\
Riverton & 6:48 & 7:05 & 7:23 & 7:51 & 8:24 & 8:38 & 8:57 & 9:28 & 9:45 & 10:03 & 10:31 & 11:04 & 11:37 \\
Palmyra & 6:50 & 7:07 & 7:25 & 7:53 & 8:26 & 8:40 & 8:59 & 9:30 & 9:47 & 10:05 & 10:33 & 11:06 & 11:39 \\
Route 73/Pennsauken & 6:53 & 7:11 & 7:29 & 7:57 & 8:29 & 8:44 & 9:03 & 9:33 & 9:51 & 10:09 & 10:37 & 11:09 & 11:43 \\
Pennsauken TC & 6:55 & 7:13 & 7:31 & 7:59 & 8:31 & 8:46 & 9:05 & 9:35 & 9:53 & 10:11 & 10:39 & 11:11 & 11:45 \\
36th Street & 7:00 & 7:17 & 7:35 & 8:03 & 8:35 & 8:50 & 9:09 & 9:40 & 9:57 & 10:15 & 10:43 & 11:15 & 11:49 \\
Walter Rand TC & 7:05 & 7:22 & 7:40 & 8:08 & 8:41 & 8:55 & 9:14 & 9:45 & 10:02 & 10:20 & 10:48 & 11:21 & 11:54 \\
Cooper St/Rutgers & 7:08 & 7:25 & 7:43 & 8:11 & 8:44 & 8:58 & 9:17 & 9:48 & 10:05 & 10:23 & 10:51 & 11:24 & 11:57 \\
Aquarium & 7:10 & 7:28 & 7:46 & 8:14 & 8:46 & 9:01 & 9:20 & 9:50 & 10:08 & 10:26 & 10:54 & 11:26 & 12:00 \\
Entertainment Center & 7:13 & 7:31 & 7:49 & 8:17 & 8:49 & 9:04 & 9:23 & 9:53 & 10:11 & 10:29 & 10:57 & 11:29 & 12:03 \\
\bottomrule\end{tabular}\end{center}
\subsection*{To Camden (read down) --- trips 14--26 of 36}
\begin{center}\footnotesize\setlength{\tabcolsep}{3pt}\begin{tabular}{@{}lrrrrrrrrrrrrr@{}}\toprule
Station & 14 & 15 & 16 & 17 & 18 & 19 & 20 & 21 & 22 & 23 & 24 & 25 & 26 \\\midrule
Trenton TC & 11:21 & 11:58 & 12:26 & 12:55 & 13:32 & 14:01 & 14:38 & 15:06 & 15:35 & 15:53 & 16:12 & 16:41 & 17:00 \\
Hamilton Avenue & 11:23 & 12:00 & 12:28 & 12:57 & 13:34 & 14:03 & 14:40 & 15:08 & 15:37 & 15:55 & 16:14 & 16:43 & 17:02 \\
Cass Street & 11:25 & 12:02 & 12:30 & 12:59 & 13:36 & 14:05 & 14:42 & 15:10 & 15:39 & 15:57 & 16:16 & 16:45 & 17:04 \\
Bordentown & 11:34 & 12:09 & 12:37 & 13:09 & 13:43 & 14:14 & 14:49 & 15:17 & 15:49 & 16:04 & 16:23 & 16:54 & 17:11 \\
Roebling & 11:39 & 12:15 & 12:43 & 13:14 & 13:49 & 14:19 & 14:55 & 15:23 & 15:54 & 16:10 & 16:29 & 16:59 & 17:17 \\
Florence & 11:42 & 12:18 & 12:46 & 13:17 & 13:52 & 14:22 & 14:58 & 15:26 & 15:57 & 16:13 & 16:32 & 17:02 & 17:20 \\
Burlington Towne Ctr & 11:48 & 12:23 & 12:51 & 13:23 & 13:57 & 14:28 & 15:03 & 15:31 & 16:03 & 16:18 & 16:37 & 17:08 & 17:25 \\
Burlington South & 11:50 & 12:25 & 12:53 & 13:25 & 13:59 & 14:30 & 15:05 & 15:33 & 16:05 & 16:20 & 16:39 & 17:10 & 17:27 \\
Beverly/Edgewater Park & 11:54 & 12:30 & 12:58 & 13:29 & 14:04 & 14:34 & 15:10 & 15:38 & 16:09 & 16:25 & 16:44 & 17:14 & 17:32 \\
Delanco & 11:58 & 12:34 & 13:02 & 13:33 & 14:08 & 14:38 & 15:14 & 15:42 & 16:13 & 16:29 & 16:48 & 17:18 & 17:36 \\
Riverside & 12:00 & 12:36 & 13:04 & 13:35 & 14:10 & 14:40 & 15:16 & 15:44 & 16:15 & 16:31 & 16:50 & 17:20 & 17:38 \\
Cinnaminson & 12:06 & 12:41 & 13:09 & 13:42 & 14:15 & 14:46 & 15:21 & 15:49 & 16:22 & 16:36 & 16:55 & 17:26 & 17:43 \\
Riverton & 12:08 & 12:43 & 13:11 & 13:44 & 14:17 & 14:48 & 15:23 & 15:51 & 16:24 & 16:38 & 16:57 & 17:28 & 17:45 \\
Palmyra & 12:10 & 12:45 & 13:13 & 13:46 & 14:19 & 14:50 & 15:25 & 15:53 & 16:26 & 16:40 & 16:59 & 17:30 & 17:47 \\
Route 73/Pennsauken & 12:13 & 12:49 & 13:17 & 13:49 & 14:23 & 14:53 & 15:29 & 15:57 & 16:29 & 16:44 & 17:03 & 17:33 & 17:51 \\
Pennsauken TC & 12:15 & 12:51 & 13:19 & 13:51 & 14:25 & 14:55 & 15:31 & 15:59 & 16:31 & 16:46 & 17:05 & 17:35 & 17:53 \\
36th Street & 12:20 & 12:55 & 13:23 & 13:55 & 14:29 & 15:00 & 15:35 & 16:03 & 16:35 & 16:50 & 17:09 & 17:40 & 17:57 \\
Walter Rand TC & 12:25 & 13:00 & 13:28 & 14:01 & 14:34 & 15:05 & 15:40 & 16:08 & 16:41 & 16:55 & 17:14 & 17:45 & 18:02 \\
Cooper St/Rutgers & 12:28 & 13:03 & 13:31 & 14:04 & 14:37 & 15:08 & 15:43 & 16:11 & 16:44 & 16:58 & 17:17 & 17:48 & 18:05 \\
Aquarium & 12:30 & 13:06 & 13:34 & 14:06 & 14:40 & 15:10 & 15:46 & 16:14 & 16:46 & 17:01 & 17:20 & 17:50 & 18:08 \\
Entertainment Center & 12:33 & 13:09 & 13:37 & 14:09 & 14:43 & 15:13 & 15:49 & 16:17 & 16:49 & 17:04 & 17:23 & 17:53 & 18:11 \\
\bottomrule\end{tabular}\end{center}
\subsection*{To Camden (read down) --- trips 27--36 of 36}
\begin{center}\footnotesize\setlength{\tabcolsep}{3pt}\begin{tabular}{@{}lrrrrrrrrrr@{}}\toprule
Station & 27 & 28 & 29 & 30 & 31 & 32 & 33 & 34 & 35 & 36 \\\midrule
Trenton TC & 17:18 & 17:46 & 18:15 & 18:33 & 18:52 & 19:21 & 19:40 & 19:58 & 20:26 & 20:55 \\
Hamilton Avenue & 17:20 & 17:48 & 18:17 & 18:35 & 18:54 & 19:23 & 19:42 & 20:00 & 20:28 & 20:57 \\
Cass Street & 17:22 & 17:50 & 18:19 & 18:37 & 18:56 & 19:25 & 19:44 & 20:02 & 20:30 & 20:59 \\
Bordentown & 17:29 & 17:57 & 18:29 & 18:44 & 19:03 & 19:34 & 19:51 & 20:09 & 20:37 & 21:09 \\
Roebling & 17:35 & 18:03 & 18:34 & 18:50 & 19:09 & 19:39 & 19:57 & 20:15 & 20:43 & 21:14 \\
Florence & 17:38 & 18:06 & 18:37 & 18:53 & 19:12 & 19:42 & 20:00 & 20:18 & 20:46 & 21:17 \\
Burlington Towne Ctr & 17:43 & 18:11 & 18:43 & 18:58 & 19:17 & 19:48 & 20:05 & 20:23 & 20:51 & 21:23 \\
Burlington South & 17:45 & 18:13 & 18:45 & 19:00 & 19:19 & 19:50 & 20:07 & 20:25 & 20:53 & 21:25 \\
Beverly/Edgewater Park & 17:50 & 18:18 & 18:49 & 19:05 & 19:24 & 19:54 & 20:12 & 20:30 & 20:58 & 21:29 \\
Delanco & 17:54 & 18:22 & 18:53 & 19:09 & 19:28 & 19:58 & 20:16 & 20:34 & 21:02 & 21:33 \\
Riverside & 17:56 & 18:24 & 18:55 & 19:11 & 19:30 & 20:00 & 20:18 & 20:36 & 21:04 & 21:35 \\
Cinnaminson & 18:01 & 18:29 & 19:02 & 19:16 & 19:35 & 20:06 & 20:23 & 20:41 & 21:09 & 21:42 \\
Riverton & 18:03 & 18:31 & 19:04 & 19:18 & 19:37 & 20:08 & 20:25 & 20:43 & 21:11 & 21:44 \\
Palmyra & 18:05 & 18:33 & 19:06 & 19:20 & 19:39 & 20:10 & 20:27 & 20:45 & 21:13 & 21:46 \\
Route 73/Pennsauken & 18:09 & 18:37 & 19:09 & 19:24 & 19:43 & 20:13 & 20:31 & 20:49 & 21:17 & 21:49 \\
Pennsauken TC & 18:11 & 18:39 & 19:11 & 19:26 & 19:45 & 20:15 & 20:33 & 20:51 & 21:19 & 21:51 \\
36th Street & 18:15 & 18:43 & 19:15 & 19:30 & 19:49 & 20:20 & 20:37 & 20:55 & 21:23 & 21:55 \\
Walter Rand TC & 18:20 & 18:48 & 19:21 & 19:35 & 19:54 & 20:25 & 20:42 & 21:00 & 21:28 & 22:01 \\
Cooper St/Rutgers & 18:23 & 18:51 & 19:24 & 19:38 & 19:57 & 20:28 & 20:45 & 21:03 & 21:31 & 22:04 \\
Aquarium & 18:26 & 18:54 & 19:26 & 19:41 & 20:00 & 20:30 & 20:48 & 21:06 & 21:34 & 22:06 \\
Entertainment Center & 18:29 & 18:57 & 19:29 & 19:44 & 20:03 & 20:33 & 20:51 & 21:09 & 21:37 & 22:09 \\
\bottomrule\end{tabular}\end{center}
\clearpage
\section{Uniform U8: every 21 min peak, 30 off-peak}\label{tt:uniform_U8}
\subsection*{To Trenton (read down) --- trips 1--13 of 37}
\begin{center}\footnotesize\setlength{\tabcolsep}{3pt}\begin{tabular}{@{}lrrrrrrrrrrrrr@{}}\toprule
Station & 1 & 2 & 3 & 4 & 5 & 6 & 7 & 8 & 9 & 10 & 11 & 12 & 13 \\\midrule
Entertainment Center & 5:45 & 6:06 & 6:27 & 6:48 & 7:09 & 7:30 & 7:51 & 8:12 & 8:33 & 8:54 & 9:15 & 9:45 & 10:15 \\
Aquarium & 5:47 & 6:08 & 6:29 & 6:50 & 7:11 & 7:32 & 7:53 & 8:14 & 8:35 & 8:56 & 9:17 & 9:47 & 10:17 \\
Cooper St/Rutgers & 5:49 & 6:10 & 6:31 & 6:52 & 7:13 & 7:34 & 7:55 & 8:16 & 8:37 & 8:58 & 9:19 & 9:49 & 10:19 \\
Walter Rand TC & 5:53 & 6:14 & 6:35 & 6:56 & 7:17 & 7:38 & 7:59 & 8:20 & 8:41 & 9:02 & 9:23 & 9:53 & 10:23 \\
36th Street & 5:58 & 6:19 & 6:40 & 7:01 & 7:22 & 7:43 & 8:04 & 8:25 & 8:46 & 9:07 & 9:28 & 9:58 & 10:28 \\
Pennsauken TC & 6:00 & 6:21 & 6:42 & 7:03 & 7:24 & 7:45 & 8:06 & 8:27 & 8:48 & 9:09 & 9:30 & 10:00 & 10:30 \\
Route 73/Pennsauken & 6:03 & 6:24 & 6:45 & 7:06 & 7:27 & 7:48 & 8:09 & 8:30 & 8:51 & 9:12 & 9:33 & 10:03 & 10:33 \\
Palmyra & 6:06 & 6:27 & 6:48 & 7:09 & 7:30 & 7:51 & 8:12 & 8:33 & 8:54 & 9:15 & 9:36 & 10:06 & 10:36 \\
Riverton & 6:08 & 6:29 & 6:50 & 7:11 & 7:32 & 7:53 & 8:14 & 8:35 & 8:56 & 9:17 & 9:38 & 10:08 & 10:38 \\
Cinnaminson & 6:10 & 6:31 & 6:52 & 7:13 & 7:34 & 7:55 & 8:16 & 8:37 & 8:58 & 9:19 & 9:40 & 10:10 & 10:40 \\
Riverside & 6:15 & 6:38 & 6:59 & 7:20 & 7:41 & 8:02 & 8:23 & 8:44 & 9:05 & 9:26 & 9:45 & 10:15 & 10:45 \\
Delanco & 6:17 & 6:40 & 7:01 & 7:22 & 7:43 & 8:04 & 8:25 & 8:46 & 9:07 & 9:28 & 9:47 & 10:17 & 10:47 \\
Beverly/Edgewater Park & 6:21 & 6:44 & 7:05 & 7:26 & 7:47 & 8:08 & 8:29 & 8:50 & 9:11 & 9:32 & 9:51 & 10:21 & 10:51 \\
Burlington South & 6:25 & 6:48 & 7:09 & 7:30 & 7:51 & 8:12 & 8:33 & 8:54 & 9:15 & 9:36 & 9:55 & 10:25 & 10:55 \\
Burlington Towne Ctr & 6:27 & 6:50 & 7:11 & 7:32 & 7:53 & 8:14 & 8:35 & 8:56 & 9:17 & 9:38 & 9:57 & 10:27 & 10:57 \\
Florence & 6:33 & 6:56 & 7:17 & 7:38 & 7:59 & 8:20 & 8:41 & 9:02 & 9:23 & 9:44 & 10:03 & 10:33 & 11:03 \\
Roebling & 6:36 & 6:59 & 7:20 & 7:41 & 8:02 & 8:23 & 8:44 & 9:05 & 9:26 & 9:47 & 10:06 & 10:36 & 11:06 \\
Bordentown & 6:41 & 7:04 & 7:25 & 7:46 & 8:07 & 8:28 & 8:49 & 9:10 & 9:31 & 9:52 & 10:11 & 10:41 & 11:11 \\
Cass Street & 6:47 & 7:10 & 7:31 & 7:52 & 8:13 & 8:34 & 8:55 & 9:16 & 9:37 & 9:58 & 10:17 & 10:47 & 11:17 \\
Hamilton Avenue & 6:49 & 7:12 & 7:33 & 7:54 & 8:15 & 8:36 & 8:57 & 9:18 & 9:39 & 10:00 & 10:19 & 10:49 & 11:19 \\
Trenton TC & 6:52 & 7:15 & 7:36 & 7:57 & 8:18 & 8:39 & 9:00 & 9:21 & 9:42 & 10:03 & 10:22 & 10:52 & 11:22 \\
\bottomrule\end{tabular}\end{center}
\subsection*{To Trenton (read down) --- trips 14--26 of 37}
\begin{center}\footnotesize\setlength{\tabcolsep}{3pt}\begin{tabular}{@{}lrrrrrrrrrrrrr@{}}\toprule
Station & 14 & 15 & 16 & 17 & 18 & 19 & 20 & 21 & 22 & 23 & 24 & 25 & 26 \\\midrule
Entertainment Center & 10:45 & 11:15 & 11:45 & 12:15 & 12:45 & 13:15 & 13:45 & 14:15 & 14:45 & 15:12 & 15:33 & 15:54 & 16:15 \\
Aquarium & 10:47 & 11:17 & 11:47 & 12:17 & 12:47 & 13:17 & 13:47 & 14:17 & 14:47 & 15:14 & 15:35 & 15:56 & 16:17 \\
Cooper St/Rutgers & 10:49 & 11:19 & 11:49 & 12:19 & 12:49 & 13:19 & 13:49 & 14:19 & 14:49 & 15:16 & 15:37 & 15:58 & 16:19 \\
Walter Rand TC & 10:53 & 11:23 & 11:53 & 12:23 & 12:53 & 13:23 & 13:53 & 14:23 & 14:53 & 15:20 & 15:41 & 16:02 & 16:23 \\
36th Street & 10:58 & 11:28 & 11:58 & 12:28 & 12:58 & 13:28 & 13:58 & 14:28 & 14:58 & 15:25 & 15:46 & 16:07 & 16:28 \\
Pennsauken TC & 11:00 & 11:30 & 12:00 & 12:30 & 13:00 & 13:30 & 14:00 & 14:30 & 15:00 & 15:27 & 15:48 & 16:09 & 16:30 \\
Route 73/Pennsauken & 11:03 & 11:33 & 12:03 & 12:33 & 13:03 & 13:33 & 14:03 & 14:33 & 15:03 & 15:30 & 15:51 & 16:12 & 16:33 \\
Palmyra & 11:06 & 11:36 & 12:06 & 12:36 & 13:06 & 13:36 & 14:06 & 14:36 & 15:06 & 15:33 & 15:54 & 16:15 & 16:36 \\
Riverton & 11:08 & 11:38 & 12:08 & 12:38 & 13:08 & 13:38 & 14:08 & 14:38 & 15:08 & 15:35 & 15:56 & 16:17 & 16:38 \\
Cinnaminson & 11:10 & 11:40 & 12:10 & 12:40 & 13:10 & 13:40 & 14:10 & 14:40 & 15:10 & 15:37 & 15:58 & 16:19 & 16:40 \\
Riverside & 11:15 & 11:45 & 12:15 & 12:45 & 13:15 & 13:45 & 14:15 & 14:45 & 15:15 & 15:44 & 16:05 & 16:26 & 16:47 \\
Delanco & 11:17 & 11:47 & 12:17 & 12:47 & 13:17 & 13:47 & 14:17 & 14:47 & 15:17 & 15:46 & 16:07 & 16:28 & 16:49 \\
Beverly/Edgewater Park & 11:21 & 11:51 & 12:21 & 12:51 & 13:21 & 13:51 & 14:21 & 14:51 & 15:21 & 15:50 & 16:11 & 16:32 & 16:53 \\
Burlington South & 11:25 & 11:55 & 12:25 & 12:55 & 13:25 & 13:55 & 14:25 & 14:55 & 15:25 & 15:54 & 16:15 & 16:36 & 16:57 \\
Burlington Towne Ctr & 11:27 & 11:57 & 12:27 & 12:57 & 13:27 & 13:57 & 14:27 & 14:57 & 15:27 & 15:56 & 16:17 & 16:38 & 16:59 \\
Florence & 11:33 & 12:03 & 12:33 & 13:03 & 13:33 & 14:03 & 14:33 & 15:03 & 15:33 & 16:02 & 16:23 & 16:44 & 17:05 \\
Roebling & 11:36 & 12:06 & 12:36 & 13:06 & 13:36 & 14:06 & 14:36 & 15:06 & 15:36 & 16:05 & 16:26 & 16:47 & 17:08 \\
Bordentown & 11:41 & 12:11 & 12:41 & 13:11 & 13:41 & 14:11 & 14:41 & 15:11 & 15:41 & 16:10 & 16:31 & 16:52 & 17:13 \\
Cass Street & 11:47 & 12:17 & 12:47 & 13:17 & 13:47 & 14:17 & 14:47 & 15:17 & 15:47 & 16:16 & 16:37 & 16:58 & 17:19 \\
Hamilton Avenue & 11:49 & 12:19 & 12:49 & 13:19 & 13:49 & 14:19 & 14:49 & 15:19 & 15:49 & 16:18 & 16:39 & 17:00 & 17:21 \\
Trenton TC & 11:52 & 12:22 & 12:52 & 13:22 & 13:52 & 14:22 & 14:52 & 15:22 & 15:52 & 16:21 & 16:42 & 17:03 & 17:24 \\
\bottomrule\end{tabular}\end{center}
\subsection*{To Trenton (read down) --- trips 27--37 of 37}
\begin{center}\footnotesize\setlength{\tabcolsep}{3pt}\begin{tabular}{@{}lrrrrrrrrrrr@{}}\toprule
Station & 27 & 28 & 29 & 30 & 31 & 32 & 33 & 34 & 35 & 36 & 37 \\\midrule
Entertainment Center & 16:36 & 16:57 & 17:18 & 17:39 & 18:00 & 18:21 & 18:42 & 19:15 & 19:45 & 20:15 & 20:45 \\
Aquarium & 16:38 & 16:59 & 17:20 & 17:41 & 18:02 & 18:23 & 18:44 & 19:17 & 19:47 & 20:17 & 20:47 \\
Cooper St/Rutgers & 16:40 & 17:01 & 17:22 & 17:43 & 18:04 & 18:25 & 18:46 & 19:19 & 19:49 & 20:19 & 20:49 \\
Walter Rand TC & 16:44 & 17:05 & 17:26 & 17:47 & 18:08 & 18:29 & 18:50 & 19:23 & 19:53 & 20:23 & 20:53 \\
36th Street & 16:49 & 17:10 & 17:31 & 17:52 & 18:13 & 18:34 & 18:55 & 19:28 & 19:58 & 20:28 & 20:58 \\
Pennsauken TC & 16:51 & 17:12 & 17:33 & 17:54 & 18:15 & 18:36 & 18:57 & 19:30 & 20:00 & 20:30 & 21:00 \\
Route 73/Pennsauken & 16:54 & 17:15 & 17:36 & 17:57 & 18:18 & 18:39 & 19:00 & 19:33 & 20:03 & 20:33 & 21:03 \\
Palmyra & 16:57 & 17:18 & 17:39 & 18:00 & 18:21 & 18:42 & 19:03 & 19:36 & 20:06 & 20:36 & 21:06 \\
Riverton & 16:59 & 17:20 & 17:41 & 18:02 & 18:23 & 18:44 & 19:05 & 19:38 & 20:08 & 20:38 & 21:08 \\
Cinnaminson & 17:01 & 17:22 & 17:43 & 18:04 & 18:25 & 18:46 & 19:07 & 19:40 & 20:10 & 20:40 & 21:10 \\
Riverside & 17:08 & 17:29 & 17:50 & 18:11 & 18:32 & 18:53 & 19:14 & 19:45 & 20:15 & 20:45 & 21:15 \\
Delanco & 17:10 & 17:31 & 17:52 & 18:13 & 18:34 & 18:55 & 19:16 & 19:47 & 20:17 & 20:47 & 21:17 \\
Beverly/Edgewater Park & 17:14 & 17:35 & 17:56 & 18:17 & 18:38 & 18:59 & 19:20 & 19:51 & 20:21 & 20:51 & 21:21 \\
Burlington South & 17:18 & 17:39 & 18:00 & 18:21 & 18:42 & 19:03 & 19:24 & 19:55 & 20:25 & 20:55 & 21:25 \\
Burlington Towne Ctr & 17:20 & 17:41 & 18:02 & 18:23 & 18:44 & 19:05 & 19:26 & 19:57 & 20:27 & 20:57 & 21:27 \\
Florence & 17:26 & 17:47 & 18:08 & 18:29 & 18:50 & 19:11 & 19:32 & 20:03 & 20:33 & 21:03 & 21:33 \\
Roebling & 17:29 & 17:50 & 18:11 & 18:32 & 18:53 & 19:14 & 19:35 & 20:06 & 20:36 & 21:06 & 21:36 \\
Bordentown & 17:34 & 17:55 & 18:16 & 18:37 & 18:58 & 19:19 & 19:40 & 20:11 & 20:41 & 21:11 & 21:41 \\
Cass Street & 17:40 & 18:01 & 18:22 & 18:43 & 19:04 & 19:25 & 19:46 & 20:17 & 20:47 & 21:17 & 21:47 \\
Hamilton Avenue & 17:42 & 18:03 & 18:24 & 18:45 & 19:06 & 19:27 & 19:48 & 20:19 & 20:49 & 21:19 & 21:49 \\
Trenton TC & 17:45 & 18:06 & 18:27 & 18:48 & 19:09 & 19:30 & 19:51 & 20:22 & 20:52 & 21:22 & 21:52 \\
\bottomrule\end{tabular}\end{center}
\subsection*{To Camden (read down) --- trips 1--13 of 37}
\begin{center}\footnotesize\setlength{\tabcolsep}{3pt}\begin{tabular}{@{}lrrrrrrrrrrrrr@{}}\toprule
Station & 1 & 2 & 3 & 4 & 5 & 6 & 7 & 8 & 9 & 10 & 11 & 12 & 13 \\\midrule
Trenton TC & 5:57 & 6:06 & 6:27 & 6:48 & 7:09 & 7:30 & 7:51 & 8:12 & 8:33 & 8:54 & 9:27 & 9:57 & 10:27 \\
Hamilton Avenue & 5:59 & 6:08 & 6:29 & 6:50 & 7:11 & 7:32 & 7:53 & 8:14 & 8:35 & 8:56 & 9:29 & 9:59 & 10:29 \\
Cass Street & 6:01 & 6:10 & 6:31 & 6:52 & 7:13 & 7:34 & 7:55 & 8:16 & 8:37 & 8:58 & 9:31 & 10:01 & 10:31 \\
Bordentown & 6:08 & 6:17 & 6:38 & 6:59 & 7:20 & 7:41 & 8:02 & 8:23 & 8:44 & 9:05 & 9:38 & 10:08 & 10:38 \\
Roebling & 6:13 & 6:22 & 6:43 & 7:04 & 7:25 & 7:46 & 8:07 & 8:28 & 8:49 & 9:10 & 9:43 & 10:13 & 10:43 \\
Florence & 6:16 & 6:25 & 6:46 & 7:07 & 7:28 & 7:49 & 8:10 & 8:31 & 8:52 & 9:13 & 9:46 & 10:16 & 10:46 \\
Burlington Towne Ctr & 6:21 & 6:30 & 6:51 & 7:12 & 7:33 & 7:54 & 8:15 & 8:36 & 8:57 & 9:18 & 9:51 & 10:21 & 10:51 \\
Burlington South & 6:23 & 6:32 & 6:53 & 7:14 & 7:35 & 7:56 & 8:17 & 8:38 & 8:59 & 9:20 & 9:53 & 10:23 & 10:53 \\
Beverly/Edgewater Park & 6:27 & 6:36 & 6:57 & 7:18 & 7:39 & 8:00 & 8:21 & 8:42 & 9:03 & 9:24 & 9:57 & 10:27 & 10:57 \\
Delanco & 6:31 & 6:40 & 7:01 & 7:22 & 7:43 & 8:04 & 8:25 & 8:46 & 9:07 & 9:28 & 10:01 & 10:31 & 11:01 \\
Riverside & 6:33 & 6:42 & 7:03 & 7:24 & 7:45 & 8:06 & 8:27 & 8:48 & 9:09 & 9:30 & 10:03 & 10:33 & 11:03 \\
Cinnaminson & 6:38 & 6:47 & 7:08 & 7:29 & 7:50 & 8:11 & 8:32 & 8:53 & 9:14 & 9:35 & 10:08 & 10:38 & 11:08 \\
Riverton & 6:40 & 6:49 & 7:10 & 7:31 & 7:52 & 8:13 & 8:34 & 8:55 & 9:16 & 9:37 & 10:10 & 10:40 & 11:10 \\
Palmyra & 6:42 & 6:51 & 7:12 & 7:33 & 7:54 & 8:15 & 8:36 & 8:57 & 9:18 & 9:39 & 10:12 & 10:42 & 11:12 \\
Route 73/Pennsauken & 6:45 & 6:54 & 7:15 & 7:36 & 7:57 & 8:18 & 8:39 & 9:00 & 9:21 & 9:42 & 10:15 & 10:45 & 11:15 \\
Pennsauken TC & 6:47 & 6:56 & 7:17 & 7:38 & 7:59 & 8:20 & 8:41 & 9:02 & 9:23 & 9:44 & 10:17 & 10:47 & 11:17 \\
36th Street & 6:51 & 7:06 & 7:27 & 7:48 & 8:09 & 8:30 & 8:51 & 9:12 & 9:33 & 9:54 & 10:21 & 10:51 & 11:21 \\
Walter Rand TC & 6:56 & 7:11 & 7:32 & 7:53 & 8:14 & 8:35 & 8:56 & 9:17 & 9:38 & 9:59 & 10:26 & 10:56 & 11:26 \\
Cooper St/Rutgers & 6:59 & 7:14 & 7:35 & 7:56 & 8:17 & 8:38 & 8:59 & 9:20 & 9:41 & 10:02 & 10:29 & 10:59 & 11:29 \\
Aquarium & 7:01 & 7:16 & 7:37 & 7:58 & 8:19 & 8:40 & 9:01 & 9:22 & 9:43 & 10:04 & 10:31 & 11:01 & 11:31 \\
Entertainment Center & 7:04 & 7:19 & 7:40 & 8:01 & 8:22 & 8:43 & 9:04 & 9:25 & 9:46 & 10:07 & 10:34 & 11:04 & 11:34 \\
\bottomrule\end{tabular}\end{center}
\subsection*{To Camden (read down) --- trips 14--26 of 37}
\begin{center}\footnotesize\setlength{\tabcolsep}{3pt}\begin{tabular}{@{}lrrrrrrrrrrrrr@{}}\toprule
Station & 14 & 15 & 16 & 17 & 18 & 19 & 20 & 21 & 22 & 23 & 24 & 25 & 26 \\\midrule
Trenton TC & 10:57 & 11:27 & 11:57 & 12:27 & 12:57 & 13:27 & 13:57 & 14:27 & 14:57 & 15:12 & 15:33 & 15:54 & 16:15 \\
Hamilton Avenue & 10:59 & 11:29 & 11:59 & 12:29 & 12:59 & 13:29 & 13:59 & 14:29 & 14:59 & 15:14 & 15:35 & 15:56 & 16:17 \\
Cass Street & 11:01 & 11:31 & 12:01 & 12:31 & 13:01 & 13:31 & 14:01 & 14:31 & 15:01 & 15:16 & 15:37 & 15:58 & 16:19 \\
Bordentown & 11:08 & 11:38 & 12:08 & 12:38 & 13:08 & 13:38 & 14:08 & 14:38 & 15:08 & 15:23 & 15:44 & 16:05 & 16:26 \\
Roebling & 11:13 & 11:43 & 12:13 & 12:43 & 13:13 & 13:43 & 14:13 & 14:43 & 15:13 & 15:28 & 15:49 & 16:10 & 16:31 \\
Florence & 11:16 & 11:46 & 12:16 & 12:46 & 13:16 & 13:46 & 14:16 & 14:46 & 15:16 & 15:31 & 15:52 & 16:13 & 16:34 \\
Burlington Towne Ctr & 11:21 & 11:51 & 12:21 & 12:51 & 13:21 & 13:51 & 14:21 & 14:51 & 15:21 & 15:36 & 15:57 & 16:18 & 16:39 \\
Burlington South & 11:23 & 11:53 & 12:23 & 12:53 & 13:23 & 13:53 & 14:23 & 14:53 & 15:23 & 15:38 & 15:59 & 16:20 & 16:41 \\
Beverly/Edgewater Park & 11:27 & 11:57 & 12:27 & 12:57 & 13:27 & 13:57 & 14:27 & 14:57 & 15:27 & 15:42 & 16:03 & 16:24 & 16:45 \\
Delanco & 11:31 & 12:01 & 12:31 & 13:01 & 13:31 & 14:01 & 14:31 & 15:01 & 15:31 & 15:46 & 16:07 & 16:28 & 16:49 \\
Riverside & 11:33 & 12:03 & 12:33 & 13:03 & 13:33 & 14:03 & 14:33 & 15:03 & 15:33 & 15:48 & 16:09 & 16:30 & 16:51 \\
Cinnaminson & 11:38 & 12:08 & 12:38 & 13:08 & 13:38 & 14:08 & 14:38 & 15:08 & 15:38 & 15:53 & 16:14 & 16:35 & 16:56 \\
Riverton & 11:40 & 12:10 & 12:40 & 13:10 & 13:40 & 14:10 & 14:40 & 15:10 & 15:40 & 15:55 & 16:16 & 16:37 & 16:58 \\
Palmyra & 11:42 & 12:12 & 12:42 & 13:12 & 13:42 & 14:12 & 14:42 & 15:12 & 15:42 & 15:57 & 16:18 & 16:39 & 17:00 \\
Route 73/Pennsauken & 11:45 & 12:15 & 12:45 & 13:15 & 13:45 & 14:15 & 14:45 & 15:15 & 15:45 & 16:00 & 16:21 & 16:42 & 17:03 \\
Pennsauken TC & 11:47 & 12:17 & 12:47 & 13:17 & 13:47 & 14:17 & 14:47 & 15:17 & 15:47 & 16:02 & 16:23 & 16:44 & 17:05 \\
36th Street & 11:51 & 12:21 & 12:51 & 13:21 & 13:51 & 14:21 & 14:51 & 15:21 & 15:51 & 16:12 & 16:33 & 16:54 & 17:15 \\
Walter Rand TC & 11:56 & 12:26 & 12:56 & 13:26 & 13:56 & 14:26 & 14:56 & 15:26 & 15:56 & 16:17 & 16:38 & 16:59 & 17:20 \\
Cooper St/Rutgers & 11:59 & 12:29 & 12:59 & 13:29 & 13:59 & 14:29 & 14:59 & 15:29 & 15:59 & 16:20 & 16:41 & 17:02 & 17:23 \\
Aquarium & 12:01 & 12:31 & 13:01 & 13:31 & 14:01 & 14:31 & 15:01 & 15:31 & 16:01 & 16:22 & 16:43 & 17:04 & 17:25 \\
Entertainment Center & 12:04 & 12:34 & 13:04 & 13:34 & 14:04 & 14:34 & 15:04 & 15:34 & 16:04 & 16:25 & 16:46 & 17:07 & 17:28 \\
\bottomrule\end{tabular}\end{center}
\subsection*{To Camden (read down) --- trips 27--37 of 37}
\begin{center}\footnotesize\setlength{\tabcolsep}{3pt}\begin{tabular}{@{}lrrrrrrrrrrr@{}}\toprule
Station & 27 & 28 & 29 & 30 & 31 & 32 & 33 & 34 & 35 & 36 & 37 \\\midrule
Trenton TC & 16:36 & 16:57 & 17:18 & 17:39 & 18:00 & 18:21 & 18:42 & 19:27 & 19:57 & 20:27 & 20:57 \\
Hamilton Avenue & 16:38 & 16:59 & 17:20 & 17:41 & 18:02 & 18:23 & 18:44 & 19:29 & 19:59 & 20:29 & 20:59 \\
Cass Street & 16:40 & 17:01 & 17:22 & 17:43 & 18:04 & 18:25 & 18:46 & 19:31 & 20:01 & 20:31 & 21:01 \\
Bordentown & 16:47 & 17:08 & 17:29 & 17:50 & 18:11 & 18:32 & 18:53 & 19:38 & 20:08 & 20:38 & 21:08 \\
Roebling & 16:52 & 17:13 & 17:34 & 17:55 & 18:16 & 18:37 & 18:58 & 19:43 & 20:13 & 20:43 & 21:13 \\
Florence & 16:55 & 17:16 & 17:37 & 17:58 & 18:19 & 18:40 & 19:01 & 19:46 & 20:16 & 20:46 & 21:16 \\
Burlington Towne Ctr & 17:00 & 17:21 & 17:42 & 18:03 & 18:24 & 18:45 & 19:06 & 19:51 & 20:21 & 20:51 & 21:21 \\
Burlington South & 17:02 & 17:23 & 17:44 & 18:05 & 18:26 & 18:47 & 19:08 & 19:53 & 20:23 & 20:53 & 21:23 \\
Beverly/Edgewater Park & 17:06 & 17:27 & 17:48 & 18:09 & 18:30 & 18:51 & 19:12 & 19:57 & 20:27 & 20:57 & 21:27 \\
Delanco & 17:10 & 17:31 & 17:52 & 18:13 & 18:34 & 18:55 & 19:16 & 20:01 & 20:31 & 21:01 & 21:31 \\
Riverside & 17:12 & 17:33 & 17:54 & 18:15 & 18:36 & 18:57 & 19:18 & 20:03 & 20:33 & 21:03 & 21:33 \\
Cinnaminson & 17:17 & 17:38 & 17:59 & 18:20 & 18:41 & 19:02 & 19:23 & 20:08 & 20:38 & 21:08 & 21:38 \\
Riverton & 17:19 & 17:40 & 18:01 & 18:22 & 18:43 & 19:04 & 19:25 & 20:10 & 20:40 & 21:10 & 21:40 \\
Palmyra & 17:21 & 17:42 & 18:03 & 18:24 & 18:45 & 19:06 & 19:27 & 20:12 & 20:42 & 21:12 & 21:42 \\
Route 73/Pennsauken & 17:24 & 17:45 & 18:06 & 18:27 & 18:48 & 19:09 & 19:30 & 20:15 & 20:45 & 21:15 & 21:45 \\
Pennsauken TC & 17:26 & 17:47 & 18:08 & 18:29 & 18:50 & 19:11 & 19:32 & 20:17 & 20:47 & 21:17 & 21:47 \\
36th Street & 17:36 & 17:57 & 18:18 & 18:39 & 19:00 & 19:21 & 19:42 & 20:21 & 20:51 & 21:21 & 21:51 \\
Walter Rand TC & 17:41 & 18:02 & 18:23 & 18:44 & 19:05 & 19:26 & 19:47 & 20:26 & 20:56 & 21:26 & 21:56 \\
Cooper St/Rutgers & 17:44 & 18:05 & 18:26 & 18:47 & 19:08 & 19:29 & 19:50 & 20:29 & 20:59 & 21:29 & 21:59 \\
Aquarium & 17:46 & 18:07 & 18:28 & 18:49 & 19:10 & 19:31 & 19:52 & 20:31 & 21:01 & 21:31 & 22:01 \\
Entertainment Center & 17:49 & 18:10 & 18:31 & 18:52 & 19:13 & 19:34 & 19:55 & 20:34 & 21:04 & 21:34 & 22:04 \\
\bottomrule\end{tabular}\end{center}
\clearpage
\section{Uniform U7: every 23 min peak, 30 off-peak}\label{tt:uniform_U7}
\subsection*{To Trenton (read down) --- trips 1--13 of 35}
\begin{center}\footnotesize\setlength{\tabcolsep}{3pt}\begin{tabular}{@{}lrrrrrrrrrrrrr@{}}\toprule
Station & 1 & 2 & 3 & 4 & 5 & 6 & 7 & 8 & 9 & 10 & 11 & 12 & 13 \\\midrule
Entertainment Center & 5:45 & 6:08 & 6:31 & 6:54 & 7:17 & 7:40 & 8:03 & 8:26 & 8:49 & 9:15 & 9:45 & 10:15 & 10:45 \\
Aquarium & 5:47 & 6:10 & 6:33 & 6:56 & 7:19 & 7:42 & 8:05 & 8:28 & 8:51 & 9:17 & 9:47 & 10:17 & 10:47 \\
Cooper St/Rutgers & 5:49 & 6:12 & 6:35 & 6:58 & 7:21 & 7:44 & 8:07 & 8:30 & 8:53 & 9:19 & 9:49 & 10:19 & 10:49 \\
Walter Rand TC & 5:53 & 6:16 & 6:39 & 7:02 & 7:25 & 7:48 & 8:11 & 8:34 & 8:57 & 9:23 & 9:53 & 10:23 & 10:53 \\
36th Street & 5:58 & 6:21 & 6:44 & 7:07 & 7:30 & 7:53 & 8:16 & 8:39 & 9:02 & 9:28 & 9:58 & 10:28 & 10:58 \\
Pennsauken TC & 6:00 & 6:23 & 6:46 & 7:09 & 7:32 & 7:55 & 8:18 & 8:41 & 9:04 & 9:30 & 10:00 & 10:30 & 11:00 \\
Route 73/Pennsauken & 6:03 & 6:26 & 6:49 & 7:12 & 7:35 & 7:58 & 8:21 & 8:44 & 9:07 & 9:33 & 10:03 & 10:33 & 11:03 \\
Palmyra & 6:06 & 6:29 & 6:52 & 7:15 & 7:38 & 8:01 & 8:24 & 8:47 & 9:10 & 9:36 & 10:06 & 10:36 & 11:06 \\
Riverton & 6:08 & 6:31 & 6:54 & 7:17 & 7:40 & 8:03 & 8:26 & 8:49 & 9:12 & 9:38 & 10:08 & 10:38 & 11:08 \\
Cinnaminson & 6:10 & 6:33 & 6:56 & 7:19 & 7:42 & 8:05 & 8:28 & 8:51 & 9:14 & 9:40 & 10:10 & 10:40 & 11:10 \\
Riverside & 6:15 & 6:38 & 7:01 & 7:24 & 7:47 & 8:10 & 8:33 & 8:56 & 9:19 & 9:45 & 10:15 & 10:45 & 11:15 \\
Delanco & 6:17 & 6:40 & 7:03 & 7:26 & 7:49 & 8:12 & 8:35 & 8:58 & 9:21 & 9:47 & 10:17 & 10:47 & 11:17 \\
Beverly/Edgewater Park & 6:21 & 6:44 & 7:07 & 7:30 & 7:53 & 8:16 & 8:39 & 9:02 & 9:25 & 9:51 & 10:21 & 10:51 & 11:21 \\
Burlington South & 6:25 & 6:48 & 7:11 & 7:34 & 7:57 & 8:20 & 8:43 & 9:06 & 9:29 & 9:55 & 10:25 & 10:55 & 11:25 \\
Burlington Towne Ctr & 6:27 & 6:50 & 7:13 & 7:36 & 7:59 & 8:22 & 8:45 & 9:08 & 9:31 & 9:57 & 10:27 & 10:57 & 11:27 \\
Florence & 6:33 & 6:56 & 7:19 & 7:42 & 8:05 & 8:28 & 8:51 & 9:14 & 9:37 & 10:03 & 10:33 & 11:03 & 11:33 \\
Roebling & 6:36 & 6:59 & 7:22 & 7:45 & 8:08 & 8:31 & 8:54 & 9:17 & 9:40 & 10:06 & 10:36 & 11:06 & 11:36 \\
Bordentown & 6:41 & 7:04 & 7:27 & 7:50 & 8:13 & 8:36 & 8:59 & 9:22 & 9:45 & 10:11 & 10:41 & 11:11 & 11:41 \\
Cass Street & 6:47 & 7:10 & 7:33 & 7:56 & 8:19 & 8:42 & 9:05 & 9:28 & 9:51 & 10:17 & 10:47 & 11:17 & 11:47 \\
Hamilton Avenue & 6:49 & 7:12 & 7:35 & 7:58 & 8:21 & 8:44 & 9:07 & 9:30 & 9:53 & 10:19 & 10:49 & 11:19 & 11:49 \\
Trenton TC & 6:52 & 7:15 & 7:38 & 8:01 & 8:24 & 8:47 & 9:10 & 9:33 & 9:56 & 10:22 & 10:52 & 11:22 & 11:52 \\
\bottomrule\end{tabular}\end{center}
\subsection*{To Trenton (read down) --- trips 14--26 of 35}
\begin{center}\footnotesize\setlength{\tabcolsep}{3pt}\begin{tabular}{@{}lrrrrrrrrrrrrr@{}}\toprule
Station & 14 & 15 & 16 & 17 & 18 & 19 & 20 & 21 & 22 & 23 & 24 & 25 & 26 \\\midrule
Entertainment Center & 11:15 & 11:45 & 12:15 & 12:45 & 13:15 & 13:45 & 14:15 & 14:45 & 15:20 & 15:43 & 16:06 & 16:29 & 16:52 \\
Aquarium & 11:17 & 11:47 & 12:17 & 12:47 & 13:17 & 13:47 & 14:17 & 14:47 & 15:22 & 15:45 & 16:08 & 16:31 & 16:54 \\
Cooper St/Rutgers & 11:19 & 11:49 & 12:19 & 12:49 & 13:19 & 13:49 & 14:19 & 14:49 & 15:24 & 15:47 & 16:10 & 16:33 & 16:56 \\
Walter Rand TC & 11:23 & 11:53 & 12:23 & 12:53 & 13:23 & 13:53 & 14:23 & 14:53 & 15:28 & 15:51 & 16:14 & 16:37 & 17:00 \\
36th Street & 11:28 & 11:58 & 12:28 & 12:58 & 13:28 & 13:58 & 14:28 & 14:58 & 15:33 & 15:56 & 16:19 & 16:42 & 17:05 \\
Pennsauken TC & 11:30 & 12:00 & 12:30 & 13:00 & 13:30 & 14:00 & 14:30 & 15:00 & 15:35 & 15:58 & 16:21 & 16:44 & 17:07 \\
Route 73/Pennsauken & 11:33 & 12:03 & 12:33 & 13:03 & 13:33 & 14:03 & 14:33 & 15:03 & 15:38 & 16:01 & 16:24 & 16:47 & 17:10 \\
Palmyra & 11:36 & 12:06 & 12:36 & 13:06 & 13:36 & 14:06 & 14:36 & 15:06 & 15:41 & 16:04 & 16:27 & 16:50 & 17:13 \\
Riverton & 11:38 & 12:08 & 12:38 & 13:08 & 13:38 & 14:08 & 14:38 & 15:08 & 15:43 & 16:06 & 16:29 & 16:52 & 17:15 \\
Cinnaminson & 11:40 & 12:10 & 12:40 & 13:10 & 13:40 & 14:10 & 14:40 & 15:10 & 15:45 & 16:08 & 16:31 & 16:54 & 17:17 \\
Riverside & 11:45 & 12:15 & 12:45 & 13:15 & 13:45 & 14:15 & 14:45 & 15:15 & 15:50 & 16:13 & 16:36 & 16:59 & 17:22 \\
Delanco & 11:47 & 12:17 & 12:47 & 13:17 & 13:47 & 14:17 & 14:47 & 15:17 & 15:52 & 16:15 & 16:38 & 17:01 & 17:24 \\
Beverly/Edgewater Park & 11:51 & 12:21 & 12:51 & 13:21 & 13:51 & 14:21 & 14:51 & 15:21 & 15:56 & 16:19 & 16:42 & 17:05 & 17:28 \\
Burlington South & 11:55 & 12:25 & 12:55 & 13:25 & 13:55 & 14:25 & 14:55 & 15:25 & 16:00 & 16:23 & 16:46 & 17:09 & 17:32 \\
Burlington Towne Ctr & 11:57 & 12:27 & 12:57 & 13:27 & 13:57 & 14:27 & 14:57 & 15:27 & 16:02 & 16:25 & 16:48 & 17:11 & 17:34 \\
Florence & 12:03 & 12:33 & 13:03 & 13:33 & 14:03 & 14:33 & 15:03 & 15:33 & 16:08 & 16:31 & 16:54 & 17:17 & 17:40 \\
Roebling & 12:06 & 12:36 & 13:06 & 13:36 & 14:06 & 14:36 & 15:06 & 15:36 & 16:11 & 16:34 & 16:57 & 17:20 & 17:43 \\
Bordentown & 12:11 & 12:41 & 13:11 & 13:41 & 14:11 & 14:41 & 15:11 & 15:41 & 16:16 & 16:39 & 17:02 & 17:25 & 17:48 \\
Cass Street & 12:17 & 12:47 & 13:17 & 13:47 & 14:17 & 14:47 & 15:17 & 15:47 & 16:22 & 16:45 & 17:08 & 17:31 & 17:54 \\
Hamilton Avenue & 12:19 & 12:49 & 13:19 & 13:49 & 14:19 & 14:49 & 15:19 & 15:49 & 16:24 & 16:47 & 17:10 & 17:33 & 17:56 \\
Trenton TC & 12:22 & 12:52 & 13:22 & 13:52 & 14:22 & 14:52 & 15:22 & 15:52 & 16:27 & 16:50 & 17:13 & 17:36 & 17:59 \\
\bottomrule\end{tabular}\end{center}
\subsection*{To Trenton (read down) --- trips 27--35 of 35}
\begin{center}\footnotesize\setlength{\tabcolsep}{3pt}\begin{tabular}{@{}lrrrrrrrrr@{}}\toprule
Station & 27 & 28 & 29 & 30 & 31 & 32 & 33 & 34 & 35 \\\midrule
Entertainment Center & 17:15 & 17:38 & 18:01 & 18:24 & 18:47 & 19:15 & 19:45 & 20:15 & 20:45 \\
Aquarium & 17:17 & 17:40 & 18:03 & 18:26 & 18:49 & 19:17 & 19:47 & 20:17 & 20:47 \\
Cooper St/Rutgers & 17:19 & 17:42 & 18:05 & 18:28 & 18:51 & 19:19 & 19:49 & 20:19 & 20:49 \\
Walter Rand TC & 17:23 & 17:46 & 18:09 & 18:32 & 18:55 & 19:23 & 19:53 & 20:23 & 20:53 \\
36th Street & 17:28 & 17:51 & 18:14 & 18:37 & 19:00 & 19:28 & 19:58 & 20:28 & 20:58 \\
Pennsauken TC & 17:30 & 17:53 & 18:16 & 18:39 & 19:02 & 19:30 & 20:00 & 20:30 & 21:00 \\
Route 73/Pennsauken & 17:33 & 17:56 & 18:19 & 18:42 & 19:05 & 19:33 & 20:03 & 20:33 & 21:03 \\
Palmyra & 17:36 & 17:59 & 18:22 & 18:45 & 19:08 & 19:36 & 20:06 & 20:36 & 21:06 \\
Riverton & 17:38 & 18:01 & 18:24 & 18:47 & 19:10 & 19:38 & 20:08 & 20:38 & 21:08 \\
Cinnaminson & 17:40 & 18:03 & 18:26 & 18:49 & 19:12 & 19:40 & 20:10 & 20:40 & 21:10 \\
Riverside & 17:45 & 18:08 & 18:31 & 18:54 & 19:17 & 19:45 & 20:15 & 20:45 & 21:15 \\
Delanco & 17:47 & 18:10 & 18:33 & 18:56 & 19:19 & 19:47 & 20:17 & 20:47 & 21:17 \\
Beverly/Edgewater Park & 17:51 & 18:14 & 18:37 & 19:00 & 19:23 & 19:51 & 20:21 & 20:51 & 21:21 \\
Burlington South & 17:55 & 18:18 & 18:41 & 19:04 & 19:27 & 19:55 & 20:25 & 20:55 & 21:25 \\
Burlington Towne Ctr & 17:57 & 18:20 & 18:43 & 19:06 & 19:29 & 19:57 & 20:27 & 20:57 & 21:27 \\
Florence & 18:03 & 18:26 & 18:49 & 19:12 & 19:35 & 20:03 & 20:33 & 21:03 & 21:33 \\
Roebling & 18:06 & 18:29 & 18:52 & 19:15 & 19:38 & 20:06 & 20:36 & 21:06 & 21:36 \\
Bordentown & 18:11 & 18:34 & 18:57 & 19:20 & 19:43 & 20:11 & 20:41 & 21:11 & 21:41 \\
Cass Street & 18:17 & 18:40 & 19:03 & 19:26 & 19:49 & 20:17 & 20:47 & 21:17 & 21:47 \\
Hamilton Avenue & 18:19 & 18:42 & 19:05 & 19:28 & 19:51 & 20:19 & 20:49 & 21:19 & 21:49 \\
Trenton TC & 18:22 & 18:45 & 19:08 & 19:31 & 19:54 & 20:22 & 20:52 & 21:22 & 21:52 \\
\bottomrule\end{tabular}\end{center}
\subsection*{To Camden (read down) --- trips 1--13 of 36}
\begin{center}\footnotesize\setlength{\tabcolsep}{3pt}\begin{tabular}{@{}lrrrrrrrrrrrrr@{}}\toprule
Station & 1 & 2 & 3 & 4 & 5 & 6 & 7 & 8 & 9 & 10 & 11 & 12 & 13 \\\midrule
Trenton TC & 5:57 & 6:14 & 6:37 & 7:00 & 7:23 & 7:46 & 8:09 & 8:32 & 8:55 & 9:27 & 9:57 & 10:27 & 10:57 \\
Hamilton Avenue & 5:59 & 6:16 & 6:39 & 7:02 & 7:25 & 7:48 & 8:11 & 8:34 & 8:57 & 9:29 & 9:59 & 10:29 & 10:59 \\
Cass Street & 6:01 & 6:18 & 6:41 & 7:04 & 7:27 & 7:50 & 8:13 & 8:36 & 8:59 & 9:31 & 10:01 & 10:31 & 11:01 \\
Bordentown & 6:08 & 6:25 & 6:48 & 7:11 & 7:34 & 7:57 & 8:20 & 8:43 & 9:06 & 9:38 & 10:08 & 10:38 & 11:08 \\
Roebling & 6:13 & 6:30 & 6:53 & 7:16 & 7:39 & 8:02 & 8:25 & 8:48 & 9:11 & 9:43 & 10:13 & 10:43 & 11:13 \\
Florence & 6:16 & 6:33 & 6:56 & 7:19 & 7:42 & 8:05 & 8:28 & 8:51 & 9:14 & 9:46 & 10:16 & 10:46 & 11:16 \\
Burlington Towne Ctr & 6:21 & 6:38 & 7:01 & 7:24 & 7:47 & 8:10 & 8:33 & 8:56 & 9:19 & 9:51 & 10:21 & 10:51 & 11:21 \\
Burlington South & 6:23 & 6:40 & 7:03 & 7:26 & 7:49 & 8:12 & 8:35 & 8:58 & 9:21 & 9:53 & 10:23 & 10:53 & 11:23 \\
Beverly/Edgewater Park & 6:27 & 6:44 & 7:07 & 7:30 & 7:53 & 8:16 & 8:39 & 9:02 & 9:25 & 9:57 & 10:27 & 10:57 & 11:27 \\
Delanco & 6:31 & 6:48 & 7:11 & 7:34 & 7:57 & 8:20 & 8:43 & 9:06 & 9:29 & 10:01 & 10:31 & 11:01 & 11:31 \\
Riverside & 6:33 & 6:50 & 7:13 & 7:36 & 7:59 & 8:22 & 8:45 & 9:08 & 9:31 & 10:03 & 10:33 & 11:03 & 11:33 \\
Cinnaminson & 6:38 & 6:58 & 7:21 & 7:44 & 8:07 & 8:30 & 8:53 & 9:16 & 9:39 & 10:08 & 10:38 & 11:08 & 11:38 \\
Riverton & 6:40 & 7:00 & 7:23 & 7:46 & 8:09 & 8:32 & 8:55 & 9:18 & 9:41 & 10:10 & 10:40 & 11:10 & 11:40 \\
Palmyra & 6:42 & 7:03 & 7:26 & 7:49 & 8:12 & 8:35 & 8:58 & 9:21 & 9:44 & 10:12 & 10:42 & 11:12 & 11:42 \\
Route 73/Pennsauken & 6:45 & 7:06 & 7:29 & 7:52 & 8:15 & 8:38 & 9:01 & 9:24 & 9:47 & 10:15 & 10:45 & 11:15 & 11:45 \\
Pennsauken TC & 6:47 & 7:08 & 7:31 & 7:54 & 8:17 & 8:40 & 9:03 & 9:26 & 9:49 & 10:17 & 10:47 & 11:17 & 11:47 \\
36th Street & 6:51 & 7:12 & 7:35 & 7:58 & 8:21 & 8:44 & 9:07 & 9:30 & 9:53 & 10:21 & 10:51 & 11:21 & 11:51 \\
Walter Rand TC & 6:56 & 7:17 & 7:40 & 8:03 & 8:26 & 8:49 & 9:12 & 9:35 & 9:58 & 10:26 & 10:56 & 11:26 & 11:56 \\
Cooper St/Rutgers & 6:59 & 7:20 & 7:43 & 8:06 & 8:29 & 8:52 & 9:15 & 9:38 & 10:01 & 10:29 & 10:59 & 11:29 & 11:59 \\
Aquarium & 7:01 & 7:22 & 7:45 & 8:08 & 8:31 & 8:54 & 9:17 & 9:40 & 10:03 & 10:31 & 11:01 & 11:31 & 12:01 \\
Entertainment Center & 7:04 & 7:25 & 7:48 & 8:11 & 8:34 & 8:57 & 9:20 & 9:43 & 10:06 & 10:34 & 11:04 & 11:34 & 12:04 \\
\bottomrule\end{tabular}\end{center}
\subsection*{To Camden (read down) --- trips 14--26 of 36}
\begin{center}\footnotesize\setlength{\tabcolsep}{3pt}\begin{tabular}{@{}lrrrrrrrrrrrrr@{}}\toprule
Station & 14 & 15 & 16 & 17 & 18 & 19 & 20 & 21 & 22 & 23 & 24 & 25 & 26 \\\midrule
Trenton TC & 11:27 & 11:57 & 12:27 & 12:57 & 13:27 & 13:57 & 14:27 & 14:57 & 15:03 & 15:26 & 15:49 & 16:12 & 16:35 \\
Hamilton Avenue & 11:29 & 11:59 & 12:29 & 12:59 & 13:29 & 13:59 & 14:29 & 14:59 & 15:05 & 15:28 & 15:51 & 16:14 & 16:37 \\
Cass Street & 11:31 & 12:01 & 12:31 & 13:01 & 13:31 & 14:01 & 14:31 & 15:01 & 15:07 & 15:30 & 15:53 & 16:16 & 16:39 \\
Bordentown & 11:38 & 12:08 & 12:38 & 13:08 & 13:38 & 14:08 & 14:38 & 15:08 & 15:14 & 15:37 & 16:00 & 16:23 & 16:46 \\
Roebling & 11:43 & 12:13 & 12:43 & 13:13 & 13:43 & 14:13 & 14:43 & 15:13 & 15:19 & 15:42 & 16:05 & 16:28 & 16:51 \\
Florence & 11:46 & 12:16 & 12:46 & 13:16 & 13:46 & 14:16 & 14:46 & 15:16 & 15:22 & 15:45 & 16:08 & 16:31 & 16:54 \\
Burlington Towne Ctr & 11:51 & 12:21 & 12:51 & 13:21 & 13:51 & 14:21 & 14:51 & 15:21 & 15:27 & 15:50 & 16:13 & 16:36 & 16:59 \\
Burlington South & 11:53 & 12:23 & 12:53 & 13:23 & 13:53 & 14:23 & 14:53 & 15:23 & 15:29 & 15:52 & 16:15 & 16:38 & 17:01 \\
Beverly/Edgewater Park & 11:57 & 12:27 & 12:57 & 13:27 & 13:57 & 14:27 & 14:57 & 15:27 & 15:33 & 15:56 & 16:19 & 16:42 & 17:05 \\
Delanco & 12:01 & 12:31 & 13:01 & 13:31 & 14:01 & 14:31 & 15:01 & 15:31 & 15:37 & 16:00 & 16:23 & 16:46 & 17:09 \\
Riverside & 12:03 & 12:33 & 13:03 & 13:33 & 14:03 & 14:33 & 15:03 & 15:33 & 15:39 & 16:02 & 16:25 & 16:48 & 17:11 \\
Cinnaminson & 12:08 & 12:38 & 13:08 & 13:38 & 14:08 & 14:38 & 15:08 & 15:38 & 15:47 & 16:10 & 16:33 & 16:56 & 17:19 \\
Riverton & 12:10 & 12:40 & 13:10 & 13:40 & 14:10 & 14:40 & 15:10 & 15:40 & 15:49 & 16:12 & 16:35 & 16:58 & 17:21 \\
Palmyra & 12:12 & 12:42 & 13:12 & 13:42 & 14:12 & 14:42 & 15:12 & 15:42 & 15:52 & 16:15 & 16:38 & 17:01 & 17:24 \\
Route 73/Pennsauken & 12:15 & 12:45 & 13:15 & 13:45 & 14:15 & 14:45 & 15:15 & 15:45 & 15:55 & 16:18 & 16:41 & 17:04 & 17:27 \\
Pennsauken TC & 12:17 & 12:47 & 13:17 & 13:47 & 14:17 & 14:47 & 15:17 & 15:47 & 15:57 & 16:20 & 16:43 & 17:06 & 17:29 \\
36th Street & 12:21 & 12:51 & 13:21 & 13:51 & 14:21 & 14:51 & 15:21 & 15:51 & 16:01 & 16:24 & 16:47 & 17:10 & 17:33 \\
Walter Rand TC & 12:26 & 12:56 & 13:26 & 13:56 & 14:26 & 14:56 & 15:26 & 15:56 & 16:06 & 16:29 & 16:52 & 17:15 & 17:38 \\
Cooper St/Rutgers & 12:29 & 12:59 & 13:29 & 13:59 & 14:29 & 14:59 & 15:29 & 15:59 & 16:09 & 16:32 & 16:55 & 17:18 & 17:41 \\
Aquarium & 12:31 & 13:01 & 13:31 & 14:01 & 14:31 & 15:01 & 15:31 & 16:01 & 16:11 & 16:34 & 16:57 & 17:20 & 17:43 \\
Entertainment Center & 12:34 & 13:04 & 13:34 & 14:04 & 14:34 & 15:04 & 15:34 & 16:04 & 16:14 & 16:37 & 17:00 & 17:23 & 17:46 \\
\bottomrule\end{tabular}\end{center}
\subsection*{To Camden (read down) --- trips 27--36 of 36}
\begin{center}\footnotesize\setlength{\tabcolsep}{3pt}\begin{tabular}{@{}lrrrrrrrrrr@{}}\toprule
Station & 27 & 28 & 29 & 30 & 31 & 32 & 33 & 34 & 35 & 36 \\\midrule
Trenton TC & 16:58 & 17:21 & 17:44 & 18:07 & 18:30 & 18:53 & 19:27 & 19:57 & 20:27 & 20:57 \\
Hamilton Avenue & 17:00 & 17:23 & 17:46 & 18:09 & 18:32 & 18:55 & 19:29 & 19:59 & 20:29 & 20:59 \\
Cass Street & 17:02 & 17:25 & 17:48 & 18:11 & 18:34 & 18:57 & 19:31 & 20:01 & 20:31 & 21:01 \\
Bordentown & 17:09 & 17:32 & 17:55 & 18:18 & 18:41 & 19:04 & 19:38 & 20:08 & 20:38 & 21:08 \\
Roebling & 17:14 & 17:37 & 18:00 & 18:23 & 18:46 & 19:09 & 19:43 & 20:13 & 20:43 & 21:13 \\
Florence & 17:17 & 17:40 & 18:03 & 18:26 & 18:49 & 19:12 & 19:46 & 20:16 & 20:46 & 21:16 \\
Burlington Towne Ctr & 17:22 & 17:45 & 18:08 & 18:31 & 18:54 & 19:17 & 19:51 & 20:21 & 20:51 & 21:21 \\
Burlington South & 17:24 & 17:47 & 18:10 & 18:33 & 18:56 & 19:19 & 19:53 & 20:23 & 20:53 & 21:23 \\
Beverly/Edgewater Park & 17:28 & 17:51 & 18:14 & 18:37 & 19:00 & 19:23 & 19:57 & 20:27 & 20:57 & 21:27 \\
Delanco & 17:32 & 17:55 & 18:18 & 18:41 & 19:04 & 19:27 & 20:01 & 20:31 & 21:01 & 21:31 \\
Riverside & 17:34 & 17:57 & 18:20 & 18:43 & 19:06 & 19:29 & 20:03 & 20:33 & 21:03 & 21:33 \\
Cinnaminson & 17:42 & 18:05 & 18:28 & 18:51 & 19:14 & 19:37 & 20:08 & 20:38 & 21:08 & 21:38 \\
Riverton & 17:44 & 18:07 & 18:30 & 18:53 & 19:16 & 19:39 & 20:10 & 20:40 & 21:10 & 21:40 \\
Palmyra & 17:47 & 18:10 & 18:33 & 18:56 & 19:19 & 19:42 & 20:12 & 20:42 & 21:12 & 21:42 \\
Route 73/Pennsauken & 17:50 & 18:13 & 18:36 & 18:59 & 19:22 & 19:45 & 20:15 & 20:45 & 21:15 & 21:45 \\
Pennsauken TC & 17:52 & 18:15 & 18:38 & 19:01 & 19:24 & 19:47 & 20:17 & 20:47 & 21:17 & 21:47 \\
36th Street & 17:56 & 18:19 & 18:42 & 19:05 & 19:28 & 19:51 & 20:21 & 20:51 & 21:21 & 21:51 \\
Walter Rand TC & 18:01 & 18:24 & 18:47 & 19:10 & 19:33 & 19:56 & 20:26 & 20:56 & 21:26 & 21:56 \\
Cooper St/Rutgers & 18:04 & 18:27 & 18:50 & 19:13 & 19:36 & 19:59 & 20:29 & 20:59 & 21:29 & 21:59 \\
Aquarium & 18:06 & 18:29 & 18:52 & 19:15 & 19:38 & 20:01 & 20:31 & 21:01 & 21:31 & 22:01 \\
Entertainment Center & 18:09 & 18:32 & 18:55 & 19:18 & 19:41 & 20:04 & 20:34 & 21:04 & 21:34 & 22:04 \\
\bottomrule\end{tabular}\end{center}
\clearpage
\section{Uniform U6: every 26 min peak, 30 off-peak}\label{tt:uniform_U6}
\subsection*{To Trenton (read down) --- trips 1--13 of 33}
\begin{center}\footnotesize\setlength{\tabcolsep}{3pt}\begin{tabular}{@{}lrrrrrrrrrrrrr@{}}\toprule
Station & 1 & 2 & 3 & 4 & 5 & 6 & 7 & 8 & 9 & 10 & 11 & 12 & 13 \\\midrule
Entertainment Center & 5:45 & 6:11 & 6:37 & 7:03 & 7:29 & 7:55 & 8:21 & 8:47 & 9:15 & 9:45 & 10:15 & 10:45 & 11:15 \\
Aquarium & 5:47 & 6:13 & 6:39 & 7:05 & 7:31 & 7:57 & 8:23 & 8:49 & 9:17 & 9:47 & 10:17 & 10:47 & 11:17 \\
Cooper St/Rutgers & 5:49 & 6:15 & 6:41 & 7:07 & 7:33 & 7:59 & 8:25 & 8:51 & 9:19 & 9:49 & 10:19 & 10:49 & 11:19 \\
Walter Rand TC & 5:53 & 6:19 & 6:45 & 7:11 & 7:37 & 8:03 & 8:29 & 8:55 & 9:23 & 9:53 & 10:23 & 10:53 & 11:23 \\
36th Street & 5:58 & 6:24 & 6:50 & 7:16 & 7:42 & 8:08 & 8:34 & 9:00 & 9:28 & 9:58 & 10:28 & 10:58 & 11:28 \\
Pennsauken TC & 6:00 & 6:26 & 6:52 & 7:18 & 7:44 & 8:10 & 8:36 & 9:02 & 9:30 & 10:00 & 10:30 & 11:00 & 11:30 \\
Route 73/Pennsauken & 6:03 & 6:29 & 6:55 & 7:21 & 7:47 & 8:13 & 8:39 & 9:05 & 9:33 & 10:03 & 10:33 & 11:03 & 11:33 \\
Palmyra & 6:06 & 6:32 & 6:58 & 7:24 & 7:50 & 8:16 & 8:42 & 9:08 & 9:36 & 10:06 & 10:36 & 11:06 & 11:36 \\
Riverton & 6:08 & 6:34 & 7:00 & 7:26 & 7:52 & 8:18 & 8:44 & 9:10 & 9:38 & 10:08 & 10:38 & 11:08 & 11:38 \\
Cinnaminson & 6:10 & 6:36 & 7:02 & 7:28 & 7:54 & 8:20 & 8:46 & 9:12 & 9:40 & 10:10 & 10:40 & 11:10 & 11:40 \\
Riverside & 6:15 & 6:41 & 7:07 & 7:33 & 7:59 & 8:25 & 8:51 & 9:17 & 9:45 & 10:15 & 10:45 & 11:15 & 11:45 \\
Delanco & 6:17 & 6:43 & 7:09 & 7:35 & 8:01 & 8:27 & 8:53 & 9:19 & 9:47 & 10:17 & 10:47 & 11:17 & 11:47 \\
Beverly/Edgewater Park & 6:21 & 6:47 & 7:13 & 7:39 & 8:05 & 8:31 & 8:57 & 9:23 & 9:51 & 10:21 & 10:51 & 11:21 & 11:51 \\
Burlington South & 6:25 & 6:51 & 7:17 & 7:43 & 8:09 & 8:35 & 9:01 & 9:27 & 9:55 & 10:25 & 10:55 & 11:25 & 11:55 \\
Burlington Towne Ctr & 6:27 & 6:53 & 7:19 & 7:45 & 8:11 & 8:37 & 9:03 & 9:29 & 9:57 & 10:27 & 10:57 & 11:27 & 11:57 \\
Florence & 6:33 & 6:59 & 7:25 & 7:51 & 8:17 & 8:43 & 9:09 & 9:35 & 10:03 & 10:33 & 11:03 & 11:33 & 12:03 \\
Roebling & 6:36 & 7:02 & 7:28 & 7:54 & 8:20 & 8:46 & 9:12 & 9:38 & 10:06 & 10:36 & 11:06 & 11:36 & 12:06 \\
Bordentown & 6:41 & 7:07 & 7:33 & 7:59 & 8:25 & 8:51 & 9:17 & 9:43 & 10:11 & 10:41 & 11:11 & 11:41 & 12:11 \\
Cass Street & 6:47 & 7:13 & 7:39 & 8:05 & 8:31 & 8:57 & 9:23 & 9:49 & 10:17 & 10:47 & 11:17 & 11:47 & 12:17 \\
Hamilton Avenue & 6:49 & 7:15 & 7:41 & 8:07 & 8:33 & 8:59 & 9:25 & 9:51 & 10:19 & 10:49 & 11:19 & 11:49 & 12:19 \\
Trenton TC & 6:52 & 7:18 & 7:44 & 8:10 & 8:36 & 9:02 & 9:28 & 9:54 & 10:22 & 10:52 & 11:22 & 11:52 & 12:22 \\
\bottomrule\end{tabular}\end{center}
\subsection*{To Trenton (read down) --- trips 14--26 of 33}
\begin{center}\footnotesize\setlength{\tabcolsep}{3pt}\begin{tabular}{@{}lrrrrrrrrrrrrr@{}}\toprule
Station & 14 & 15 & 16 & 17 & 18 & 19 & 20 & 21 & 22 & 23 & 24 & 25 & 26 \\\midrule
Entertainment Center & 11:45 & 12:15 & 12:45 & 13:15 & 13:45 & 14:15 & 14:45 & 15:17 & 15:43 & 16:09 & 16:35 & 17:01 & 17:27 \\
Aquarium & 11:47 & 12:17 & 12:47 & 13:17 & 13:47 & 14:17 & 14:47 & 15:19 & 15:45 & 16:11 & 16:37 & 17:03 & 17:29 \\
Cooper St/Rutgers & 11:49 & 12:19 & 12:49 & 13:19 & 13:49 & 14:19 & 14:49 & 15:21 & 15:47 & 16:13 & 16:39 & 17:05 & 17:31 \\
Walter Rand TC & 11:53 & 12:23 & 12:53 & 13:23 & 13:53 & 14:23 & 14:53 & 15:25 & 15:51 & 16:17 & 16:43 & 17:09 & 17:35 \\
36th Street & 11:58 & 12:28 & 12:58 & 13:28 & 13:58 & 14:28 & 14:58 & 15:30 & 15:56 & 16:22 & 16:48 & 17:14 & 17:40 \\
Pennsauken TC & 12:00 & 12:30 & 13:00 & 13:30 & 14:00 & 14:30 & 15:00 & 15:32 & 15:58 & 16:24 & 16:50 & 17:16 & 17:42 \\
Route 73/Pennsauken & 12:03 & 12:33 & 13:03 & 13:33 & 14:03 & 14:33 & 15:03 & 15:35 & 16:01 & 16:27 & 16:53 & 17:19 & 17:45 \\
Palmyra & 12:06 & 12:36 & 13:06 & 13:36 & 14:06 & 14:36 & 15:06 & 15:38 & 16:04 & 16:30 & 16:56 & 17:22 & 17:48 \\
Riverton & 12:08 & 12:38 & 13:08 & 13:38 & 14:08 & 14:38 & 15:08 & 15:40 & 16:06 & 16:32 & 16:58 & 17:24 & 17:50 \\
Cinnaminson & 12:10 & 12:40 & 13:10 & 13:40 & 14:10 & 14:40 & 15:10 & 15:42 & 16:08 & 16:34 & 17:00 & 17:26 & 17:52 \\
Riverside & 12:15 & 12:45 & 13:15 & 13:45 & 14:15 & 14:45 & 15:15 & 15:47 & 16:13 & 16:39 & 17:05 & 17:31 & 17:57 \\
Delanco & 12:17 & 12:47 & 13:17 & 13:47 & 14:17 & 14:47 & 15:17 & 15:49 & 16:15 & 16:41 & 17:07 & 17:33 & 17:59 \\
Beverly/Edgewater Park & 12:21 & 12:51 & 13:21 & 13:51 & 14:21 & 14:51 & 15:21 & 15:53 & 16:19 & 16:45 & 17:11 & 17:37 & 18:03 \\
Burlington South & 12:25 & 12:55 & 13:25 & 13:55 & 14:25 & 14:55 & 15:25 & 15:57 & 16:23 & 16:49 & 17:15 & 17:41 & 18:07 \\
Burlington Towne Ctr & 12:27 & 12:57 & 13:27 & 13:57 & 14:27 & 14:57 & 15:27 & 15:59 & 16:25 & 16:51 & 17:17 & 17:43 & 18:09 \\
Florence & 12:33 & 13:03 & 13:33 & 14:03 & 14:33 & 15:03 & 15:33 & 16:05 & 16:31 & 16:57 & 17:23 & 17:49 & 18:15 \\
Roebling & 12:36 & 13:06 & 13:36 & 14:06 & 14:36 & 15:06 & 15:36 & 16:08 & 16:34 & 17:00 & 17:26 & 17:52 & 18:18 \\
Bordentown & 12:41 & 13:11 & 13:41 & 14:11 & 14:41 & 15:11 & 15:41 & 16:13 & 16:39 & 17:05 & 17:31 & 17:57 & 18:23 \\
Cass Street & 12:47 & 13:17 & 13:47 & 14:17 & 14:47 & 15:17 & 15:47 & 16:19 & 16:45 & 17:11 & 17:37 & 18:03 & 18:29 \\
Hamilton Avenue & 12:49 & 13:19 & 13:49 & 14:19 & 14:49 & 15:19 & 15:49 & 16:21 & 16:47 & 17:13 & 17:39 & 18:05 & 18:31 \\
Trenton TC & 12:52 & 13:22 & 13:52 & 14:22 & 14:52 & 15:22 & 15:52 & 16:24 & 16:50 & 17:16 & 17:42 & 18:08 & 18:34 \\
\bottomrule\end{tabular}\end{center}
\subsection*{To Trenton (read down) --- trips 27--33 of 33}
\begin{center}\footnotesize\setlength{\tabcolsep}{3pt}\begin{tabular}{@{}lrrrrrrr@{}}\toprule
Station & 27 & 28 & 29 & 30 & 31 & 32 & 33 \\\midrule
Entertainment Center & 17:53 & 18:19 & 18:45 & 19:15 & 19:45 & 20:15 & 20:45 \\
Aquarium & 17:55 & 18:21 & 18:47 & 19:17 & 19:47 & 20:17 & 20:47 \\
Cooper St/Rutgers & 17:57 & 18:23 & 18:49 & 19:19 & 19:49 & 20:19 & 20:49 \\
Walter Rand TC & 18:01 & 18:27 & 18:53 & 19:23 & 19:53 & 20:23 & 20:53 \\
36th Street & 18:06 & 18:32 & 18:58 & 19:28 & 19:58 & 20:28 & 20:58 \\
Pennsauken TC & 18:08 & 18:34 & 19:00 & 19:30 & 20:00 & 20:30 & 21:00 \\
Route 73/Pennsauken & 18:11 & 18:37 & 19:03 & 19:33 & 20:03 & 20:33 & 21:03 \\
Palmyra & 18:14 & 18:40 & 19:06 & 19:36 & 20:06 & 20:36 & 21:06 \\
Riverton & 18:16 & 18:42 & 19:08 & 19:38 & 20:08 & 20:38 & 21:08 \\
Cinnaminson & 18:18 & 18:44 & 19:10 & 19:40 & 20:10 & 20:40 & 21:10 \\
Riverside & 18:23 & 18:49 & 19:15 & 19:45 & 20:15 & 20:45 & 21:15 \\
Delanco & 18:25 & 18:51 & 19:17 & 19:47 & 20:17 & 20:47 & 21:17 \\
Beverly/Edgewater Park & 18:29 & 18:55 & 19:21 & 19:51 & 20:21 & 20:51 & 21:21 \\
Burlington South & 18:33 & 18:59 & 19:25 & 19:55 & 20:25 & 20:55 & 21:25 \\
Burlington Towne Ctr & 18:35 & 19:01 & 19:27 & 19:57 & 20:27 & 20:57 & 21:27 \\
Florence & 18:41 & 19:07 & 19:33 & 20:03 & 20:33 & 21:03 & 21:33 \\
Roebling & 18:44 & 19:10 & 19:36 & 20:06 & 20:36 & 21:06 & 21:36 \\
Bordentown & 18:49 & 19:15 & 19:41 & 20:11 & 20:41 & 21:11 & 21:41 \\
Cass Street & 18:55 & 19:21 & 19:47 & 20:17 & 20:47 & 21:17 & 21:47 \\
Hamilton Avenue & 18:57 & 19:23 & 19:49 & 20:19 & 20:49 & 21:19 & 21:49 \\
Trenton TC & 19:00 & 19:26 & 19:52 & 20:22 & 20:52 & 21:22 & 21:52 \\
\bottomrule\end{tabular}\end{center}
\subsection*{To Camden (read down) --- trips 1--13 of 33}
\begin{center}\footnotesize\setlength{\tabcolsep}{3pt}\begin{tabular}{@{}lrrrrrrrrrrrrr@{}}\toprule
Station & 1 & 2 & 3 & 4 & 5 & 6 & 7 & 8 & 9 & 10 & 11 & 12 & 13 \\\midrule
Trenton TC & 5:57 & 6:16 & 6:42 & 7:08 & 7:34 & 8:00 & 8:26 & 8:52 & 9:27 & 9:57 & 10:27 & 10:57 & 11:27 \\
Hamilton Avenue & 5:59 & 6:18 & 6:44 & 7:10 & 7:36 & 8:02 & 8:28 & 8:54 & 9:29 & 9:59 & 10:29 & 10:59 & 11:29 \\
Cass Street & 6:01 & 6:20 & 6:46 & 7:12 & 7:38 & 8:04 & 8:30 & 8:56 & 9:31 & 10:01 & 10:31 & 11:01 & 11:31 \\
Bordentown & 6:08 & 6:27 & 6:53 & 7:19 & 7:45 & 8:11 & 8:37 & 9:03 & 9:38 & 10:08 & 10:38 & 11:08 & 11:38 \\
Roebling & 6:13 & 6:32 & 6:58 & 7:24 & 7:50 & 8:16 & 8:42 & 9:08 & 9:43 & 10:13 & 10:43 & 11:13 & 11:43 \\
Florence & 6:16 & 6:35 & 7:01 & 7:27 & 7:53 & 8:19 & 8:45 & 9:11 & 9:46 & 10:16 & 10:46 & 11:16 & 11:46 \\
Burlington Towne Ctr & 6:21 & 6:40 & 7:06 & 7:32 & 7:58 & 8:24 & 8:50 & 9:16 & 9:51 & 10:21 & 10:51 & 11:21 & 11:51 \\
Burlington South & 6:23 & 6:42 & 7:08 & 7:34 & 8:00 & 8:26 & 8:52 & 9:18 & 9:53 & 10:23 & 10:53 & 11:23 & 11:53 \\
Beverly/Edgewater Park & 6:27 & 6:46 & 7:12 & 7:38 & 8:04 & 8:30 & 8:56 & 9:22 & 9:57 & 10:27 & 10:57 & 11:27 & 11:57 \\
Delanco & 6:31 & 6:50 & 7:16 & 7:42 & 8:08 & 8:34 & 9:00 & 9:26 & 10:01 & 10:31 & 11:01 & 11:31 & 12:01 \\
Riverside & 6:33 & 6:52 & 7:18 & 7:44 & 8:10 & 8:36 & 9:02 & 9:28 & 10:03 & 10:33 & 11:03 & 11:33 & 12:03 \\
Cinnaminson & 6:38 & 6:57 & 7:23 & 7:49 & 8:15 & 8:41 & 9:07 & 9:33 & 10:08 & 10:38 & 11:08 & 11:38 & 12:08 \\
Riverton & 6:40 & 6:59 & 7:25 & 7:51 & 8:17 & 8:43 & 9:09 & 9:35 & 10:10 & 10:40 & 11:10 & 11:40 & 12:10 \\
Palmyra & 6:42 & 7:01 & 7:27 & 7:53 & 8:19 & 8:45 & 9:11 & 9:37 & 10:12 & 10:42 & 11:12 & 11:42 & 12:12 \\
Route 73/Pennsauken & 6:45 & 7:04 & 7:30 & 7:56 & 8:22 & 8:48 & 9:14 & 9:40 & 10:15 & 10:45 & 11:15 & 11:45 & 12:15 \\
Pennsauken TC & 6:47 & 7:06 & 7:32 & 7:58 & 8:24 & 8:50 & 9:16 & 9:42 & 10:17 & 10:47 & 11:17 & 11:47 & 12:17 \\
36th Street & 6:51 & 7:10 & 7:36 & 8:02 & 8:28 & 8:54 & 9:20 & 9:46 & 10:21 & 10:51 & 11:21 & 11:51 & 12:21 \\
Walter Rand TC & 6:56 & 7:15 & 7:41 & 8:07 & 8:33 & 8:59 & 9:25 & 9:51 & 10:26 & 10:56 & 11:26 & 11:56 & 12:26 \\
Cooper St/Rutgers & 6:59 & 7:18 & 7:44 & 8:10 & 8:36 & 9:02 & 9:28 & 9:54 & 10:29 & 10:59 & 11:29 & 11:59 & 12:29 \\
Aquarium & 7:01 & 7:20 & 7:46 & 8:12 & 8:38 & 9:04 & 9:30 & 9:56 & 10:31 & 11:01 & 11:31 & 12:01 & 12:31 \\
Entertainment Center & 7:04 & 7:23 & 7:49 & 8:15 & 8:41 & 9:07 & 9:33 & 9:59 & 10:34 & 11:04 & 11:34 & 12:04 & 12:34 \\
\bottomrule\end{tabular}\end{center}
\subsection*{To Camden (read down) --- trips 14--26 of 33}
\begin{center}\footnotesize\setlength{\tabcolsep}{3pt}\begin{tabular}{@{}lrrrrrrrrrrrrr@{}}\toprule
Station & 14 & 15 & 16 & 17 & 18 & 19 & 20 & 21 & 22 & 23 & 24 & 25 & 26 \\\midrule
Trenton TC & 11:57 & 12:27 & 12:57 & 13:27 & 13:57 & 14:27 & 14:57 & 15:22 & 15:48 & 16:14 & 16:40 & 17:06 & 17:32 \\
Hamilton Avenue & 11:59 & 12:29 & 12:59 & 13:29 & 13:59 & 14:29 & 14:59 & 15:24 & 15:50 & 16:16 & 16:42 & 17:08 & 17:34 \\
Cass Street & 12:01 & 12:31 & 13:01 & 13:31 & 14:01 & 14:31 & 15:01 & 15:26 & 15:52 & 16:18 & 16:44 & 17:10 & 17:36 \\
Bordentown & 12:08 & 12:38 & 13:08 & 13:38 & 14:08 & 14:38 & 15:08 & 15:33 & 15:59 & 16:25 & 16:51 & 17:17 & 17:43 \\
Roebling & 12:13 & 12:43 & 13:13 & 13:43 & 14:13 & 14:43 & 15:13 & 15:38 & 16:04 & 16:30 & 16:56 & 17:22 & 17:48 \\
Florence & 12:16 & 12:46 & 13:16 & 13:46 & 14:16 & 14:46 & 15:16 & 15:41 & 16:07 & 16:33 & 16:59 & 17:25 & 17:51 \\
Burlington Towne Ctr & 12:21 & 12:51 & 13:21 & 13:51 & 14:21 & 14:51 & 15:21 & 15:46 & 16:12 & 16:38 & 17:04 & 17:30 & 17:56 \\
Burlington South & 12:23 & 12:53 & 13:23 & 13:53 & 14:23 & 14:53 & 15:23 & 15:48 & 16:14 & 16:40 & 17:06 & 17:32 & 17:58 \\
Beverly/Edgewater Park & 12:27 & 12:57 & 13:27 & 13:57 & 14:27 & 14:57 & 15:27 & 15:52 & 16:18 & 16:44 & 17:10 & 17:36 & 18:02 \\
Delanco & 12:31 & 13:01 & 13:31 & 14:01 & 14:31 & 15:01 & 15:31 & 15:56 & 16:22 & 16:48 & 17:14 & 17:40 & 18:06 \\
Riverside & 12:33 & 13:03 & 13:33 & 14:03 & 14:33 & 15:03 & 15:33 & 15:58 & 16:24 & 16:50 & 17:16 & 17:42 & 18:08 \\
Cinnaminson & 12:38 & 13:08 & 13:38 & 14:08 & 14:38 & 15:08 & 15:38 & 16:03 & 16:29 & 16:55 & 17:21 & 17:47 & 18:13 \\
Riverton & 12:40 & 13:10 & 13:40 & 14:10 & 14:40 & 15:10 & 15:40 & 16:05 & 16:31 & 16:57 & 17:23 & 17:49 & 18:15 \\
Palmyra & 12:42 & 13:12 & 13:42 & 14:12 & 14:42 & 15:12 & 15:42 & 16:07 & 16:33 & 16:59 & 17:25 & 17:51 & 18:17 \\
Route 73/Pennsauken & 12:45 & 13:15 & 13:45 & 14:15 & 14:45 & 15:15 & 15:45 & 16:10 & 16:36 & 17:02 & 17:28 & 17:54 & 18:20 \\
Pennsauken TC & 12:47 & 13:17 & 13:47 & 14:17 & 14:47 & 15:17 & 15:47 & 16:12 & 16:38 & 17:04 & 17:30 & 17:56 & 18:22 \\
36th Street & 12:51 & 13:21 & 13:51 & 14:21 & 14:51 & 15:21 & 15:51 & 16:16 & 16:42 & 17:08 & 17:34 & 18:00 & 18:26 \\
Walter Rand TC & 12:56 & 13:26 & 13:56 & 14:26 & 14:56 & 15:26 & 15:56 & 16:21 & 16:47 & 17:13 & 17:39 & 18:05 & 18:31 \\
Cooper St/Rutgers & 12:59 & 13:29 & 13:59 & 14:29 & 14:59 & 15:29 & 15:59 & 16:24 & 16:50 & 17:16 & 17:42 & 18:08 & 18:34 \\
Aquarium & 13:01 & 13:31 & 14:01 & 14:31 & 15:01 & 15:31 & 16:01 & 16:26 & 16:52 & 17:18 & 17:44 & 18:10 & 18:36 \\
Entertainment Center & 13:04 & 13:34 & 14:04 & 14:34 & 15:04 & 15:34 & 16:04 & 16:29 & 16:55 & 17:21 & 17:47 & 18:13 & 18:39 \\
\bottomrule\end{tabular}\end{center}
\subsection*{To Camden (read down) --- trips 27--33 of 33}
\begin{center}\footnotesize\setlength{\tabcolsep}{3pt}\begin{tabular}{@{}lrrrrrrr@{}}\toprule
Station & 27 & 28 & 29 & 30 & 31 & 32 & 33 \\\midrule
Trenton TC & 17:58 & 18:24 & 18:50 & 19:27 & 19:57 & 20:27 & 20:57 \\
Hamilton Avenue & 18:00 & 18:26 & 18:52 & 19:29 & 19:59 & 20:29 & 20:59 \\
Cass Street & 18:02 & 18:28 & 18:54 & 19:31 & 20:01 & 20:31 & 21:01 \\
Bordentown & 18:09 & 18:35 & 19:01 & 19:38 & 20:08 & 20:38 & 21:08 \\
Roebling & 18:14 & 18:40 & 19:06 & 19:43 & 20:13 & 20:43 & 21:13 \\
Florence & 18:17 & 18:43 & 19:09 & 19:46 & 20:16 & 20:46 & 21:16 \\
Burlington Towne Ctr & 18:22 & 18:48 & 19:14 & 19:51 & 20:21 & 20:51 & 21:21 \\
Burlington South & 18:24 & 18:50 & 19:16 & 19:53 & 20:23 & 20:53 & 21:23 \\
Beverly/Edgewater Park & 18:28 & 18:54 & 19:20 & 19:57 & 20:27 & 20:57 & 21:27 \\
Delanco & 18:32 & 18:58 & 19:24 & 20:01 & 20:31 & 21:01 & 21:31 \\
Riverside & 18:34 & 19:00 & 19:26 & 20:03 & 20:33 & 21:03 & 21:33 \\
Cinnaminson & 18:39 & 19:05 & 19:31 & 20:08 & 20:38 & 21:08 & 21:38 \\
Riverton & 18:41 & 19:07 & 19:33 & 20:10 & 20:40 & 21:10 & 21:40 \\
Palmyra & 18:43 & 19:09 & 19:35 & 20:12 & 20:42 & 21:12 & 21:42 \\
Route 73/Pennsauken & 18:46 & 19:12 & 19:38 & 20:15 & 20:45 & 21:15 & 21:45 \\
Pennsauken TC & 18:48 & 19:14 & 19:40 & 20:17 & 20:47 & 21:17 & 21:47 \\
36th Street & 18:52 & 19:18 & 19:44 & 20:21 & 20:51 & 21:21 & 21:51 \\
Walter Rand TC & 18:57 & 19:23 & 19:49 & 20:26 & 20:56 & 21:26 & 21:56 \\
Cooper St/Rutgers & 19:00 & 19:26 & 19:52 & 20:29 & 20:59 & 21:29 & 21:59 \\
Aquarium & 19:02 & 19:28 & 19:54 & 20:31 & 21:01 & 21:31 & 22:01 \\
Entertainment Center & 19:05 & 19:31 & 19:57 & 20:34 & 21:04 & 21:34 & 22:04 \\
\bottomrule\end{tabular}\end{center}
\clearpage
\section{Uniform U5: every 29 min peak, 30 off-peak}\label{tt:uniform_U5}
\subsection*{To Trenton (read down) --- trips 1--13 of 31}
\begin{center}\footnotesize\setlength{\tabcolsep}{3pt}\begin{tabular}{@{}lrrrrrrrrrrrrr@{}}\toprule
Station & 1 & 2 & 3 & 4 & 5 & 6 & 7 & 8 & 9 & 10 & 11 & 12 & 13 \\\midrule
Entertainment Center & 5:45 & 6:14 & 6:43 & 7:12 & 7:41 & 8:10 & 8:39 & 9:15 & 9:45 & 10:15 & 10:45 & 11:15 & 11:45 \\
Aquarium & 5:47 & 6:16 & 6:45 & 7:14 & 7:43 & 8:12 & 8:41 & 9:17 & 9:47 & 10:17 & 10:47 & 11:17 & 11:47 \\
Cooper St/Rutgers & 5:49 & 6:18 & 6:47 & 7:16 & 7:45 & 8:14 & 8:43 & 9:19 & 9:49 & 10:19 & 10:49 & 11:19 & 11:49 \\
Walter Rand TC & 5:53 & 6:22 & 6:51 & 7:20 & 7:49 & 8:18 & 8:47 & 9:23 & 9:53 & 10:23 & 10:53 & 11:23 & 11:53 \\
36th Street & 5:58 & 6:27 & 6:56 & 7:25 & 7:54 & 8:23 & 8:52 & 9:28 & 9:58 & 10:28 & 10:58 & 11:28 & 11:58 \\
Pennsauken TC & 6:00 & 6:29 & 6:58 & 7:27 & 7:56 & 8:25 & 8:54 & 9:30 & 10:00 & 10:30 & 11:00 & 11:30 & 12:00 \\
Route 73/Pennsauken & 6:03 & 6:32 & 7:01 & 7:30 & 7:59 & 8:28 & 8:57 & 9:33 & 10:03 & 10:33 & 11:03 & 11:33 & 12:03 \\
Palmyra & 6:06 & 6:35 & 7:04 & 7:33 & 8:02 & 8:31 & 9:00 & 9:36 & 10:06 & 10:36 & 11:06 & 11:36 & 12:06 \\
Riverton & 6:08 & 6:37 & 7:06 & 7:35 & 8:04 & 8:33 & 9:02 & 9:38 & 10:08 & 10:38 & 11:08 & 11:38 & 12:08 \\
Cinnaminson & 6:10 & 6:39 & 7:08 & 7:37 & 8:06 & 8:35 & 9:04 & 9:40 & 10:10 & 10:40 & 11:10 & 11:40 & 12:10 \\
Riverside & 6:15 & 6:44 & 7:13 & 7:42 & 8:11 & 8:40 & 9:09 & 9:45 & 10:15 & 10:45 & 11:15 & 11:45 & 12:15 \\
Delanco & 6:17 & 6:46 & 7:15 & 7:44 & 8:13 & 8:42 & 9:11 & 9:47 & 10:17 & 10:47 & 11:17 & 11:47 & 12:17 \\
Beverly/Edgewater Park & 6:21 & 6:50 & 7:19 & 7:48 & 8:17 & 8:46 & 9:15 & 9:51 & 10:21 & 10:51 & 11:21 & 11:51 & 12:21 \\
Burlington South & 6:25 & 6:54 & 7:23 & 7:52 & 8:21 & 8:50 & 9:19 & 9:55 & 10:25 & 10:55 & 11:25 & 11:55 & 12:25 \\
Burlington Towne Ctr & 6:27 & 6:56 & 7:25 & 7:54 & 8:23 & 8:52 & 9:21 & 9:57 & 10:27 & 10:57 & 11:27 & 11:57 & 12:27 \\
Florence & 6:33 & 7:02 & 7:31 & 8:00 & 8:29 & 8:58 & 9:27 & 10:03 & 10:33 & 11:03 & 11:33 & 12:03 & 12:33 \\
Roebling & 6:36 & 7:05 & 7:34 & 8:03 & 8:32 & 9:01 & 9:30 & 10:06 & 10:36 & 11:06 & 11:36 & 12:06 & 12:36 \\
Bordentown & 6:41 & 7:10 & 7:39 & 8:08 & 8:37 & 9:06 & 9:35 & 10:11 & 10:41 & 11:11 & 11:41 & 12:11 & 12:41 \\
Cass Street & 6:47 & 7:16 & 7:45 & 8:14 & 8:43 & 9:12 & 9:41 & 10:17 & 10:47 & 11:17 & 11:47 & 12:17 & 12:47 \\
Hamilton Avenue & 6:49 & 7:18 & 7:47 & 8:16 & 8:45 & 9:14 & 9:43 & 10:19 & 10:49 & 11:19 & 11:49 & 12:19 & 12:49 \\
Trenton TC & 6:52 & 7:21 & 7:50 & 8:19 & 8:48 & 9:17 & 9:46 & 10:22 & 10:52 & 11:22 & 11:52 & 12:22 & 12:52 \\
\bottomrule\end{tabular}\end{center}
\subsection*{To Trenton (read down) --- trips 14--26 of 31}
\begin{center}\footnotesize\setlength{\tabcolsep}{3pt}\begin{tabular}{@{}lrrrrrrrrrrrrr@{}}\toprule
Station & 14 & 15 & 16 & 17 & 18 & 19 & 20 & 21 & 22 & 23 & 24 & 25 & 26 \\\midrule
Entertainment Center & 12:15 & 12:45 & 13:15 & 13:45 & 14:15 & 14:45 & 15:25 & 15:54 & 16:23 & 16:52 & 17:21 & 17:50 & 18:19 \\
Aquarium & 12:17 & 12:47 & 13:17 & 13:47 & 14:17 & 14:47 & 15:27 & 15:56 & 16:25 & 16:54 & 17:23 & 17:52 & 18:21 \\
Cooper St/Rutgers & 12:19 & 12:49 & 13:19 & 13:49 & 14:19 & 14:49 & 15:29 & 15:58 & 16:27 & 16:56 & 17:25 & 17:54 & 18:23 \\
Walter Rand TC & 12:23 & 12:53 & 13:23 & 13:53 & 14:23 & 14:53 & 15:33 & 16:02 & 16:31 & 17:00 & 17:29 & 17:58 & 18:27 \\
36th Street & 12:28 & 12:58 & 13:28 & 13:58 & 14:28 & 14:58 & 15:38 & 16:07 & 16:36 & 17:05 & 17:34 & 18:03 & 18:32 \\
Pennsauken TC & 12:30 & 13:00 & 13:30 & 14:00 & 14:30 & 15:00 & 15:40 & 16:09 & 16:38 & 17:07 & 17:36 & 18:05 & 18:34 \\
Route 73/Pennsauken & 12:33 & 13:03 & 13:33 & 14:03 & 14:33 & 15:03 & 15:43 & 16:12 & 16:41 & 17:10 & 17:39 & 18:08 & 18:37 \\
Palmyra & 12:36 & 13:06 & 13:36 & 14:06 & 14:36 & 15:06 & 15:46 & 16:15 & 16:44 & 17:13 & 17:42 & 18:11 & 18:40 \\
Riverton & 12:38 & 13:08 & 13:38 & 14:08 & 14:38 & 15:08 & 15:48 & 16:17 & 16:46 & 17:15 & 17:44 & 18:13 & 18:42 \\
Cinnaminson & 12:40 & 13:10 & 13:40 & 14:10 & 14:40 & 15:10 & 15:50 & 16:19 & 16:48 & 17:17 & 17:46 & 18:15 & 18:44 \\
Riverside & 12:45 & 13:15 & 13:45 & 14:15 & 14:45 & 15:15 & 15:55 & 16:24 & 16:53 & 17:22 & 17:51 & 18:20 & 18:49 \\
Delanco & 12:47 & 13:17 & 13:47 & 14:17 & 14:47 & 15:17 & 15:57 & 16:26 & 16:55 & 17:24 & 17:53 & 18:22 & 18:51 \\
Beverly/Edgewater Park & 12:51 & 13:21 & 13:51 & 14:21 & 14:51 & 15:21 & 16:01 & 16:30 & 16:59 & 17:28 & 17:57 & 18:26 & 18:55 \\
Burlington South & 12:55 & 13:25 & 13:55 & 14:25 & 14:55 & 15:25 & 16:05 & 16:34 & 17:03 & 17:32 & 18:01 & 18:30 & 18:59 \\
Burlington Towne Ctr & 12:57 & 13:27 & 13:57 & 14:27 & 14:57 & 15:27 & 16:07 & 16:36 & 17:05 & 17:34 & 18:03 & 18:32 & 19:01 \\
Florence & 13:03 & 13:33 & 14:03 & 14:33 & 15:03 & 15:33 & 16:13 & 16:42 & 17:11 & 17:40 & 18:09 & 18:38 & 19:07 \\
Roebling & 13:06 & 13:36 & 14:06 & 14:36 & 15:06 & 15:36 & 16:16 & 16:45 & 17:14 & 17:43 & 18:12 & 18:41 & 19:10 \\
Bordentown & 13:11 & 13:41 & 14:11 & 14:41 & 15:11 & 15:41 & 16:21 & 16:50 & 17:19 & 17:48 & 18:17 & 18:46 & 19:15 \\
Cass Street & 13:17 & 13:47 & 14:17 & 14:47 & 15:17 & 15:47 & 16:27 & 16:56 & 17:25 & 17:54 & 18:23 & 18:52 & 19:21 \\
Hamilton Avenue & 13:19 & 13:49 & 14:19 & 14:49 & 15:19 & 15:49 & 16:29 & 16:58 & 17:27 & 17:56 & 18:25 & 18:54 & 19:23 \\
Trenton TC & 13:22 & 13:52 & 14:22 & 14:52 & 15:22 & 15:52 & 16:32 & 17:01 & 17:30 & 17:59 & 18:28 & 18:57 & 19:26 \\
\bottomrule\end{tabular}\end{center}
\subsection*{To Trenton (read down) --- trips 27--31 of 31}
\begin{center}\footnotesize\setlength{\tabcolsep}{3pt}\begin{tabular}{@{}lrrrrr@{}}\toprule
Station & 27 & 28 & 29 & 30 & 31 \\\midrule
Entertainment Center & 18:48 & 19:15 & 19:45 & 20:15 & 20:45 \\
Aquarium & 18:50 & 19:17 & 19:47 & 20:17 & 20:47 \\
Cooper St/Rutgers & 18:52 & 19:19 & 19:49 & 20:19 & 20:49 \\
Walter Rand TC & 18:56 & 19:23 & 19:53 & 20:23 & 20:53 \\
36th Street & 19:01 & 19:28 & 19:58 & 20:28 & 20:58 \\
Pennsauken TC & 19:03 & 19:30 & 20:00 & 20:30 & 21:00 \\
Route 73/Pennsauken & 19:06 & 19:33 & 20:03 & 20:33 & 21:03 \\
Palmyra & 19:09 & 19:36 & 20:06 & 20:36 & 21:06 \\
Riverton & 19:11 & 19:38 & 20:08 & 20:38 & 21:08 \\
Cinnaminson & 19:13 & 19:40 & 20:10 & 20:40 & 21:10 \\
Riverside & 19:18 & 19:45 & 20:15 & 20:45 & 21:15 \\
Delanco & 19:20 & 19:47 & 20:17 & 20:47 & 21:17 \\
Beverly/Edgewater Park & 19:24 & 19:51 & 20:21 & 20:51 & 21:21 \\
Burlington South & 19:28 & 19:55 & 20:25 & 20:55 & 21:25 \\
Burlington Towne Ctr & 19:30 & 19:57 & 20:27 & 20:57 & 21:27 \\
Florence & 19:36 & 20:03 & 20:33 & 21:03 & 21:33 \\
Roebling & 19:39 & 20:06 & 20:36 & 21:06 & 21:36 \\
Bordentown & 19:44 & 20:11 & 20:41 & 21:11 & 21:41 \\
Cass Street & 19:50 & 20:17 & 20:47 & 21:17 & 21:47 \\
Hamilton Avenue & 19:52 & 20:19 & 20:49 & 21:19 & 21:49 \\
Trenton TC & 19:55 & 20:22 & 20:52 & 21:22 & 21:52 \\
\bottomrule\end{tabular}\end{center}
\subsection*{To Camden (read down) --- trips 1--13 of 32}
\begin{center}\footnotesize\setlength{\tabcolsep}{3pt}\begin{tabular}{@{}lrrrrrrrrrrrrr@{}}\toprule
Station & 1 & 2 & 3 & 4 & 5 & 6 & 7 & 8 & 9 & 10 & 11 & 12 & 13 \\\midrule
Trenton TC & 5:57 & 6:00 & 6:29 & 6:58 & 7:27 & 7:56 & 8:25 & 8:54 & 9:27 & 9:57 & 10:27 & 10:57 & 11:27 \\
Hamilton Avenue & 5:59 & 6:02 & 6:31 & 7:00 & 7:29 & 7:58 & 8:27 & 8:56 & 9:29 & 9:59 & 10:29 & 10:59 & 11:29 \\
Cass Street & 6:01 & 6:04 & 6:33 & 7:02 & 7:31 & 8:00 & 8:29 & 8:58 & 9:31 & 10:01 & 10:31 & 11:01 & 11:31 \\
Bordentown & 6:08 & 6:11 & 6:40 & 7:09 & 7:38 & 8:07 & 8:36 & 9:05 & 9:38 & 10:08 & 10:38 & 11:08 & 11:38 \\
Roebling & 6:13 & 6:16 & 6:45 & 7:14 & 7:43 & 8:12 & 8:41 & 9:10 & 9:43 & 10:13 & 10:43 & 11:13 & 11:43 \\
Florence & 6:16 & 6:19 & 6:48 & 7:17 & 7:46 & 8:15 & 8:44 & 9:13 & 9:46 & 10:16 & 10:46 & 11:16 & 11:46 \\
Burlington Towne Ctr & 6:21 & 6:24 & 6:53 & 7:22 & 7:51 & 8:20 & 8:49 & 9:18 & 9:51 & 10:21 & 10:51 & 11:21 & 11:51 \\
Burlington South & 6:23 & 6:26 & 6:55 & 7:24 & 7:53 & 8:22 & 8:51 & 9:20 & 9:53 & 10:23 & 10:53 & 11:23 & 11:53 \\
Beverly/Edgewater Park & 6:27 & 6:30 & 6:59 & 7:28 & 7:57 & 8:26 & 8:55 & 9:24 & 9:57 & 10:27 & 10:57 & 11:27 & 11:57 \\
Delanco & 6:31 & 6:34 & 7:03 & 7:32 & 8:01 & 8:30 & 8:59 & 9:28 & 10:01 & 10:31 & 11:01 & 11:31 & 12:01 \\
Riverside & 6:33 & 6:36 & 7:05 & 7:34 & 8:03 & 8:32 & 9:01 & 9:30 & 10:03 & 10:33 & 11:03 & 11:33 & 12:03 \\
Cinnaminson & 6:38 & 6:41 & 7:10 & 7:39 & 8:08 & 8:37 & 9:06 & 9:35 & 10:08 & 10:38 & 11:08 & 11:38 & 12:08 \\
Riverton & 6:40 & 6:43 & 7:12 & 7:41 & 8:10 & 8:39 & 9:08 & 9:37 & 10:10 & 10:40 & 11:10 & 11:40 & 12:10 \\
Palmyra & 6:42 & 6:45 & 7:14 & 7:43 & 8:12 & 8:41 & 9:10 & 9:39 & 10:12 & 10:42 & 11:12 & 11:42 & 12:12 \\
Route 73/Pennsauken & 6:45 & 6:48 & 7:17 & 7:46 & 8:15 & 8:44 & 9:13 & 9:42 & 10:15 & 10:45 & 11:15 & 11:45 & 12:15 \\
Pennsauken TC & 6:47 & 6:50 & 7:19 & 7:48 & 8:17 & 8:46 & 9:15 & 9:44 & 10:17 & 10:47 & 11:17 & 11:47 & 12:17 \\
36th Street & 6:51 & 6:54 & 7:23 & 7:52 & 8:21 & 8:50 & 9:19 & 9:48 & 10:21 & 10:51 & 11:21 & 11:51 & 12:21 \\
Walter Rand TC & 6:56 & 6:59 & 7:28 & 7:57 & 8:26 & 8:55 & 9:24 & 9:53 & 10:26 & 10:56 & 11:26 & 11:56 & 12:26 \\
Cooper St/Rutgers & 6:59 & 7:02 & 7:31 & 8:00 & 8:29 & 8:58 & 9:27 & 9:56 & 10:29 & 10:59 & 11:29 & 11:59 & 12:29 \\
Aquarium & 7:01 & 7:04 & 7:33 & 8:02 & 8:31 & 9:00 & 9:29 & 9:58 & 10:31 & 11:01 & 11:31 & 12:01 & 12:31 \\
Entertainment Center & 7:04 & 7:07 & 7:36 & 8:05 & 8:34 & 9:03 & 9:32 & 10:01 & 10:34 & 11:04 & 11:34 & 12:04 & 12:34 \\
\bottomrule\end{tabular}\end{center}
\subsection*{To Camden (read down) --- trips 14--26 of 32}
\begin{center}\footnotesize\setlength{\tabcolsep}{3pt}\begin{tabular}{@{}lrrrrrrrrrrrrr@{}}\toprule
Station & 14 & 15 & 16 & 17 & 18 & 19 & 20 & 21 & 22 & 23 & 24 & 25 & 26 \\\midrule
Trenton TC & 11:57 & 12:27 & 12:57 & 13:27 & 13:57 & 14:27 & 14:57 & 15:11 & 15:40 & 16:09 & 16:38 & 17:07 & 17:36 \\
Hamilton Avenue & 11:59 & 12:29 & 12:59 & 13:29 & 13:59 & 14:29 & 14:59 & 15:13 & 15:42 & 16:11 & 16:40 & 17:09 & 17:38 \\
Cass Street & 12:01 & 12:31 & 13:01 & 13:31 & 14:01 & 14:31 & 15:01 & 15:15 & 15:44 & 16:13 & 16:42 & 17:11 & 17:40 \\
Bordentown & 12:08 & 12:38 & 13:08 & 13:38 & 14:08 & 14:38 & 15:08 & 15:22 & 15:51 & 16:20 & 16:49 & 17:18 & 17:47 \\
Roebling & 12:13 & 12:43 & 13:13 & 13:43 & 14:13 & 14:43 & 15:13 & 15:27 & 15:56 & 16:25 & 16:54 & 17:23 & 17:52 \\
Florence & 12:16 & 12:46 & 13:16 & 13:46 & 14:16 & 14:46 & 15:16 & 15:30 & 15:59 & 16:28 & 16:57 & 17:26 & 17:55 \\
Burlington Towne Ctr & 12:21 & 12:51 & 13:21 & 13:51 & 14:21 & 14:51 & 15:21 & 15:35 & 16:04 & 16:33 & 17:02 & 17:31 & 18:00 \\
Burlington South & 12:23 & 12:53 & 13:23 & 13:53 & 14:23 & 14:53 & 15:23 & 15:37 & 16:06 & 16:35 & 17:04 & 17:33 & 18:02 \\
Beverly/Edgewater Park & 12:27 & 12:57 & 13:27 & 13:57 & 14:27 & 14:57 & 15:27 & 15:41 & 16:10 & 16:39 & 17:08 & 17:37 & 18:06 \\
Delanco & 12:31 & 13:01 & 13:31 & 14:01 & 14:31 & 15:01 & 15:31 & 15:45 & 16:14 & 16:43 & 17:12 & 17:41 & 18:10 \\
Riverside & 12:33 & 13:03 & 13:33 & 14:03 & 14:33 & 15:03 & 15:33 & 15:47 & 16:16 & 16:45 & 17:14 & 17:43 & 18:12 \\
Cinnaminson & 12:38 & 13:08 & 13:38 & 14:08 & 14:38 & 15:08 & 15:38 & 15:52 & 16:21 & 16:50 & 17:19 & 17:48 & 18:17 \\
Riverton & 12:40 & 13:10 & 13:40 & 14:10 & 14:40 & 15:10 & 15:40 & 15:54 & 16:23 & 16:52 & 17:21 & 17:50 & 18:19 \\
Palmyra & 12:42 & 13:12 & 13:42 & 14:12 & 14:42 & 15:12 & 15:42 & 15:56 & 16:25 & 16:54 & 17:23 & 17:52 & 18:21 \\
Route 73/Pennsauken & 12:45 & 13:15 & 13:45 & 14:15 & 14:45 & 15:15 & 15:45 & 15:59 & 16:28 & 16:57 & 17:26 & 17:55 & 18:24 \\
Pennsauken TC & 12:47 & 13:17 & 13:47 & 14:17 & 14:47 & 15:17 & 15:47 & 16:01 & 16:30 & 16:59 & 17:28 & 17:57 & 18:26 \\
36th Street & 12:51 & 13:21 & 13:51 & 14:21 & 14:51 & 15:21 & 15:51 & 16:05 & 16:34 & 17:03 & 17:32 & 18:01 & 18:30 \\
Walter Rand TC & 12:56 & 13:26 & 13:56 & 14:26 & 14:56 & 15:26 & 15:56 & 16:10 & 16:39 & 17:08 & 17:37 & 18:06 & 18:35 \\
Cooper St/Rutgers & 12:59 & 13:29 & 13:59 & 14:29 & 14:59 & 15:29 & 15:59 & 16:13 & 16:42 & 17:11 & 17:40 & 18:09 & 18:38 \\
Aquarium & 13:01 & 13:31 & 14:01 & 14:31 & 15:01 & 15:31 & 16:01 & 16:15 & 16:44 & 17:13 & 17:42 & 18:11 & 18:40 \\
Entertainment Center & 13:04 & 13:34 & 14:04 & 14:34 & 15:04 & 15:34 & 16:04 & 16:18 & 16:47 & 17:16 & 17:45 & 18:14 & 18:43 \\
\bottomrule\end{tabular}\end{center}
\subsection*{To Camden (read down) --- trips 27--32 of 32}
\begin{center}\footnotesize\setlength{\tabcolsep}{3pt}\begin{tabular}{@{}lrrrrrr@{}}\toprule
Station & 27 & 28 & 29 & 30 & 31 & 32 \\\midrule
Trenton TC & 18:05 & 18:34 & 19:27 & 19:57 & 20:27 & 20:57 \\
Hamilton Avenue & 18:07 & 18:36 & 19:29 & 19:59 & 20:29 & 20:59 \\
Cass Street & 18:09 & 18:38 & 19:31 & 20:01 & 20:31 & 21:01 \\
Bordentown & 18:16 & 18:45 & 19:38 & 20:08 & 20:38 & 21:08 \\
Roebling & 18:21 & 18:50 & 19:43 & 20:13 & 20:43 & 21:13 \\
Florence & 18:24 & 18:53 & 19:46 & 20:16 & 20:46 & 21:16 \\
Burlington Towne Ctr & 18:29 & 18:58 & 19:51 & 20:21 & 20:51 & 21:21 \\
Burlington South & 18:31 & 19:00 & 19:53 & 20:23 & 20:53 & 21:23 \\
Beverly/Edgewater Park & 18:35 & 19:04 & 19:57 & 20:27 & 20:57 & 21:27 \\
Delanco & 18:39 & 19:08 & 20:01 & 20:31 & 21:01 & 21:31 \\
Riverside & 18:41 & 19:10 & 20:03 & 20:33 & 21:03 & 21:33 \\
Cinnaminson & 18:46 & 19:15 & 20:08 & 20:38 & 21:08 & 21:38 \\
Riverton & 18:48 & 19:17 & 20:10 & 20:40 & 21:10 & 21:40 \\
Palmyra & 18:50 & 19:19 & 20:12 & 20:42 & 21:12 & 21:42 \\
Route 73/Pennsauken & 18:53 & 19:22 & 20:15 & 20:45 & 21:15 & 21:45 \\
Pennsauken TC & 18:55 & 19:24 & 20:17 & 20:47 & 21:17 & 21:47 \\
36th Street & 18:59 & 19:28 & 20:21 & 20:51 & 21:21 & 21:51 \\
Walter Rand TC & 19:04 & 19:33 & 20:26 & 20:56 & 21:26 & 21:56 \\
Cooper St/Rutgers & 19:07 & 19:36 & 20:29 & 20:59 & 21:29 & 21:59 \\
Aquarium & 19:09 & 19:38 & 20:31 & 21:01 & 21:31 & 22:01 \\
Entertainment Center & 19:12 & 19:41 & 20:34 & 21:04 & 21:34 & 22:04 \\
\bottomrule\end{tabular}\end{center}
\clearpage
\section{Uniform U4: every 38 min peak, 38 off-peak}\label{tt:uniform_U4}
\subsection*{To Trenton (read down) --- trips 1--13 of 25}
\begin{center}\footnotesize\setlength{\tabcolsep}{3pt}\begin{tabular}{@{}lrrrrrrrrrrrrr@{}}\toprule
Station & 1 & 2 & 3 & 4 & 5 & 6 & 7 & 8 & 9 & 10 & 11 & 12 & 13 \\\midrule
Entertainment Center & 5:45 & 6:23 & 7:01 & 7:39 & 8:17 & 8:55 & 9:33 & 10:11 & 10:49 & 11:27 & 12:05 & 12:43 & 13:21 \\
Aquarium & 5:47 & 6:25 & 7:03 & 7:41 & 8:19 & 8:57 & 9:35 & 10:13 & 10:51 & 11:29 & 12:07 & 12:45 & 13:23 \\
Cooper St/Rutgers & 5:49 & 6:27 & 7:05 & 7:43 & 8:21 & 8:59 & 9:37 & 10:15 & 10:53 & 11:31 & 12:09 & 12:47 & 13:25 \\
Walter Rand TC & 5:53 & 6:31 & 7:09 & 7:47 & 8:25 & 9:03 & 9:41 & 10:19 & 10:57 & 11:35 & 12:13 & 12:51 & 13:29 \\
36th Street & 5:58 & 6:36 & 7:14 & 7:52 & 8:30 & 9:08 & 9:46 & 10:24 & 11:02 & 11:40 & 12:18 & 12:56 & 13:34 \\
Pennsauken TC & 6:00 & 6:38 & 7:16 & 7:54 & 8:32 & 9:10 & 9:48 & 10:26 & 11:04 & 11:42 & 12:20 & 12:58 & 13:36 \\
Route 73/Pennsauken & 6:03 & 6:41 & 7:19 & 7:57 & 8:35 & 9:13 & 9:51 & 10:29 & 11:07 & 11:45 & 12:23 & 13:01 & 13:39 \\
Palmyra & 6:06 & 6:44 & 7:22 & 8:00 & 8:38 & 9:16 & 9:54 & 10:32 & 11:10 & 11:48 & 12:26 & 13:04 & 13:42 \\
Riverton & 6:08 & 6:46 & 7:24 & 8:02 & 8:40 & 9:18 & 9:56 & 10:34 & 11:12 & 11:50 & 12:28 & 13:06 & 13:44 \\
Cinnaminson & 6:10 & 6:48 & 7:26 & 8:04 & 8:42 & 9:20 & 9:58 & 10:36 & 11:14 & 11:52 & 12:30 & 13:08 & 13:46 \\
Riverside & 6:15 & 6:53 & 7:31 & 8:09 & 8:47 & 9:25 & 10:03 & 10:41 & 11:19 & 11:57 & 12:35 & 13:13 & 13:51 \\
Delanco & 6:17 & 6:55 & 7:33 & 8:11 & 8:49 & 9:27 & 10:05 & 10:43 & 11:21 & 11:59 & 12:37 & 13:15 & 13:53 \\
Beverly/Edgewater Park & 6:21 & 6:59 & 7:37 & 8:15 & 8:53 & 9:31 & 10:09 & 10:47 & 11:25 & 12:03 & 12:41 & 13:19 & 13:57 \\
Burlington South & 6:25 & 7:03 & 7:41 & 8:19 & 8:57 & 9:35 & 10:13 & 10:51 & 11:29 & 12:07 & 12:45 & 13:23 & 14:01 \\
Burlington Towne Ctr & 6:27 & 7:05 & 7:43 & 8:21 & 8:59 & 9:37 & 10:15 & 10:53 & 11:31 & 12:09 & 12:47 & 13:25 & 14:03 \\
Florence & 6:33 & 7:11 & 7:49 & 8:27 & 9:05 & 9:43 & 10:21 & 10:59 & 11:37 & 12:15 & 12:53 & 13:31 & 14:09 \\
Roebling & 6:36 & 7:14 & 7:52 & 8:30 & 9:08 & 9:46 & 10:24 & 11:02 & 11:40 & 12:18 & 12:56 & 13:34 & 14:12 \\
Bordentown & 6:41 & 7:19 & 7:57 & 8:35 & 9:13 & 9:51 & 10:29 & 11:07 & 11:45 & 12:23 & 13:01 & 13:39 & 14:17 \\
Cass Street & 6:47 & 7:25 & 8:03 & 8:41 & 9:19 & 9:57 & 10:35 & 11:13 & 11:51 & 12:29 & 13:07 & 13:45 & 14:23 \\
Hamilton Avenue & 6:49 & 7:27 & 8:05 & 8:43 & 9:21 & 9:59 & 10:37 & 11:15 & 11:53 & 12:31 & 13:09 & 13:47 & 14:25 \\
Trenton TC & 6:52 & 7:30 & 8:08 & 8:46 & 9:24 & 10:02 & 10:40 & 11:18 & 11:56 & 12:34 & 13:12 & 13:50 & 14:28 \\
\bottomrule\end{tabular}\end{center}
\subsection*{To Trenton (read down) --- trips 14--25 of 25}
\begin{center}\footnotesize\setlength{\tabcolsep}{3pt}\begin{tabular}{@{}lrrrrrrrrrrrr@{}}\toprule
Station & 14 & 15 & 16 & 17 & 18 & 19 & 20 & 21 & 22 & 23 & 24 & 25 \\\midrule
Entertainment Center & 13:59 & 14:37 & 15:15 & 15:53 & 16:31 & 17:09 & 17:47 & 18:25 & 19:03 & 19:41 & 20:19 & 20:57 \\
Aquarium & 14:01 & 14:39 & 15:17 & 15:55 & 16:33 & 17:11 & 17:49 & 18:27 & 19:05 & 19:43 & 20:21 & 20:59 \\
Cooper St/Rutgers & 14:03 & 14:41 & 15:19 & 15:57 & 16:35 & 17:13 & 17:51 & 18:29 & 19:07 & 19:45 & 20:23 & 21:01 \\
Walter Rand TC & 14:07 & 14:45 & 15:23 & 16:01 & 16:39 & 17:17 & 17:55 & 18:33 & 19:11 & 19:49 & 20:27 & 21:05 \\
36th Street & 14:12 & 14:50 & 15:28 & 16:06 & 16:44 & 17:22 & 18:00 & 18:38 & 19:16 & 19:54 & 20:32 & 21:10 \\
Pennsauken TC & 14:14 & 14:52 & 15:30 & 16:08 & 16:46 & 17:24 & 18:02 & 18:40 & 19:18 & 19:56 & 20:34 & 21:12 \\
Route 73/Pennsauken & 14:17 & 14:55 & 15:33 & 16:11 & 16:49 & 17:27 & 18:05 & 18:43 & 19:21 & 19:59 & 20:37 & 21:15 \\
Palmyra & 14:20 & 14:58 & 15:36 & 16:14 & 16:52 & 17:30 & 18:08 & 18:46 & 19:24 & 20:02 & 20:40 & 21:18 \\
Riverton & 14:22 & 15:00 & 15:38 & 16:16 & 16:54 & 17:32 & 18:10 & 18:48 & 19:26 & 20:04 & 20:42 & 21:20 \\
Cinnaminson & 14:24 & 15:02 & 15:40 & 16:18 & 16:56 & 17:34 & 18:12 & 18:50 & 19:28 & 20:06 & 20:44 & 21:22 \\
Riverside & 14:29 & 15:07 & 15:45 & 16:23 & 17:01 & 17:39 & 18:17 & 18:55 & 19:33 & 20:11 & 20:49 & 21:27 \\
Delanco & 14:31 & 15:09 & 15:47 & 16:25 & 17:03 & 17:41 & 18:19 & 18:57 & 19:35 & 20:13 & 20:51 & 21:29 \\
Beverly/Edgewater Park & 14:35 & 15:13 & 15:51 & 16:29 & 17:07 & 17:45 & 18:23 & 19:01 & 19:39 & 20:17 & 20:55 & 21:33 \\
Burlington South & 14:39 & 15:17 & 15:55 & 16:33 & 17:11 & 17:49 & 18:27 & 19:05 & 19:43 & 20:21 & 20:59 & 21:37 \\
Burlington Towne Ctr & 14:41 & 15:19 & 15:57 & 16:35 & 17:13 & 17:51 & 18:29 & 19:07 & 19:45 & 20:23 & 21:01 & 21:39 \\
Florence & 14:47 & 15:25 & 16:03 & 16:41 & 17:19 & 17:57 & 18:35 & 19:13 & 19:51 & 20:29 & 21:07 & 21:45 \\
Roebling & 14:50 & 15:28 & 16:06 & 16:44 & 17:22 & 18:00 & 18:38 & 19:16 & 19:54 & 20:32 & 21:10 & 21:48 \\
Bordentown & 14:55 & 15:33 & 16:11 & 16:49 & 17:27 & 18:05 & 18:43 & 19:21 & 19:59 & 20:37 & 21:15 & 21:53 \\
Cass Street & 15:01 & 15:39 & 16:17 & 16:55 & 17:33 & 18:11 & 18:49 & 19:27 & 20:05 & 20:43 & 21:21 & 21:59 \\
Hamilton Avenue & 15:03 & 15:41 & 16:19 & 16:57 & 17:35 & 18:13 & 18:51 & 19:29 & 20:07 & 20:45 & 21:23 & 22:01 \\
Trenton TC & 15:06 & 15:44 & 16:22 & 17:00 & 17:38 & 18:16 & 18:54 & 19:32 & 20:10 & 20:48 & 21:26 & 22:04 \\
\bottomrule\end{tabular}\end{center}
\subsection*{To Camden (read down) --- trips 1--13 of 25}
\begin{center}\footnotesize\setlength{\tabcolsep}{3pt}\begin{tabular}{@{}lrrrrrrrrrrrrr@{}}\toprule
Station & 1 & 2 & 3 & 4 & 5 & 6 & 7 & 8 & 9 & 10 & 11 & 12 & 13 \\\midrule
Trenton TC & 5:47 & 6:25 & 7:03 & 7:41 & 8:19 & 8:57 & 9:35 & 10:13 & 10:51 & 11:29 & 12:07 & 12:45 & 13:23 \\
Hamilton Avenue & 5:49 & 6:27 & 7:05 & 7:43 & 8:21 & 8:59 & 9:37 & 10:15 & 10:53 & 11:31 & 12:09 & 12:47 & 13:25 \\
Cass Street & 5:51 & 6:29 & 7:07 & 7:45 & 8:23 & 9:01 & 9:39 & 10:17 & 10:55 & 11:33 & 12:11 & 12:49 & 13:27 \\
Bordentown & 5:58 & 6:36 & 7:14 & 7:52 & 8:30 & 9:08 & 9:46 & 10:24 & 11:02 & 11:40 & 12:18 & 12:56 & 13:34 \\
Roebling & 6:03 & 6:41 & 7:19 & 7:57 & 8:35 & 9:13 & 9:51 & 10:29 & 11:07 & 11:45 & 12:23 & 13:01 & 13:39 \\
Florence & 6:06 & 6:44 & 7:22 & 8:00 & 8:38 & 9:16 & 9:54 & 10:32 & 11:10 & 11:48 & 12:26 & 13:04 & 13:42 \\
Burlington Towne Ctr & 6:11 & 6:49 & 7:27 & 8:05 & 8:43 & 9:21 & 9:59 & 10:37 & 11:15 & 11:53 & 12:31 & 13:09 & 13:47 \\
Burlington South & 6:13 & 6:51 & 7:29 & 8:07 & 8:45 & 9:23 & 10:01 & 10:39 & 11:17 & 11:55 & 12:33 & 13:11 & 13:49 \\
Beverly/Edgewater Park & 6:17 & 6:55 & 7:33 & 8:11 & 8:49 & 9:27 & 10:05 & 10:43 & 11:21 & 11:59 & 12:37 & 13:15 & 13:53 \\
Delanco & 6:21 & 6:59 & 7:37 & 8:15 & 8:53 & 9:31 & 10:09 & 10:47 & 11:25 & 12:03 & 12:41 & 13:19 & 13:57 \\
Riverside & 6:23 & 7:01 & 7:39 & 8:17 & 8:55 & 9:33 & 10:11 & 10:49 & 11:27 & 12:05 & 12:43 & 13:21 & 13:59 \\
Cinnaminson & 6:28 & 7:06 & 7:44 & 8:22 & 9:00 & 9:38 & 10:16 & 10:54 & 11:32 & 12:10 & 12:48 & 13:26 & 14:04 \\
Riverton & 6:30 & 7:08 & 7:46 & 8:24 & 9:02 & 9:40 & 10:18 & 10:56 & 11:34 & 12:12 & 12:50 & 13:28 & 14:06 \\
Palmyra & 6:32 & 7:10 & 7:48 & 8:26 & 9:04 & 9:42 & 10:20 & 10:58 & 11:36 & 12:14 & 12:52 & 13:30 & 14:08 \\
Route 73/Pennsauken & 6:35 & 7:13 & 7:51 & 8:29 & 9:07 & 9:45 & 10:23 & 11:01 & 11:39 & 12:17 & 12:55 & 13:33 & 14:11 \\
Pennsauken TC & 6:37 & 7:15 & 7:53 & 8:31 & 9:09 & 9:47 & 10:25 & 11:03 & 11:41 & 12:19 & 12:57 & 13:35 & 14:13 \\
36th Street & 6:41 & 7:19 & 7:57 & 8:35 & 9:13 & 9:51 & 10:29 & 11:07 & 11:45 & 12:23 & 13:01 & 13:39 & 14:17 \\
Walter Rand TC & 6:46 & 7:24 & 8:02 & 8:40 & 9:18 & 9:56 & 10:34 & 11:12 & 11:50 & 12:28 & 13:06 & 13:44 & 14:22 \\
Cooper St/Rutgers & 6:49 & 7:27 & 8:05 & 8:43 & 9:21 & 9:59 & 10:37 & 11:15 & 11:53 & 12:31 & 13:09 & 13:47 & 14:25 \\
Aquarium & 6:51 & 7:29 & 8:07 & 8:45 & 9:23 & 10:01 & 10:39 & 11:17 & 11:55 & 12:33 & 13:11 & 13:49 & 14:27 \\
Entertainment Center & 6:54 & 7:32 & 8:10 & 8:48 & 9:26 & 10:04 & 10:42 & 11:20 & 11:58 & 12:36 & 13:14 & 13:52 & 14:30 \\
\bottomrule\end{tabular}\end{center}
\subsection*{To Camden (read down) --- trips 14--25 of 25}
\begin{center}\footnotesize\setlength{\tabcolsep}{3pt}\begin{tabular}{@{}lrrrrrrrrrrrr@{}}\toprule
Station & 14 & 15 & 16 & 17 & 18 & 19 & 20 & 21 & 22 & 23 & 24 & 25 \\\midrule
Trenton TC & 14:01 & 14:39 & 15:17 & 15:55 & 16:33 & 17:11 & 17:49 & 18:27 & 19:05 & 19:43 & 20:21 & 20:59 \\
Hamilton Avenue & 14:03 & 14:41 & 15:19 & 15:57 & 16:35 & 17:13 & 17:51 & 18:29 & 19:07 & 19:45 & 20:23 & 21:01 \\
Cass Street & 14:05 & 14:43 & 15:21 & 15:59 & 16:37 & 17:15 & 17:53 & 18:31 & 19:09 & 19:47 & 20:25 & 21:03 \\
Bordentown & 14:12 & 14:50 & 15:28 & 16:06 & 16:44 & 17:22 & 18:00 & 18:38 & 19:16 & 19:54 & 20:32 & 21:10 \\
Roebling & 14:17 & 14:55 & 15:33 & 16:11 & 16:49 & 17:27 & 18:05 & 18:43 & 19:21 & 19:59 & 20:37 & 21:15 \\
Florence & 14:20 & 14:58 & 15:36 & 16:14 & 16:52 & 17:30 & 18:08 & 18:46 & 19:24 & 20:02 & 20:40 & 21:18 \\
Burlington Towne Ctr & 14:25 & 15:03 & 15:41 & 16:19 & 16:57 & 17:35 & 18:13 & 18:51 & 19:29 & 20:07 & 20:45 & 21:23 \\
Burlington South & 14:27 & 15:05 & 15:43 & 16:21 & 16:59 & 17:37 & 18:15 & 18:53 & 19:31 & 20:09 & 20:47 & 21:25 \\
Beverly/Edgewater Park & 14:31 & 15:09 & 15:47 & 16:25 & 17:03 & 17:41 & 18:19 & 18:57 & 19:35 & 20:13 & 20:51 & 21:29 \\
Delanco & 14:35 & 15:13 & 15:51 & 16:29 & 17:07 & 17:45 & 18:23 & 19:01 & 19:39 & 20:17 & 20:55 & 21:33 \\
Riverside & 14:37 & 15:15 & 15:53 & 16:31 & 17:09 & 17:47 & 18:25 & 19:03 & 19:41 & 20:19 & 20:57 & 21:35 \\
Cinnaminson & 14:42 & 15:20 & 15:58 & 16:36 & 17:14 & 17:52 & 18:30 & 19:08 & 19:46 & 20:24 & 21:02 & 21:40 \\
Riverton & 14:44 & 15:22 & 16:00 & 16:38 & 17:16 & 17:54 & 18:32 & 19:10 & 19:48 & 20:26 & 21:04 & 21:42 \\
Palmyra & 14:46 & 15:24 & 16:02 & 16:40 & 17:18 & 17:56 & 18:34 & 19:12 & 19:50 & 20:28 & 21:06 & 21:44 \\
Route 73/Pennsauken & 14:49 & 15:27 & 16:05 & 16:43 & 17:21 & 17:59 & 18:37 & 19:15 & 19:53 & 20:31 & 21:09 & 21:47 \\
Pennsauken TC & 14:51 & 15:29 & 16:07 & 16:45 & 17:23 & 18:01 & 18:39 & 19:17 & 19:55 & 20:33 & 21:11 & 21:49 \\
36th Street & 14:55 & 15:33 & 16:11 & 16:49 & 17:27 & 18:05 & 18:43 & 19:21 & 19:59 & 20:37 & 21:15 & 21:53 \\
Walter Rand TC & 15:00 & 15:38 & 16:16 & 16:54 & 17:32 & 18:10 & 18:48 & 19:26 & 20:04 & 20:42 & 21:20 & 21:58 \\
Cooper St/Rutgers & 15:03 & 15:41 & 16:19 & 16:57 & 17:35 & 18:13 & 18:51 & 19:29 & 20:07 & 20:45 & 21:23 & 22:01 \\
Aquarium & 15:05 & 15:43 & 16:21 & 16:59 & 17:37 & 18:15 & 18:53 & 19:31 & 20:09 & 20:47 & 21:25 & 22:03 \\
Entertainment Center & 15:08 & 15:46 & 16:24 & 17:02 & 17:40 & 18:18 & 18:56 & 19:34 & 20:12 & 20:50 & 21:28 & 22:06 \\
\bottomrule\end{tabular}\end{center}
\end{document}
